\documentclass[12pt,a4paper]{article}
% ============================================================
%  Preamble dung chung cho de cuong NCKH sinh vien - RelyFetal
%  Bien dich bang: xelatex (2 lan) de co muc luc va tham chieu cheo
% ============================================================
\usepackage{fontspec}
\usepackage{polyglossia}
\setdefaultlanguage{vietnamese}
\setotherlanguage{english}
\setmainfont{Times New Roman}[Ligatures=TeX]
\setsansfont{Arial}
\setmonofont{Consolas}[Scale=0.88]

\usepackage[a4paper,top=2.5cm,bottom=2.5cm,left=3cm,right=2cm]{geometry}
\usepackage{amsmath,amssymb}
\usepackage{booktabs}
\usepackage{longtable}
\usepackage{array}
\usepackage{tabularx}
\usepackage{multirow}
\usepackage{makecell}
\usepackage{graphicx}
\usepackage{xcolor}
\usepackage{colortbl}
\usepackage{enumitem}
\usepackage{fancyhdr}
\usepackage{titlesec}
\usepackage{caption}
\usepackage{float}
\usepackage[hidelinks,bookmarksnumbered]{hyperref}
\usepackage{bookmark}
\usepackage{setspace}
\usepackage{ragged2e}
\usepackage{tcolorbox}
\tcbuselibrary{skins,breakable}

% ---------- mau ----------
\definecolor{tdtblue}{HTML}{1F3864}
\definecolor{tdtgray}{HTML}{5A5A5A}
\definecolor{rowbg}{HTML}{EFF3F8}
\definecolor{okgreen}{HTML}{1E6B33}
\definecolor{warnred}{HTML}{9B1C1C}
\definecolor{boxbg}{HTML}{F5F7FA}

% ---------- giai bay trang ----------
\onehalfspacing
\setlength{\parindent}{1.2em}
\setlength{\parskip}{3pt}
\sloppy
\hyphenpenalty=10000
\exhyphenpenalty=10000

% ---------- tieu de muc ----------
\titleformat{\section}{\normalfont\Large\bfseries\color{tdtblue}}{\thesection.}{0.6em}{}
\titleformat{\subsection}{\normalfont\large\bfseries\color{tdtblue}}{\thesubsection.}{0.6em}{}
\titleformat{\subsubsection}{\normalfont\normalsize\bfseries}{\thesubsubsection.}{0.6em}{}
\titlespacing*{\section}{0pt}{14pt}{7pt}
\titlespacing*{\subsection}{0pt}{11pt}{5pt}
\titlespacing*{\subsubsection}{0pt}{9pt}{4pt}

% ---------- dau trang chan trang ----------
\pagestyle{fancy}
\fancyhf{}
\renewcommand{\headrulewidth}{0.4pt}
\fancyhead[L]{\footnotesize\color{tdtgray} Đề cương NCKH sinh viên --- RelyFetal}
\fancyhead[R]{\footnotesize\color{tdtgray} Trường Đại học Tôn Đức Thắng}
\fancyfoot[C]{\thepage}

% ---------- bang ----------
\renewcommand{\arraystretch}{1.28}
\captionsetup[table]{font=small,labelfont=bf,skip=5pt}
\captionsetup[figure]{font=small,labelfont=bf,skip=5pt}
\newcolumntype{L}[1]{>{\RaggedRight\arraybackslash}p{#1}}
\newcolumntype{C}[1]{>{\Centering\arraybackslash}p{#1}}
\newcolumntype{R}[1]{>{\RaggedLeft\arraybackslash}p{#1}}

% ---------- hop ghi chu ----------
\newtcolorbox{ghichu}[1][]{
  enhanced,breakable,colback=boxbg,colframe=tdtblue,
  boxrule=0.5pt,left=8pt,right=8pt,top=6pt,bottom=6pt,
  fonttitle=\bfseries\small,coltitle=white,
  attach boxed title to top left={yshift=-2mm,xshift=6mm},
  boxed title style={colback=tdtblue,boxrule=0pt,arc=1pt},#1}

\newtcolorbox{canhbao}[1][]{
  enhanced,breakable,colback=white,colframe=warnred,
  boxrule=0.7pt,left=8pt,right=8pt,top=6pt,bottom=6pt,
  fonttitle=\bfseries\small,coltitle=white,
  attach boxed title to top left={yshift=-2mm,xshift=6mm},
  boxed title style={colback=warnred,boxrule=0pt,arc=1pt},#1}

% ---------- lenh tat ----------
\newcommand{\tot}[1]{\textcolor{okgreen}{\textbf{#1}}}
\newcommand{\xau}[1]{\textcolor{warnred}{\textbf{#1}}}
\newcommand{\mh}{\textsc{FetalQRS-TCN}}
\newcommand{\hethong}{\textsc{RelyFetal}}

\fancyhead[L]{\footnotesize\color{tdtgray} Báo cáo đọc 30 công trình}
\begin{document}
\begin{titlepage}\centering\setstretch{1.0}
\vspace*{1.5cm}
{\large TRƯỜNG ĐẠI HỌC TÔN ĐỨC THẮNG}\\[1.6cm]
\rule{\textwidth}{0.8pt}\\[0.7cm]
{\LARGE\bfseries\color{tdtblue} Báo cáo đọc 30 công trình\\[6pt] về điện tim thai và phân tích dữ liệu topo\par}\vspace{0.6cm}
\rule{\textwidth}{0.8pt}\\[1.2cm]
{\large Họ làm gì, bằng phương pháp nào,\\ và mô hình của nhóm đứng ở đâu so với họ\par}
\vspace{2.2cm}
\begin{tabular}{@{}r@{\hspace{1.2em}}l@{}}
\textbf{Thuộc đề tài} & RelyFetal --- Đề cương NCKH sinh viên 2026 \\[3pt]
\textbf{Người thực hiện} & Ngô Bình Minh \\[3pt]
\textbf{Số công trình} & 30, đọc toàn văn \\[3pt]
\textbf{Ngày} & 09/09/2026 \\
\end{tabular}
\vfill
\begin{ghichu}[title={Cách đọc báo cáo này}]\footnotesize\justifying
Mỗi công trình được trình bày theo chín mục giống nhau để dễ so sánh chéo. Mục quan trọng nhất là \textbf{Giao thức đánh giá} --- đây là chỗ quyết định một con số công bố có so sánh được với con số của nhóm hay không. Nhiều công trình có F1 rất cao nhưng đạt được bằng cách chia dữ liệu ngẫu nhiên trong cùng một bệnh nhân, nới dung sai, gộp nhiều kênh, hoặc tinh chỉnh trên chính tập đích. Những chỗ bài báo \textbf{không nêu} thông tin cũng được ghi rõ, vì đó là dấu hiệu công trình không tái lập được.
\end{ghichu}
\end{titlepage}
\pagenumbering{roman}\tableofcontents\clearpage\pagenumbering{arabic}
\section{Bảng tóm tắt: ai so được với nhóm}
Trong 30 công trình, chỉ một số ít có giao thức đối chiếu trực tiếp được. Bảng này là bản đồ nhanh trước khi đi vào chi tiết từng bài.

\begin{longtable}{@{}L{1.0cm} L{3.6cm} L{6.6cm} C{3.2cm}@{}}
\toprule \textbf{Mã} & \textbf{Công trình} & \textbf{Kết quả công bố} & \textbf{So với nhóm} \\ \midrule \endfirsthead
\toprule \textbf{Mã} & \textbf{Công trình} & \textbf{Kết quả công bố} & \textbf{So với nhóm} \\ \midrule \endhead
\bottomrule \endlastfoot
\multicolumn{4}{@{}l}{\cellcolor{rowbg}\textbf{Học sâu dò QRS thai}}\\[1pt]
p01 & {\scriptsize Zhong và cs. 2018, Physiol Meas} & {\scriptsize CNN đề xuất đạt Precision 75,33\%, Recall 80,54\%, F-measure (F1) 77,85\% và Accuracy 77,38\%; so sánh accuracy bốn bộ phân lớp (Bảng 1): KNN 68,76\%, Naive Bayes 55,52\%, SVM 70,65\%, CNN 77,38\%. Bước SQI ...} & {\scriptsize\color{tdtblue}\textbf{so được một phần}} \\[2pt]
p09 & {\scriptsize Xu và cs. 2026, Sensors} & {\scriptsize Bộ phân lớp IMF đạt AUC trung bình 0,9282 ± 0,0189 qua tách theo chủ thể trên 20 đối tượng và 3 lần chạy (cao nhất r08 0,9840 ± 0,0035 và Case 12 0,9647; thấp nhất r07 0,8550 ± 0,0161 và Case 9 ...} & {\scriptsize\color{warnred}\textbf{không so trực tiếp được}} \\[2pt]
p10 & {\scriptsize Esmaeili Alidash và cs. 2025, Scientific ...} & {\scriptsize Phương án 1 (25 bản ghi, hàm MAPE, 100 epoch, ngưỡng 0,5) cho kết quả cao nhất: Accuracy 98,36 ± 0,07; Sensitivity 98,79 ± 0,13; Specificity 96,85 ± 0,43; Precision 99,09 ± 0,12; F1 96,35 ± 0,16; ...} & {\scriptsize\color{warnred}\textbf{không so trực tiếp được}} \\[2pt]
p12 & {\scriptsize Wahbah và cs. 2024, Front Physiol} & {\scriptsize Subject-dependent 5-fold (có rò rỉ): accuracy từng fold 92,5 / 93,4 / 93,5 / 93,2 / 92,6, trung bình 93,0\% với F1 = 0,96; sau khi thêm FECGPP đạt accuracy 94,2\%, F1 = 0,97 và độ nhạy tăng từ 81\% lên ...} & {\scriptsize\color{warnred}\textbf{không so trực tiếp được}} \\[2pt]
\multicolumn{4}{@{}l}{\cellcolor{rowbg}\textbf{Tách nguồn và tái tạo dạng sóng}}\\[1pt]
p03 & {\scriptsize Mohebbian và cs. 2022, IEEE JBHI} & {\scriptsize Chất lượng tái tạo tín hiệu trên A\&D FECG (Bảng I, tách theo chủ thể): R-squared trung bình 98\% [CI 95\%: 97–99\%], theo từng đối tượng R-squared = 0,99 / 0,98 / 0,96 / 0,98 / 0,98; ICC = 0,99 / 0,99 ...} & {\scriptsize\color{okgreen}\textbf{so sánh trực tiếp được}} \\[2pt]
p04 & {\scriptsize Basak và cs. 2024, Expert Syst Appl} & {\scriptsize Chất lượng tái tạo fECG theo hệ số tương quan Pearson (Bảng 1, trung bình 5 fold): ResNet 13 khối + Basic 0,882 (tốt nhất), ResNet 13 khối + Self 0,871, ResNet 9 khối + Basic 0,864, ResNet 9 khối + ...} & {\scriptsize\color{warnred}\textbf{không so trực tiếp được}} \\[2pt]
p05 & {\scriptsize Shokouhmand và Tavassolian 2023, IEEE Trans ...} & {\scriptsize Trên FECGSYNDB (tách theo chủ thể, cùng phân phối), F1 trung bình đạt 99,03\% (±0,39) và SI-SNR +14,92 dB (±1,02), tốt nhất là chủ thể 8 với F1 99,35\% và SI-SNR +15,94 dB, yếu nhất là chủ thể 2 với ...} & {\scriptsize\color{tdtblue}\textbf{so được một phần}} \\[2pt]
p06 & {\scriptsize Chen và cs. 2025, Sensors} & {\scriptsize Về chất lượng tái tạo tín hiệu, FECGSYN đạt MSE 0,096 ± 0,016 và MAE 0,082 ± 0,006 (theo kịch bản: C0 MSE 0,079 / C1 0,083 / C2 0,081 / C3 0,122 / C4 0,114 / C5 0,097), còn ADFECGDB đạt MSE 0,025 ± ...} & {\scriptsize\color{warnred}\textbf{không so trực tiếp được}} \\[2pt]
p07 & {\scriptsize Asadi và cs. 2025, IEEE Access} & {\scriptsize SCTD-2Net đánh giá zero-shot đạt Se 97,19\% / PPV 98,80\% / F1 97,97\% trên PCDB (80 kênh) và Se 95,17\% / PPV 97,94\% / F1 96,52\% trên ADFECGDB (17 kênh), vượt SCTD-1Net ở cả ba chỉ số, rõ nhất ở độ ...} & {\scriptsize\color{tdtblue}\textbf{so được một phần}} \\[2pt]
p11 & {\scriptsize Orvas và cs. 2025, arXiv preprint} & {\scriptsize Trên 100 đối tượng mô phỏng (Bảng II), CUNet đạt PRD 26,2 ± 11,6; PCC 95,9 ± 5,7; F-score 99,8 ± 0,2; SE 99,6 ± 0,4; HRerr 0,40 ± 0,81 bpm, so với UNet F-score 99,3 ± 1,6 / PCC 77,4; AttUNet 98,4 ± ...} & {\scriptsize\color{okgreen}\textbf{so sánh trực tiếp được}} \\[2pt]
\multicolumn{4}{@{}l}{\cellcolor{rowbg}\textbf{Xử lý tín hiệu cổ điển}}\\[1pt]
p16 & {\scriptsize Niknazar và cs. 2013, IEEE Trans Biomed Eng} & {\scriptsize Mọi con số định lượng đều là mức cải thiện SNR/SIR trên dữ liệu tổng hợp chứ không phải độ chính xác phát hiện đỉnh: khi hỗn hợp gần như sạch (công suất nhiễu -30 dB), mức cải thiện SNR tối thiểu ...} & {\scriptsize\color{warnred}\textbf{không so trực tiếp được}} \\[2pt]
p17 & {\scriptsize Castillo và cs. 2018, PLoS ONE} & {\scriptsize Trên ADFECGDB với dung sai 50 ms: toàn bộ 17 tín hiệu (10.839 fQRS, 10.000 TP, 839 FN, 413 FP) đạt Se 92,26\%, PPV 96,03\%, Acc 88,87\%, F1 94,11\%; riêng Nhóm 1 gồm 11 tín hiệu sạch (6.990 fQRS, 175 FN ...} & {\scriptsize\color{okgreen}\textbf{so sánh trực tiếp được}} \\[2pt]
p19 & {\scriptsize Sulas và cs. 2020, Math Biosci Eng} & {\scriptsize Trên dữ liệu thật (trung vị), độ chính xác phát hiện fQRS của SR so với MR là 0,62 so với 0,85 ở kênh ngang, 0,69 so với 0,92 ở kênh dọc, 0,61 so với 0,89 ở kênh chéo và 0,70 so với 0,86 ở kênh đơn ...} & {\scriptsize\color{warnred}\textbf{không so trực tiếp được}} \\[2pt]
\multicolumn{4}{@{}l}{\cellcolor{rowbg}\textbf{Đối chuẩn, bộ dữ liệu và thử thách}}\\[1pt]
p13 & {\scriptsize Behar 2014, Luận án Tiến sĩ} & {\scriptsize Chương 5: bSQI đơn lẻ tốt nhất với Ac = 0,969 (trên pcaSQI 0,950, sSQI 0,892, kSQI 0,879, qSQI 0,766, pSQI 0,752, basSQI 0,626), kết hợp 6 chỉ số cho Ac kiểm thử 0,985, nhưng trên MIMIC II hoàn toàn ...} & {\scriptsize\color{tdtblue}\textbf{so được một phần}} \\[2pt]
p14 & {\scriptsize Behar và cs. 2014, Physiol Meas} & {\scriptsize Trên set-a với fb = 10 Hz, thứ hạng F1 là FUSE-SMOOTH 96,0\% (Se 95,9, PPV 96,0, HRE 19,1, RRE 6,3), FUSE 95,0, TSpca-ICA 93,0, TSc-ICA 92,6, ICA-TSpca-ICA 92,4, TS-ICA 92,0, TSm-ICA 91,2, ICA-TSpca ...} & {\scriptsize\color{tdtblue}\textbf{so được một phần}} \\[2pt]
p15 & {\scriptsize Andreotti và cs. 2016, Physiol Meas} & {\scriptsize Thí nghiệm 1 (số kênh cho BSS): với 8 kênh, F1 trung bình đạt 97,46\% (BSSica-JADE), 97,40\% (BSSpca) và 97,22\% (FAST-ICA); khi Nch = 16, bước giảm chiều bằng PCA nâng F1 trung bình của ICA từ 95,3\% ...} & {\scriptsize\color{tdtblue}\textbf{so được một phần}} \\[2pt]
p20 & {\scriptsize Clifford và cs. 2014, Physiol Meas} & {\scriptsize 53 đội quốc tế tham gia, sinh ra 208 bộ nhãn và 93 bài dự thi mã nguồn mở, với điểm tốt nhất từng sự kiện là E1 = 179,439 (nhịp/phút)\textasciicircum{}2, E2 = 20,793 ms, E4 = 18,083 (nhịp/phút)\textasciicircum{}2 và E5 = 4,337 ms; ...} & {\scriptsize\color{tdtblue}\textbf{so được một phần}} \\[2pt]
p21 & {\scriptsize Matonia và cs. 2020, Scientific Data} & {\scriptsize Phần Technical Validation trích lại kết quả của chính nhóm trên bộ dữ liệu này, tất cả với dung sai ±40 ms: toàn bộ chuỗi xử lý tín hiệu bụng đạt PI tốt nhất 81,4\% nhưng khác biệt rất lớn giữa các ...} & {\scriptsize\color{tdtblue}\textbf{so được một phần}} \\[2pt]
\multicolumn{4}{@{}l}{\cellcolor{rowbg}\textbf{Lý thuyết phân tích dữ liệu topo}}\\[1pt]
p22 & {\scriptsize Perea và Harer 2015, Found Comput Math} & {\scriptsize Thí nghiệm 1 xếp hạng 10 tín hiệu theo điểm tuần hoàn: SW1PerS cho điểm từ 0,9488 (cosin thuần) giảm dần qua 0,8810 (ba tín hiệu), 0,8132, 0,7455, 0,6777, 0,6770 xuống 0,4066 (hai tín hiệu cuối), ...} & {\scriptsize\color{warnred}\textbf{không so trực tiếp được}} \\[2pt]
p23 & {\scriptsize Adams và cs. 2017, JMLR} & {\scriptsize Bảng 1 (ma trận khoảng cách 150×150, phân cụm K-medoids, độ chính xác và thời gian, mức nhiễu 0,05 và 0,1): sơ đồ gốc PD H0 với L1 đạt 96,0\% trong 37.346 s và 96,0\% trong 42.613 s, PD H0 L2 đạt ...} & {\scriptsize\color{warnred}\textbf{không so trực tiếp được}} \\[2pt]
p24 & {\scriptsize Carrière và cs. 2020, AISTATS} & {\scriptsize Bảng 1 (ORBIT, độ chính xác ± độ lệch chuẩn): trên ORBIT5K, Persistence Scale Space Kernel đạt 72,38 (±2,4), Persistence Weighted Gaussian Kernel 76,63 (±0,7), Sliced Wasserstein Kernel 83,6 (±0,9), ...} & {\scriptsize\color{warnred}\textbf{không so trực tiếp được}} \\[2pt]
p25 & {\scriptsize Som và cs. 2020, CVPRW} & {\scriptsize F1 có trọng số trên GENEactiv với khung 250 bước: MLP-PI 46,45 ± 0,32; MLP-Signal PI-Net 49,47 ± 0,69; MLP-SF 41,70 ± 0,41; CNN 1D 53,56 ± 0,31; CNN 1D + PI 54,28 ± 0,23; CNN 1D + Signal PI-Net ...} & {\scriptsize\color{warnred}\textbf{không so trực tiếp được}} \\[2pt]
p26 & {\scriptsize Turkes và cs. 2022, NeurIPS} & {\scriptsize Đếm số lỗ hổng, độ chính xác trên dữ liệu gốc: PH 1,00; PH simple 0,94; ML 0,67; NN nông 0,53; NN sâu 0,50; PointNet 1,00 — nhưng khi tịnh tiến, PointNet sụp còn 0,20 trong khi PH giữ 1,00, khi xoay ...} & {\scriptsize\color{warnred}\textbf{không so trực tiếp được}} \\[2pt]
\multicolumn{4}{@{}l}{\cellcolor{rowbg}\textbf{Nhúng trễ và chọn tham số}}\\[1pt]
p27 & {\scriptsize Trần và cs. 2019, arXiv} & {\scriptsize Độ chính xác phân loại (\%) của delay-variant so với single-delay: ECG200 90,0 so với 87,0 (1NN Euclid và 1NN DTW cùng 88,0); ECG5000 93,6 so với 92,1; ECGFiveDays 99,9 so với 92,0; TwoLeadECG 99,4 ...} & {\scriptsize\color{tdtblue}\textbf{so được một phần}} \\[2pt]
p28 & {\scriptsize Tan và cs. 2023, arXiv:2302.03447} & {\scriptsize Về phát hiện thang thời gian: với chuỗi tổng sóng sin (kỳ vọng τ = 250, 50, 8), thông tin tương hỗ chỉ tìm được thành phần tần số cao nhất, PECUZAL chỉ ra một độ trễ sai τ = 341 và MDOP ra τ = (241, ...} & {\scriptsize\color{warnred}\textbf{không so trực tiếp được}} \\[2pt]
\multicolumn{4}{@{}l}{\cellcolor{rowbg}\textbf{Topo cho tín hiệu tim}}\\[1pt]
p29 & {\scriptsize Ren và cs. 2023, Frontiers in Neuroscience} & {\scriptsize Trên bộ 12 chuyển đạo, SubLevel-Set tốt nhất với mAcc 98,04 ± 0,11 | F1 98,40 ± 0,06 | Se 97,15 ± 0,11 | Sp 98,93 ± 0,23 (Acc 97,70), còn Vietoris-Rips giảm dần theo cửa sổ từ mAcc 95,16 ± 0,20 ở 1 ...} & {\scriptsize\color{tdtblue}\textbf{so được một phần}} \\[2pt]
p30 & {\scriptsize Dindin và cs. 2019, arXiv:1906.05795} & {\scriptsize Mọi con số đều là độ chính xác có trọng số chứ không phải F1: phát hiện loạn nhịp nhị phân đạt 98\% trên tập kiểm định và 90\% trên bệnh nhân mới, còn phân loại 13 lớp đạt 97,3\% kiểm định và 80,5\% ...} & {\scriptsize\color{tdtblue}\textbf{so được một phần}} \\[2pt]
p31 & {\scriptsize Chung và cs. 2020} & {\scriptsize Bài toán Thức vs Ngủ, huấn luyện trên CGMH-training (Bảng 1): trên CGMH-validation đạt SE 70,9 ± 16,0\%, SP 78,9 ± 5,4\%, Acc 75,8 ± 4,4\%, PR 38,1 ± 19,6\%, F1 0,452 ± 0,140, AUC 0,824 ± 0,090, Kappa ...} & {\scriptsize\color{tdtblue}\textbf{so được một phần}} \\[2pt]
p32 & {\scriptsize Dominguez-Monterroza và cs. 2025, PLoS One} & {\scriptsize Chỉ số HRV cổ điển (Bảng 2, trung vị [khoảng tứ phân vị]): Mean RR tăng đơn điệu từ 404 ms [395-417] ở sơ sinh lên 674 ms [669-680] ở thiếu niên, SDNN từ 62,9 ms [57-66] ở sơ sinh xuống thấp nhất ...} & {\scriptsize\color{warnred}\textbf{không so trực tiếp được}} \\[2pt]
p33 & {\scriptsize Karan và Kaygun 2021} & {\scriptsize WESAD ba lớp (Bảng 2, tách theo chủ thể, cửa sổ 60 giây): gộp tất cả kênh với SVC và dùng tất cả sơ đồ đạt 81,35\% độ chính xác, F1 73,44\% (bản gốc WESAD 79,57\%); riêng từng kênh thì ECG 76,27\% ...} & {\scriptsize\color{tdtblue}\textbf{so được một phần}} \\[2pt]
\end{longtable}
\clearpage
\section{Học sâu dò QRS thai}
Nhóm công trình dùng mạng nơ-ron để dò trực tiếp vị trí nhịp thai. Đây là họ gần nhất với mô hình của nhóm, nên cũng là nhóm đáng so nhất.

\subsection{Zhong và cs. 2018, Physiol Meas \normalfont\small (p01)}
\noindent{\footnotesize\itshape Zhong W, Liao L, Guo X, Wang G. "A deep learning approach for fetal QRS complex detection." Physiological Measurement (IOP Publishing / IPEM), vol. 39, no. 4, art. 045004, 9 trang, 2018. DOI: 10.1088/1361-6579/aab297. Nhan bai 15/12/2017, sua 22/02/2018, chap nhan 27/02/2018, xuat ban 20/04/2018. Sun Yat-sen University, Guangzhou, Trung Quoc. Ban doc: PDF nha xuat ban (10 trang, gom 1 trang bia).}\par\vspace{2pt}
\noindent{\small\color{tdtblue}\textbf{Kết luận: so được một phần}}\par\vspace{4pt}
\noindent \textbf{Họ làm gì.} Đây là công trình đầu tiên (theo tuyên bố của chính tác giả) dùng mạng nơ-ron tích chập (CNN) để dò phức bộ QRS thai nhi. Câu hỏi khoa học chính là có thể dò fQRS trực tiếp từ tín hiệu ECG bụng thô, bỏ hẳn bước khử ECG mẹ của quy trình hai bước cổ điển (lọc thích nghi hoặc trừ mẫu rồi dò đỉnh trên phần dư) hay không; bài cũng khảo sát ảnh hưởng của hàm kích hoạt và của bước đánh giá chỉ số chất lượng tín hiệu (SQI).\par\vspace{3pt}
\noindent \textbf{Phương pháp.} Dùng set-a của PhysioNet/CinC Challenge 2013, loại 7 bản ghi có chú thích sai (a33, a38, a47, a52, a54, a71, a74) còn 68 bản ghi, rồi lọc kênh bằng Sample Entropy (m = 2, r = 1,5, N = 500; giữ kênh có SampEn trung bình < 1,5, nếu một bản ghi có ít hơn 2 kênh đạt thì lấy 2 kênh có SampEn nhỏ nhất và nhỏ nhì). Tín hiệu được cắt thành cửa sổ 100 ms (kích thước khung chọn theo dung sai chuẩn ±50 ms), gán nhãn nhị phân có/không có đỉnh R thai nhi bên trong, cân bằng lớp một cách nhân tạo bằng bước trượt 10 ms cho vùng có fQRS và 35 ms cho vùng không có fQRS, sau đó chuẩn hóa min-max về [0,1]. Dữ liệu được đưa vào CNN 1D ba lớp huấn luyện bằng SGD (lr = 10\textasciicircum{}-2, mất mát cross-entropy, chọn mô hình tốt nhất theo tập validation) và so sánh với KNN, Naive Bayes, SVM nhân RBF Gauss chạy trên MATLAB; không có bước hậu xử lý lấy đỉnh, không có thời gian trơ, không có bước khớp nhịp, đầu ra cuối cùng là nhãn lớp của từng cửa sổ 100 ms.\par\vspace{3pt}
\noindent \textbf{Mô hình.} CNN 1D nông gồm 3 khối tích chập (Conv1 → BatchNorm → ReLU; Conv2 → Dropout → MaxPool → ReLU; Conv3 → Dropout → MaxPool → ReLU) với 64/128/256 bộ lọc, độ dài bộ lọc bằng 3 cho cả ba lớp (tác giả nói đã thử bộ lọc dài hơn nhưng không thấy cải thiện đáng kể), kích thước max-pooling bằng 2, nối 3 lớp fully-connected và Softmax 2 lớp; đầu vào là chuỗi 100 mẫu đơn kênh (100 ms ở 1000 Hz). Bài không nêu số tham số huấn luyện được, bài không nêu kích thước checkpoint, bài không nêu độ trễ suy luận, bài không nêu cấu hình phần cứng dùng để huấn luyện, và bài không nêu số epoch, batch size hay hệ số dropout.\par\vspace{3pt}
\noindent \textbf{Dữ liệu.} PhysioNet/Computing in Cardiology Challenge 2013 database (PCDB), set-a: 75 bản ghi ECG bụng a01–a75, mỗi bản ghi 4 kênh, tần số lấy mẫu 1000 Hz; sau khi loại 7 bản ghi chú thích sai còn 68 bản ghi. Chia cố định thành tập kiểm thử 7 bản ghi a01–a07 (3.924 phức bộ fQRS), tập validation 6 bản ghi a08–a13 (3.348 fQRS) và tập huấn luyện 55 bản ghi còn lại (31.056 fQRS); sau khi áp dụng bước SQI, số fQRS đổi thành 26.639 (huấn luyện) / 3.082 (validation) / 3.538 (kiểm thử), ba tập được tác giả ghi rõ là loại trừ lẫn nhau, không dùng dữ liệu tổng hợp và không dùng ADFECGDB.\par\vspace{3pt}
\noindent \textbf{Giao thức đánh giá.} Chia cố định theo bản ghi một lần duy nhất, không phải tách theo bản ghi (leave-one-record-out), không phải tách theo chủ thể, cũng không phải k-fold, nên không có độ lệch chuẩn, không có khoảng tin cậy và không có kết quả từng bản ghi; điểm tốt là việc chia được thực hiện theo bản ghi nên không có rò rỉ dữ liệu giữa các bản ghi (mạnh hơn p03 và p04 ở điểm này). Đơn vị phân tích là cửa sổ 100 ms chứ không phải nhịp tim: dung sai ±50 ms chỉ được dùng để chọn kích thước khung 100 ms và hoàn toàn không được dùng để khớp nhịp dò được với nhịp thật, vì vậy TP/FP/FN của Zhong là TP/FP/FN của bài toán phân loại nhị phân cửa sổ, không phải của bài toán dò nhịp. Ngoài ra tập kiểm thử đã bị cân bằng lớp nhân tạo (bước trượt 10 ms cho lớp dương so với 35 ms cho lớp âm) nên tỷ lệ tiên nghiệm dương/âm trong tập kiểm thử không bằng tỷ lệ thực tế khi chạy trên dòng tín hiệu liên tục.\par\vspace{3pt}
\noindent \textbf{Kết quả chính.} CNN đề xuất đạt Precision 75,33\%, Recall 80,54\%, F-measure (F1) 77,85\% và Accuracy 77,38\%; so sánh accuracy bốn bộ phân lớp (Bảng 1): KNN 68,76\%, Naive Bayes 55,52\%, SVM 70,65\%, CNN 77,38\%. Bước SQI chỉ nâng accuracy từ 76,52\% lên 77,38\%, tức hơn 0,86 điểm (Bảng 2), trong khi hàm kích hoạt ảnh hưởng mạnh hơn nhiều: Sigmoid 52,37\%, Tanh 71,92\%, ReLU 77,38\% (Bảng 3). Về lọc, việc dùng lọc chắn dải kết hợp Butterworth dải thông 0,5–100 Hz (ghép một lọc thông thấp bậc 2 và một lọc thông cao bậc 2, lọc pha không) chỉ "cải thiện nhẹ" accuracy của CNN; bài không đưa con số cụ thể cho phần này mà chỉ vẽ ở Hình 5.\par\vspace{3pt}
\noindent \textbf{Điểm yếu.} Nghiêm trọng nhất, chỉ số F1 77,85\% không phải là F1 dò nhịp theo chuẩn CinC 2013 mà là F1 của bài toán phân loại cửa sổ 100 ms trên tập đã cân bằng lớp nhân tạo, trong khi toàn bộ văn liệu về sau (kể cả Mohebbian 2022 và Basak 2024 trong bảng so sánh của họ) đều trích con số này như F1 dò fQRS — một lỗi trích dẫn hệ thống của cả lĩnh vực. Không có hậu xử lý lấy đỉnh và không áp thời gian trơ nên một phức bộ QRS thật có thể sinh ra nhiều cửa sổ dương liên tiếp mà không bị tính là FP; chỉ có một lần chia cố định với 7 bản ghi kiểm thử, không lặp lại, không có độ lệch chuẩn, không có kiểm định thống kê; và bài không báo cáo bất kỳ chi phí tính toán nào (số tham số, độ trễ, kích thước mô hình, phần cứng) nên không tái lập được về mặt kiến trúc chính xác. Cửa sổ ngữ cảnh chỉ 100 ms, ngắn hơn rất nhiều chu kỳ tim thai (khoảng 430 ms) nên mạng không thể nhìn thấy thông tin nhịp điệu; dải lọc 0,5–100 Hz là dải rộng nhất nên giữ lại nhiều sóng T mẹ và nhiễu cơ; tuyên bố "đơn kênh" chỉ đúng ở mức đầu vào của từng cửa sổ vì thực tế mỗi bản ghi đóng góp ít nhất 2 kênh vào tập dữ liệu, còn bài không nêu rõ lúc kiểm thử lấy kết quả của một kênh hay gộp nhiều kênh của cùng một bản ghi; và bước SQI chỉ mang lại 0,86 điểm, gần như không đáng kể, nhưng bài vẫn kết luận là tốt hơn.\par\vspace{3pt}
\noindent \textbf{\color{tdtblue}So với mô hình của nhóm.} Không so sánh trực tiếp được ở mức con số vì đơn vị phân tích khác hẳn: Zhong tính F1 trên cửa sổ 100 ms đã cân bằng lớp, còn FetalQRS-TCN tính F1 trên nhịp tim với khớp ±50 ms; dù vậy bài vẫn là mốc đối chiếu quan trọng theo ba hướng. Về kiến trúc, Zhong dùng CNN 1D ba lớp với bộ lọc dài 3 nên tổng trường tiếp nhận chỉ vài chục mili giây trên đầu vào 100 ms, trong khi mô hình của nhóm có trường tiếp nhận 1.516 ms và khảo sát trường tiếp nhận của nhóm (172 ms → 98,73; 748 ms → 99,25; 1.516 ms → 99,50; Wilcoxon p = 0,0000) cho thấy ngữ cảnh 100 ms nằm dưới cả điểm thấp nhất nhóm đã đo; về dải lọc, dải 0,5–100 Hz mà họ dùng chính là dải cho kết quả thấp nhất trong khảo sát 8 dải của nhóm (72,17 so với 91,06 của dải 10–60 Hz, chênh 18,89 điểm), đây là bằng chứng trực tiếp và rất mạnh cho phát hiện cốt lõi của nhóm rằng công trình học sâu nền tảng của lĩnh vực đã chọn dải lọc dưới tối ưu. Về quy mô thì không so được vì Zhong không công bố số tham số; suy luận (không phải dữ kiện) là với 3 lớp tích chập 64/128/256 bộ lọc và 3 lớp fully-connected trên vector độ dài 100, mô hình của Zhong gần như chắc chắn lớn hơn 113.481 tham số của nhóm nhiều lần do khối dense, nhưng không thể khẳng định vì bài không nêu kích thước các lớp dense.\par\vspace{3pt}
\noindent \textbf{\color{tdtblue}Nhóm rút ra gì.} Áp dụng: đưa Zhong 2018 vào bài báo của nhóm như bằng chứng nguồn gốc cho luận điểm "tiền xử lý quan trọng hơn kiến trúc" (họ dùng 0,5–100 Hz, dải mà khảo sát của nhóm cho thấy mất 18,89 điểm F1 so với 10–60 Hz), học cách họ liệt kê rõ bản ghi của từng tập (a01–a07 kiểm thử, a08–a13 validation, 55 bản ghi còn lại huấn luyện) để nhóm cũng liệt kê bản ghi của từng fold tách theo bản ghi, và coi ý tưởng dùng SampEn làm chỉ số chất lượng tín hiệu là tiền lệ trực tiếp cho hướng fSQI mới nhưng phải nhắm mục tiêu cao hơn nhiều vì SQI của Zhong chỉ thêm 0,86 điểm accuracy và không có quyền từ chối trả lời, chỉ dùng để lọc kênh lúc chuẩn bị dữ liệu. Tránh: tuyệt đối không dùng khung phân loại cửa sổ nhị phân mà giữ nguyên sơ đồ sequence-to-sequence + heatmap Gaussian + lấy đỉnh + khớp nhịp ±50 ms, không cân bằng lớp bằng cách đổi bước trượt rồi báo cáo precision trên tập đã cân bằng mà phải báo cáo trên dòng tín hiệu liên tục thật, và trong bảng so sánh phải ghi chú rõ 77,85\% của Zhong là F1 mức cửa sổ chứ không phải F1 khớp nhịp, nếu không nhóm sẽ lặp lại lỗi trích dẫn của cả lĩnh vực và có thể bị phản biện ở vòng review Q1.\par\vspace{3pt}
\vspace{6pt}
\subsection{Xu và cs. 2026, Sensors \normalfont\small (p09)}
\noindent{\footnotesize\itshape Xu, Xiaojian; Zhang, Yifan; Rao, Yufei; Xu, Yinru; Gao, Yang; Tu, Huating. "A Deep Learning-Guided Ensemble Empirical Mode Decomposition Method for Single-Channel Fetal Electrocardiogram Extraction." Sensors (MDPI), 2026, tap 26, so 7, bai so 2037. DOI: 10.3390/s26072037. Nhan bai 29/01/2026, chap nhan 23/03/2026, xuat ban 25/03/2026. Bien tap vien: Pawel Plawiak. Open access CC BY. 24 trang.}\par\vspace{2pt}
\noindent{\small\color{warnred}\textbf{Kết luận: không so trực tiếp được}}\par\vspace{4pt}
\noindent \textbf{Họ làm gì.} Bài giải bài toán tách tín hiệu ECG thai nhi từ ECG bụng đơn kênh (không phải dò đỉnh trực tiếp): EEMD phân rã tín hiệu thành 8 thành phần IMF, còn việc chọn IMF nào chứa thông tin thai nhi - trước đây làm thủ công hoặc bằng luật cứng theo ngưỡng năng lượng và tần số nên rất kém ổn định - được thay bằng một CNN 1 chiều (IMFClassifier1D) chấm điểm xác suất "liên quan đến thai nhi" cho từng IMF. Toàn bộ đường ống, gọi là CNN-2xEEMD, gồm: lọc, khử ECG mẹ bằng mẫu trung bình, EEMD lần 1, CNN chọn 2 IMF điểm cao nhất, EEMD lần 2 trên phần dư, CNN chọn 1 IMF điểm cao nhất, rồi lọc 10-60 Hz cho kết quả cuối.\par\vspace{3pt}
\noindent \textbf{Phương pháp.} Tiền xử lý: đồng bộ tần số lấy mẫu về 1000 Hz (dữ liệu mô phỏng 500 Hz được nâng bằng lọc polyphase, ADFECGDB vốn đã 1000 Hz), lọc notch IIR 50 Hz và Butterworth bậc 4 dải thông 0,1-100 Hz kiểu zero-phase; riêng dữ liệu mô phỏng được tăng nhiễu EMG (nhiễu trắng Gauss lọc 20-150 Hz, biên độ bằng 0,1 lần độ lệch chuẩn tín hiệu gốc) và trôi đường nền (tổng hai sin 0,2 Hz và 0,33 Hz, biên độ 0,18 lần độ lệch chuẩn) để ép SNR về đúng 10 dB. Khử ECG mẹ: dò đỉnh R mẹ bằng NeuroKit2 với phương án dự phòng là ngưỡng biên độ kèm ràng buộc RR tối thiểu 500 ms (tương ứng nhịp mẹ tối đa 120 bpm), dựng mẫu mẹ bằng trung bình đồng pha rồi trừ khỏi tín hiệu gốc; EEMD lần 1 chạy trên phần dư với ensemble 30, noise\_width 0,2, tối đa 8 IMF, cửa sổ 10 giây (10.000 mẫu ở 1000 Hz), CNN cộng 2 IMF điểm cao nhất thành ước lượng fECG ban đầu, sau đó lấy hiệu giữa bản tối ưu và bản ban đầu làm phần dư cho EEMD lần 2 và CNN chọn thêm 1 IMF, đầu ra cuối qua Butterworth bậc 4 zero-phase 10-60 Hz. Nhãn huấn luyện cho CNN sinh tự động theo quy tắc (công thức 3): một IMF được gán nhãn 1 nếu biên độ đỉnh-đỉnh cục bộ lớn nhất trong cửa sổ ±50 mẫu quanh đỉnh R thai vượt đồng thời hai ngưỡng 1,4σ và 0,75 lần biên độ đỉnh-đỉnh toàn cục, lấy trung vị qua các nhịp để chống nhiễu xung.\par\vspace{3pt}
\noindent \textbf{Mô hình.} IMFClassifier1D là CNN 1 chiều nhận đầu vào 1 kênh dài 10.000 mẫu (10 giây ở 1000 Hz), gồm 5 khối conv-BatchNorm-ReLU-pooling với kích thước nhân giảm dần [31, 21, 15, 11, 7] và số kênh đặc trưng hình kim tự tháp [32, 64, 128, 256, 512] (conv đầu tiên stride 2, max pooling cho lớp 1-4 và adaptive average pooling cho lớp 5), tiếp theo là 3 lớp fully-connected thu về 1024, 256 rồi 1 logit; tổng cộng 5,91 triệu tham số, 0,08 GFLOPs, độ trễ suy luận 2,24 ms. Bài đối chiếu với Transformer (d=128, 2 lớp: 5,90 M / 12,50 GFLOPs / 35 ms / AUC 0,9380) và Bi-LSTM (h=128, 2 lớp: 2,10 M / 1,25 GFLOPs / 12 ms / AUC 0,9320).\par\vspace{3pt}
\noindent \textbf{Dữ liệu.} Ba nguồn: 15 ca mô phỏng sinh bởi máy mô phỏng phần cứng H2-3000F (Xuzhou Mingsheng, Trung Quốc) với nhịp mẹ 60-120 bpm, nhịp thai 110-160 bpm, tỷ lệ biên độ thai/mẹ từ 1:1,17 đến 1:16 và SNR ép về 10 dB; 5 bản ghi lâm sàng ADFECGDB (PhysioNet, 1000 Hz: r01, r04, r07, r08, r10) chỉ lấy một kênh bụng làm đầu vào, còn điện đạo thai trực tiếp chỉ dùng để sinh nhãn huấn luyện chứ không dùng khi suy luận; hai tập kiểm thử ngoài là DaISy (KU Leuven, 250 Hz, mỗi bản ghi chỉ dài 10 giây, 5 kênh bụng, không có chú thích đỉnh R thai sẵn có nên nhóm tác giả tự chọn kênh 2 và kênh 3 sau "phân tích sơ bộ", tự đánh dấu từng nhịp bằng tay và loại bỏ các mẫu mơ hồ) và NIFEA DB (2 ca NR14 và ARR01). Tổng cộng 1.505 mẫu IMF từ 20 đối tượng, trong đó 744 dương tính (49,44\%).\par\vspace{3pt}
\noindent \textbf{Giao thức đánh giá.} Tách theo chủ thể chỉ được áp cho bộ phân lớp IMF chứ không cho bước dò đỉnh R đầu-cuối: 20 đối tượng (15 mô phỏng và 5 lâm sàng), mỗi fold giữ 1 đối tượng làm validation và huấn luyện trên 19 đối tượng còn lại, 60 epoch, chạy 3 lần với seed khác nhau cho ma trận AUC 3 × 20, chỉ số báo cáo là AUC ở mức IMF. Toàn bộ các con số F1 dò đỉnh R thai (0,8372-0,9565) đo trên DaISy và NIFEA là tập ngoài, không có kiểm định chéo nào, chỉ là 4 trường hợp đơn lẻ; điểm cộng là họ không tinh chỉnh trên hai tập này. Dung sai khớp nhịp: bài không nêu - toàn văn không có câu nào định nghĩa một đỉnh được coi là đúng nếu nằm trong bao nhiêu ms, con số ±50 mẫu chỉ xuất hiện trong quy tắc gán nhãn IMF chứ không phải tiêu chí khớp đỉnh khi tính F1; công thức F1/Precision/Recall và cách tính cột "Accuracy" ở Bảng 8 cũng không được nêu.\par\vspace{3pt}
\noindent \textbf{Kết quả chính.} Bộ phân lớp IMF đạt AUC trung bình 0,9282 ± 0,0189 qua tách theo chủ thể trên 20 đối tượng và 3 lần chạy (cao nhất r08 0,9840 ± 0,0035 và Case 12 0,9647; thấp nhất r07 0,8550 ± 0,0161 và Case 9 0,8840), còn cấu hình mạng tốt nhất trong khảo sát loại bỏ đạt AUC 0,9445 (Bảng 4, 5,91 M tham số). Dò đỉnh R thai đầu-cuối (Bảng 7 và Bảng 8): DaISy kênh 2 đạt Precision 0,9167 / Recall 1,0000 / F1 0,9565 / CC 0,96 / SNRimp 15,88 dB; DaISy kênh 3 đạt 0,9524 / 0,9091 / F1 0,9302 / CC 0,94 / SNRimp 13,39 dB; NIFEA NR14 đạt 0,9130 / 0,8750 / F1 0,8936 và ARR01 đạt 1,0000 / 0,7200 / F1 0,8372. So ba chiến lược chọn IMF trên DaISy, luật cứng Auto-2xEEMD chỉ đạt F1 0,8293 và 0,7000 trong khi CNN-2xEEMD đạt 0,9565 và 0,9302, ngang với chọn thủ công (Wilcoxon bắt cặp: CNN so với Auto p = 0,002 và p = 0,003; CNN so với Manual p = 0,147 và p = 0,089, tức không khác biệt có ý nghĩa); cấu hình chọn IMF N1=2, N2=1 cho CC 0,94, F1 0,93 và SNRimp 14,8 dB trên 15 ca mô phỏng, tốt hơn N1=1 (CC 0,81, p=0,001) và N1=3 (CC 0,85, p=0,002).\par\vspace{3pt}
\noindent \textbf{Điểm yếu.} Lỗ hổng tái lập nghiêm trọng nhất là bài không nêu dung sai khớp nhịp tính bằng ms, khiến mọi con số F1 không thể đối chiếu với bất kỳ công trình nào khác; bài cũng không viết công thức F1, cột "Accuracy" ở Bảng 8 ghi 0,0172 và 0,0152 cho DaISy là vô nghĩa và không định nghĩa được, còn Hình 11 ghi độ rộng QRS 28 ms với sóng T chỉ cách S 16 ms - theo suy luận của tôi là bất khả thi về sinh lý (QRS thai bình thường khoảng 50-60 ms) nên rất có thể đó là dao động của bộ lọc 10-60 Hz bị đọc nhầm thành sóng P và T. Cỡ mẫu đánh giá quá nhỏ: mỗi bản ghi DaISy chỉ dài 10 giây, từ Recall 1,0000 và Precision 0,9167 suy ra TP = 22 và FP = 2, tức F1 0,9565 được tính trên khoảng 22 nhịp tim nên chỉ một nhịp sai đã làm F1 đổi 4 điểm; tệ hơn, nhãn chuẩn trên DaISy do chính nhóm tác giả tự đánh dấu bằng mắt rồi loại bỏ các mẫu mơ hồ (tức loại đúng những nhịp khó, làm F1 tăng giả tạo), còn kênh 2 và kênh 3 được chọn sau khi đã nhìn dữ liệu vì chúng có cạnh lên/xuống sắc hơn và RR ổn định hơn. Tách theo chủ thể chỉ áp cho AUC mức IMF chứ không cho F1 dò đỉnh nên hai con số này không liên quan trực tiếp, 15/20 đối tượng huấn luyện là mô phỏng từ máy phần cứng và chỉ 5 đối tượng là thật, con số 0,08 GFLOPs và 2,24 ms chỉ tính riêng CNN chứ không tính EEMD (30 ensemble × 8 IMF × 2 giai đoạn trên cửa sổ 10 giây, phần đắt nhất của cả đường ống) nên tuyên bố chạy được trên STM32H7 Cortex-M7 chưa được chứng minh (bài tự thừa nhận chưa thử trên phần cứng thật), và Bảng 6 so sánh với 1D-CycleGAN, R2W-Net, DSS chỉ lấy số từ bài gốc chứ không chạy lại trên cùng giao thức.\par\vspace{3pt}
\noindent \textbf{\color{tdtblue}So với mô hình của nhóm.} Không so sánh trực tiếp được: họ có dùng ADFECGDB nhưng chỉ để tính AUC mức IMF (0,9282) và không bao giờ báo cáo F1 dò đỉnh R thai trên bộ này, nên không có con số nào đối chiếu được với macro F1 97,43 (tách theo bản ghi, 4 kênh) hay 99,18 (kênh tốt nhất) của mô hình của nhóm; hơn nữa họ không nêu dung sai ms nên dù có số cũng không biết họ đo ở ±50 ms hay rộng hơn. Về quy mô đánh giá và chi phí tính toán thì mô hình của nhóm mạnh hơn hẳn: nhóm đánh giá tách theo bản ghi 20 lần trên 5 sản phụ với hàng chục nghìn nhịp, còn họ đo F1 trên 4 trường hợp đơn lẻ mà 2 trường hợp chỉ dài 10 giây (khoảng 22 nhịp); nhóm dùng 113.481 tham số / 0,48 MB / 4,35 ms cho cửa sổ 4 giây, còn họ dùng 5,91 triệu tham số (gấp 52 lần) và con số 2,24 ms của họ chưa tính phần EEMD đắt đỏ. Họ hơn nhóm ở hai điểm: có thử nghiệm trên NIFEA (loạn nhịp thai) mà nhóm chưa có, và có so sánh thống kê với baseline thủ công.\par\vspace{3pt}
\noindent \textbf{\color{tdtblue}Nhóm rút ra gì.} Nên trích dẫn bài này làm bằng chứng độc lập cho luận điểm "tiền xử lý quan trọng hơn kiến trúc": họ khởi đầu bằng dải thông rộng 0,1-100 Hz nhưng đầu ra cuối cùng lại là Butterworth bậc 4 zero-phase 10-60 Hz, đúng dải mà nhóm chọn sau khi quét 8 dải (10-60 Hz thắng với 91,06), và họ cũng ghi nhận năng lượng QRS thai tập trung khoảng 8-45 Hz - khác biệt là một bài Sensors 2026 tìm ra dải này bằng tay còn nhóm chứng minh định lượng bằng khảo sát 8 dải có kiểm định; đồng thời nên học cách họ báo cáo tham số / GFLOPs / độ trễ và cách họ dùng Wilcoxon bắt cặp (giống việc nhóm so RF 1516 ms với 172 ms, p = 0,0000), nhưng phải tính cả chi phí tiền xử lý khử ECG mẹ chứ không chỉ chi phí mạng. Cần tránh ba lỗi của họ: báo cáo F1 trên đoạn 10 giây (khoảng 22 nhịp) rồi gọi đó là tập kiểm thử độc lập; tự đánh dấu nhãn chuẩn bằng mắt rồi loại bỏ các nhịp mơ hồ (đây là rò rỉ ngược, phải giữ nguyên mọi nhịp); và quên nêu dung sai ±50 ms trong phần phương pháp - chính điều này khiến kết quả của họ không thể đối chiếu được.\par\vspace{3pt}
\vspace{6pt}
\subsection{Esmaeili Alidash và cs. 2025, Scientific Reports \normalfont\small (p10)}
\noindent{\footnotesize\itshape Esmaeili Alidash, Kimia; Danandeh Hesar, Hamed. "Deep learning-based approach for accurate detection of fetal QRS complexes in abdominal ECG signals." Scientific Reports (Nature Portfolio), 2025, tap 15, bai so 39260. DOI: 10.1038/s41598-025-22999-9. Department of Biomedical Engineering, Sahand University of Technology, Tabriz, Iran. Nhan bai 15/12/2024, chap nhan 03/10/2025. Open access CC ...}\par\vspace{2pt}
\noindent{\small\color{warnred}\textbf{Kết luận: không so trực tiếp được}}\par\vspace{4pt}
\noindent \textbf{Họ làm gì.} Bài dò QRS thai nhi trực tiếp từ ECG bụng mà không khử ECG mẹ, và không coi đây là bài toán tìm vị trí đỉnh mà là phân lớp nhị phân từng ô 100 ms: tín hiệu được cắt thành cửa sổ 1 giây, mỗi cửa sổ chia thành 10 ô 100 ms, mạng xuất 10 bit với bit thứ k bằng 1 nếu ô thứ k có một QRS thai. Họ biện minh cho con số 100 ms bằng lập luận rằng chuẩn cổ điển coi một đỉnh là đúng nếu nằm trong ±50 ms so với chú thích nên một khung 100 ms là tương đương, rồi chạy 5 "phương án" khác nhau về số bản ghi, bước tiền xử lý, hàm mất mát và tỷ lệ chia dữ liệu, báo cáo phương án tốt nhất và nhấn mạnh rằng "chỉ cần 20 tín hiệu để huấn luyện".\par\vspace{3pt}
\noindent \textbf{Phương pháp.} Loại bỏ 15 trong 75 bản ghi của Set A - 7 bản ghi do chú thích sai (a33, a38, a47, a52, a54, a71, a74) và 8 bản ghi do thiếu dữ liệu (a01, a02, a07, a09, a11, a16, a18, a43) - còn lại 60 bản ghi. Tiền xử lý gồm khử trôi đường nền bằng Butterworth thông cao bậc 2 tần số cắt 5 Hz chạy cả hai chiều (cách thứ hai là ước lượng đường nền bằng Butterworth thông thấp bậc 1 tần số cắt 3,17 Hz rồi trừ đi), khử nhiễu xung bằng lọc trung vị cửa sổ 60 ms (một chỗ khác ghi 0,26 giây) rồi thay đoạn hỏng vượt ngưỡng bằng giá trị ước lượng, lọc notch khử nhiễu lưới điện (bài không nêu tần số notch là 50 hay 60 Hz), và chuẩn hóa về [-1, 1] hoặc [0, 1] với min/max tính riêng cho từng tín hiệu; riêng phương án 4 dùng thêm DWT với wavelet Biorthogonal Bior 4.4, phân rã 10 mức và chỉ giữ hệ số chi tiết mức 1-4 để tái tạo. Tín hiệu được cắt thành đoạn 1 giây với bước trượt 100 ms (chồng lấp 90\%, dùng làm tăng cường dữ liệu), mỗi ô 100 ms gán nhãn 0 hoặc 1 theo có hay không có QRS thai, huấn luyện bằng Adam với learning rate 0,00003, hàm mất mát MSE hoặc MAPE (còn thử Huber, LogCosh, MSLE, Cosine similarity), 30-100 epoch, ngưỡng trên đầu ra sigmoid là 0,4 hoặc 0,5, và mỗi kết quả là trung bình của 5 lần chạy độc lập.\par\vspace{3pt}
\noindent \textbf{Mô hình.} CNN 1 chiều nhận đầu vào 4 (kênh) × 1000 (mẫu), tức đoạn 1 giây ở 1 kHz, gồm 5 lớp tích chập (lớp 1 và 2 dùng stride 2), 7 lớp batch normalization, 3 lớp dropout và 3 lớp fully-connected, kích hoạt ReLU sau conv và sigmoid ở các lớp dense, đầu ra 10 neuron ứng với 10 dự đoán nhị phân cho 10 ô 100 ms; số tham số là 31.407.738 (ghi rõ trong Bảng 13). Bài không nêu kích thước nhân, số kênh từng lớp, tỷ lệ dropout, batch size, kích thước mô hình (MB), FLOPs, độ trễ suy luận, và cũng không nêu trường tiếp nhận.\par\vspace{3pt}
\noindent \textbf{Dữ liệu.} PhysioNet/Computing in Cardiology Challenge 2013, Set A (bài gọi nhầm là "NI-FECGDB", link dữ liệu họ cung cấp là physionet.org/content/challenge-2013/1.0.0): mỗi bản ghi dài 1 phút, 4 kênh ECG bụng, lấy mẫu 1 kHz tức 60.000 mẫu mỗi kênh, chú thích đỉnh R thai lấy từ điện cực da đầu thai nhi do ban tổ chức Challenge cung cấp, và mỗi bản ghi sinh 600 dự đoán (60.000 mẫu chia cho 100 mẫu mỗi ô). Từ 75 bản ghi Set A loại 15 còn 60; chế độ 1 dùng 25 bản ghi (20 huấn luyện, 5 kiểm thử), chế độ 2 dùng 60 bản ghi (48 huấn luyện, 12 kiểm thử).\par\vspace{3pt}
\noindent \textbf{Giao thức đánh giá.} Không phải tách theo chủ thể cũng không phải tách theo bản ghi, mà là chia ngẫu nhiên 80-20\% (hoặc 33-67\% ở phương án 3), và bài tự mâu thuẫn: Bảng 10 ghi "25 signals (20 training, 5 testing)" và "60 signals (48 training, 12 testing)" - đọc theo hướng này thì chia ở mức bản ghi, chấp nhận được - nhưng trang 13 (mục "Novel approach for labeling data") lại viết "prior to partitioning the data into test and validation sets, the data is first shuffled, and then split before being fed into the proposed network", tức xáo trộn rồi mới chia, mà kết hợp với cửa sổ chồng lấp 90\% thì các cửa sổ của cùng một bệnh nhân nằm ở cả tập huấn luyện lẫn tập kiểm thử, tức rò rỉ dữ liệu nghiêm trọng, và bài không làm rõ được điều này. Họ còn gọi lẫn lộn tập thứ hai là "test and validation sets" và Bảng 14 ghi "final validation accuracy", nên rất có thể con số báo cáo đo trên tập validation vốn dùng để chọn epoch, ngưỡng và hàm mất mát chứ không phải tập kiểm thử độc lập. Về dung sai, họ viện dẫn chuẩn ±50 ms (trích dẫn Zhong và cs. 2018) rồi áp đặt cơ chế khung 100 ms, nhưng đơn vị phân tích không phải là nhịp mà là ô 100 ms và các chỉ số bao gồm cả ô âm, nên accuracy/specificity của họ không thể so với F1 đo theo nhịp của bất kỳ bài nào; họ chỉ dùng một bộ dữ liệu duy nhất nên không có bước tinh chỉnh trên tập đích.\par\vspace{3pt}
\noindent \textbf{Kết quả chính.} Phương án 1 (25 bản ghi, hàm MAPE, 100 epoch, ngưỡng 0,5) cho kết quả cao nhất: Accuracy 98,36 ± 0,07; Sensitivity 98,79 ± 0,13; Specificity 96,85 ± 0,43; Precision 99,09 ± 0,12; F1 96,35 ± 0,16; MSE 0,01 ± 0. Phương án 5 (60 bản ghi, MSE, 100 epoch, ngưỡng 0,4) - con số được đưa lên tóm tắt - đạt Accuracy 96,79 ± 0,87; Sensitivity 97,91 ± 1,01; Specificity 92,79 ± 1,50; Precision 97,88 ± 0,43; F1 93,02 ± 1,36; còn phương án 2 đạt Acc 92,56 / F1 83,67, phương án 3 (chia 33-67\%) đạt Acc 94,54 / F1 87,54 / Se 96,65 / Precision 96,25, và phương án 4 (DWT) đạt Acc 95,69 / F1 90,98. Bảng 13 đối chiếu với các công trình trước trên cùng bộ dữ liệu (Zhong 2018 Acc 77,38 / F1 77,85; Zhong 2019 RCED-Net F1 93,62; Vo 2020 OctConv F1 91,1; Ting 2021 Acc 95,20; Ghonchi 2022 dual-attention Acc 96,26 / F1 98,02; Lee 2022 W-net F1 96,91; Krupa 2022 Acc 91,32 / F1 90,12), còn Bảng 14 cho thấy learning rate 0,00003 đạt 96,79\%, 0,0001 đạt 95,12\%, 0,0003 đạt 93,45\% và 0,001 không hội tụ.\par\vspace{3pt}
\noindent \textbf{Điểm yếu.} Mâu thuẫn số học nghiêm trọng (suy luận của tôi, đã kiểm chứng bằng tính toán): từ Sensitivity, Specificity và Precision họ báo cáo, giải ngược ra tỷ lệ lớp dương là 77,6\% ở phương án 1 và 77,3\% ở phương án 5, tính lại Accuracy từ tỷ lệ đó cho 0,9836 và 0,9675 - khớp chính xác với con số bài ghi (0,9836 và 0,9679) - nhưng với 10 ô mỗi giây và nhịp thai 110-160 bpm thì chỉ khoảng 20-27\% số ô có QRS thai, không thể là 77\%, nên các chỉ số này được tính với lớp dương là ô KHÔNG có QRS thai, tức "Sensitivity 97,91\%" trong tóm tắt thực chất là độ nhạy của việc nhận ra ô không có nhịp thai, vô nghĩa về mặt lâm sàng; kiểm chứng thêm, F1 báo cáo (96,35 và 93,02) không bằng trung bình điều hòa của Precision và Sensitivity họ ghi (phải là 98,94 và 97,89) nhưng khớp chính xác nếu tính cho lớp thiểu số, nên con số có ý nghĩa thật sự của bài là F1 96,35\% (25 bản ghi) và F1 93,02\% (60 bản ghi). Giao thức có rò rỉ dữ liệu tiềm tàng (cửa sổ chồng lấp 90\% cộng với câu "xáo trộn rồi mới chia" - nếu xáo trộn ở mức cửa sổ thì mọi con số trong bài đều vô giá trị) và không có tập kiểm thử độc lập rõ ràng trong khi epoch, ngưỡng và hàm mất mát đều được chọn theo chính con số báo cáo; ngoài ra đầu vào là 4 (kênh) × 1000 (mẫu) nên đây là mô hình đa kênh chứ không phải đơn kênh, mà Bảng 1 của chính họ lại phân loại các bài khác thành đơn kênh/đa kênh khiến người đọc dễ nhầm. Mô hình có 31.407.738 tham số nhưng được gọi là "lightweight", "low-complexity", "suitable for real-time deployment" hàng chục lần mà không kèm FLOPs, độ trễ hay kích thước file; bài cũng không nêu tần số notch, batch size, kích thước nhân và số kênh từng lớp nên không tái lập được, ghi nhầm năm của chính họ là "2023" trong một bài xuất bản 2025, trộn lẫn hai phương án trong tóm tắt (96,79\% là của phương án 5 dùng 48 bản ghi huấn luyện, còn "chỉ 20 tín hiệu" là phương án 1 đạt 98,36\%), và chỉ có độ lệch qua 5 lần khởi tạo chứ không có độ lệch theo bệnh nhân nên không biết mô hình ổn định thế nào giữa các sản phụ.\par\vspace{3pt}
\noindent \textbf{\color{tdtblue}So với mô hình của nhóm.} Không so sánh trực tiếp được, và mô hình của nhóm mạnh hơn rõ rệt ở mọi trục quan trọng: đơn vị phân tích khác hẳn (họ đánh giá trên ô 100 ms có tính cả ô âm, nhóm đánh giá trên nhịp với dung sai ±50 ms, nên 96,79\% accuracy của họ và macro F1 97,43 của nhóm là hai đại lượng khác loại, đặt cạnh nhau là nguy hiểm), họ dùng 4 kênh đồng thời còn nhóm dùng 1 kênh nên bài toán của nhóm khó hơn hẳn, họ chia ngẫu nhiên có xáo trộn với cửa sổ chồng lấp 90\% còn nhóm dùng tách theo bản ghi nghiêm ngặt, và họ cần 31.407.738 tham số mà không báo cáo độ trễ còn nhóm chỉ cần 113.481 tham số / 0,48 MB / 4,35 ms. Về bộ dữ liệu, họ đánh giá trong bộ CinC 2013 (chia ngẫu nhiên trên chính bộ đó, có khả năng rò rỉ) còn nhóm báo cáo CinC 2013 là kết quả xuyên bộ dữ liệu không tinh chỉnh (77,34 ± 28,54), tức con số của nhóm là zero-shot còn con số 93-96\% của họ là in-domain, hai thứ này không so được. Nếu bị chất vấn vì sao nhóm chỉ đạt 77,34 trên CinC 2013 trong khi Esmaeili Alidash đạt 96,79, câu trả lời ngắn gọn là: họ huấn luyện trên chính CinC 2013 còn nhóm thì không, họ dùng 4 kênh còn nhóm dùng 1 kênh, và họ đếm ô 100 ms còn nhóm đếm nhịp.\par\vspace{3pt}
\noindent \textbf{\color{tdtblue}Nhóm rút ra gì.} Bài này là ví dụ để nhóm viết một đoạn về "tại sao giao thức đánh giá quan trọng hơn con số" trong phần Related Work hoặc Discussion, nêu đích danh rằng một bài Scientific Reports 2025 báo cáo 98,36\% accuracy nhưng đầu vào là 4 kênh chứ không phải 1 kênh, cửa sổ chồng lấp 90\% rồi xáo trộn trước khi chia, và đơn vị là ô 100 ms chứ không phải nhịp; đồng thời nên chỉ rõ rằng dùng khung 100 ms để xấp xỉ dung sai ±50 ms là một lỗi thiết kế, vì một QRS nằm sát biên ô sẽ rơi vào ô sai dù chỉ lệch 1 ms trong khi một QRS lệch 49 ms vẫn được tính đúng nếu cùng ô, còn cách của nhóm (heatmap Gaussian sigma 12 ms, lấy đỉnh rồi khớp ±50 ms trên trục thời gian liên tục) là đúng về nguyên tắc. Nhóm bắt buộc phải nêu rõ đơn vị phân tích, lớp nào là lớp dương và công thức F1 (chính vì bài này không làm vậy mà số liệu của họ không đọc được), không bao giờ xáo trộn cửa sổ chồng lấp rồi mới chia train/test mà giữ nguyên tách theo bản ghi và tách theo chủ thể, không trộn lẫn số liệu của hai thí nghiệm khác nhau trong một câu ở tóm tắt, và chỉ dùng từ "nhẹ" khi kèm số tham số và độ trễ - với 113.481 tham số nhóm hoàn toàn có quyền dùng từ đó kèm bảng so sánh trực tiếp với 31,4 triệu tham số của bài này.\par\vspace{3pt}
\vspace{6pt}
\subsection{Wahbah và cs. 2024, Front Physiol \normalfont\small (p12)}
\noindent{\footnotesize\itshape Wahbah M, Zitouni MS, Al Sakaji R, Funamoto K, Widatalla N, Krishnan A, Kimura Y, Khandoker AH (2024). "A deep learning framework for noninvasive fetal ECG signal extraction." Frontiers in Physiology, 15:1329313. DOI: 10.3389/fphys.2024.1329313. Nhan 28/10/2023, chap nhan 22/03/2024, xuat ban 22/04/2024. Giay phep CC BY. File PDF 13 trang, doc toan van.}\par\vspace{2pt}
\noindent{\small\color{warnred}\textbf{Kết luận: không so trực tiếp được}}\par\vspace{4pt}
\noindent \textbf{Họ làm gì.} Bài báo xây dựng một khung học sâu phát hiện trực tiếp đỉnh R của thai nhi từ tín hiệu ECG bụng 12 kênh, không tách fECG ra trước. Tín hiệu bụng được cắt thành các đoạn 65 ms rồi phân loại nhị phân từng đoạn: Class 1 nếu đoạn chứa một đỉnh R thai nhi, Class 0 nếu không.\par\vspace{3pt}
\noindent \textbf{Phương pháp.} Tín hiệu bụng lưỡng cực từ 12 điện cực được thu ở fs = 1 kHz, độ phân giải 16 bit, lọc dải thông 0,05–100 Hz kèm notch 50 Hz hoặc 60 Hz, mỗi sản phụ chỉ lấy 1 phút. Nhãn chuẩn được tạo theo phương pháp FRPL: dùng thuật toán BSSR (Sato 2007) tách fECG khỏi hỗn hợp bụng rồi chạy một chương trình MATLAB R2022b tự viết để dò đỉnh R; tín hiệu được chuẩn hoá z-score bằng mean/std của tập huấn luyện rồi đưa vào mạng BiLSTM với trọng số lớp w1 = 3,17 và w0 = 0,59 để bù mất cân bằng. Hai khối hậu xử lý FECGPP (giảm báo động giả) và FHRE (làm sạch nhịp bằng SD-ROM kết hợp lọc thích nghi trên cửa sổ 10 giây) hoàn tất pipeline trước khi tính nhịp tim thai và đối chiếu bằng tương quan Pearson và Bland-Altman.\par\vspace{3pt}
\noindent \textbf{Mô hình.} Kiến trúc gồm 3 lớp BiLSTM, 2 lớp dropout (0,5) xen giữa, 1 lớp fully connected 2 đầu ra, softmax và lớp phân loại, với đầu vào chuỗi 12 đặc trưng tương ứng 12 kênh bụng. Số tham số tổng: bài không nêu; ngoài ra phần văn bản ghi 100/50/20 hidden units trong khi Table 1 ngụ ý 200/100/50 nên hai chỗ mâu thuẫn nhau, còn kích thước checkpoint, độ trễ suy luận và phần cứng huấn luyện đều bài không nêu.\par\vspace{3pt}
\noindent \textbf{Dữ liệu.} 70 sản phụ có thai không ghi nhận bất thường, từ ba cơ sở: Children's National Hospital (Hoa Kỳ, 12 mẫu = 17\%), Kanagawa Children's Medical Center (Nhật, 13 mẫu = 19\%) và Tohoku University Hospital (Nhật, 45 mẫu = 64\%), với 12 kênh bụng, fs = 1 kHz, 16 bit, mỗi người chỉ 1 phút tín hiệu; nhân khẩu học của 26 đối tượng: tuổi thai 30,3 ± 5,6 tuần, cân nặng thai 1.518 ± 794 g, tuổi mẹ 31,9 ± 5,9, BMI 23,9 ± 4,4 kg/m². Dữ liệu chia 75\% huấn luyện / 25\% kiểm thử, nhưng kiểm thử tách theo chủ thể chỉ thực hiện trên 26 đối tượng vì 44 người còn lại không được công bố thông tin theo IRB nên luôn nằm trong tập huấn luyện; dữ liệu bị ràng buộc bởi IRB Tohoku và IRB Children's National nên không công khai và không tái lập được.\par\vspace{3pt}
\noindent \textbf{Giao thức đánh giá.} Bài chạy hai giao thức và báo cáo cả hai: (A) subject-dependent với 5-fold cross-validation, là rò rỉ dữ liệu vì cùng một bệnh nhân xuất hiện ở cả tập huấn luyện lẫn kiểm thử — bài tự nhận điều này qua cách đặt tên nhưng vẫn đưa con số 94,2\% lên Abstract và Conclusion; (B) tách theo chủ thể trên 26 đối tượng, là giao thức sạch. Dung sai khớp nhịp: bài không nêu và thực chất không có — bài không dùng cơ chế khớp đỉnh với dung sai ±50 ms theo chuẩn CinC 2013 mà biến bài toán thành phân loại nhị phân đoạn 65 ms, tức một bài toán khác và dễ hơn vì không đòi hỏi đỉnh nằm đúng vị trí. Đơn vị phân tích là đoạn 65 ms, metric chính là accuracy trên tập mất cân bằng, nhãn chuẩn là đầu ra thuật toán BSSR chứ không phải ECG da đầu thai nhi, và bài không đánh giá xuyên bộ dữ liệu nào.\par\vspace{3pt}
\noindent \textbf{Kết quả chính.} Subject-dependent 5-fold (có rò rỉ): accuracy từng fold 92,5 / 93,4 / 93,5 / 93,2 / 92,6, trung bình 93,0\% với F1 = 0,96; sau khi thêm FECGPP đạt accuracy 94,2\%, F1 = 0,97 và độ nhạy tăng từ 81\% lên 85,8\%. Tách theo chủ thể trên 26 đối tượng: accuracy trung bình 88,8 ± 6,4, thấp nhất là đối tượng \#5 (68,5\%, BMI 33,0, có biến chứng mẹ) và cao nhất là \#16 (97,1\%); phân tầng cho 86,7\% ở giai đoạn hình thành vernix caseosa 28–34 tuần, 88\% khi tuổi thai < 30 tuần, 93\% khi > 34 tuần, 90\% ở 16 đối tượng mẹ khoẻ mạnh BMI bình thường, 85\% ở 5 đối tượng BMI > 26 kg/m², 87\% khi có biến chứng mẹ và 91\% khi không. Nhịp tim thai chỉ đạt tương quan Pearson r = 0,56 (p < 0,05) với Bland-Altman bias = −5,49 (bài viết sai đơn vị là "weeks", đúng ra phải là bpm) và giới hạn đồng thuận 18,05 / −29,03, tức biên độ khoảng 47 bpm; số tham số, độ trễ và kích thước mô hình đều bài không nêu.\par\vspace{3pt}
\noindent \textbf{Điểm yếu.} Mô hình dùng đồng thời 12 kênh bụng nên không thể so với mô hình đơn kênh, và nhãn chuẩn là đầu ra thuật toán BSSR chứ không phải ECG da đầu, nghĩa là mô hình chỉ học bắt chước một thuật toán khác mà không ai biết BSSR chính xác đến đâu — đây là điểm yếu nghiêm trọng nhất về mặt khoa học. Bài không có dung sai khớp nhịp, dùng accuracy trên tập mất cân bằng làm metric chính (từ w1 = 3,17 suy ngược ra lớp đỉnh R chỉ chiếm khoảng 15,8\% số đoạn), và F1 = 0,96/0,97 chắc chắn không phải F1 của lớp đỉnh R vì với Se = 0,858 thì giải 2·Se·PPV/(Se+PPV) = 0,97 cho PPV = 1,116 > 1, vô lý — nên con số này không được phép đem so với F1 97,43 của nhóm. Ngoài ra còn rò rỉ dữ liệu ở thí nghiệm 5-fold, mâu thuẫn kiến trúc giữa văn bản và Table 1, chú thích Table 3 ngược nghĩa, sai đơn vị Bland-Altman, dải lọc 0,05–100 Hz rộng nhất mà không hề khảo sát, chi phí tính toán bài không nêu, chỉ 1 phút tín hiệu mỗi người, và dữ liệu riêng tư bị khoá bởi IRB nên không tái lập được.\par\vspace{3pt}
\noindent \textbf{\color{tdtblue}So với mô hình của nhóm.} Không so sánh trực tiếp được vì bốn lý do độc lập: họ dùng 12 kênh đồng thời còn nhóm dùng 1 kênh; họ phân loại nhị phân đoạn 65 ms còn nhóm định vị đỉnh với dung sai ±50 ms nên bài toán của nhóm khó hơn; nhãn chuẩn của họ là đầu ra thuật toán BSSR còn nhóm dùng ECG da đầu của ADFECGDB là tiêu chuẩn vàng thực sự; metric của họ là accuracy trên tập mất cân bằng chứ không phải macro F1 của lớp đỉnh. Nếu buộc phải đặt cạnh nhau, kết quả tách theo chủ thể 88,8 ± 6,4 của họ (accuracy, 12 kênh, đoạn 65 ms, nhãn chuẩn thuật toán) so với 97,43 ± 4,37 tách theo bản ghi của nhóm (macro F1, 1 kênh, ±50 ms, nhãn chuẩn ECG da đầu) cho thấy mô hình của nhóm ở vị thế mạnh hơn về mọi mặt, nhưng tuyệt đối không được viết "chúng tôi vượt 88,8\% của Wahbah 2024" trong bài báo vì đó là so sánh sai thang đo và reviewer Q1 sẽ bắt lỗi ngay. Điểm duy nhất họ hơn nhóm là cỡ mẫu lớn hơn nhiều (26 đối tượng so với 5 sản phụ ADFECGDB) và trải rộng tuổi thai 20–39 tuần trong khi nhóm chỉ có 38–41 tuần chuyển dạ.\par\vspace{3pt}
\noindent \textbf{\color{tdtblue}Nhóm rút ra gì.} Nên sao chép cách phân tầng kết quả của họ (theo tuổi thai, BMI mẹ, có hay không biến chứng mẹ), bổ sung phân tích Bland-Altman kèm tương quan Pearson cho nhịp tim thai vì nhóm sẽ dễ dàng vượt xa r = 0,56, thêm hậu xử lý lọc nhịp kiểu SD-ROM lên chuỗi RR, và làm nổi bật việc nhóm báo cáo 113.481 tham số, 0,48 MB, 4,35 ms trong khi bài Q1 năm 2024 này không báo cáo gì về chi phí tính toán. Phải tránh việc dùng accuracy làm metric chính trên bài toán mất cân bằng — luôn báo cáo Se, PPV, F1 của lớp đỉnh R kèm dung sai ±50 ms ghi rõ ràng — không đặt con số có rò rỉ lên Abstract, và nếu trích dẫn bài này thì phải ghi kèm "12 kênh, phân loại đoạn 65 ms, nhãn chuẩn từ thuật toán BSSR" ngay bên cạnh con số.\par\vspace{3pt}
\vspace{6pt}
\clearpage
\section{Tách nguồn và tái tạo dạng sóng}
Nhóm công trình đặt bài toán khác: thay vì dò vị trí nhịp, họ tái tạo lại toàn bộ dạng sóng điện tim thai rồi mới dò. Cách này cho thêm thông tin hình thái nhưng khó hơn và tốn hơn.

\subsection{Mohebbian và cs. 2022, IEEE JBHI \normalfont\small (p03)}
\noindent{\footnotesize\itshape Mohebbian MR, Vedaei SS, Wahid KA, Dinh A, Marateb HR, Tavakolian K. "Fetal ECG Extraction From Maternal ECG Using Attention-Based CycleGAN." IEEE Journal of Biomedical and Health Informatics (JBHI), vol. 26, no. 2, pp. 515-526, 2022. DOI: 10.1109/JBHI.2021.3111873. University of Saskatchewan (Canada), University of Isfahan (Iran), UPC BarcelonaTech (Tay Ban Nha), University of North Dakota (Hoa ...}\par\vspace{2pt}
\noindent{\small\color{okgreen}\textbf{Kết luận: so sánh trực tiếp được}}\par\vspace{4pt}
\noindent \textbf{Họ làm gì.} Bài coi việc tách ECG thai nhi khỏi ECG bụng mẹ là bài toán chuyển đổi tín hiệu sang tín hiệu, dùng CycleGAN một chiều với hai bộ sinh (một ánh xạ ECG bụng → ECG thai nhi, một ánh xạ ngược lại) cùng mất mát nhất quán chu kỳ; hai điểm mới là một lớp chú ý (attention) đóng vai trò mặt nạ che vùng tín hiệu, xuất phát từ quan sát rằng nhân tích chập khi khuếch đại sóng R thai nhi thì khuếch đại luôn sóng R mẹ, và việc dùng hàm kích hoạt sin thay LeakyReLU vì sin giữ chi tiết tín hiệu tốt hơn. Mục tiêu chính của họ là tái tạo toàn bộ dạng sóng fECG (cả sóng P, sóng T và đoạn ST), còn dò fQRS chỉ là một ứng dụng phụ để chứng minh năng lực, điều chính tác giả nói rõ ở phần Discussion.\par\vspace{3pt}
\noindent \textbf{Phương pháp.} Cắt bỏ 10 giây đầu và 10 giây cuối mỗi bản ghi (nhiễu nhất), hạ tần số lấy mẫu về 200 Hz, lọc dải thông 1–100 Hz rồi làm mượt Savitzky-Golay, cắt cửa sổ chữ nhật độ dài M = 200 mẫu (1,0 giây) và chuẩn hóa z-score từng cửa sổ. Bộ sinh gồm lớp MultiHead self-attention (2 đầu, số chiều nhúng bằng kích thước cửa sổ) → chuẩn hóa z-score đầu ra attention → cắt ngưỡng [0,8; 1] để tạo mặt nạ thay cho softmax thông thường → nhân từng phần tử với tín hiệu gốc → 4 lớp Conv1D với hàm kích hoạt sin, huấn luyện 50 epoch, mất mát log(cosh), tối ưu Nesterov Adam, hệ số lambda cho mất mát chu kỳ bằng 4. Sau khi tái tạo fECG, vị trí R được lấy bằng bộ dò đỉnh Pan \& Tompkins áp lên tín hiệu tái tạo; điểm cực kỳ quan trọng là với hai bộ NI-FECG và CinC Challenge (chỉ có chú thích vị trí R chứ không có dạng sóng fECG thật), tác giả tự sinh tín hiệu fECG đích bằng wavelet Daubechies dựa trên chính chuỗi R-R đáp án, nghĩa là tín hiệu mục tiêu mà mạng được dạy để tạo ra chính là tín hiệu dựng từ đáp án.\par\vspace{3pt}
\noindent \textbf{Mô hình.} CycleGAN 1D có chú ý: bộ sinh gồm MultiHead self-attention (2 đầu, số chiều nhúng bằng kích thước cửa sổ) → chuẩn hóa và cắt ngưỡng [0,8; 1] → nhân phần tử → Conv1D 13 bộ lọc kernel 3, Conv1D 7 bộ lọc kernel 5, Conv1D 5 bộ lọc kernel 13, Conv1D 1 bộ lọc kernel 3, tất cả với hàm sin và các lớp cuối dùng stride 2 cùng kernel lớn để làm mượt; bộ phân biệt gồm Conv1D 12/7/3/1 bộ lọc kernel 3 với sin rồi Dense 32 và softmax, ngoài ra có một lớp Conv1D (1 bộ lọc, kernel 31) ở đường tiền xử lý. Bài không nêu tổng số tham số huấn luyện được, bài không nêu kích thước mô hình (MB) và bài không nêu độ trễ cho một cửa sổ đơn lẻ; chỉ nêu thời gian huấn luyện tổng 10.400 giây (208 giây/epoch × 50 epoch) và thời gian xử lý 1.000 mẫu là 0,14 giây trên máy Core-i9, RAM 16 GB, GPU NVIDIA GeForce GTX 1060 6 GB, cùng phân tích độ nhạy cho thấy lambda tối ưu nằm trong khoảng 2 < lambda < 10 và việc bỏ một lớp Conv1D sau attention (từ 4 xuống 3 lớp) làm R-squared trung bình giảm 5\%.\par\vspace{3pt}
\noindent \textbf{Dữ liệu.} Bốn nguồn: A\&D FECG (PhysioNet Abdominal and Direct FECG Database, chính là ADFECGDB nhóm đang dùng) với 5 sản phụ, 5 bản ghi (r01, r07, r10, r04, r08), tuổi thai 38–41 tuần, mỗi bản ghi 5 phút, 4 kênh bụng và 1 kênh fECG da đầu, 1000 Hz, 16 bit, đã lọc 0–100 Hz và lọc điện lưới lúc thu; NI-FECG database (nifecgdb) với 55 bản ghi ECG bụng đa kênh, 1000 Hz, 16 bit, lọc 0–100 Hz, mà chính bài ghi rõ đều lấy từ một sản phụ duy nhất ở tuần thai 21–40 với vị trí điện cực thay đổi giữa các lần đo, trong đó họ chọn 14 bản ghi (154, 192, 244, 274, 290, 323, 368, 444, 597, 733, 746, 811, 826, 906). Ngoài ra có CinC Challenge 2013 set-a với 75 bản ghi, 1000 Hz, 4 kênh, loại a33, a38, a52, a54, a71, a74 còn 69 bản ghi (lưu ý họ không loại a47, khác với Zhong 2018), và dữ liệu tổng hợp bằng bộ công cụ FECGSYN gồm 24 cặp ECG mẹ - thai với nhịp thai 115/125/135/145/155/160 bpm và nhịp mẹ 65/77/89/100 bpm.\par\vspace{3pt}
\noindent \textbf{Giao thức đánh giá.} Bài dùng ba giao thức có chất lượng rất khác nhau: trên A\&D FECG là tách theo chủ thể thật sự, lặp 5 lần cho 5 đối tượng — giao thức sạch và là giao thức duy nhất trong bài có thể tin cậy hoàn toàn; trên NI-FECG và CinC Challenge là "4-fold cross-validation" nhưng bài không nêu cách chia fold theo bản ghi hay theo cửa sổ ngẫu nhiên, không nêu danh sách bản ghi của từng fold, không nêu có đảm bảo các cửa sổ của cùng một bản ghi không nằm cả ở tập huấn luyện lẫn tập kiểm thử hay không; và ngoại kiểm là huấn luyện trên toàn bộ A\&D FECG rồi kiểm thử trên NI-FECG (14 tín hiệu) và CinC Challenge (69 tín hiệu), đây là giao thức chặt chẽ nhất trong bài và là con số nên dùng để so sánh. Riêng với NI-FECG, vì cả 14 bản ghi đều từ cùng một sản phụ theo đúng mô tả của chính bài, 4-fold cross-validation trên đó chắc chắn là rò rỉ dữ liệu theo chủ thể: mô hình được huấn luyện và kiểm thử trên cùng một người. Dung sai khớp nhịp là 30 ms (6 mẫu ở 200 Hz), chặt hơn chuẩn CinC 50 ms và đây là điểm cộng cho họ; bộ dò đỉnh là Pan \& Tompkins áp lên fECG tái tạo; đơn vị phân tích là nhịp tim (đúng chuẩn); có báo cáo khoảng tin cậy 95\% theo chuẩn STARD và TRIPOD; Bảng III ghi "1 Abdominal" cho phương pháp đề xuất nhưng bài không nêu chọn kênh bụng nào và theo tiêu chí gì.\par\vspace{3pt}
\noindent \textbf{Kết quả chính.} Chất lượng tái tạo tín hiệu trên A\&D FECG (Bảng I, tách theo chủ thể): R-squared trung bình 98\% [CI 95\%: 97–99\%], theo từng đối tượng R-squared = 0,99 / 0,98 / 0,96 / 0,98 / 0,98; ICC = 0,99 / 0,99 / 0,98 / 0,99 / 0,99; WEDD = 7,0\% / 5,7\% / 9,2\% / 6,3\% / 7,2\% (xếp loại từ "very good" đến "good"); kiểm định t bắt cặp không có khác biệt có ý nghĩa (p > 0,1). Dò fQRS (Bảng II): A\&D FECG (tách theo chủ thể, 5 tín hiệu) F1 = 99,7 ± 0,4\% [CI 95\%: 97,8–99,9], PPV 99,6 ± 0,7\%, Se 99,4 ± 0,6\%; NI-FECG (4-fold, 14 tín hiệu) F1 = 99,6 ± 1,3\% [98,2–99,9]; CinC Challenge (4-fold, 69 tín hiệu) F1 = 99,3 ± 1,3\% [95,3–99,9], PPV 99,1 ± 1,0\%, Se 99,2 ± 0,6\%; ngoại kiểm khi huấn luyện trên A\&D FECG cho F1 = 97,9\% [96,5–98,4] trên NI-FECG và F1 = 94,7\% [92,6–96,5] với PPV 93,9\%, Se 93,9\% trên CinC Challenge. Nghiên cứu loại bỏ cho thấy attention là thành phần quyết định: bỏ attention và thay bằng Conv1D + LeakyReLU thì R-squared tụt xuống 0,79 / 0,80 / 0,77 / 0,82 / 0,81, còn thay LeakyReLU bằng sin mà vẫn không có attention chỉ lên 0,80 / 0,81 / 0,78 / 0,83 / 0,81, tức attention đóng góp khoảng 17 điểm R-squared còn hàm sin chỉ khoảng 1 điểm; các bản ghi CinC kém nhất là a29, a38, a40, a42, a52, a53, a54, a56, a61, a65, và khi thêm nhiễu EMG thì R-squared của đối tượng 1 giảm dần từ khoảng 100\% ở 40 dB xuống khoảng 70\% ở 10 dB.\par\vspace{3pt}
\noindent \textbf{Điểm yếu.} Nghiêm trọng nhất, với NI-FECG và CinC Challenge, tín hiệu fECG đích dùng để huấn luyện được tự sinh từ chuỗi R-R đáp án bằng wavelet Daubechies: mạng được dạy để tạo ra một tín hiệu mà các đỉnh nằm đúng vị trí đáp án, sau đó lại lấy đỉnh trên chính tín hiệu đó bằng Pan-Tompkins rồi so với đáp án, nên nếu có bất kỳ rò rỉ dữ liệu nào giữa fold huấn luyện và fold kiểm thử (điều bài không chứng minh là không có) thì con số 99,3\% sẽ hoàn toàn vô nghĩa. Bài không nêu cách chia 4 fold theo bản ghi hay theo cửa sổ ngẫu nhiên — đây là thông tin bắt buộc phải có, và việc thiếu nó làm các con số 99,3\% và 99,6\% không tái lập được và không kiểm chứng được; bộ NI-FECG chỉ có một sản phụ nên 4-fold cross-validation trên đó chắc chắn là rò rỉ theo chủ thể và 99,6\% không nói lên khả năng tổng quát hóa; bài cũng không nêu chọn kênh bụng nào trong 4 kênh, nếu chọn hậu nghiệm theo kết quả tốt nhất thì con số bị thổi phồng. Ngoài ra bài không nêu tổng số tham số, không nêu kích thước mô hình và không nêu độ trễ một cửa sổ (chính tác giả thừa nhận ở phần hạn chế rằng độ phức tạp mô hình là gánh nặng cho hệ thống nhúng), bộ A\&D FECG chỉ có 5 đối tượng nên cỡ mẫu quá nhỏ để tuyên bố tổng quát hóa, và phương pháp bắt buộc phải có fECG da đầu làm nhãn huấn luyện (hoặc phải giả lập nó) — một ràng buộc lâm sàng rất nặng vì phần lớn dữ liệu thực chỉ có chú thích vị trí R.\par\vspace{3pt}
\noindent \textbf{\color{tdtblue}So với mô hình của nhóm.} Trên cùng bộ ADFECGDB và cùng kiểu giao thức bỏ một đối tượng, họ đạt F1 99,7 ± 0,4\% ở dung sai 30 ms, còn nhóm đạt 99,18\% (Se 99,59, PPV 98,79) trên kênh lâm sàng tốt nhất và 97,43 ± 4,37\% trên toàn bộ 4 kênh ở dung sai 50 ms; vì mỗi bản ghi ADFECGDB là một sản phụ nên tách theo bản ghi của nhóm tương đương tách theo chủ thể của họ, do đó họ nhỉnh hơn nhóm (khoảng 0,5 điểm trên kênh tốt nhất, khoảng 2,3 điểm trên trung bình 4 kênh) và còn ở dung sai chặt hơn, nên nhóm không được tuyên bố vượt họ trên bộ này; điểm giảm nhẹ duy nhất là họ không nói rõ chọn kênh nào, còn nhóm báo cáo cả 20 đánh giá tách theo bản ghi trên đủ 4 kênh. Ngược lại, khi chuyển sang CinC 2013 không tinh chỉnh, họ đạt 94,7\% [92,6–96,5] trên 69 bản ghi trong khi nhóm chỉ đạt 77,34 ± 28,54 \% (lưỡng cực: 4 bản ghi F1 = 100, 3 bản ghi 24–49\%) — đây là khoảng cách thật khoảng 17 điểm mà nhóm phải thừa nhận và giải quyết, không thể viện cớ dung sai vì dung sai của họ (30 ms) chặt hơn của nhóm (50 ms); giả thuyết chưa kiểm chứng về nguyên nhân gồm việc họ huấn luyện mạng tái tạo dạng sóng fECG đầy đủ từ ECG da đầu nên biểu diễn giàu thông tin hình thái hơn và phụ thuộc biên độ hơn, họ hạ tần số về 200 Hz và dùng cửa sổ 1 giây khác cấu hình của nhóm, và việc bài không nêu cách chọn kênh trên CinC khiến một phần khoảng cách có thể đến từ chọn kênh thuận lợi. Về chi phí thì không so được số tham số vì bài không nêu; về thông lượng, họ xử lý 1.000 mẫu (tức 5 giây tín hiệu ở 200 Hz) trong 0,14 giây trên GPU GTX 1060, khoảng 36 lần thời gian thực, còn nhóm xử lý 4 giây tín hiệu trong 4,35 ms trên CPU, khoảng 920 lần thời gian thực — suy luận là mô hình của nhóm nhanh hơn khoảng 25 lần và còn chạy trên CPU thay vì GPU, đây là lợi thế rõ ràng cần nhấn mạnh nhưng phải ghi chú rằng hai phép đo thực hiện trên hai hệ thống khác nhau.\par\vspace{3pt}
\noindent \textbf{\color{tdtblue}Nhóm rút ra gì.} Áp dụng ngay: bổ sung báo cáo kết quả ở dung sai 30 ms bên cạnh 50 ms cho tất cả thực nghiệm để đặt cạnh trực tiếp con số 99,7\% và 94,7\% của họ (nếu chỉ báo cáo 50 ms, phản biện Q1 sẽ nói nhóm chọn dung sai để đẹp số), bổ sung khoảng tin cậy 95\% cho macro F1 theo chuẩn STARD/TRIPOD thay vì chỉ có trung bình ± độ lệch chuẩn, đưa con số 94,7\% xuyên bộ dữ liệu của họ vào bảng so sánh và giải thích thẳng thắn vì sao 77,34 \% của nhóm thấp hơn, thử ý tưởng chú ý làm mặt nạ với cắt ngưỡng [0,8; 1] vì thành phần này mang lại gần như toàn bộ hiệu quả của họ (+17 điểm R-squared) với chi phí tham số rất thấp, và làm thực nghiệm nhiễu EMG ở 40/30/20/10 dB vì liên quan trực tiếp đến hướng fSQI mới. Tránh tuyệt đối: không bao giờ tạo tín hiệu đích tổng hợp từ chính đáp án rồi huấn luyện và kiểm thử trên cùng bộ dữ liệu (lỗ hổng lớn nhất của bài này), phải viết rõ cách chia fold, danh sách bản ghi từng fold và câu khẳng định "không có cửa sổ nào của bản ghi kiểm thử xuất hiện trong tập huấn luyện" — chính vì thiếu câu này mà con số 99,3\% của họ không thể tin cậy hoàn toàn, và nếu trích dẫn bài này thì phải phân biệt rõ ba giao thức (tách theo chủ thể / 4-fold không rõ cách chia / ngoại kiểm) thay vì trích chung một con số 99,7\%.\par\vspace{3pt}
\vspace{6pt}
\subsection{Basak và cs. 2024, Expert Syst Appl \normalfont\small (p04)}
\noindent{\footnotesize\itshape Basak P, Sakib AHMN, Chowdhury MEH, Al-Emadi N, Yalcin HC, Pedersen S, Mahmud S, Kiranyaz S, Al-Maadeed S. "A novel deep learning technique for morphology preserved fetal ECG extraction using 1D-CycleGAN." Expert Systems with Applications (Elsevier), vol. 235, art. 121196, 2024. DOI: 10.1016/j.eswa.2023.121196. University of Dhaka (Bangladesh) va Qatar University (Qatar). Tai tro: Qatar ...}\par\vspace{2pt}
\noindent{\small\color{warnred}\textbf{Kết luận: không so trực tiếp được}}\par\vspace{4pt}
\noindent \textbf{Họ làm gì.} Bài cũng dùng CycleGAN một chiều để chuyển ECG bụng mẹ thành ECG thai nhi, nhưng mục tiêu khác hẳn các bài chỉ dò đỉnh: họ muốn giữ nguyên hình thái tín hiệu thai nhi, tức giữ được cả sóng P, sóng T, đoạn ST, khoảng PR và QT, vì những thành phần này mới cho phép chẩn đoán bệnh tim thai chứ không chỉ nhịp tim. Đóng góp chính gồm một hàm mất mát có trọng số bốn thành phần (mất mát L1 đối kháng, mất mát phổ, mất mát thời gian và mất mát công suất), một khối tiền xử lý rất công phu để khử trôi đường cơ sở, khử trễ pha và khử sóng dư ở hai đầu đoạn, cùng việc gộp hai bộ dữ liệu thực khác nhau để mô hình bớt phụ thuộc vào thiết bị đo và vị trí điện cực; dò fQRS và ước lượng biến thiên nhịp tim (HRV) là các ứng dụng phía sau, chạy trên tín hiệu đã tái tạo.\par\vspace{3pt}
\noindent \textbf{Phương pháp.} Lấy mẫu lại tất cả tín hiệu (cả mECG và fECG) về 512 Hz, lọc chắn dải quanh 50 Hz để khử nhiễu điện lưới, lọc thông dải Butterworth bậc cao với tần số cắt 0,1–70 Hz, khử trôi đường cơ sở bằng cách khớp một đa thức vào đường cơ sở rồi trừ đi (bậc 8 cho mECG và bậc 36 cho fECG, chọn bằng thử sai), và lọc trung bình trượt cửa sổ 4 mẫu cho mECG, 10 mẫu cho fECG để làm mượt. Tiếp đó cắt đoạn cửa sổ 512 mẫu (1,0 giây), sửa lại đường cơ sở cục bộ cho từng đoạn, chuẩn hóa min-max về [0,1] từng đoạn, khử trễ pha bằng kỹ thuật "áp dụng - lật - áp dụng lại" (chính là lọc pha không) và khử sóng dư đầu/cuối bằng cửa sổ chồng lấn, tạo ma trận N × M × C với C = 4 cho mECG và C = 1 cho fECG. Huấn luyện CycleGAN 1D với bộ sinh nhận đầu vào 4 kênh và tạo ra fECG 1 kênh; sau khi tái tạo, dò fQRS bằng thuật toán Engelse \& Zeelenberg (EngZee) sửa đổi và tính các chỉ số HRV gồm muy\_RR, muy\_HR, sigma\_HR, RMSSD, SDNN/RMSSD và PNN50.\par\vspace{3pt}
\noindent \textbf{Mô hình.} CycleGAN 1D với bộ sinh kiểu ResNet: cấu hình tốt nhất là ResNet 13 khối gồm 3 lớp Conv1D hạ mẫu kích thước 16, 32, 64 (mỗi lớp kèm BatchNorm + Dense + ReLU) → n khối ResNet (mỗi khối 2 đơn vị, mỗi đơn vị gồm ReflectionPad2d, 64 lớp Conv1d và BatchNorm, xen giữa có một lớp fully-connected và một lớp dropout, cộng với kết nối tắt) → 3 lớp ConvTranspose1D nâng mẫu kích thước 16, 32, 64, kết hợp bộ phân biệt PatchGAN 3×3 "Basic"; các biến thể được so sánh gồm bộ sinh Unet 128, Unet 256, ResNet 9 khối, ResNet 13 khối, Self-FPN và bộ phân biệt Basic (PatchGAN 3×3) hoặc Self (PatchGAN 1×1 dùng lớp 1D Self-ONN). Số tham số được nêu rõ: 0,337 triệu tham số huấn luyện được cho mỗi bộ sinh, 0,044 triệu cho mỗi bộ phân biệt, tổng 0,762 triệu cho toàn mạng, với tốc độ học 0,00001 (tốt nhất trong ba mức 1e-4, 1e-5, 1e-6), 150 epoch (100 epoch thiếu khớp, 300 epoch quá khớp) và trọng số mất mát p = 2 (phổ), q = 4 (thời gian), r = 1 (công suất); bài không nêu kích thước checkpoint tính bằng MB và bài không nêu độ trễ cho một cửa sổ 1 giây, chỉ nêu thông lượng tổng.\par\vspace{3pt}
\noindent \textbf{Dữ liệu.} Hai bộ dữ liệu thực được gộp lại: bộ 1 là Abdominal and Direct fECG Database của PhysioNet (chính là ADFECGDB) với 5 sản phụ, 5 bản ghi, tuổi thai 38–41 tuần, mỗi bản ghi 5 phút, 4 kênh ECG bụng cùng 1 kênh fECG da đầu đồng thời và chú thích vị trí đỉnh R do chuyên gia gán, 1000 Hz, 16 bit, băng thông 1–150 Hz; bộ 2 là "Fetal electrocardiograms, direct and abdominal with reference heartbeat annotations" (Matonia 2020, bộ Silesia) gồm phần 2a tín hiệu chuyển dạ với 12 sản phụ, 12 bản ghi 5 phút, 4 kênh bụng lấy mẫu 500 Hz, có fECG da đầu đồng thời lấy mẫu 1000 Hz và chú thích đỉnh R, và phần 2b tín hiệu thai kỳ với 10 sản phụ, 10 bản ghi 20 phút, 4 kênh bụng 500 Hz, không có fECG da đầu mà chỉ có chú thích vị trí đỉnh R. Như vậy tổng số bản ghi có dạng sóng fECG da đầu (bắt buộc để huấn luyện phương pháp này) chỉ là 5 + 12 = 17; bài không nêu rõ có dùng bộ 2b hay không và nếu có thì lấy nhãn dạng sóng ở đâu, bài không nêu tổng số đoạn cửa sổ N sau khi cắt, và bài không nêu có kiểm tra trùng lặp giữa bộ 1 và bộ 2a hay không dù cả hai đều do cùng nhóm Silesia thu thập theo cùng giao thức.\par\vspace{3pt}
\noindent \textbf{Giao thức đánh giá.} Kiểm định chéo 5-fold trên dữ liệu gộp của cả hai bộ, mô tả nguyên văn chỉ gồm một câu duy nhất "The datasets were split using the 5-Fold cross-validation technique.": bài không nêu fold được chia theo chủ thể, theo bản ghi hay theo đoạn ngẫu nhiên, không nêu danh sách bản ghi của từng fold, không nêu có đảm bảo các cửa sổ của cùng một bản ghi không nằm cả ở tập huấn luyện lẫn tập kiểm thử hay không — đây là khoảng trống nghiêm trọng nhất của bài. Dung sai khớp nhịp là 31,25 ms (16 mẫu ở 512 Hz), nguyên văn "any R-peak within 31.25 ms of the ground truth is considered to be correct", chặt hơn chuẩn CinC 50 ms; bộ dò đỉnh là EngZee sửa đổi áp lên fECG tái tạo; đơn vị phân tích là nhịp tim (đúng chuẩn); đầu vào là bốn kênh bụng (C = 4) nên đây không phải là phương pháp đơn kênh. Không có tinh chỉnh trên tập đích vì không có tập đích riêng, tất cả đều nằm trong 5-fold gộp; không có ngoại kiểm trên bộ dữ liệu thứ ba; không báo cáo khoảng tin cậy 95\%, chỉ báo cáo giá trị từng fold và trung bình.\par\vspace{3pt}
\noindent \textbf{Kết quả chính.} Chất lượng tái tạo fECG theo hệ số tương quan Pearson (Bảng 1, trung bình 5 fold): ResNet 13 khối + Basic 0,882 (tốt nhất), ResNet 13 khối + Self 0,871, ResNet 9 khối + Basic 0,864, ResNet 9 khối + Self 0,855, Unet 256 + Basic 0,780, Unet 256 + Self 0,770, Self-FPN + Basic 0,762, Unet 128 + Basic 0,751, Self-FPN + Self 0,738, Unet 128 + Self 0,732; mô hình tốt nhất theo từng fold (Bảng 2) có RMSE 0,106 / 0,109 / 0,103 / 0,102 / 0,115, MAE 0,069 / 0,073 / 0,068 / 0,069 / 0,077, PCC 0,89 / 0,88 / 0,89 / 0,89 / 0,87 và tương quan phổ 0,90 / 0,88 / 0,90 / 0,90 / 0,89, tóm tắt công bố là PCC trung bình 88,4\% và tương quan phổ trung bình 89,4\%. Dò fQRS (Bảng 3) đạt trung bình Precision 97,6\%, Recall 94,8\%, F1 96,4\% và Accuracy 92,6\%, với các fold lần lượt là 0,98/0,95/0,96/0,93; 0,97/0,95/0,96/0,92; 0,98/0,95/0,97/0,93; 0,97/0,95/0,96/0,92; 0,98/0,94/0,97/0,92. Về HRV (Bảng 5, trung bình): muy\_RR thật 473,188 ms so với tái tạo 474,464 ms (sai số 0,27\%), muy\_HR 128,52 bpm so với 128,852 bpm (0,25\%), sigma\_HR 14,142 bpm so với 18,244 bpm (sai số 29\%, rất lớn), SDNN/RMSSD 0,700 so với 0,696 (0,57\%), PNN50 40,448\% so với 41,386\% (2,32\%); nghiên cứu loại bỏ cho thấy hàm mất mát đề xuất làm tăng PCC gần 10 điểm, còn chi phí tính toán là 0,762 triệu tham số tổng trên máy ảo đám mây Intel Xeon 4 nhân 8 luồng 2,0 GHz, RAM 25 GB, GPU NVIDIA Tesla T4 16 GB, tách được 1.740 giây tín hiệu fECG trong khoảng 73 giây (tương đương khoảng 23,8 lần thời gian thực).\par\vspace{3pt}
\noindent \textbf{Điểm yếu.} Nghiêm trọng nhất, bài chỉ nói "5-fold cross-validation" mà không nêu cách chia fold, trong khi có một dấu hiệu suy luận rất mạnh cho thấy fold được chia ngẫu nhiên theo đoạn chứ không theo chủ thể: ở Bảng 5, thống kê của chính dữ liệu thật (nhãn chuẩn) gần như y hệt nhau ở cả 5 fold, với muy\_RR = 473,19 / 473,18 / 473,19 / 473,18 / 473,20 ms và muy\_HR = 128,52 / 128,52 / 128,51 / 128,52 / 128,53 bpm, mà nếu mỗi fold kiểm thử chứa những thai nhi khác nhau thì nhịp tim thật trung bình phải khác nhau rõ rệt vì mỗi thai nhi có nhịp riêng — đây là rò rỉ dữ liệu (xin nhắc lại: suy luận từ con số trong bài, không phải điều bài tuyên bố), và kết hợp với việc họ dùng cửa sổ chồng lấn ở bước tiền xử lý thì khả năng các đoạn gần như trùng nhau nằm cả hai bên là rất cao. Đây cũng không phải phương pháp đơn kênh vì đầu vào bộ sinh là 4 kênh bụng nên mọi so sánh với phương pháp đơn kênh đều không công bằng; bài không nêu cách xử lý bộ 2b (10 bản ghi không có fECG da đầu) trong khi phương pháp bắt buộc phải có fECG da đầu làm nhãn; và công thức Accuracy in trong bài (phương trình 20) có chứa TN vốn không xác định được trong bài toán dò đỉnh, kiểm tra số học cho thấy từ P = 0,976 và R = 0,948 ta có TP/(TP+FP+FN) = 1/(1/0,948 + 1/0,976 − 1) = 0,9264, đúng bằng 0,926 mà họ báo cáo, nên "accuracy" của họ thực chất là chỉ số TP/(TP+FP+FN) kiểu CinC chứ không phải công thức họ in ra. Sai số sigma\_HR lên đến 29\% cho thấy biến thiên nhịp tim tái tạo không đáng tin trong khi đây mới là thứ quan trọng nhất để phát hiện suy thai (họ có thừa nhận một câu nhưng không phân tích thêm); ngoài ra không có ngoại kiểm trên bộ dữ liệu thứ ba nên không biết mô hình tổng quát hóa thế nào, không có khoảng tin cậy 95\% và không có kiểm định thống kê giữa các cấu hình, bài không nêu độ trễ cho một cửa sổ đơn lẻ mà chỉ nêu thông lượng tổng, và một giả thuyết cần kiểm tra là bộ 1 (ADFECGDB) với bộ 2a (Silesia chuyển dạ) đều do cùng nhóm Silesian University of Technology thu thập theo cùng giao thức, nên nếu có bản ghi trùng nhau thì việc gộp hai bộ không tăng tính độc lập như họ tuyên bố — khả năng này bài không hề đề cập.\par\vspace{3pt}
\noindent \textbf{\color{tdtblue}So với mô hình của nhóm.} Không so sánh trực tiếp được vì ba lý do độc lập mà mỗi lý do đều đủ phá vỡ phép so sánh: họ dùng 4 kênh bụng làm đầu vào bộ sinh còn FetalQRS-TCN dùng đơn kênh, nên 96,4\% của họ là kết quả của một bài toán dễ hơn hẳn; giao thức 5-fold không nêu cách nhóm và bằng chứng số học ở Bảng 5 (nhịp thật giống hệt nhau ở cả 5 fold) chỉ rất mạnh tới việc chia ngẫu nhiên theo đoạn trên dữ liệu gộp, tức là có rò rỉ dữ liệu, trong khi con số 97,43 ± 4,37\% của nhóm đến từ giao thức tách theo bản ghi chặt chẽ hơn nhiều; và tập dữ liệu khác nhau vì họ gộp ADFECGDB với bộ Silesia còn nhóm chỉ dùng ADFECGDB, cùng CinC 2013 làm kiểm thử xuyên bộ. Điểm duy nhất đối chiếu được một cách có ý nghĩa là chi phí tính toán và ở đây nhóm thắng rõ ràng: mô hình của nhóm có 113.481 tham số (0,113 triệu), checkpoint 0,48 MB, 4,35 ms cho một cửa sổ 4 giây trên CPU, so với 0,762 triệu tham số toàn mạng của họ (gấp 6,7 lần) mà riêng một bộ sinh đã là 0,337 triệu (gấp 3,0 lần), và thông lượng 1.740 giây tín hiệu trong 73 giây trên GPU Tesla T4 (khoảng 23,8 lần thời gian thực) so với khoảng 920 lần thời gian thực của nhóm trên CPU — tức nhóm nhỏ hơn khoảng 6,7 lần, nhanh hơn khoảng 39 lần và không cần GPU. Về dải lọc, họ dùng 0,1–70 Hz; khảo sát 8 dải của nhóm không có chính xác dải này nhưng hai dải gần nhất là 0,5–100 Hz (72,17) và 1–45 (80,06) đều thấp hơn 10–60 Hz (91,06) từ 11 đến 19 điểm, từ đó có thể suy luận rằng việc họ phải dùng hàm mất mát bốn thành phần và bộ sinh gấp 6,7 lần tham số một phần là để bù cho dải lọc quá rộng — nhưng đây chỉ là giả thuyết, không thể khẳng định vì họ không làm khảo sát dải lọc.\par\vspace{3pt}
\noindent \textbf{\color{tdtblue}Nhóm rút ra gì.} Áp dụng: dùng con số 0,762 triệu tham số của họ làm mốc đối chiếu chính cho luận điểm hiệu quả tính toán của nhóm (FetalQRS-TCN chỉ bằng 1/3 một bộ sinh của họ và 1/6,7 toàn mạng mà vẫn không cần fECG da đầu làm nhãn), học cách họ báo cáo chi phí tính toán với phần cứng cụ thể (Xeon 4 nhân 8 luồng 2,0 GHz, 25 GB RAM, Tesla T4 16 GB), số tham số tách riêng từng thành phần và thông lượng (1.740 giây tín hiệu / 73 giây) rồi bổ sung thêm độ trễ một cửa sổ 4,35 ms và kích thước checkpoint 0,48 MB là hai thứ họ không có, báo cáo thêm chỉ số TP/(TP+FP+FN) để có điểm đối chiếu với cả họ lẫn bảng xếp hạng CinC 2013, thử hàm mất mát có thành phần phổ nếu mở rộng sang giữ hình thái (mất mát phổ + thời gian + công suất làm tăng PCC gần 10 điểm ở bài của họ), và đối chiếu các chi tiết tiền xử lý như khử trôi đường cơ sở bằng khớp đa thức bậc cao (bậc 8) cùng kỹ thuật "áp dụng - lật - áp dụng lại" chính là lọc pha không nhóm đã dùng. Tránh: tuyệt đối không viết "5-fold cross-validation" mà không nói rõ cách nhóm — bài của nhóm phải có câu nguyên văn kiểu "Chúng tôi dùng tách theo bản ghi; không có cửa sổ nào của bản ghi kiểm thử xuất hiện trong tập huấn luyện; danh sách bản ghi của từng fold được liệt kê ở Bảng X" — tránh công bố các thống kê nhãn chuẩn mà không kiểm tra chúng, vì chính việc nhịp thật giống hệt nhau ở cả 5 fold đã tố cáo giao thức của họ, nên nhóm nên chủ động in bảng nhịp tim thật của từng fold tách theo bản ghi để chứng minh các fold thật sự khác nhau, và tránh tuyên bố "đơn kênh" nếu thực tế gộp nhiều kênh, ngược lại phải nhấn mạnh nhóm thật sự đơn kênh trong khi Basak dùng 4 kênh vì đây là lợi thế lâm sàng lớn (ít điện cực hơn, dễ triển khai hơn).\par\vspace{3pt}
\vspace{6pt}
\subsection{Shokouhmand và Tavassolian 2023, IEEE Trans Biomed Eng \normalfont\small (p05)}
\noindent{\footnotesize\itshape Arash Shokouhmand, Negar Tavassolian, "Fetal Electrocardiogram Extraction Using Dual-Path Source Separation of Single-Channel Non-Invasive Abdominal Recordings", IEEE Transactions on Biomedical Engineering, vol. 70, no. 1, pp. 283-295, January 2023. DOI: 10.1109/TBME.2022.3189617. (Nhan 09/05/2022, chinh sua 18/06/2022, chap nhan 06/07/2022, xuat ban 11/07/2022. Tai tro: NSF Award 1855394. ...}\par\vspace{2pt}
\noindent{\small\color{tdtblue}\textbf{Kết luận: so được một phần}}\par\vspace{4pt}
\noindent \textbf{Họ làm gì.} Bài giải bài toán tách nguồn mù từ một kênh ECG bụng duy nhất, đồng thời tách tín hiệu trộn AECG thành ECG mẹ (MECG) và ECG thai (FECG), bằng cách chuyển kiến trúc Dual-Path RNN của Luo 2020 từ lĩnh vực tách giọng nói sang tín hiệu ECG. Nhóm tác giả không dò đỉnh trực tiếp mà tái tạo lại dạng sóng FECG rồi chạy Pan-Tompkins cải biên trên tín hiệu tái tạo, với lập luận rằng bài toán một kênh là bài toán thiếu xác định nên ICA/PCA truyền thống (đòi hỏi số cảm biến lớn hơn hoặc bằng số nguồn) không dùng được.\par\vspace{3pt}
\noindent \textbf{Phương pháp.} Tiền xử lý gồm lọc thông cao Butterworth zero-phase 2 Hz để khử trôi đường cơ sở, lọc notch IIR bậc 2 tại 50 Hz hoặc 60 Hz với hệ số phẩm chất Q = 25 chạy cả chiều thuận và chiều ngược, lấy mẫu lại NIFECGC và FECGSYNDB về 500 Hz bằng FFT, cắt cửa sổ 4 giây (2.000 mẫu tại 500 Hz) chồng lấn 90\% và chuẩn hoá z-score. Khối khử nhiễu là mạng sinh đối kháng bình phương tối thiểu có điều kiện (LSCGAN) với bộ sinh Conv1d 16 bộ lọc → khối DP-LSTM → tích chập theo chiều sâu và bộ phân biệt gồm các khối residual 1D, hàm mất mát sinh có thêm số hạng L1 với lambda = 0,01; khối tách nguồn dùng Inception trích đặc trưng thưa nhiều kích thước bộ lọc → chuỗi khối DP-LSTM → lớp Conv2D với 2M bộ lọc và sigmoid sinh ma trận mặt nạ kích thước K × 2M × L, tách đôi thành mặt nạ thai và mặt nạ mẹ rồi đưa về miền thời gian, huấn luyện bằng hàm mất mát −SI-SNR với Permutation Invariant Training. Ba kiểu tăng cường dữ liệu được dùng là jittering (công suất nhiễu thấp hơn năng lượng phức bộ QRS thai 3 dB, kèm nhiễu đặc thù thai kỳ), cắt ngẫu nhiên với độ rộng lấy từ phân phối Beta và che thời gian–tần số dưới 50 Hz; cấu hình cuối gồm 3 khối DP-LSTM, 6 bộ lọc trích đặc trưng, đoạn con L = 100 ms, huấn luyện bằng Adam, lr 0,001 (giảm hệ số 0,98 sau mỗi 5 epoch không cải thiện), 200 epoch, batch 32 trên GPU NVIDIA RTX 2070.\par\vspace{3pt}
\noindent \textbf{Mô hình.} DPSS (Dual-Path Source Separation) gồm khối khử nhiễu LSCGAN (bộ sinh DP-LSTM và bộ phân biệt ResBlock) ghép với khối tách nguồn dựa trên mặt nạ (Inception và 3 khối DP-LSTM). Số tham số, kích thước mô hình (MB) và số FLOPs: bài không nêu; bài chỉ báo cáo thời gian xử lý.\par\vspace{3pt}
\noindent \textbf{Dữ liệu.} Ba bộ dữ liệu: FECGSYNDB tổng hợp (10 thai phụ mô phỏng, 32 kênh bụng, 145,8 giờ, mỗi bản ghi 5 phút, 250 Hz, nhiễu đặc thù thai kỳ trộn vào để tạo SNR trong khoảng −9 đến +2 dB), ADFECGDB phiên bản mở rộng Silesia (4 kênh bụng, 500 Hz, 10 ca mang thai và 12 ca chuyển dạ, tổng 22 ca, tuổi thai 38–42 tuần) và NIFECGC tức CinC 2013 Set A (75 bản ghi 4 kênh, 1 kHz, 60 giây mỗi bản ghi, loại 7 bản ghi a33, a38, a47, a52, a54, a71, a74 do chú thích sai, còn 69 bản ghi). Cần lưu ý bộ thứ hai không phải ADFECGDB cổ điển 5 bản ghi 1 kHz mà nhóm ta dùng; số đoạn kiểm thử lần lượt là 190.720 đoạn trên FECGSYNDB (trên 3.432.960 đoạn huấn luyện sau tăng cường), 2.999 và 749 đoạn/kênh trên ADFECGDB mang thai và chuyển dạ, 149 đoạn/kênh trên NIFECGC.\par\vspace{3pt}
\noindent \textbf{Giao thức đánh giá.} Mô hình được huấn luyện bằng xác thực chéo tách theo chủ thể chỉ trên FECGSYNDB (9 chủ thể tổng hợp để huấn luyện, 1 chủ thể giữ lại kiểm thử, lặp 10 lần), sau đó đem kiểm thử zero-shot trực tiếp trên ADFECGDB (mang thai và chuyển dạ) cùng NIFECGC mà không tinh chỉnh; đây là giao thức sạch nhất trong ba bài vì không một mẫu thực tế nào được dùng để huấn luyện nên không có rò rỉ dữ liệu từ tập thực tế. Dung sai là 50 ms theo chuẩn CinC 2013 (Behar 2016), trùng với dung sai của nhóm ta; đầu vào một kênh nhưng là kênh được chọn thủ công, cụ thể là kênh "có thành phần FECG rõ nhất" của từng chủ thể trong NIFECGC và ADFECGDB, riêng FECGSYNDB dùng cả 32 kênh. Đơn vị phân tích là kênh/bản ghi rồi lấy trung bình; ngoài F1 còn báo cáo độ chính xác FHR (tỷ lệ đoạn có FHR ước lượng nằm trong ±5 nhịp/phút so với tham chiếu) và các chỉ số phân cụm (purity, Jaccard, Davies-Bouldin) trên bản đồ mặt nạ sau khi giảm chiều bằng t-SNE, còn kiểm định thống kê so sánh với các phương pháp khác thì bài không nêu (không có p-value, không có khoảng tin cậy).\par\vspace{3pt}
\noindent \textbf{Kết quả chính.} Trên FECGSYNDB (tách theo chủ thể, cùng phân phối), F1 trung bình đạt 99,03\% (±0,39) và SI-SNR +14,92 dB (±1,02), tốt nhất là chủ thể 8 với F1 99,35\% và SI-SNR +15,94 dB, yếu nhất là chủ thể 2 với F1 98,24\%. Kiểm thử zero-shot cho F1 97,7\% và độ chính xác FHR 88,61\% trên ADFECGDB 22 ca (riêng tập chuyển dạ đạt F1 98,08\% và FHR 89,22\%), F1 95,3\% và độ chính xác FHR 83,86\% trên NIFECGC 69 ca, trung bình trên dữ liệu thật là Se 86,23\% / PPV 95,75\% / Acc 97,29\% / F1 93,43\%; phân cụm cho purity trung bình 0,9750, Jaccard 0,9705 và Davies-Bouldin 0,7429. Nghiên cứu loại bỏ trên FECGSYNDB cho thấy khử nhiễu nâng F1 từ 61,76\% lên 84,21\% (+22,45 điểm), trích đặc trưng thưa 6 bộ lọc nâng lên 94,68\% và tăng cường dữ liệu nâng lên 99,03\% (+4,35 điểm); quét cửa sổ cho thấy 4 giây là tốt nhất (98,37\%, giảm còn 90,39\% ở 10 giây) và đoạn con 100 ms tốt nhất (99,03\% so với 97,11\% ở 50 ms); xử lý một đoạn 4 giây hết 0,52 s trên CPU Intel 3,6 GHz, 1,2 s trên Broadcom BCM2711 (Raspberry Pi), 1,8 s trên Qualcomm Snapdragon S1 và 1,0 s trên Apple S7, đều dưới 4 s nên nhóm tác giả kết luận là chạy được thời gian thực.\par\vspace{3pt}
\noindent \textbf{Điểm yếu.} Bài không nêu số tham số, kích thước mô hình và số FLOPs nên không thể tái lập chính xác kiến trúc cũng như không so sánh được ngân sách tính toán, đồng thời tồn tại nhiều mâu thuẫn nội tại: mục IV-C ghi 6 bộ lọc độ dài "1, 3, 5, 7, 9 và 11" trong khi mục IV-E ghi "1, 3, 5, 7, 9 và 13"; dòng trung bình Bảng II ghi Se 86,23\% và PPV 95,75\% nhưng F1 tương ứng phải khoảng 90,72\% chứ không phải 93,43\%, và trung bình của 97,7\% với 95,3\% cũng không thể ra 93,43\%; Hình 9(a) ghi F1 tốt nhất 98,37\% tại 4 giây trong khi Bảng I ghi 99,03\% cho cùng cấu hình. Việc chọn thủ công kênh "có thành phần FECG rõ nhất" cho từng chủ thể tức là đã dùng thông tin nhãn để chọn kênh, một dạng rò rỉ nhẹ (chọn kênh kiểu oracle) không khả thi trên lâm sàng — nhóm tác giả có đề xuất dùng chỉ số chất lượng tín hiệu Bayes ngây thơ của Andreotti 2017 để tự động hoá nhưng không thực hiện trong bài; ngoài ra việc gọi bộ 22 bản ghi Silesia là "ADFECGDB" gây nhầm lẫn với ADFECGDB cổ điển 5 bản ghi, khiến mọi so sánh với các công trình dùng 5 bản ghi đều không cùng tập dữ liệu (chính họ thừa nhận điều này ở mục IV-F). Bài cũng không báo cáo jitter/sai số định vị đỉnh (ms), không có khoảng tin cậy hay kiểm định thống kê khi so với đối thủ; do chỉ lọc thông cao 2 Hz và không lọc thông thấp nên dải thông thực tế 2–250 Hz rất rộng, giữ lại nhiều sóng T mẹ và nhiễu cơ và phải bù bằng mạng GAN khử nhiễu rất nặng, và chính nhóm tác giả thừa nhận hai chế độ hỏng: GAN có thể khuếch đại artifact giống QRS thai (sinh dương tính giả) hoặc làm mượt mất hoàn toàn đoạn tín hiệu nhiễu mạnh (sinh âm tính giả).\par\vspace{3pt}
\noindent \textbf{\color{tdtblue}So với mô hình của nhóm.} Không so sánh trực tiếp được hoàn toàn, nhưng đây là công trình gần nhất với giao thức của nhóm ta và là mốc so sánh đáng tin nhất trong ba bài: cùng dung sai 50 ms, cùng cửa sổ 4 giây (họ tìm ra 4 giây là tối ưu bằng thực nghiệm, trùng khớp với kết quả quét trường tiếp nhận của nhóm ta là bão hoà quanh 1,5 s), cùng định hướng một kênh và cùng cách chọn kênh tốt nhất. Hai con số không thể đặt cạnh nhau vì khác tập dữ liệu và khác nguồn huấn luyện: 97,7\% của họ đo trên 22 bản ghi Silesia mở rộng sau khi huấn luyện 100\% trên FECGSYNDB (145,8 giờ) rồi kiểm thử zero-shot, còn 99,18\% của nhóm ta đo trên 5 bản ghi ADFECGDB cổ điển với tách theo bản ghi; họ hơn nhóm ta rõ rệt về khả năng tổng quát hoá xuyên bộ dữ liệu khi đạt 95,3\% trên CinC 2013 (69 bản ghi) zero-shot so với 77,34 ± 28,54 của nhóm ta, tức nhóm ta kém hơn khoảng 18 điểm và đây là điểm yếu cần thừa nhận thẳng thắn. Nhiệm vụ cũng khác vì họ tái tạo dạng sóng rồi mới dò đỉnh (hai bước) nên buộc phải giữ dải tần rộng từ 2 Hz trở lên để bảo toàn sóng P và sóng T, còn nhóm ta dò đỉnh trực tiếp sequence-to-sequence nên được phép lọc chặt 10–60 Hz; về chi phí thì mô hình của nhóm thắng tuyệt đối vì 0,52 s cho một cửa sổ 4 giây trên CPU 3,6 GHz là chậm hơn 4,35 ms của nhóm ta khoảng 120 lần.\par\vspace{3pt}
\noindent \textbf{\color{tdtblue}Nhóm rút ra gì.} Nên chạy ngay thí nghiệm ưu tiên số một là huấn luyện chính FetalQRS-TCN trên FECGSYNDB (145,8 giờ, 250 Hz, đúng tần số nhóm ta đang dùng nên không cần lấy mẫu lại) rồi đo trên ADFECGDB và CinC 2013 mà không tinh chỉnh, kèm ba kiểu tăng cường dữ liệu của họ (jittering với công suất nhiễu thấp hơn năng lượng QRS thai 3 dB, cắt ngẫu nhiên theo phân phối Beta, che thời gian–tần số chỉ trên các ô dưới 50 Hz) vốn mang lại +4,35 điểm F1, lớn hơn khoảng 10 lần toàn bộ khoảng cách giữa 16 kiến trúc nhóm ta đã so (+0,41 điểm); đồng thời bổ sung độ chính xác FHR (±5 nhịp/phút) bên cạnh F1, làm bảng thời gian chạy trên nhiều loại chip như Hình 11 của họ, và trích dẫn đề xuất dùng chỉ số chất lượng tín hiệu để chọn kênh mà họ nêu nhưng không thực hiện làm động lực cho hướng fSQI-TDA. Cần tránh ba lỗi của họ: không nêu số tham số (lý do khiến công trình không tái lập được), để các con số mâu thuẫn nội tại trong bảng (Se 86,23 và PPV 95,75 không thể ra F1 93,43 — phải tự kiểm tra lại mọi bảng bằng công thức trước khi nộp), và gọi bộ 22 bản ghi Silesia là "ADFECGDB" — nhóm ta phải ghi rõ số bản ghi và ID bản ghi trong mọi bảng.\par\vspace{3pt}
\vspace{6pt}
\subsection{Chen và cs. 2025, Sensors \normalfont\small (p06)}
\noindent{\footnotesize\itshape Lin Chen, Shuicai Wu, Zhuhuang Zhou, "Fetal ECG Signal Extraction from Maternal Abdominal ECG Signals Using Attention R2W-Net", Sensors, vol. 25, no. 3, art. no. 601, 2025. DOI: 10.3390/s25030601. (MDPI, open access CC BY. Nhan 22/10/2024, chinh sua 18/01/2025, chap nhan 19/01/2025, xuat ban 21/01/2025. College of Chemistry and Life Sciences, Beijing University of Technology, Trung Quoc. Bai ...}\par\vspace{2pt}
\noindent{\small\color{warnred}\textbf{Kết luận: không so trực tiếp được}}\par\vspace{4pt}
\noindent \textbf{Họ làm gì.} Bài xây dựng một mạng hình chữ W gồm hai mạng Attention R2U-Net song song, trong đó nhánh phải tái tạo ECG mẹ, nhánh trái tái tạo ECG thai, và đặc trưng thai được lấy bằng cách trừ đặc trưng mẹ (qua hàm tanh) ngay trong bộ mã hoá; nhóm tác giả khẳng định mô hình xử lý đầu vào một kênh. Khác biệt kỹ thuật so với W-Net gốc là thay lớp tích chập tiến truyền bằng khối tích chập hồi quy có dư (Recurrent Residual Convolution - RRC) và thay kết nối tắt bằng cổng chú ý; mục tiêu không chỉ là dò đỉnh QRS thai mà còn giữ dạng sóng (sóng P và sóng T) để bác sĩ đọc được hình thái ECG thai.\par\vspace{3pt}
\noindent \textbf{Phương pháp.} Tiền xử lý gồm lọc thông thấp 100 Hz và lọc thông cao 3 Hz (lý do nêu trong bài: đỉnh QRS của FECG nằm trong khoảng 10–15 Hz và băng thông tối thiểu cần 100 Hz), lấy mẫu lại tất cả về 250 Hz, cắt đoạn 1.024 mẫu (khoảng 4,096 giây) chồng lấn 24 mẫu ở hai đầu để bảo đảm liên tục và chuẩn hoá Z-score; loại lọc (Butterworth hay Chebyshev), bậc lọc và việc có zero-phase hay không: bài không nêu, còn lọc notch tần số lưới điện thì bài không nêu rõ (chỉ nói nhiễu đường dây điện "đã được loại bỏ hiệu quả" sau khi lọc thông dải). Kiến trúc gồm hai mạng Attention R2U-Net ghép chữ W, mỗi mạng 5 tầng với bộ mã hoá 5 khối RRC và 4 lớp max-pooling, bộ giải mã đối xứng 4 lớp upsampling và 5 khối RRC; khối RRC dùng T = 2 (một lớp tích chập thường và hai lớp hồi quy luân phiên qua 2 bước thời gian) cộng kết nối dư; kích thước nhân là 5 cho mạng trích MECG và 3 cho mạng trích FECG; cổng chú ý dùng chú ý cộng, qua ReLU rồi Sigmoid để sinh hệ số alpha cho từng điểm mẫu. Huấn luyện bằng Python 3.8, PyTorch và MONAI trên 2 CPU Intel Xeon Gold 6132 @2,60 GHz, GPU NVIDIA TITAN RTX 24 GB, RAM 128 GB, với hàm mất mát MAE (có thí nghiệm loại bỏ chứng minh MAE tốt hơn MSE), Adam, lr 0,0001, batch 32, 30 epoch, padding zero, stride 1 và kích hoạt Leaky ReLU; đỉnh được dò bằng thuật toán Pan-Tompkins trên tín hiệu FECG đã tái tạo.\par\vspace{3pt}
\noindent \textbf{Mô hình.} Attention R2W-Net gồm hai mạng Attention R2U-Net 1D ghép chữ W, dùng khối tích chập hồi quy có dư (RRC) và cổng chú ý, với số kênh các tầng ghi trong bài là "64, 128, 256, 512, 256 và 1024". Số tham số, kích thước mô hình (MB) và số FLOPs: bài không nêu; chi phí duy nhất được báo cáo là 0,18 giây để trích FECG từ một đoạn maECG 4 giây (so với CycleGAN 0,20 s, PA2NET 0,28 s, DPSS 0,46 s, W-NETR 0,34 s, đo trong cùng môi trường nhưng không nêu đo trên CPU hay GPU nào).\par\vspace{3pt}
\noindent \textbf{Dữ liệu.} Bốn bộ dữ liệu: FECGSYNDB dùng để huấn luyện (10 đối tượng mô phỏng, 32 kênh bụng và 2 kênh ngực, 7 kịch bản sinh lý Baseline và C0–C5 trong đó C5 là song thai, 5 mức nhiễu cộng thêm 0, 3, 6, 9, 12 dB, mỗi kịch bản mô phỏng 5 lần, tổng 1.750 tín hiệu, mỗi tín hiệu 5 phút tại 250 Hz, nhưng chỉ dùng 4 kênh là kênh 11, 19, 22 và 25) và ADFECGDB (5 thai phụ 38–41 tuần, 4 kênh bụng và 1 kênh điện cực da đầu thai nhi, mỗi bản ghi 5 phút tại 1 kHz). Ngoài ra có PCDB tức CinC 2013 Set A (75 bản ghi 4 kênh, 1 kHz, 60 giây mỗi bản ghi, chú thích đỉnh R thai do chuyên gia đánh dấu tay) và NIFECGDB (55 bản ghi đa kênh gồm 3–4 kênh bụng và 2 kênh ngực, 1 kHz), trong đó chọn 14 tập con 154, 192, 244, 274, 290, 323, 368, 444, 597, 733, 746, 811, 826 và 906 để so sánh với Ghonchi 2022.\par\vspace{3pt}
\noindent \textbf{Giao thức đánh giá.} Đây là điểm yếu nghiêm trọng nhất của bài: trên FECGSYNDB chỉ có một lần chia duy nhất — dữ liệu của đối tượng thứ 10 làm tập kiểm thử, các đối tượng còn lại làm tập huấn luyện — tức không phải tách theo chủ thể đầy đủ (chỉ 1 fold, không lặp 10 lần) nên không có độ lệch chuẩn xuyên đối tượng và con số có thể là may mắn của một fold; trên ADFECGDB, bài ghi nguyên văn rằng mô hình "được tinh chỉnh bằng một tập con của ADFECGDB, phần dữ liệu còn lại được dùng để kiểm thử" nhưng Bảng 5 lại báo cáo kết quả trên cả 5 bản ghi R01, R04, R07, R08, R10 tức toàn bộ ADFECGDB, mà không nêu bản ghi nào dùng để tinh chỉnh, bản ghi nào để kiểm thử hay tỷ lệ chia bao nhiêu; với PCDB và NIFECGDB, bài không nêu rõ có tinh chỉnh hay không và cũng không nêu cách chia. Suy luận (không phải dữ kiện trong bài): nếu tập tinh chỉnh và tập kiểm thử cùng nằm trong 5 bản ghi đó thì đây là chia ngẫu nhiên trong cùng bệnh nhân, tức rò rỉ dữ liệu điển hình, và con số 99,17\% không còn ý nghĩa. Dung sai là 31,25 ms — con số phi tiêu chuẩn, chặt hơn chuẩn CinC 2013 (50 ms) và chặt hơn nhiều so với AAMI EC57 (150 ms), bằng 1000/32 ms nên có vẻ xuất phát từ 8 mẫu tại 256 Hz chứ không từ chuẩn nào, nhưng vì là ngưỡng chặt hơn nên nó không làm đẹp số liệu, đây là điểm cộng hiếm hoi của bài; đầu vào một kênh (có chạy riêng trên từng kênh trong 4 kênh của ADFECGDB), đơn vị phân tích là bản ghi, và có dùng t-test so sánh MSE/MAE với các mô hình khác (tất cả p < 0,05) nhưng không có kiểm định thống kê cho chỉ số F1 dò đỉnh.\par\vspace{3pt}
\noindent \textbf{Kết quả chính.} Về chất lượng tái tạo tín hiệu, FECGSYN đạt MSE 0,096 ± 0,016 và MAE 0,082 ± 0,006 (theo kịch bản: C0 MSE 0,079 / C1 0,083 / C2 0,081 / C3 0,122 / C4 0,114 / C5 0,097), còn ADFECGDB đạt MSE 0,025 ± 0,004, MAE 0,013 ± 0,003 và SNR 8,03 ± 0,67 dB (R01 SNR 9,08; R04 7,17; R07 7,94; R08 8,41; R10 7,53). Về dò đỉnh fQRS ở dung sai 31,25 ms, ADFECGDB đạt Se 99,24\% / PPV 99,11\% / F1 99,17\% (R01 99,40; R04 99,18; R07 99,24; R08 99,08; R10 98,96), PCDB đạt Se 98,23\% / PPV 97,74\% / F1 98,03\%, NIFECGDB đạt Se 96,76\% / PPV 97,39\% / F1 97,08\% và FECGSYNDB đạt Se 97,43\% / PPV 97,98\% / F1 97,56\%. Nghiên cứu loại bỏ trên ADFECGDB (Bảng 7) cho F1 lần lượt 96,41\% với tích chập tiến truyền, 98,29\% với tích chập có dư, 98,50\% với tích chập hồi quy và 99,17\% với khối RRC đề xuất, tức toàn bộ phần thưởng của kiến trúc RRC so với tích chập thường chỉ là +2,76 điểm và so với tích chập có dư chỉ +0,88 điểm; số lớp tích chập tối ưu là 5 (4 lớp thì thiếu, từ 6 lớp trở lên thì quá khớp), hàm mất mát MAE tốt hơn MSE, và chi phí là 0,18 giây cho một đoạn 4 giây.\par\vspace{3pt}
\noindent \textbf{Điểm yếu.} Rò rỉ dữ liệu tiềm ẩn nghiêm trọng trên ADFECGDB: mô hình được tinh chỉnh trên một phần ADFECGDB rồi báo cáo kết quả trên cả 5 bản ghi của chính bộ đó mà không ghi rõ cách chia, nên con số 99,17\% không được coi là kết quả ngoài mẫu — chính bài p07 (Asadi 2025) đã chỉ đích danh điều này khi viết rằng "một phần bộ ADFECGDB đã được dùng để huấn luyện trong các mô hình này, điều nhiều khả năng đã góp phần làm tăng hiệu năng của chúng trên bộ dữ liệu này", và Attention R2W-Net được nêu tên trong nhóm đó. Bài không nêu số tham số, kích thước mô hình hay số FLOPs (với 5 tầng, số kênh lên tới 1024 và nhân đôi vì kiến trúc chữ W, ước lượng — suy luận, không phải dữ kiện — mô hình ở thang hàng chục triệu tham số, lớn hơn FetalQRS-TCN ít nhất 2–3 bậc độ lớn); dãy số kênh "64, 128, 256, 512, 256 và 1024" giảm rồi lại tăng nên không thể là một bộ mã hoá hợp lệ, nhiều khả năng là lỗi đánh máy của "64, 128, 256, 512, 1024" và khiến kiến trúc không tái lập được; công thức (11) in F1 = TP/(TP+FP+FN) thực chất là chỉ số Jaccard chứ không phải F1 (với Se 99,24 và PPV 99,11 thì TP/(TP+FP+FN) = 98,36\% trong khi trung bình điều hoà đúng bằng 99,17\%, tức họ tính đúng nhưng in sai công thức — một lỗi biên tập cho thấy bài không được soát kỹ). Bảng so sánh (Bảng 6) chứa các con số bất khả thi về mặt toán học — W-NETR trên ADFECGDB ghi Se 99,45 / PPV 99,50 / F1 99,88 trong khi F1 không thể lớn hơn cả Se lẫn PPV, CycleGAN ghi 99,40 / 99,60 / 99,70 trên ADFECGDB và 96,80 / 97,20 / 97,90 trên PCDB — nghĩa là số liệu được chép lại mà không kiểm tra, nên mọi con số trong bảng này phải coi là không đáng tin; ngoài ra so sánh chồng bộ dữ liệu không đồng nhất về dung sai (các bài được so dùng 50 ms, họ dùng 31,25 ms) mà bài không thảo luận, chỉ 1 fold trên FECGSYNDB nên không có độ lệch chuẩn xuyên đối tượng cho F1, không báo cáo jitter/sai số định vị đỉnh, và bài được gán thể loại "Study Protocol" dù nội dung là một bài thực nghiệm đầy đủ, dấu hiệu quy trình biên tập lỏng lẻo.\par\vspace{3pt}
\noindent \textbf{\color{tdtblue}So với mô hình của nhóm.} Không so sánh trực tiếp được, vì con số 99,17\% trên ADFECGDB của họ được tạo ra sau khi tinh chỉnh trên chính bộ ADFECGDB, còn 97,43 ± 4,37 (toàn bộ 4 kênh) và 99,18 (kênh tốt nhất) của nhóm ta là kết quả tách theo bản ghi, trong đó mọi bản ghi kiểm thử đều chưa từng xuất hiện khi huấn luyện; đặt hai con số cạnh nhau là so một kết quả có thể có rò rỉ với một kết quả sạch, về nguyên tắc không được làm. Thêm ba khác biệt: dung sai của họ là 31,25 ms so với 50 ms của nhóm ta, tức thang đo chặt hơn 1,6 lần nên nếu đo lại FetalQRS-TCN ở 31,25 ms thì số của nhóm ta sẽ giảm và đây là thí nghiệm nên chạy để báo cáo trung thực; mục tiêu của họ là tái tạo dạng sóng giữ được sóng P và sóng T nên buộc phải dùng dải lọc rộng 3–100 Hz, còn nhóm ta chỉ cần dò đỉnh nên được phép lọc 10–60 Hz, vì vậy dải lọc của họ không mâu thuẫn với phát hiện 10–60 Hz của nhóm ta mà là một bài toán khác; về chi phí, 0,18 s cho một cửa sổ 4 giây so với 4,35 ms của nhóm ta là chậm hơn khoảng 41 lần, và mô hình của họ tuy không nêu số tham số nhưng chắc chắn lớn hơn 113.481 tham số của nhóm ta rất nhiều. Kết luận cho bài báo của nhóm: nên trích dẫn bài này như một ví dụ về "con số cao nhưng giao thức không minh bạch", không đặt vào cột so sánh chính, hoặc nếu đặt thì phải ghi chú rõ là có tinh chỉnh trên tập đích.\par\vspace{3pt}
\noindent \textbf{\color{tdtblue}Nhóm rút ra gì.} Nên học ba điểm: báo cáo kết quả theo từng bản ghi như Bảng 2 và Bảng 5 của họ (R01, R04, R07, R08, R10 riêng từng dòng) thay vì chỉ báo cáo macro F1 97,43 ± 4,37 vì độ lệch chuẩn 4,37 đang giấu đi bản ghi yếu; báo cáo F1 ở cả ba mức dung sai — 50 ms (chuẩn CinC, để so sánh), 31,25 ms (để đối chiếu với bài này) và 150 ms (AAMI EC57) — thành một bảng độ nhạy theo dung sai, vì với jitter chỉ 3,05 ms của FetalQRS-TCN thì nhóm ta gần như chắc chắn giữ được F1 cao ngay cả ở 31,25 ms, một lập luận mạnh mà không bài nào trong ba bài đưa ra được; và trích Bảng 7 của họ (tiến truyền 96,41 → có dư 98,29 → hồi quy 98,50 → RRC 99,17) làm bằng chứng độc lập cho luận điểm đổi kiến trúc chỉ ăn vài điểm phần trăm trong khi đổi dải lọc ăn 11,00 điểm, đồng thời cân nhắc dùng MAE thay MSE nếu nhóm ta làm nhánh phụ tái tạo dạng sóng. Cần tránh tuyệt đối: tinh chỉnh trên một phần bộ dữ liệu rồi báo cáo trên toàn bộ bộ đó (đúng lỗi mà README cũ của repo fECG-TDA đã mắc, 91\% ảo so với 59\% tách theo chủ thể thật, nhóm ta đã trả giá một lần rồi), chép bảng so sánh từ bài khác mà không tự tính lại F1 từ Se và PPV cho mỗi dòng (nếu F1 lớn hơn max(Se, PPV) thì con số đó sai và phải ghi chú hoặc loại bỏ), in sai công thức F1 (phải là 2·Se·PPV/(Se+PPV) chứ không phải TP/(TP+FP+FN)), và bỏ sót số tham số, kích thước checkpoint cùng độ trễ — vốn đúng là thế mạnh của nhóm ta (113.481 tham số, 0,48 MB, 4,35 ms) và là thứ mà cả p05 lẫn p06 đều thiếu.\par\vspace{3pt}
\vspace{6pt}
\subsection{Asadi và cs. 2025, IEEE Access \normalfont\small (p07)}
\noindent{\footnotesize\itshape Amirreza Asadi, Aboozar Ghaffari, Zahra Hatami, "U-Net-Based Deep Learning Framework for Single-Channel Fetal ECG Extraction Using Time-Delay Representation", IEEE Access, vol. 13, pp. 206367-206379+ (bai dai 14 trang), 2025. DOI: 10.1109/ACCESS.2025.3639272. (Nhan 05/11/2025, chap nhan 26/11/2025, xuat ban 03/12/2025. Biomedical Engineering Department, Iran University of Science and Technology, ...}\par\vspace{2pt}
\noindent{\small\color{tdtblue}\textbf{Kết luận: so được một phần}}\par\vspace{4pt}
\noindent \textbf{Họ làm gì.} Bài chuyển tín hiệu một kênh thành ảnh hai chiều bằng phép nhúng trễ (time-delay embedding, dựa trên định lý Takens) rồi dùng CNN 2D để tách ECG thai, với lập luận rằng tích chập 1D không bắt được cấu trúc thời gian phức tạp của tín hiệu bụng, còn mạng 2D học được quan hệ giữa các lát cắt trễ nhau vài mẫu. Hai biến thể được đề xuất là SCTD-1Net (trừ trực tiếp FECG từ ma trận trễ của AECG) và SCTD-2Net (hai giai đoạn: U-Net 1D ước lượng MECG trước, sau đó xếp xen kẽ các hàng vector trễ của AECG và của MECG vào một ma trận để mạng 2D tự học phép trừ); đây là bài quan trọng nhất với nhóm ta vì nó là đối thủ trực tiếp của hướng biểu diễn 2D cho tín hiệu 1D mà nhóm ta đã bác bỏ qua thí nghiệm Persistence Image (27,42\%), và là bài duy nhất trong ba bài báo cáo đầy đủ số tham số, dung lượng cùng FLOPs đồng thời không tinh chỉnh trên tập đích.\par\vspace{3pt}
\noindent \textbf{Phương pháp.} Tiền xử lý gồm hạ tần số lấy mẫu về 250 Hz, lọc thông cao 3 Hz để khử trôi đường cơ sở rồi lọc thông thấp 100 Hz để khử nhiễu tần cao và nhiễu điện cơ (loại lọc, bậc lọc, zero-phase và lọc notch 50/60 Hz: bài không nêu), cắt đoạn dài 1.024 mẫu tương đương khoảng 4 giây tại 250 Hz và chuẩn hoá min-max theo công thức (x − min)/(max − min) rồi trừ trung bình để đưa về quanh 0. Giai đoạn 1 trích MECG bằng U-Net 1D với đường thu hẹp 4 mức AveragePooling1D kích thước 2, số kênh 32/64/128/256/512, đường mở rộng đối xứng 4 mức deconvolution, mỗi mức 3 lớp tích chập nhân 7 kèm BatchNorm, LeakyReLU (alpha = 0,01) và dropout 0,1, kết nối tắt giữa hai đường và lớp Conv1D nhân 1 ở cuối; tiếp đó phép nhúng trễ tạo K = N − L + 1 vector trễ độ dài L, trong SCTD-2Net dùng L = 4, N = 1.019, K = 1.016 với các hàng ma trận trễ của AECG và MECG xếp xen kẽ thành ma trận 1.016 × 8, còn SCTD-1Net dùng L = 8, N = 1.023, K = 1.016 và chỉ dùng ma trận trễ của AECG. Giai đoạn 2 là Ursa-Net (CNN 2D đặt theo chòm sao Đại Hùng) với đường thu hẹp 4 mức AveragePooling2D kích thước gộp (1,2) ở mức 1–3 và (2,1) ở mức 4, mỗi mức 3 lớp tích chập 2D nhân (3,7), số kênh 16/32/64/128/256, đường mở rộng đối xứng có deconvolution và kết nối tắt, cùng khối tái tạo gồm 3 lớp tích chập nối tiếp kèm BatchNorm, LeakyReLU và AveragePooling2D kích thước (2,1) đưa ma trận từ 1.016 × 8 về 1.016 × 1 trước lớp Conv2D nhân 1 sinh FECG; hai mạng được huấn luyện độc lập với cùng tham số (hàm mất mát MSE, Adam, lr 0,001, batch 64, 60 epoch) và đỉnh được dò bằng Pan-Tompkins trên FECG đã tái tạo.\par\vspace{3pt}
\noindent \textbf{Mô hình.} SCTD-Net gồm SCTD-1Net (chỉ Ursa-Net 2D) và SCTD-2Net (U-Net 1D trích MECG kết hợp Ursa-Net 2D trích FECG). Bảng 3 báo cáo đầy đủ: SCTD-1Net có 7,68 triệu tham số, dung lượng 92,78 MB và 3,91 GFLOPs mỗi lần suy luận; SCTD-2Net có 9,22 triệu tham số, 111,79 MB và 4,63 GFLOPs — đây là bài duy nhất trong ba bài báo cáo đủ ba chỉ số này, tuy độ trễ thực tế (ms trên CPU cụ thể) thì bài không nêu.\par\vspace{3pt}
\noindent \textbf{Dữ liệu.} Huấn luyện hoàn toàn trên dữ liệu tổng hợp sinh bằng bộ công cụ FECGSYN (Behar/Andreotti), mỗi mô phỏng có 34 kênh gồm 2 kênh MECG tham chiếu và 32 kênh AECG, dùng 3 kịch bản sinh lý (case 0 cơ bản kèm nhiễu, case 1 thai cử động kèm nhiễu, case 2 tăng/giảm nhịp mẹ hoặc nhịp thai kèm nhiễu) với các nguồn nhiễu cơ MA, trôi đường cơ sở BW và nhiễu chuyển động điện cực EM; tổng 55.000 mẫu huấn luyện chia theo chỉ số SINR thành nhóm chất lượng cao 20.000 mẫu (SINR\_f trong [−5; 0] dB, SINR\_m trong [−5; 5] dB), nhóm trung bình 25.000 mẫu (SINR\_f trong [−15; −5], SINR\_m trong [0; 10]) và nhóm thấp 10.000 mẫu (SINR\_f trong [−25; −15], SINR\_m trong [10; 15]). Kiểm thử chỉ trên dữ liệu thật gồm PCDB tức CinC 2013 Set A (75 bản ghi a01–a75, 4 kênh mỗi bản ghi thành 300 kênh, 1 kHz, 60 giây mỗi bản ghi, đánh giá trên đúng 80 kênh AECG, cùng tập con với AECG-DecompNet, RCED-Net, GAN, W-Net và SCTD-ICA) và ADFECGDB (5 thai phụ r01, r04, r07, r08, r10, tuổi thai 38–41 tuần, mỗi bản ghi 5 phút tại 1 kHz, 4 kênh bụng và 1 kênh điện cực da đầu thai làm tham chiếu, loại 3 kênh quá nhiễu r01-Abdomen1, r07-Abdomen1, r10-Abdomen3 theo tiền lệ của RCED-Net, còn lại 17 kênh).\par\vspace{3pt}
\noindent \textbf{Giao thức đánh giá.} Đây là giao thức sạch nhất trong ba bài và cũng là giao thức duy nhất có thể đối chiếu công bằng: huấn luyện 100\% trên dữ liệu tổng hợp FECGSYN (55.000 mẫu), không dùng một mẫu thật nào, rồi đánh giá zero-shot trên PCDB và ADFECGDB hoàn toàn không tinh chỉnh — nhóm tác giả nói rõ SCTD-2Net "chỉ được huấn luyện trên bộ FECGSYN và được đánh giá bằng các bộ dữ liệu thực tế, bảo đảm một đánh giá thực tế hơn về khả năng tổng quát hoá", và công khai phê bình các bài khác (TCGAN, CBLS-CycleGAN, Improved CycleGAN, DAW-Net, Attention R2W-Net tức p06, và dual-branch Transformer-CNN) rằng "một phần bộ ADFECGDB đã được dùng để huấn luyện trong các mô hình này, điều nhiều khả năng đã góp phần làm tăng hiệu năng của chúng trên bộ dữ liệu này". Dung sai là 50 ms theo chuẩn CinC 2013, trùng khớp với nhóm ta, và công thức F1 = 2·PPV·Se/(PPV+Se) được in đúng. Đầu vào một kênh (SCTD-2Net dùng thêm MECG do chính mạng ước lượng ra từ kênh đó chứ không phải kênh thứ hai đo được), đơn vị phân tích là kênh (80 kênh PCDB, 17 kênh ADFECGDB) chứ không phải bản ghi — khác biệt quan trọng với nhóm ta; bài không có kiểm định thống kê, không có p-value hay khoảng tin cậy khi so với các phương pháp khác, và không có xác thực chéo trên dữ liệu thật (không cần, vì không huấn luyện trên dữ liệu thật).\par\vspace{3pt}
\noindent \textbf{Kết quả chính.} SCTD-2Net đánh giá zero-shot đạt Se 97,19\% / PPV 98,80\% / F1 97,97\% trên PCDB (80 kênh) và Se 95,17\% / PPV 97,94\% / F1 96,52\% trên ADFECGDB (17 kênh), vượt SCTD-1Net ở cả ba chỉ số, rõ nhất ở độ nhạy với mức cao hơn khoảng 1,10 điểm Se trên PCDB và 1,56 điểm Se trên ADFECGDB (con số tuyệt đối của SCTD-1Net trong Bảng 2 là ảnh nên không trích xuất được; suy luận từ độ lệch: Se khoảng 96,09\% trên PCDB và khoảng 93,61\% trên ADFECGDB). So với AECG-DecompNet, SCTD-2Net cao hơn 3,88 điểm Se, 1,58 điểm PPV và 2,74 điểm F1; so với FastICA-AEFLN-NWT (dùng 3–4 kênh, chỉ đánh giá trên 16 bản ghi) cao hơn 1,77 điểm F1 và làm trên 80 bản ghi; còn so với Improved CycleGAN và Attention R2W-Net (p06) thì Se và F1 thấp hơn nhưng PPV cao hơn lần lượt 0,56 và 1,06 điểm, kèm lưu ý rõ ràng rằng các bản ghi được dùng có thể không giống nhau. Về chi phí tính toán (Bảng 3), SCTD-1Net có 7,68 triệu tham số / 92,78 MB / 3,91 GFLOPs và SCTD-2Net có 9,22 triệu tham số / 111,79 MB / 4,63 GFLOPs, mức tăng độ phức tạp từ 1Net lên 2Net được nhận xét là khiêm tốn và được bù lại bằng độ chính xác; Bảng 4 so sánh với 7 phương pháp truyền thống và 11 phương pháp học sâu là ảnh nên không trích xuất được số liệu từng dòng.\par\vspace{3pt}
\noindent \textbf{Điểm yếu.} Kết quả tuyệt đối trên ADFECGDB (F1 96,52\%) thấp hơn đáng kể các bài có tinh chỉnh — chính nhóm tác giả thừa nhận thua TCGAN, CBLS-CycleGAN, Improved CycleGAN, DAW-Net, Attention R2W-Net và dual-branch Transformer-CNN trên bộ này, tuy lập luận (đúng) rằng đó là do khác biệt quy trình huấn luyện chứ không phải do mô hình kém; ngoài ra mô hình rất nặng, 9,22 triệu tham số và 111,79 MB cho một bài toán một kênh nên không thể triển khai trên thiết bị đeo, và bài báo cáo FLOPs nhưng không đo độ trễ thực tế (ms) trên bất kỳ CPU nào nên không biết có chạy được thời gian thực hay không. Phần nhúng trễ vừa có lỗi số vừa thiếu biện minh: đoạn mô tả ghi tín hiệu AECG có N = 1.019 và vector trễ độ dài L = 4 cho K = 1.016 (đúng), trong khi phần chuẩn bị dữ liệu trước đó lại ghi đoạn dài 1.024 mẫu, ba con số 1.024 / 1.023 / 1.019 xuất hiện lẫn lộn mà không giải thích cách chuyển đổi nên bước này không tái lập chính xác được; độ trễ giữa các vector chỉ là 1 mẫu (4 ms tại 250 Hz) nên toàn bộ ma trận 2D chỉ bao phủ 8 × 4 = 32 ms lịch sử, một cửa sổ cực ngắn so với chu kỳ tim thai khoảng 430 ms, và tuy viện dẫn định lý Takens họ không chọn độ trễ theo bất kỳ tiêu chí nào của Takens (không dùng thông tin tương hỗ, không dùng hàm tự tương quan, không quét L) — chính họ cũng thừa nhận trong chú thích Hình 4 rằng "vì độ trễ giữa các vector liên tiếp nhỏ nên khác biệt giữa chúng trông rất tinh vi". Suy luận (không phải dữ kiện trong bài): biểu diễn 2D ở đây không mang thông tin hình học động lực thật mà gần như tương đương với việc cho mạng một đầu vào 8 kênh gồm tín hiệu gốc, các bản dịch chuyển 1–3 mẫu và MECG ước lượng, tức một dạng "2 kênh cộng trễ" chứ không phải nhúng Takens đúng nghĩa; bài cũng không có nghiên cứu loại bỏ cho L, cho độ trễ hay cho cách xếp xen kẽ — ba siêu tham số quan trọng nhất của đóng góp chính — không có kiểm định thống kê, không có độ lệch chuẩn trong phần chữ (Bảng 2 có ghi "mean (std)" nhưng là ảnh), không báo cáo jitter/sai số định vị đỉnh, và không nêu lọc notch tần số lưới điện dù dữ liệu thật thu ở Ba Lan (50 Hz) và dải thông kéo đến 100 Hz.\par\vspace{3pt}
\noindent \textbf{\color{tdtblue}So với mô hình của nhóm.} Đây là mốc so sánh tốt nhất trong ba bài và đối chiếu được một phần, vì bốn tham số trùng khớp hoàn toàn với FetalQRS-TCN: dung sai 50 ms, đầu vào một kênh, cửa sổ khoảng 4 giây và tần số làm việc 250 Hz. Khác biệt nằm ở nguồn huấn luyện và đơn vị phân tích: họ huấn luyện 100\% trên dữ liệu tổng hợp rồi kiểm thử zero-shot còn nhóm ta huấn luyện tách theo bản ghi trong nội bộ ADFECGDB, nên 96,52\% của họ và 97,43 ± 4,37 của nhóm ta không cùng nghĩa (nhóm ta có lợi thế phân phối, họ có lợi thế tổng quát hoá); họ tính theo kênh (17 kênh ADFECGDB sau khi loại 3 kênh nhiễu, 80 kênh PCDB) còn nhóm ta tính theo 20 đánh giá tách theo bản ghi trên 4 kênh, nên muốn so công bằng thì nhóm ta phải loại đúng 3 kênh r01-Abdomen1, r07-Abdomen1, r10-Abdomen3 giống họ và báo cáo macro F1 trên 17 kênh còn lại — thí nghiệm này rẻ và nên làm ngay. Trên CinC 2013 họ đạt 97,97\% (80 kênh, zero-shot từ dữ liệu tổng hợp) so với 77,34 ± 28,54 của nhóm ta, tức nhóm ta kém hơn khoảng 20 điểm, nhưng đó là so sánh về nguồn huấn luyện (55.000 mẫu tổng hợp trải rộng 3 mức chất lượng SINR so với 5 bản ghi) chứ không phải về kiến trúc; về chi phí, mô hình của nhóm nhỏ hơn 81 lần về tham số và 233 lần về dung lượng (113.481 tham số / 0,48 MB so với 9,22 triệu tham số / 111,79 MB / 4,63 GFLOPs), nên nếu FetalQRS-TCN đạt F1 tương đương sau khi huấn luyện trên dữ liệu tổng hợp thì đó là một luận điểm công bố rất mạnh; và về mặt khoa học, bài này là đối chứng trực tiếp cho kết quả âm tính của nhóm ta — họ dùng biểu diễn 2D là ma trận trễ, một ánh xạ tuyến tính khả nghịch giữ nguyên biên độ và mốc thời gian tuyệt đối, nên thành công, còn nhóm ta dùng biểu diễn 2D là Persistence Image, một ánh xạ mất mát thông tin bỏ đi thứ tự thời gian, nên thất bại ở mức 27,42\%, đó chính là cơ chế giải thích vì sao một cái chạy được và một cái không.\par\vspace{3pt}
\noindent \textbf{\color{tdtblue}Nhóm rút ra gì.} Thí nghiệm ưu tiên số một là huấn luyện FetalQRS-TCN trên dữ liệu tổng hợp FECGSYN theo đúng công thức của bài này — 55.000 đoạn chia theo SINR thành ba nhóm (20.000 chất lượng cao với SINR\_f trong [−5; 0], 25.000 trung bình với SINR\_f trong [−15; −5], 10.000 thấp với SINR\_f trong [−25; −15]), các kịch bản sinh lý và nhiễu MA/BW/EM, trong đó nhóm ta có thể dùng cả 7 kịch bản của FECGSYNDB thay vì 3 (thêm co thắt tử cung C3, ngoại tâm thu C4, song thai C5) để tăng tính mạnh mẽ và tạo một khác biệt bảo vệ được — rồi đo zero-shot trên ADFECGDB và CinC 2013, vì độ phân tầng theo SINR là mấu chốt ép mô hình học cả vùng tín hiệu xấu chứ không chỉ vùng đẹp. Song song đó nên sao chép giao thức loại 3 kênh r01-Abdomen1, r07-Abdomen1, r10-Abdomen3 để có một cột kết quả trên 17 kênh đặt cạnh 96,52\% của họ một cách hợp lệ, bổ sung GFLOPs vào bảng chi phí bên cạnh 113.481 tham số / 0,48 MB / 4,35 ms để đối chiếu với 3,91 và 4,63 GFLOPs của họ, trích dẫn nguyên văn câu phê bình của họ về việc các mô hình khác huấn luyện trên một phần ADFECGDB, và chỉ ra rằng khung nhúng trễ 32 ms của SCTD-Net thấp hơn gần 50 lần ngưỡng bão hoà khoảng 1,5 giây mà nhóm ta đo được (172 ms 98,73 → 748 ms 99,25 → 1516 ms 99,50 → 3052 ms 99,52, Wilcoxon p = 0,0000) mà vẫn đạt 96–98\%; đồng thời cần tránh việc gán mác "nhúng Takens" cho một phép xếp trễ 1 mẫu mà không chứng minh (nếu viết về TDA thì phải nêu rõ tiêu chí chọn độ trễ bằng thông tin tương hỗ hoặc hàm tự tương quan, đúng như bài p28 Tan 2023), tránh mô hình 9 triệu tham số và 111 MB vốn ngược với định vị của nhóm, và tránh để lẫn lộn ba con số độ dài đoạn 1.024 / 1.023 / 1.019 như họ vì mọi con số trong bài phải nhất quán từ đầu đến cuối.\par\vspace{3pt}
\vspace{6pt}
\subsection{Orvas và cs. 2025, arXiv preprint \normalfont\small (p11)}
\noindent{\footnotesize\itshape Orvas, Iulia; Radu, Andrei; Galli, Alessandra; Neacsu, Ana; Peri, Elisabetta. "A Complex UNet Approach for Non-Invasive Fetal ECG Extraction Using Single-Channel Dry Textile Electrodes." Preprint arXiv:2506.22457v1 [eess.SP], 16/06/2025, ghi ro "preprint submitted to IEEE Journal of Biomedical and Health Informatics". Don vi: SIGMA Laboratory, CAMPUS Institute, National University of Science and ...}\par\vspace{2pt}
\noindent{\small\color{okgreen}\textbf{Kết luận: so sánh trực tiếp được}}\par\vspace{4pt}
\noindent \textbf{Họ làm gì.} Bài tách tín hiệu ECG thai nhi từ ECG bụng đơn kênh thu bằng điện cực vải khô - loại không dùng gel, thoải mái khi đeo lâu tại nhà nhưng nhiễu hơn hẳn điện cực ướt Ag/AgCl - và vì không có bộ dữ liệu công khai nào dùng điện cực khô, họ tự tạo bộ dữ liệu mô phỏng bằng cách đo nhiễu thật từ 9 phụ nữ không mang thai, khử ECG người lớn để còn lại nhiễu thuần túy, mô hình hóa đặc tính phổ của nhiễu đó rồi cộng vào ECG mẹ và ECG thai tổng hợp sinh bởi fecgsyn. Mô hình đề xuất là Complex UNet (CUNet): biến tín hiệu 1 chiều thành phổ đồ bằng STFT, xử lý cả phần thực lẫn phần ảo bằng một UNet giá trị phức thay vì chỉ xử lý biên độ như các bài trước, rồi ISTFT về miền thời gian, với lập luận cốt lõi rằng thông tin pha quan trọng cho việc định vị thời gian của đỉnh R nên vứt bỏ pha là một mất mát.\par\vspace{3pt}
\noindent \textbf{Phương pháp.} Đo nhiễu từ 9 bản ghi của phụ nữ không mang thai, mỗi bản ghi dài 180 giây với 4 điện cực vải khô quanh rốn, lấy đạo trình lưỡng cực để giảm nhiễu đồng pha, lọc notch IIR bậc 10 băng thông 1 Hz tại tần số trung tâm 50 Hz rồi thông dải IIR bậc 10 từ 1-100 Hz, sau đó khử ECG người lớn bằng thuật toán của Galli và cs. (nhận diện QRS kết hợp BSS) để phần còn lại là nhiễu thuần. Nhiễu được mô hình hóa gồm nhiễu hồng trong dải 0 Hz đến một giá trị ngẫu nhiên giữa 9 và 12 Hz, nhiễu trắng đến một giá trị ngẫu nhiên giữa 60 và 90 Hz, và nhiễu hỗn hợp Gauss (nhiễu nền sigma0 = 1 trộn với nhiễu xung sigma1 = 10, xác suất chọn nhiễu xung p = 0,1), trộn trong miền tần số theo trọng số nhiễu hồng 2, nhiễu trắng 0,2 và hỗn hợp Gauss 0,15 rồi IFFT về miền thời gian; từ đó sinh 10.100 bản ghi tổng hợp dài 1 phút gồm mECG và fECG sạch từ bộ sinh fecgsyn, biên độ fECG tối đa 20 microvolt, cộng nhiễu mô hình với SNR từ 5 đến 20 dB. Kiến trúc là STFT → CUNet (các khối CDown/CUp với hệ số lấy mẫu 2 mỗi khối và skip connection kiểu UNet) → ISTFT, có lớp đặc biệt gọi là Diagonal layer thực hiện dịch pha trong mặt phẳng phức được đặt làm lớp đầu và lớp cuối, hàm kích hoạt phức là tổ hợp học được của ba hàm Complex ReLU, Georgiou-Katsageras GK(z) = z/(1+|z|) và GroupSort GS(z) = min(Re,Im) + i·max(Re,Im) với trọng số mu1 + mu2 + mu3 = 1 học trong lúc huấn luyện; kết quả được so với EKF, EKS, trừ mẫu bằng SVD, UNet và AttUNet.\par\vspace{3pt}
\noindent \textbf{Mô hình.} Complex UNet (CUNet) là một UNet giá trị phức chạy trên phổ đồ STFT phức. Bài không nêu số tham số, không nêu FLOPs, không nêu độ trễ suy luận, không nêu kích thước mô hình, không nêu số lớp m, kích thước nhân k, số kênh từng tầng hay số tầng encoder/decoder cụ thể (chỉ nói "hệ số lấy mẫu 2 mỗi khối"), không nêu tham số STFT (độ dài cửa sổ, hop length, n\_fft), không nêu tần số lấy mẫu của dữ liệu mô phỏng, và không nêu hàm mất mát, số epoch, batch size, optimizer hay learning rate; về chi phí họ chỉ nói định tính trong phần Limitations rằng huấn luyện CUNet tốn kém hơn UNet và AttUNet, phù hợp với GPU, còn muốn đưa vào thiết bị đeo thì cần tối ưu thêm.\par\vspace{3pt}
\noindent \textbf{Dữ liệu.} Ba nguồn: in-silico (dùng cho huấn luyện và kiểm thử chính) gồm 10.100 bản ghi tổng hợp dài 1 phút sinh bởi fecgsyn (Andreotti và cs. 2016) với SNR 5-20 dB và biên độ fECG tối đa 20 microvolt, trong đó 100 bản ghi giữ làm tập kiểm thử và 10.000 bản ghi còn lại để huấn luyện; in-vivo thứ nhất là 75 bản ghi bụng dài 1 phút của PhysioNet/CinC Challenge 2013 thu bằng điện cực ướt với nhiều thiết lập thiết bị khác nhau, chú thích fQRS lấy từ điện cực da đầu thai. In-vivo thứ hai là bộ điện cực khô riêng gồm 4 tình nguyện viên mang thai 37-39 tuần, mỗi bản ghi dài 21-24 phút với 4 điện cực vải khô, nghiên cứu được duyệt bởi Ủy ban Y đức Maxima Medical Centre, Veldhoven, Hà Lan (METC N19.061, 06/08/2019), nhưng bộ này chỉ được đánh giá bằng mắt và không có con số nào.\par\vspace{3pt}
\noindent \textbf{Giao thức đánh giá.} Không phải tách theo chủ thể cũng không phải tách theo bản ghi, nhưng cũng không cần, vì tập huấn luyện là mô phỏng còn tập đánh giá là dữ liệu thật, tức tách hoàn toàn về nguồn gốc: huấn luyện trên 10.000 bản ghi mô phỏng rồi đánh giá trên 100 bản ghi mô phỏng không dùng để huấn luyện, 75 bản ghi thật của CinC 2013 hoàn toàn chưa từng thấy, và 4 bản ghi điện cực khô thật, hoàn toàn không tinh chỉnh trên tập đích (bài nói rõ mô hình tách được fECG trên một bộ dữ liệu hoàn toàn mới mà không cần tinh chỉnh thêm). Dung sai khớp nhịp được nêu rõ - một đỉnh R được coi là dò đúng nếu vị trí ước lượng nằm trong vòng 50 ms so với chú thích tham chiếu, dẫn chuẩn AAMI EC38-1994 - tức ±50 ms, đúng chuẩn CinC 2013 và giống hệt dung sai của nhóm; đơn vị phân tích là nhịp chứ không phải ô thời gian, và công thức cũng viết rõ: SE = TP/(TP+FN), F-score = 2TP/(2TP+FN+FP). Điểm yếu của giao thức là bài không nêu dùng kênh nào của CinC 2013 (mỗi bản ghi có 4 kênh, họ chỉ ghi "single-channel estimates" mà không nói chọn kênh nào hay chọn như thế nào), không nêu có loại bản ghi nào khỏi 75 bản ghi hay không, và không nêu tiền xử lý nào được áp khi suy luận trên CinC 2013 (dải lọc 1-100 Hz và notch 50 Hz chỉ được mô tả cho bước đo nhiễu từ phụ nữ không mang thai).\par\vspace{3pt}
\noindent \textbf{Kết quả chính.} Trên 100 đối tượng mô phỏng (Bảng II), CUNet đạt PRD 26,2 ± 11,6; PCC 95,9 ± 5,7; F-score 99,8 ± 0,2; SE 99,6 ± 0,4; HRerr 0,40 ± 0,81 bpm, so với UNet F-score 99,3 ± 1,6 / PCC 77,4; AttUNet 98,4 ± 10,1 / PCC 77,4; SVD 92,0 ± 32,4 / PCC 75,9; EKF 88,7 ± 27,0 / PCC 65,0; EKS 88,7 ± 28,0 / PCC 66,2. Theo mức SNR (Bảng I), F-score gần như không đổi qua mọi mức (99,84 / 99,81 / 99,79 / 99,79 / 99,83 cho các dải SNR từ [-25,-20) đến [-5,0)) trong khi PRD xấu đi mạnh (80,39 → 21,89) và PCC biến động lớn (64,65 → 97,80). Trên 75 bản ghi thật của CinC 2013 theo chế độ zero-shot không tinh chỉnh (Bảng III), CUNet đạt F-score 77,8 ± 18,6; SE 72,1 ± 22,5; HRerr 14,9 ± 17,8 bpm, so với EKF F-score 80,0 ± 20,0 / SE 75,5 / HRerr 15,4; EKS 77,8 ± 20,9 / SE 75,2 / HRerr 15,3; SVD 75,6 ± 26,3 / SE 74,5 / HRerr 13,6; UNet 54,3 ± 21,0 / SE 43,4 / HRerr 23,2; AttUNet 51,6 ± 22,4 / SE 41,0 / HRerr 22,0 - tức các mô hình học sâu giá trị thực sụp đổ khi chuyển từ mô phỏng sang dữ liệu thật (99,3 → 54,3) còn CUNet giữ được 77,8, trong khi bộ điện cực khô chỉ được đánh giá trực quan (Hình 8) và không có con số nào.\par\vspace{3pt}
\noindent \textbf{Điểm yếu.} Mất đối xứng nghiêm trọng giữa mô phỏng và thực tế: F-score tụt từ 99,8 trên mô phỏng xuống 77,8 trên CinC 2013 thật, mất 22 điểm, mà bản thân bài không nhấn mạnh dù dữ liệu của chính họ cho thấy rõ; Bảng I cũng đáng ngờ khi F-score gần như bất biến 99,79-99,84 trên toàn dải SNR từ -25 dB đến 0 dB trong khi PCC rơi từ 97,80 xuống 64,65 và PRD tăng từ 21,89 lên 80,39 - suy luận của tôi là dữ liệu fecgsyn có nhịp R hoàn toàn xác định và tuần hoàn nên mạng có thể học "nhịp điệu" của bộ sinh thay vì học hình dạng QRS, đúng kiểu ngộ nhận mà nhóm đã tự phát hiện khi làm việc với cửa sổ ngắn. Bài không báo cáo bất kỳ chi phí tính toán nào (không số tham số, không FLOPs, không độ trễ, không kích thước mô hình) dù đặt mục tiêu theo dõi tại nhà và thiết bị đeo, và bài tự thừa nhận CUNet đắt hơn UNet/AttUNet; bài cũng không tái lập được vì thiếu tham số STFT, tần số lấy mẫu, số lớp, kích thước nhân, hàm mất mát, optimizer, learning rate, epoch và batch size, trong khi mã nguồn chỉ hứa công bố khi bài được chấp nhận. Đóng góp trung tâm của bài - điện cực vải khô - lại không được định lượng: chỉ 4 tình nguyện viên, đều ở tuần 37-39 (thai lớn, tín hiệu mạnh nhất, không có thai sớm), và chỉ đánh giá bằng mắt; ngoài ra bài không nêu dùng kênh nào của CinC 2013 nên nếu họ chạy cả 4 kênh rồi lấy kênh tốt nhất thì con số 77,8 sẽ bị thổi phồng, đây vẫn là preprint chưa qua bình duyệt và chưa có DOI, còn độ lệch chuẩn ±18,6 trên CinC 2013 là rất lớn, cho thấy kết quả phân cực mạnh giữa các bản ghi giống hệt hiện tượng nhóm quan sát được.\par\vspace{3pt}
\noindent \textbf{\color{tdtblue}So với mô hình của nhóm.} Đây là bài duy nhất trong ba bài so sánh trực tiếp được với mô hình của nhóm, và kết quả gần như ngang nhau: CUNet đạt F-score 77,8 ± 18,6 trên CinC 2013 zero-shot không tinh chỉnh, còn nhóm đạt 77,34 ± 28,54 trên CinC 2013 xuyên bộ dữ liệu không tinh chỉnh, chênh lệch 0,28 điểm tức thực tế là ngang nhau, với các điều kiện khớp nhau về dung sai ±50 ms, đơn vị phân tích là nhịp, công thức F-score = 2TP/(2TP+FN+FP), đầu vào đơn kênh và không tinh chỉnh trên tập đích. Khác biệt cần lưu ý là họ huấn luyện trên 10.000 bản ghi mô phỏng còn nhóm huấn luyện trên dữ liệu thật của ADFECGDB, tức hai hướng xuất phát khác nhau nhưng về cùng một điểm. Nhóm có hai lợi thế rõ - báo cáo đầy đủ chi phí tính toán (113.481 tham số, 0,48 MB, 4,35 ms trên CPU) trong khi họ không báo cáo gì, và có kết quả tách theo bản ghi trong bộ trên ADFECGDB (97,43 và 99,18) trong khi họ không có con số nào trên dữ liệu thật ngoài CinC 2013 - nhưng độ lệch chuẩn của nhóm (28,54) lớn hơn của họ (18,6), tức kết quả của nhóm phân cực mạnh hơn và đây là điểm yếu nhóm cần thừa nhận, còn họ hơn nhóm ở đóng góp về điện cực khô, mô hình giá trị phức và so sánh với 5 phương pháp baseline trên cùng tập.\par\vspace{3pt}
\noindent \textbf{\color{tdtblue}Nhóm rút ra gì.} Đây là mốc chuẩn chính xác nhất mà nhóm cần: hãy đưa con số 77,8 ± 18,6 của Orvas 2025 vào bảng so sánh của nhóm ngay cạnh 77,34 ± 28,54, kèm ghi chú rõ ràng là cùng ±50 ms, cùng đơn kênh, cùng zero-shot trên CinC 2013 - điều này biến con số 77,34 từ chỗ "có vẻ thấp" thành "ngang bằng trình độ hiện tại năm 2025", đồng thời nên bổ sung chỉ số HRerr (bpm) vì đó là sai số nhịp tim thai mà bác sĩ thực sự quan tâm và cho phép so trực tiếp với 14,9 ± 17,8 bpm của họ, và nên chạy ít nhất một baseline trừ mẫu trên CinC 2013 vì trên dữ liệu thật khó thì EKF (phương pháp cổ điển) đạt 80,0, cao hơn cả CUNet 77,8, tức học sâu chưa vượt được phương pháp cổ điển. Việc UNet và AttUNet sụp từ 99,3 xuống 54,3 khi chuyển từ mô phỏng sang thật là bằng chứng bên ngoài cho luận điểm của nhóm rằng số liệu trên dữ liệu mô phỏng không nói lên điều gì, còn cách họ tính độ lệch chuẩn theo từng bản ghi làm lộ ra sự phân cực - nhóm nên dùng đúng cách trình bày này và coi hiện tượng 4 bản ghi F1 = 100 cùng 3 bản ghi 24-49\% là một phát hiện chứ không phải một thất bại, rồi nói thẳng sang giải pháp chỉ số chất lượng tín hiệu thai có quyền từ chối trả lời; đồng thời tránh ba lỗi của họ: không báo cáo chi phí tính toán, lấy số liệu mô phỏng làm kết quả chính, và không nêu rõ dùng kênh nào (nhóm đã làm đúng với giao thức DPSS).\par\vspace{3pt}
\vspace{6pt}
\clearpage
\section{Xử lý tín hiệu cổ điển}
Các phương pháp không dùng học máy. Quan trọng vì Power-MF năm 2024 cho thấy hướng cổ điển vẫn cạnh tranh được với học sâu.

\subsection{Niknazar và cs. 2013, IEEE Trans Biomed Eng \normalfont\small (p16)}
\noindent{\footnotesize\itshape Niknazar, Mohammad; Rivet, Bertrand; Jutten, Christian. "Fetal ECG Extraction by Extended State Kalman Filtering Based on Single-Channel Recordings." IEEE Transactions on Biomedical Engineering, tap 60, so 5, nam 2013, trang 1345-1352. DOI: 10.1109/TBME.2012.2234456. (Ban PDF trong thu muc la ban luu tru HAL, ma hal-00831928.)}\par\vspace{2pt}
\noindent{\small\color{warnred}\textbf{Kết luận: không so trực tiếp được}}\par\vspace{4pt}
\noindent \textbf{Họ làm gì.} Bài giải quyết bài toán tách ECG thai khỏi tín hiệu bụng khi chỉ có một điện cực duy nhất, bằng cách mở rộng vector trạng thái của bộ lọc Kalman để mô hình hóa đồng thời cả ECG mẹ và ECG thai trong cùng một hệ phương trình (par-EKF, tức song song) thay vì tách tuần tự (seq-EKF: ước lượng ECG mẹ trước, trừ đi, rồi lấy ECG thai từ phần dư). Lý do là ở cách tuần tự, ECG thai bị coi là nhiễu Gauss khi đang ước lượng ECG mẹ — giả thiết sai khiến bộ lọc thất bại khi phức bộ QRS của mẹ và thai trùng nhau hoàn toàn về thời gian — còn cách song song không mắc lỗi đó vì mỗi ECG có một số hạng riêng trong mô hình; phương pháp cũng được áp dụng cho từ trường tim (MCG) song thai.\par\vspace{3pt}
\noindent \textbf{Phương pháp.} Nền tảng là mô hình động ECG tổng hợp của McSharry 2003 và Sameni 2007, trong đó mỗi ECG được biểu diễn bằng tổng của 5 hàm Gauss ứng với các sóng P, Q, R, S, T, mỗi hàm có 3 tham số (biên độ đỉnh, độ rộng, tâm pha), còn vector trạng thái của một ECG gồm pha và biên độ. Với N tín hiệu ECG trộn trong một quan sát duy nhất, vector trạng thái mở rộng thành pha và biên độ của cả N nguồn, và phương trình đo buộc bộ lọc phải giải thích tín hiệu đo được như tổng của N ECG thay vì một ECG cộng nhiễu Gauss; thuật toán dùng bộ lọc Kalman mở rộng (EKF), và nếu chấp nhận một độ trễ nhỏ thì chạy thêm bước làm mượt lùi thành EKS, vốn phi nhân quả nên không chạy trực tuyến được. Bước ước lượng tham số là cốt tử và cũng là điểm yếu chí mạng: phải dò được đỉnh R của cả mẹ lẫn thai trước (đỉnh R thai khó dò trực tiếp nên phải chạy seq-EKF để có ước lượng thô của fECG rồi mới dò trên đó), sau đó nén tuyến tính mỗi nhịp về khoảng [0, 2π), lấy trung bình các nhịp đã nén pha và khớp 15 tham số Gauss bằng bình phương tối thiểu phi tuyến.\par\vspace{3pt}
\noindent \textbf{Mô hình.} Đây không phải mạng nơ-ron mà là bộ lọc Kalman mở rộng dựa trên mô hình động phi tuyến; bài không nêu tổng số tham số một cách tường minh, cũng không nêu độ trễ, thời gian chạy hay kích thước mô hình. Suy luận đơn giản từ mô tả của bài: mỗi ECG cần 5 hàm Gauss × 3 tham số = 15 tham số hình dạng cộng 2 biến trạng thái (pha, biên độ), nên với trường hợp mẹ và một thai (N = 2) thì vector trạng thái mở rộng có 4 chiều và khoảng 30 tham số hình dạng phải khớp lại cho từng bản ghi.\par\vspace{3pt}
\noindent \textbf{Dữ liệu.} Bốn nguồn dữ liệu nhưng không nguồn nào được đánh giá định lượng về độ chính xác phát hiện đỉnh: dữ liệu tổng hợp theo mô hình lưỡng cực ba chiều của Sameni 2007 (500 Hz, 20.000 mẫu tức 40 giây mỗi tín hiệu, 100 lần thử cho mỗi giá trị tham số khảo sát, nhịp mẹ 1,1 Hz, nhịp thai 2 Hz, tỉ biên độ fECG/mECG = 0,3, dao động ngẫu nhiên 5\% về biên độ và thời lượng) dùng cho toàn bộ phần định lượng; và DaISy chỉ với duy nhất một bản ghi của một sản phụ, 8 kênh (5 bụng, 3 ngực), 250 Hz, dài 10 giây. Hai nguồn còn lại chỉ dùng để minh họa: Non-Invasive Fetal ECG Database của PhysioNet (55 bản ghi từ một đối tượng duy nhất, tuần thai 21-40, 2 kênh ngực và 3-4 kênh bụng, 1 kHz, 16 bit, máy đo lọc sẵn 0,01-100 Hz kèm notch 50 Hz) chỉ trình bày 20 giây đầu của 4 bản ghi dưới dạng hình vẽ, và dữ liệu MCG song thai (208 kênh, 1025 Hz, 30 phút) chỉ minh họa vài kênh trong 8 giây.\par\vspace{3pt}
\noindent \textbf{Giao thức đánh giá.} Bài không nêu dung sai khớp nhịp vì chưa bao giờ đo độ chính xác phát hiện đỉnh: không có Se, không có PPV, không có F1 và không có MAE trên bất kỳ bộ dữ liệu nào; cũng không có chia huấn luyện/kiểm thử dưới bất kỳ hình thức nào (không tách theo chủ thể, không tách theo bản ghi, không chia ngẫu nhiên) vì đây là phương pháp dựa trên mô hình chứ không học từ dữ liệu có nhãn. Tuy nhiên cần ghi rõ một dạng rò rỉ dữ liệu riêng: 15 tham số hình dạng của mỗi ECG được khớp lại trên chính bản ghi đang xử lý và việc khớp này cần vị trí đỉnh R thai, tức mô hình được hiệu chỉnh riêng cho từng bản ghi trước khi chạy; trên dữ liệu tổng hợp chỉ số duy nhất là mức cải thiện SNR và SIR (dB) tính theo nguyên lý trực giao, trung bình trên 100 lần thử mỗi điểm tham số, còn trên dữ liệu thật việc đánh giá hoàn toàn định tính bằng mắt thường trên hình vẽ, ngoại trừ một con số duy nhất là cosine 0,92 giữa hai không gian con của ước lượng par-EKS và ước lượng piCA trên kênh 1 của DaISy. Tuyên bố đơn kênh bị phá vỡ trong các tình huống khó: chuẩn đối chiếu piCA dùng cả 8 kênh của DaISy, và với MCG song thai chính tác giả thừa nhận nhãn chuẩn được lấy bằng cách dùng nhiều cảm biến với Fast-ICA; ngoài ra bài không nêu bất kỳ bộ lọc dải thông tiền xử lý nào của riêng họ, vì con số 0,01-100 Hz kèm notch 50 Hz chỉ là mô tả thiết bị ghi của NIFECGDB.\par\vspace{3pt}
\noindent \textbf{Kết quả chính.} Mọi con số định lượng đều là mức cải thiện SNR/SIR trên dữ liệu tổng hợp chứ không phải độ chính xác phát hiện đỉnh: khi hỗn hợp gần như sạch (công suất nhiễu -30 dB), mức cải thiện SNR tối thiểu của fECG là 40 dB, và ngay cả với hỗn hợp rất nhiễu (công suất nhiễu 20 dB) mức cải thiện vẫn trên 20 dB. Với tỉ biên độ fECG/mECG chạy từ 0,1 đến 1,0, cải thiện SIR của fECG luôn trên 30 dB ở cả EKF lẫn EKS (EKS tốt hơn chút ít); khi khảo sát nhịp thai từ 0,3 Hz đến 3,6 Hz, đường cong có ba cực tiểu địa phương sâu tại tỉ số nhịp bằng 1, 2 và 3 do tần số cơ bản của mẹ và thai là bội số của nhau, nhưng ngay tại các điểm xấu nhất này cải thiện SIR vẫn trên 20 dB. EKF tốt hơn khi tín hiệu tương đối sạch (bước làm mượt lùi gây lọc quá tay) còn EKS vượt trội rõ khi nhiễu mạnh; trên dữ liệu thật con số duy nhất là cosine 0,92 nói trên, phần còn lại là định tính — seq-EKS làm hỏng tín hiệu thai trong khoảng t = 6 s đến 7 s của DaISy nơi sóng mẹ và thai trùng nhau trong khi par-EKS thì không, và trên MCG song thai seq-EKS không phân biệt được hai thai khi chúng chồng lấn còn par-EKS thì thành công.\par\vspace{3pt}
\noindent \textbf{Điểm yếu.} Bài không có một con số F1/Se/PPV nào nên không thể tái lập về mặt số học và không thể đặt vào bảng so sánh dò đỉnh, đồng thời bài không nêu dung sai khớp nhịp vì không đo khớp nhịp. Tồn tại một vòng lặp con gà - quả trứng: par-EKS cần vị trí đỉnh R thai làm đầu vào để khớp tham số mô hình, trong khi chính tác giả viết rằng phần cốt tử của par-EKS đề xuất là việc dò đỉnh R — nghĩa là đây không phải bộ dò đỉnh mà là bộ tách/tăng cường tín hiệu, bắt buộc phải có một bộ dò đỉnh khác đi trước; tuyên bố đơn kênh cũng bị phá vỡ ở ca song thai khi nhãn chuẩn được lấy bằng Fast-ICA trên nhiều cảm biến (họ có gợi ý thay bằng cảm biến âm thanh nhưng không thực hiện). Đánh giá trên dữ liệu thật hoàn toàn là nhìn hình với 10 giây DaISy, 20 giây của 4 bản ghi NIFECGDB (chỉ một đối tượng duy nhất, không nói được gì về khả năng tổng quát hóa giữa các bệnh nhân) và 8 giây MCG; mô hình 5 hàm Gauss không chịu được biến thiên quá lớn như loạn nhịp theo chính thừa nhận của tác giả, bài không nêu bộ lọc tiền xử lý riêng và không báo cáo chi phí tính toán dù EKF/EKS trên vector trạng thái mở rộng không hề rẻ.\par\vspace{3pt}
\noindent \textbf{\color{tdtblue}So với mô hình của nhóm.} Không so sánh được, và lý do rất quan trọng về mặt khái niệm vì đây không phải cùng một bài toán: FetalQRS-TCN là bộ dò đỉnh (đầu vào là tín hiệu, đầu ra là vị trí fQRS, đo bằng F1 với dung sai 50 ms) còn par-EKS là bộ tách/tăng cường tín hiệu cần sẵn vị trí đỉnh R thai làm đầu vào, nên mô hình của nhóm có thể dùng làm đầu vào cho par-EKS chứ không cạnh tranh với nó; nếu phản biện hỏi vì sao không so với Niknazar 2013 thì câu trả lời chuẩn là bài đó không báo cáo chỉ số phát hiện đỉnh nào, chỉ báo cáo cải thiện SNR/SIR trên dữ liệu tổng hợp 40 giây và đánh giá định tính trên 10-20 giây dữ liệu thật. Nhóm mạnh hơn rõ ràng ở ba điểm: giao thức tách theo bản ghi trên 5 sản phụ chuyển dạ với 20 đánh giá, so với 1 bản ghi DaISy và 1 đối tượng NIFECGDB của họ; báo cáo đầy đủ chi phí tính toán (113.481 tham số, 0,48 MB, 4,35 ms) trong khi họ không báo cáo gì; và không cần biết trước vị trí đỉnh R thai trong khi họ cần. Điểm cần thận trọng: par-EKS xử lý được trường hợp QRS mẹ và thai trùng nhau hoàn toàn về thời gian — ca khó kinh điển mà bài này lấy làm trọng tâm — nên nhóm nên kiểm tra riêng hiệu năng của FetalQRS-TCN trên các nhịp thai rơi vào cửa sổ QRS mẹ và báo cáo con số đó.\par\vspace{3pt}
\noindent \textbf{\color{tdtblue}Nhóm rút ra gì.} Dùng bài này làm ví dụ điển hình trong phần tổng quan về việc phần lớn công trình đơn kênh không báo cáo chỉ số phát hiện đỉnh chuẩn hóa, với cách viết bám sát văn bản gốc rằng Niknazar và cộng sự (2013) báo cáo mức cải thiện SNR trên 20 dB trên dữ liệu tổng hợp nhưng không báo cáo Se/PPV/F1 trên dữ liệu lâm sàng, khiến kết quả của họ không thể đối chiếu trực tiếp với các bộ dò đỉnh; đồng thời khai thác hướng mở rộng ghép nối tầng là dùng FetalQRS-TCN để dò vị trí fQRS rồi nạp các vị trí đó vào par-EKS để tái tạo hình thái sóng ECG thai, qua đó khắc phục đúng điểm yếu con gà - quả trứng của họ, và có thể nhắc ngắn gọn rằng bộ lọc Kalman mở rộng của họ là cơ sở lý thuyết cho phương pháp TSekf. Phải tránh trích con số 40 dB hay 20 dB vào bảng so sánh F1 vì đó là đơn vị hoàn toàn khác (dB cải thiện SNR so với phần trăm F1), tránh gọi đây là phương pháp đơn kênh mà không ghi chú rằng ca song thai họ phải dùng Fast-ICA đa kênh để lấy nhãn chuẩn, và phải rút kinh nghiệm nêu đủ dải lọc, dung sai, số kênh, giao thức chia, số tham số và độ trễ để bài của nhóm không rơi vào trạng thái không tái lập được như bài này.\par\vspace{3pt}
\vspace{6pt}
\subsection{Castillo và cs. 2018, PLoS ONE \normalfont\small (p17)}
\noindent{\footnotesize\itshape Castillo, Encarnacion; Morales, Diego P.; Garcia, Antonio; Parrilla, Luis; Ruiz, Victor U.; Alvarez-Bermejo, Jose A. "A clustering-based method for single-channel fetal heart rate monitoring." PLoS ONE, tap 13, so 6, nam 2018, bai e0199308, 22 trang. DOI: 10.1371/journal.pone.0199308. Nhan 06/06/2017, chap nhan 05/06/2018, dang 22/06/2018.}\par\vspace{2pt}
\noindent{\small\color{okgreen}\textbf{Kết luận: so sánh trực tiếp được}}\par\vspace{4pt}
\noindent \textbf{Họ làm gì.} Bài đo nhịp tim thai từ một kênh ECG bụng duy nhất, với điểm độc đáo là không khử ECG mẹ: thay vì tách fECG khỏi hỗn hợp theo cách kinh điển, họ dò trực tiếp phức bộ QRS thai trên tín hiệu bụng đã khử nhiễu. Vấn đề họ nhắm tới là các phương pháp đặt ngưỡng không tự động hóa được vì ngưỡng biên độ phụ thuộc tuần thai, vị trí điện cực và vị trí tương đối của tim mẹ với tim thai ngay cả khi đã chuẩn hóa tín hiệu, nên họ chuyển thành bài toán gom cụm không giám sát: mỗi cặp cực đại địa phương theo sau bởi cực tiểu địa phương (ứng viên cặp R-S) được trích đặc trưng rồi k-medoids++ chia thành đúng ba cụm — RS của mẹ, RS của thai và sóng khác/nhiễu — nhờ đó loại bỏ hoàn toàn ngưỡng thủ công.\par\vspace{3pt}
\noindent \textbf{Phương pháp.} Giai đoạn 1 khử nhiễu bằng wavelet (kế thừa từ Castillo 2013): phân tích tới L = 7 mức, thay hệ số xấp xỉ ở mức 7 bằng vector toàn 0 để loại trôi đường nền và thành phần một chiều, áp ngưỡng mềm phổ quát với tái tỉ lệ đơn lên các hệ số chi tiết từ mức 1 đến M = 3 rồi tái tạo tín hiệu, dùng hàm wavelet db6. Giai đoạn 2 là đóng góp mới: quét toàn tín hiệu tìm các cặp cực đại - cực tiểu, tự động chọn đặc trưng bằng cách chia dải từ 0 đến biên độ max-min lớn nhất (MA) thành 50 khoảng rồi xét phân bố đã chuẩn hóa và làm mượt — nếu phát hiện hai cực đại cách nhau hơn 35\% của MA thì chọn đặc trưng biên độ, còn nếu cách nhau dưới 35\% hoặc chỉ có một cực đại thì chọn đặc trưng biên độ nhân số mẫu (mượn thêm thông tin thời gian vì khoảng RS của mẹ dài hơn của thai) — sau đó gom cụm bằng hàm kmedoids của MATLAB với thuật toán k-medoids++, số cụm bằng 3, khoảng cách Euclid bình phương và 20 lần lặp với medoid khởi tạo mới, trong đó cụm RS thai là cụm có giá trị trung vị nằm giữa hai cụm kia, rồi áp thêm giới hạn biên độ và khoảng thời gian để cải thiện phân loại. Giai đoạn 3 sửa FP/FN dựa trên giới hạn sinh lý của khoảng RR (kế thừa từ Castillo 2013); cửa sổ xử lý là 50.000 mẫu tức 50 giây ở 1 kHz, chọn sau khi khảo sát từ 10.000 đến 60.000 mẫu không thấy khác biệt đáng kể.\par\vspace{3pt}
\noindent \textbf{Mô hình.} Đây không phải mô hình học sâu và không có tham số huấn luyện theo nghĩa mạng nơ-ron, mà là chuỗi xử lý tín hiệu kết hợp gom cụm không giám sát, cài đặt bằng MATLAB ở mức chứng minh khái niệm; bài không nêu số tham số, không nêu độ trễ, không nêu kích thước mô hình và không nêu thời gian chạy. Bù lại các siêu tham số được nêu đầy đủ nên rất dễ tái lập: db6, L = 7, M = 3, ngưỡng phổ quát, ngưỡng mềm, k = 3, 20 lần lặp khởi tạo, khoảng cách Euclid bình phương, 50 khoảng chia phân bố biên độ, ngưỡng 35\% của MA và cửa sổ 50.000 mẫu.\par\vspace{3pt}
\noindent \textbf{Dữ liệu.} Hai bộ dữ liệu lâm sàng thật công khai trên PhysioNet: tập huấn luyện là ADFECGDB (5 sản phụ chuyển dạ, mỗi bản ghi gồm 4 tín hiệu vi sai từ bụng dài 5 phút cùng ECG thai trực tiếp từ da đầu thai làm tham chiếu, 1 ksps, 16 bit, băng thông tín hiệu 1-150 Hz) — chính là bộ dữ liệu nhóm đang dùng — trong đó bác sĩ tim mạch đã loại 3 tín hiệu (r04 Ab-1, r07 Ab-1, r10 Ab-3), còn lại 17 tín hiệu chia thành Nhóm 1 (11 tín hiệu không có vấn đề rõ ràng) và Nhóm 2 (6 tín hiệu nhiễu, nhiễu loạn hoặc bão hòa), tổng 10.839 phức bộ fQRS. Tập kiểm thử là Challenge 2013 Training Set A (bản ghi 1 phút, 1 ksps, 4 kênh bụng không xâm lấn, có nhãn vị trí sóng R thai), từ đó bác sĩ tim mạch chọn ra 64 tín hiệu-kênh thuộc 26 bản ghi riêng biệt với 9.103 phức bộ fQRS; tổng thời lượng là 55 phút cho ADFECGDB và 64 phút cho Challenge Set A, cộng lại 119 phút.\par\vspace{3pt}
\noindent \textbf{Giao thức đánh giá.} Dung sai khớp nhịp là 50 ms (nguyên văn: mọi ứng viên lệch quá 50 ms so với nhãn tham chiếu đều bị coi là FP), trùng khớp hoàn toàn với dung sai của nhóm, và có chia chéo bộ dữ liệu theo hướng tốt: ADFECGDB dùng để dò các siêu tham số (ngưỡng 35\%, 50 khoảng, 20 lần lặp, kích thước cửa sổ) còn Challenge 2013 Set A làm tập kiểm thử, với đúng một bộ tham số duy nhất áp cho cả hai nên không có tinh chỉnh trên tập đích. Phương pháp thực sự đơn kênh — mỗi tín hiệu-kênh được xử lý độc lập, không gộp kênh, không dùng kênh ngực tham chiếu — nhưng có ba vấn đề nghiêm trọng về giao thức: không có tách theo chủ thể hay tách theo bản ghi nên con số trên ADFECGDB là con số trong mẫu; bác sĩ tim mạch lọc bỏ bản ghi bằng tay (3/20 kênh ADFECGDB và chỉ giữ 26/75 bản ghi Challenge Set A, tức bỏ khoảng hai phần ba bộ dữ liệu); và họ báo cáo thêm con số cho kênh tốt nhất của mỗi bản ghi, tuy có đặt cạnh con số đầy đủ nên vẫn trung thực. Bài không nêu bộ lọc dải thông theo nghĩa cổ điển mà chỉ nêu tham số wavelet (băng thông 1-150 Hz là của thiết bị ghi ADFECGDB, không phải lựa chọn xử lý của tác giả); suy luận từ mô tả thì việc triệt tiêu hệ số xấp xỉ mức 7 ở 1 kHz tương đương bỏ thành phần dưới khoảng 4 Hz, còn cửa sổ xử lý dài 50 giây so với 4 giây của nhóm và chi phí tính toán thì bài không nêu.\par\vspace{3pt}
\noindent \textbf{Kết quả chính.} Trên ADFECGDB với dung sai 50 ms: toàn bộ 17 tín hiệu (10.839 fQRS, 10.000 TP, 839 FN, 413 FP) đạt Se 92,26\%, PPV 96,03\%, Acc 88,87\%, F1 94,11\%; riêng Nhóm 1 gồm 11 tín hiệu sạch (6.990 fQRS, 175 FN tức 2,50\%, 98 FP tức 1,40\%) đạt Se 97,50\%, PPV 98,58\%, Acc 96,15\%, F1 98,04\%; riêng kênh tốt nhất mỗi bản ghi (5 tín hiệu, 3.182 fQRS) đạt Se 98,40\%, PPV 98,86\%, Acc 97,30\%, F1 98,63\%, với hai tín hiệu tệ nhất là r08 Ab-2 (F1 76,91\%) và r01 Ab-2 (F1 79,80\%). Trên Challenge 2013 Set A không tinh chỉnh: toàn bộ 64 tín hiệu-kênh (9.103 fQRS, 227 FN, 123 FP) đạt Se 97,51\%, PPV 98,63\%, Acc 96,21\%, F1 98,07\%, còn kênh tốt nhất mỗi bản ghi (26 tín hiệu, 3.627 fQRS, 75 FN, 32 FP) đạt Se 97,93\%, PPV 99,11\%, Acc 97,08\%, F1 98,52\%, trong đó 16 tín hiệu-kênh đạt F1 tuyệt đối 100\% và tệ nhất là a67 Ab-4 (91,53\%), a70 Ab-1 (91,37\%), a65 Ab-2 (92,53\%), a55 Ab-3 (92,75\%). Bảng so sánh với các phương pháp khác do chính tác giả tổng hợp từ báo cáo gốc của những nhóm khác chứ không chạy lại: STVD 89,90, EKF 81,10, TVD+TSPCA 87,40, TSPCA 86,70 (trên 68 bản ghi 1 phút); TSLpca 95,04/89,30, LMS 95,40/87,90, RLS 95,90/88,20, ESNa 97,20/90,20 (14 bản ghi 1 phút so với 11 bản ghi 5 phút); TA 97,30/85,40 và EKS 93,00/86,30 (75 bản ghi 1 phút so với 24 bản ghi 20 phút); EKS+DE+ANFIS 95,12/91,82 và EKF+DE+ANFIS 94,03/90,32; SDV+SW 99,61 nhưng chỉ trên 1 bản ghi 5 phút; CEEMDAN 98,6 và 100 nhưng chỉ trên vài bản ghi dài 30-40 giây; cùng phương pháp đặt ngưỡng trước đó của chính họ (THR) 96,76/98,24/97,50.\par\vspace{3pt}
\noindent \textbf{Điểm yếu.} Điểm yếu lớn nhất là việc lọc bản ghi bằng bác sĩ: chỉ dùng 26/75 bản ghi Challenge Set A (khoảng 35\%) và loại 3/20 kênh ADFECGDB, mà chính các bản ghi bị loại mới là các bản ghi khó, nên F1 98,07\% không phải hiệu năng trên toàn bộ Set A mà là hiệu năng trên một tập con đã được làm sạch thủ công — bài có ghi rõ điều này nhưng phần tóm tắt và bảng so sánh lại không mang theo cảnh báo đó. Không có tách theo chủ thể hay tách theo bản ghi nên con số ADFECGDB là con số trong mẫu, và con số tổng bộ 94,11\% (vốn vẫn đã loại 3 tín hiệu) bị lu mờ trước con số 98,04\% được nhắc trong kết luận, trong khi hiệu năng thật trên cả 20 kênh không bao giờ được báo cáo. Ngoài ra cửa sổ 50 giây khiến không thể theo dõi thời gian thực với độ trễ thấp (với bản ghi Challenge chỉ dài 60 giây thì cửa sổ này gần bằng cả bản ghi), bài không đo jitter nên chỉ biết đỉnh dò được nằm trong ±50 ms chứ không biết độ chính xác về thời điểm, không báo cáo chi phí tính toán, bảng so sánh trộn kết quả của các tác giả khác trên những tập dữ liệu rất khác nhau (từ 1 đến 150 bản ghi, từ 30 giây đến 20 phút), và phương pháp dựa vào giả thiết khoảng RS của mẹ dài hơn đáng kể khoảng RS của thai — giả thiết có thể vô hiệu với thai gần đủ tháng hoặc nhịp mẹ nhanh.\par\vspace{3pt}
\noindent \textbf{\color{tdtblue}So với mô hình của nhóm.} Đây là bài so sánh trực tiếp được nhất trong ba bài vì cùng bộ dữ liệu ADFECGDB, cùng dung sai 50 ms, cùng công thức F1 và thực sự đơn kênh: trên ADFECGDB với tất cả các kênh, Castillo đạt 94,11\% (17/20 tín hiệu, trong mẫu) so với 97,43 ± 4,37 của FetalQRS-TCN (tách theo bản ghi, 20 đánh giá trên đủ 4 kênh) nên nhóm cao hơn khoảng 3,3 điểm và giao thức của nhóm nghiêm ngặt hơn, còn ở kênh tốt nhất là 98,63\% so với 99,18 của nhóm (giao thức DPSS), tức nhóm cao hơn 0,55 điểm. Ngược lại ở trục xuyên bộ dữ liệu sang CinC 2013 không tinh chỉnh, họ đạt 98,07\% còn nhóm chỉ đạt 77,34 ± 28,54 — con số nhóm phải đối mặt trung thực — và có ba lý do giải thích: họ chỉ chạy trên 26/75 bản ghi do bác sĩ chọn nên đã loại hết bản ghi khó; phương pháp của họ không giám sát và không có trọng số học được nên không chịu dịch chuyển miền, trong khi mô hình học sâu của nhóm học đặc trưng của 5 sản phụ ADFECGDB nên tụt mạnh khi sang bộ dữ liệu khác thiết bị, khác tuần thai, khác tư thế (đây là điểm yếu cơ cấu của học sâu chứ không phải lỗi cài đặt); và họ có giai đoạn sửa FP/FN theo giới hạn sinh lý của chuỗi RR trong khi hậu xử lý của nhóm hiện chỉ có lấy đỉnh cùng thời gian trơ 250 ms. Nhóm mạnh hơn rõ ràng ở bốn điểm: jitter 3,05 ms trong khi họ không đo jitter; giao thức tách theo bản ghi mà họ không có; báo cáo đầy đủ chi phí tính toán (113.481 tham số, 0,48 MB, 4,35 ms cho cửa sổ 4 giây) trong khi họ không báo cáo gì; và cửa sổ 4 giây so với 50 giây của họ, tức nhóm có thể theo dõi gần thời gian thực còn họ thì không.\par\vspace{3pt}
\noindent \textbf{\color{tdtblue}Nhóm rút ra gì.} Nên đưa bài này vào bảng so sánh chính với chú thích đầy đủ (Castillo 2018, ADFECGDB, đơn kênh, dung sai 50 ms, F1 98,04\% trên 11 tín hiệu do bác sĩ chọn và 94,11\% trên 17 tín hiệu, không tách theo bản ghi) vì đây là baseline cổ điển mạnh nhất và sát nhất với bài toán của nhóm, đồng thời bắt chước cách báo cáo ba tầng của họ (toàn bộ, nhóm sạch, kênh tốt nhất) kèm thêm cột tỉ lệ đoạn được chấp nhận, thêm bước sửa FP/FN theo chuỗi RR vào hậu xử lý hiện chỉ có ngưỡng và thời gian trơ 250 ms, và kiểm tra riêng hai tín hiệu mà họ thất bại nặng là r08 Ab-2 (F1 76,91\%) và r01 Ab-2 (F1 79,80\%). Đây cũng là luận cứ mạnh nhất cho hướng fSQI vì họ phải nhờ một bác sĩ tim mạch loại bỏ tín hiệu xấu bằng tay (3/20 kênh ADFECGDB và 49/75 bản ghi CinC) trong khi nhóm tự động hóa bước này bằng chỉ số chất lượng tín hiệu thai kèm quyền từ chối trả lời và báo cáo rõ độ phủ đi kèm độ chính xác; tuyệt đối không bắt chước việc để chuyên gia lọc bản ghi rồi báo cáo như hiệu năng tổng quát (nếu có lọc thì phải bằng tiêu chí tự động và phải báo cáo cả tỉ lệ bị loại lẫn kết quả không lọc), và không được giấu việc thua ở trục xuyên bộ dữ liệu (77,34 so với 98,07).\par\vspace{3pt}
\vspace{6pt}
\subsection{Sulas và cs. 2020, Math Biosci Eng \normalfont\small (p19)}
\noindent{\footnotesize\itshape Eleonora Sulas, Monica Urru, Roberto Tumbarello, Luigi Raffo, Danilo Pani. "Systematic analysis of single- and multi-reference adaptive filters for non-invasive fetal electrocardiography". Mathematical Biosciences and Engineering (MBE), Volume 17, Issue 1, trang 286-308, 2020. DOI: 10.3934/mbe.2020016. Nhan bai 10/06/2019, chap nhan 26/09/2019, xuat ban 09/10/2019. Dai hoc Cagliari (Y) + Benh ...}\par\vspace{2pt}
\noindent{\small\color{warnred}\textbf{Kết luận: không so trực tiếp được}}\par\vspace{4pt}
\noindent \textbf{Họ làm gì.} Bài không đề xuất mô hình mới mà trả lời một câu hỏi kỹ thuật hẹp: khi dùng bộ lọc thích nghi để xóa ECG mẹ khỏi tín hiệu bụng, một điện cực ngực tham chiếu có đủ hay bắt buộc phải dùng nhiều điện cực ngực. Nhóm tác giả so sánh phiên bản đơn tham chiếu (SR, 1 kênh ngực) và đa tham chiếu (MR, 3 kênh ngực không đồng phẳng) của cùng bộ lọc QRD-RLS, và kết luận SR luôn thua MR một cách có ý nghĩa thống kê, không đủ tin cậy ngay cả khi chỉ cần đo nhịp tim thai, nên khuyến cáo không dùng SR trong mọi trường hợp.\par\vspace{3pt}
\noindent \textbf{Phương pháp.} Tín hiệu được thu bằng máy Porti7 (TMSi) ở 2048 Hz, độ phân giải 22 bit, và chỉ qua duy nhất một bước tiền xử lý là lọc thông cao 1 Hz bằng FIR equiripple bậc 1124 chạy hai chiều (pha bằng không), không có dải thông và không có notch 50 Hz trước khi lọc thích nghi. Bộ lọc QRD-RLS chạy ở cấu hình khử nhiễu thích nghi với đầu vào d(k) là kênh bụng, tham chiếu x(k) là kênh ngực, đầu ra y(k) xấp xỉ thành phần mẹ và sai số e(k) chính là fECG, dùng phân rã QR (phép quay Givens / Householder) thay cho nghịch đảo ma trận để tránh mất ổn định số học; SR dùng 1 kênh ngực (kênh 1m, vuông góc mặt phẳng dọc) còn MR dùng cả 3 kênh ngực (1m, 2m, 3m) theo kiểu đệ quy lần lượt từng kênh. Nhóm đo ba chỉ số là SIR, độ suy giảm mECG (dB) và độ chính xác phát hiện fQRS bằng thuật toán jqrs của Behar, sau đó kiểm định Lilliefors để xét phân phối chuẩn rồi dùng Wilcoxon signed rank ở mức ý nghĩa p < 0,05.\par\vspace{3pt}
\noindent \textbf{Mô hình.} Bộ lọc thích nghi QRD-RLS, không phải mạng nơ-ron và không dùng học sâu, với hệ số quên lambda = 0,999 và số tạp N = 20 cho mỗi kênh, nghĩa là MR có 60 hệ số còn SR có 20 hệ số. Bài không nêu số tham số của bất kỳ mô hình học máy nào vì không dùng học máy; toàn bộ xử lý chạy trên Matlab.\par\vspace{3pt}
\noindent \textbf{Dữ liệu.} Bộ dữ liệu thật là bộ riêng tư của dự án NInFEA (Đại học Cagliari): 28 tín hiệu đa kênh từ 20 thai phụ khỏe mạnh, tuổi thai 21 đến 27 tuần, tổng 112 kênh bụng, mỗi bản ghi dài cố định 10 giây, 2048 Hz, 22 bit, cấu hình 6 điện cực cho 3 kênh lưỡng cực ở ngực và 3 điện cực bụng cho 3 kênh lưỡng cực cùng 1 kênh đơn cực, với nhãn chuẩn fQRS xác định nhờ tín hiệu Doppler xung thu đồng thời bằng máy siêu âm Philips iE33. Bộ tổng hợp sinh bởi FECGSYN ở 2048 Hz, mỗi tín hiệu 10 giây, SNR thai/mẹ = -18 dB, 5 mức SNR mẹ/nhiễu (3, 5, 9, 12, 15 dB), nhịp mẹ cố định 90 bpm, nhịp thai cố định 150 bpm, 20 vị trí tim thai cho mỗi mức, cho 100 đối tượng ảo nhân 4 kênh bụng bằng 400 kênh tổng hợp, có thêm nhiễu chuyển động, dịch chuyển điện cực và trôi đường nền.\par\vspace{3pt}
\noindent \textbf{Giao thức đánh giá.} Không có chia train/test, không tách theo chủ thể, không tách theo bản ghi và không k-fold vì đây không phải bài học máy; tuy nhiên lambda, N và việc chọn bộ phát hiện jqrs đều được tối ưu trong "giai đoạn huấn luyện" trên chính bộ dữ liệu thật mà sau đó bài báo cáo kết quả, tức rò rỉ dữ liệu ở cấp siêu tham số và không có tập kiểm tra tách riêng. Về dung sai khớp nhịp thì bài không nêu: không có con số ms nào được định nghĩa làm ngưỡng khớp giữa đỉnh phát hiện và nhãn chuẩn, còn hai con số 40 ms và 100 ms xuất hiện trong bài chỉ là bề rộng cửa sổ mẫu của fQRS và mQRS khi tính SIR, nên không thể tái lập con số độ chính xác của họ. Đơn vị phân tích là từng kênh bụng riêng biệt (112 kênh thật, 400 kênh tổng hợp) gộp lại theo loại kênh ngang, dọc, chéo và đơn cực, với công thức bất thường Acc = (TP+TN)/(TP+TN+FP+FN) trong khi bài không định nghĩa TN cho bài toán phát hiện đỉnh.\par\vspace{3pt}
\noindent \textbf{Kết quả chính.} Trên dữ liệu thật (trung vị), độ chính xác phát hiện fQRS của SR so với MR là 0,62 so với 0,85 ở kênh ngang, 0,69 so với 0,92 ở kênh dọc, 0,61 so với 0,89 ở kênh chéo và 0,70 so với 0,86 ở kênh đơn cực, dẫn đến kết luận SR chỉ đạt "68\% trung bình, sigma = 0,0097". Độ suy giảm mECG (dB) cũng nghiêng hẳn về MR (kênh ngang 15,02 so với 25,75; dọc 12,08 so với 33,79; chéo 13,99 so với 35,09; đơn cực 12,16 so với 31,50), SIR kênh ngang đi từ -8,5 dB gốc lên -0,87 dB với SR và +1,79 dB với MR, và mọi kiểm định Wilcoxon trên dữ liệu thật đều cho p = 0,0000. Trên dữ liệu tổng hợp, SIR và suy giảm mECG vẫn khác biệt có ý nghĩa (p = 0,0000) nhưng độ chính xác phát hiện fQRS không còn khác nhau (p = 0,5800 tổng thể; từng kênh 0,4514 / 0,2055 / 0,4115 / 0,8280) và đều rất cao (SR 0,90-0,93, MR 0,92-0,96), tức ngược hẳn với dữ liệu thật.\par\vspace{3pt}
\noindent \textbf{Điểm yếu.} Bài không nêu dung sai khớp nhịp bằng ms nên không thể tái lập con số độ chính xác, lại dùng công thức Acc có chứa TN mà không định nghĩa TN cho bài toán phát hiện đỉnh, khiến kết quả khó diễn giải và khó đối chiếu với F1/Se/PPV của các công trình khác. Siêu tham số lambda, N và lựa chọn bộ phát hiện jqrs đều được tinh chỉnh trên chính bộ dữ liệu thật được báo cáo, trong khi bộ dữ liệu đó lại riêng tư, chỉ 20 thai phụ, tuổi thai 21-27 tuần và mỗi bản ghi chỉ 10 giây, nên không tái lập độc lập được và không suy rộng cho theo dõi liên tục hay giai đoạn chuyển dạ. Tiền xử lý chỉ có lọc thông cao 1 Hz, giữ nguyên dải thông tới khoảng 550 Hz và không notch 50 Hz nên giữ lại toàn bộ sóng P, T của mẹ và nhiễu đường điện, chính tác giả thừa nhận đầu ra "vẫn còn nhiễu dư và nhiễu đường điện"; bài không báo cáo bất kỳ chi phí tính toán nào, và kết quả trên bộ tổng hợp mâu thuẫn với bộ thật ở một số kết luận (kênh chéo tốt nhất trên tổng hợp, kênh ngang tốt nhất trên thật).\par\vspace{3pt}
\noindent \textbf{\color{tdtblue}So với mô hình của nhóm.} Không so sánh trực tiếp được, vì bốn lý do độc lập: họ dùng NInFEA riêng tư với 20 thai phụ tuần 21-27 và nhãn chuẩn từ Doppler xung còn nhóm dùng ADFECGDB với 5 sản phụ chuyển dạ và nhãn chuẩn từ điện cực da đầu thai nhi; độ đo của họ là Acc = (TP+TN)/(TP+TN+FP+FN) không rõ cách đếm TN trong khi nhóm báo cáo macro F1, Se và PPV; họ không nêu dung sai khớp nhịp nên không biết khoan dung hơn hay chặt hơn mức ±50 ms của nhóm; và ràng buộc phần cứng khác hẳn vì cấu hình "đơn tham chiếu" của họ vẫn cần tối thiểu 3-4 điện cực (1 điện cực bụng cộng 1 kênh ngực lưỡng cực) còn MR cần 6 điện cực ngực cộng điện cực bụng, trong khi FetalQRS-TCN chỉ dùng 1 kênh bụng và không có điện cực ngực nào. Về mặt luận điểm (đây là suy luận, không phải đối chiếu số học), bài này là bằng chứng mạnh nhất trong y văn rằng tuyến lọc thích nghi với ít kênh là ngõ cụt, khi SR chỉ đạt 68\% độ chính xác trung bình trên dữ liệu thật, còn mô hình của nhóm đạt 97,43\% macro F1 (tách theo bản ghi, toàn bộ kênh) và 99,18\% ở kênh lâm sàng tốt nhất với số kênh ít hơn. Điểm mạnh của nhóm không nằm ở chỗ khử mẹ tốt hơn mà ở chỗ không đòi hỏi khử mẹ phải hoàn hảo, vì mạng TCN học cách phát hiện đỉnh trên tín hiệu dư còn lẫn tạp mẹ; ngoài ra có một giả thuyết cần kiểm chứng là một phần con số SR thấp của họ đến từ tiền xử lý (chỉ lọc thông cao 1 Hz và giữ dải tới khoảng 550 Hz, tương ứng vùng xấu nhất trong khảo sát của nhóm, nơi dải 0,5-100 Hz chỉ đạt 72,17 điểm F1, thấp hơn dải 10-60 Hz tới 12,4-18,9 điểm), song không thể khẳng định vì họ không chạy thí nghiệm đổi dải lọc.\par\vspace{3pt}
\noindent \textbf{\color{tdtblue}Nhóm rút ra gì.} Trích dẫn bài này làm mốc đối chứng cho luận điểm "1 kênh cộng học sâu" trong phần Related Work (lọc thích nghi đơn tham chiếu chỉ đạt khoảng 68\% độ chính xác trên dữ liệu thật dù đã thêm điện cực ngực, trong khi nhóm đạt 97,43\% F1 mà không cần điện cực ngực nào), dùng bề rộng QRS thai 40 ms và QRS mẹ 100 ms mà họ dẫn từ Taylor 2003 và Surawicz 2009 để biện minh cho sigma = 12 ms của nhãn Gaussian (khoảng ±2 sigma bằng 48 ms, xấp xỉ bề rộng fQRS 40 ms), và giữ nguyên Wilcoxon signed rank với p < 0,05 làm phương pháp thống kê chung, có thể trích dẫn bài này làm tiền lệ cho thí nghiệm trường tiếp nhận. Cần tránh ba lỗi của họ: phải nêu rõ dung sai ±50 ms ngay trong phần Methods và lặp lại trong caption bảng kết quả, không dùng công thức Acc = (TP+TN)/(...) cho bài toán phát hiện đỉnh mà giữ F1/Se/PPV, và không tinh chỉnh siêu tham số trên chính tập báo cáo kết quả (giao thức tách theo bản ghi của nhóm đã tránh được, phải nói rõ điều đó); đồng thời rút bài học rằng không được dùng FECGSYN làm bằng chứng chính cho hiệu năng phát hiện vì trên đó SR và MR không khác biệt (p = 0,58) trong khi trên dữ liệu thật khác biệt cực mạnh (p = 0,0000), chỉ nên dùng để stress-test và trình bày riêng kèm hạn chế này.\par\vspace{3pt}
\vspace{6pt}
\clearpage
\section{Đối chuẩn, bộ dữ liệu và thử thách}
Các công trình định nghĩa luật chơi của cả lĩnh vực: bộ dữ liệu nào, chấm điểm thế nào, dung sai bao nhiêu.

\subsection{Behar 2014, Luận án Tiến sĩ \normalfont\small (p13)}
\noindent{\footnotesize\itshape Behar J (2014). "Extraction of Clinical Information from the Non-Invasive Fetal Electrocardiogram." Luan an Tien si (DPhil), Balliol College, University of Oxford. Nguoi huong dan: Prof. Gari D. Clifford (dong huong dan: Dr Julien Oster). Michaelmas 2014, 249 trang. LUU Y VE TRICH DAN: ten file trong repo ghi "arxiv" va (o ban p14) ghi arXiv:1606.01093, NHUNG BEN TRONG FILE PDF KHONG HE CO dau ...}\par\vspace{2pt}
\noindent{\small\color{tdtblue}\textbf{Kết luận: so được một phần}}\par\vspace{4pt}
\noindent \textbf{Họ làm gì.} Đây là luận án Tiến sĩ 249 trang, tài liệu nền tảng của lĩnh vực fECG không xâm lấn và là nguồn gốc của chuẩn đánh giá mà nhóm đang dùng. Tác giả làm bốn việc lớn: đồng tổ chức PhysioNet/CinC Challenge 2013 và định nghĩa chuẩn đánh giá (Se/PPV/F1 với dung sai ±50 ms), xây bộ mô phỏng fecgsyn, xây bộ chỉ số chất lượng tín hiệu, và so sánh công bằng một loạt phương pháp tách fECG trên cùng dữ liệu và cùng thiết lập.\par\vspace{3pt}
\noindent \textbf{Phương pháp.} Chương 4 tạo bốn bộ dữ liệu mới (CRDB, CDB, AQTDB, RQTDB) cùng giao diện gán nhãn CrowdLabel và bộ mô phỏng fecgsyn (hơn 2.500 dòng mã, GPL, trên PhysioNet); Chương 5 tính 7 chỉ số chất lượng tín hiệu (kSQI, sSQI, pSQI, basSQI, bSQI, qSQI, pcaSQI) trên cửa sổ 10 giây rồi phân loại tốt/xấu bằng SVM nhân Gaussian. Chương 6 dò fQRS khi có kênh ngực tham chiếu: tiền lọc bằng hai bộ Butterworth zero-phase nối tiếp (thông cao bậc 3, thông thấp bậc 5) và chuẩn hoá bằng tanh, khử ECG mẹ bằng bốn họ phương pháp gồm trừ mẫu (TSc, TSm, TSlp, TSpca), LMS, RLS và Echo State Network, rồi dò đỉnh trên phần dư bằng bộ dò kiểu Pan-Tompkins với thời gian trơ 150 ms. Toàn bộ tham số kể cả hai tần số cắt fb (1–49 Hz, bước 3) và fh (30–120 Hz, bước 5) được quét lưới vét cạn; Chương 8 phục hồi khoảng QT thai nhi bằng Dual Extended Kalman Filter và Chương 9 đưa ra khuyến nghị về bộ dữ liệu và chuẩn đánh giá tương lai.\par\vspace{3pt}
\noindent \textbf{Mô hình.} Luận án không dùng mạng học sâu; các mô hình có tham số gồm ESN với bể chứa M = 90 nơ ron, độ thưa ψ = 20\%, bán kính phổ ρ = 0,4, hệ số rò rỉ α = 0,4 và input scaling γ = 1, LMS với N = 20 và μ = 0,1, RLS với N = 20 và λ = 0,999, cùng SVM nhân Gaussian cho chỉ số chất lượng tín hiệu dùng 7 đặc trưng với C = 25 và gamma = 1. Số tham số tổng, kích thước mô hình tính bằng MB và độ trễ đều bài không nêu — chỉ có nhận xét định tính rằng LMS ít tốn kém tính toán hơn nên phù hợp cho thiết bị giá rẻ — nhưng mã nguồn fecgsyn, mã tách nguồn CinC 2013, rs-toolbox và bộ dò jqrs đều được mở trên PhysioNet.\par\vspace{3pt}
\noindent \textbf{Dữ liệu.} CRDB gồm 262 phút từ 9 sản phụ với 37.452 nhãn fQRS tham chiếu ở fs = 1 kHz, chia thành CRDB1 làm tập huấn luyện (14 bản ghi công khai từ NIFECGDB, tuổi thai 22–40 tuần, 42 phút, 6.444 nhãn gán thủ công) và CRDB2 làm tập kiểm thử (11 bản ghi 5 phút từ 8 sản phụ thuộc một bộ dữ liệu thương mại riêng tư, 220 phút, 31.008 nhãn lấy từ ECG da đầu). CDB là bộ CinC 2013 với 447 bản ghi 1 phút, 4 kênh bụng, fs = 1 kHz, gộp từ 5 nguồn (ADFECGDB 25, NIFECGDB 14, AFECGDB 20 bản ghi nhân tạo, PNIFECGDB 340, PSFECGDB 48) và chia set-a 75 / set-b 100 / set-c 272; ngoài ra có AQTDB 220 phút mô phỏng, RQTDB 105 phút từ 22 sản phụ đủ tháng chuyển dạ do 3 bác sĩ tim mạch nhi gán nhãn, và ba bộ ECG người lớn cho chương chỉ số chất lượng tín hiệu (CinC 2011, MIT-BIH Arrhythmia, MIMIC II).\par\vspace{3pt}
\noindent \textbf{Giao thức đánh giá.} Dung sai khớp nhịp là ±50 ms, ghi rõ tại trang 59 và nhắc lại ở mục 9.2.2 — đây chính là nguồn gốc chuẩn mà nhóm đang dùng — với Se = TP/(TP+FN), PPV = TP/(TP+FP), F1 = 2TP/(2TP+FN+FP) theo hướng dẫn ANSI/AAMI, kèm chỉ số HRm là tỷ lệ phần trăm thời điểm nhịp tim thai ước lượng nằm trong ±5 bpm so với nhịp tham chiếu. Chương 6 huấn luyện và tối ưu trên CRDB1 rồi kiểm thử trên CRDB2 khác phần cứng, khác giao thức thu, khác giai đoạn thai kỳ và khác đối tượng, tức kiểm thử xuyên bộ dữ liệu thật sự, tương đương hoặc khắt khe hơn tách theo bản ghi của nhóm; Chương 7 tối ưu trên set-a nhưng cũng chỉ báo cáo F1 trên chính set-a; Chương 5 kiểm thử cuối cùng trên MIMIC II hoàn toàn độc lập; Chương 8 tách theo hình thái VCG giữa huấn luyện (1–4) và kiểm thử (5–9). Chương 6 xử lý từng kênh bụng độc lập nhưng bắt buộc phải có 1 kênh ngực mẹ tham chiếu còn Chương 7 dùng 4 kênh bụng; dải lọc là biến được quét tường minh, cho fb = 20 Hz và fh = 95 Hz ở Chương 6 nhưng fb = 10 Hz và fh = 99 Hz ở Chương 7, và không có tinh chỉnh trên tập đích ở bất kỳ chỗ nào.\par\vspace{3pt}
\noindent \textbf{Kết quả chính.} Chương 5: bSQI đơn lẻ tốt nhất với Ac = 0,969 (trên pcaSQI 0,950, sSQI 0,892, kSQI 0,879, qSQI 0,766, pSQI 0,752, basSQI 0,626), kết hợp 6 chỉ số cho Ac kiểm thử 0,985, nhưng trên MIMIC II hoàn toàn độc lập thì tụt còn Ac = 89,3\%, Se = 74,4\%, Sp = 93,5\%, tệ nhất ở vô tâm thu (Ac = 0,683) và nhịp chậm xoang (Ac = 0,708 với Se = 0). Chương 6: trên CRDB1 huấn luyện, ESNna đạt F1 = 97,2\% (Se 97,2, PPV 97,3), trên ESNa 97,0, RLS 95,9, TSpca 95,4, LMS 95,4, TSlp 92,0, TSc 90,8 và TSm 90,1; trên CRDB2 kiểm thử độc lập, ESNa dẫn đầu với F1 = 90,2\% (Se 91,4, PPV 88,9, HRm 78,7) trên TSpca 89,3, RLS 88,2, LMS 87,9 và ESNna 87,9, tức tất cả tụt khoảng 7 điểm khi đổi phần cứng và đối tượng. Loại bỏ đoạn có bSQI < 0,8 trên kênh ngực cải thiện F1 trên CRDB2 thêm +0,95 đến +1,11 điểm trong khi chỉ loại 3,6\% tổng độ dài tín hiệu; Chương 8 đạt sai số tuyệt đối và RMSE khoảng QT thai (7,6 ms; 12,7 ms) trên dữ liệu nhân tạo và (13 ms; 17,3 ms) trên dữ liệu thật với EKFD; chi phí tính toán bài không nêu ở bất kỳ chương nào.\par\vspace{3pt}
\noindent \textbf{Điểm yếu.} Chương 7 báo cáo F1 trong mẫu vì mọi tham số đều quét trên set-a rồi F1 cuối cùng cũng tính trên set-a, còn tập giấu chỉ được chấm bằng điểm Challenge kiểu RMS; nhãn chuẩn của CinC 2013 lại không đồng nhất (một phần từ ECG da đầu, một phần gán tay, một phần suy từ kênh bụng mẹ, một phần từ tham số mô phỏng) và chính tác giả thừa nhận một số lỗi nhãn vẫn còn tồn tại, bên cạnh việc loại 7 bản ghi khỏi set-a bằng quan sát mắt thường sau khi đã nhìn dữ liệu. Không có bất kỳ con số chi phí tính toán nào trong 249 trang; chương chỉ số chất lượng tín hiệu hoàn toàn dùng ECG người lớn và chính tác giả chỉ rõ bSQI thất bại trên tín hiệu bụng vì bộ dò QRS bắt nhầm cả MQRS lẫn FQRS nên ông chỉ dùng bSQI trên kênh ngực; CRDB2 và RQTDB đều là dữ liệu thương mại riêng tư nên phần quan trọng nhất của Chương 6 và Chương 8 không tái lập được. Dải lọc tối ưu không ổn định giữa hai chương (fb = 20 Hz so với 10 Hz), Chương 6 bắt buộc phải có kênh ngực mẹ nên không so được với các phương pháp chỉ dùng tín hiệu bụng, CinC 2013 chỉ có 4 kênh và 1 phút mỗi bản ghi mà không có ca bệnh lý nào, và nhãn CRDB1 gán tay nên chỉ gán được khi fECG đã nhìn thấy bằng mắt — chính tác giả cảnh báo cách này tạo thiên lệch.\par\vspace{3pt}
\noindent \textbf{\color{tdtblue}So với mô hình của nhóm.} Điểm chung quan trọng nhất là cùng dung sai ±50 ms và cùng công thức F1 = 2TP/(2TP+FN+FP), nên con số của Behar và của nhóm nằm trên cùng một thang đo — điều mà rất ít bài trong danh sách đạt được. Tuy vậy chỉ so sánh được một phần: Chương 6 dùng 1 kênh bụng nhưng cộng thêm 1 kênh ngực mẹ còn Chương 7 dùng 4 kênh bụng nên đều nhiều thông tin hơn nhóm; Behar không bao giờ báo cáo F1 riêng cho ADFECGDB vì 25 bản ghi ADFECGDB bị trộn vào 447 bản ghi CinC 2013 và không tách ra, nên con số 97,43 ± 4,37 tách theo bản ghi của nhóm không có đối thủ trực tiếp trong luận án; điểm đối chiếu khả dĩ nhất là ESNa đạt F1 = 90,2\% trên CRDB2 với sự trợ giúp của kênh ngực, so với 97,43\% của nhóm không cần kênh ngực. Cảnh báo rò rỉ dữ liệu cần kiểm chứng ngay: Bảng 4.2 ghi CinC 2013 chứa 25 bản ghi từ ADFECGDB, đúng bằng 5 sản phụ × 5 phút, nên nếu bất kỳ bản ghi nào trong số đó rơi vào set-a thì con số xuyên bộ dữ liệu 77,34 ± 28,54 của nhóm không phải xuyên bộ dữ liệu thật sự — điều này giải thích trực tiếp hiện tượng lưỡng cực nhóm đã quan sát (4 bản ghi F1 = 100 và 3 bản ghi 24–49\%).\par\vspace{3pt}
\noindent \textbf{\color{tdtblue}Nhóm rút ra gì.} Việc phải làm ngay là kiểm tra chồng lấn giữa ADFECGDB (r01, r04, r07, r08, r10) và CinC 2013 set-a bằng tương quan chéo tín hiệu hoặc chuỗi RR tham chiếu trước khi công bố con số 77,34, vì nếu không kiểm tra thì reviewer Q1 có thể tự phát hiện và bài sẽ bị từ chối vì rò rỉ dữ liệu. Nên trích dẫn Behar làm căn cứ cho dung sai ±50 ms (AAMI EC57 cho phép 150 ms nhưng đó là ngưỡng cho ECG người lớn, còn nhịp thai nhanh gấp đôi), báo cáo thêm HRm (±5 bpm) và đường cơ sở CONST-HR bên cạnh Se/PPV/F1, áp dụng hậu xử lý làm mượt chuỗi RR của Chương 7 (+1,0 điểm F1), dùng fecgsyn để stress test ca bệnh lý, đưa con số 113.481 tham số / 0,48 MB / 4,35 ms vào Abstract như một đóng góp, và chính thức hoá luận điểm mới cho fSQI dựa trên việc chính Behar chỉ ra bSQI không áp dụng được cho tín hiệu bụng.\par\vspace{3pt}
\vspace{6pt}
\subsection{Behar và cs. 2014, Physiol Meas \normalfont\small (p14)}
\noindent{\footnotesize\itshape Behar J (2014), Chuong 7 "Fetal QRS detection with no maternal chest reference", trang 114-137 trong luan an Tien si "Extraction of Clinical Information from the Non-Invasive Fetal Electrocardiogram", Balliol College, University of Oxford, Michaelmas 2014. Ban tap chi tuong ung (ghi trong danh muc cong bo cua chinh luan an, muc [9] trang 10): Behar J, Oster J, Clifford GD. "Combining and ...}\par\vspace{2pt}
\noindent{\small\color{tdtblue}\textbf{Kết luận: so được một phần}}\par\vspace{4pt}
\noindent \textbf{Họ làm gì.} Chương này so sánh công bằng 15 phương pháp tách fECG trên cùng một bộ dữ liệu, cùng bộ tiền xử lý, cùng bộ dò đỉnh và cùng bộ thống kê, với cấu hình chỉ có 4 kênh bụng, không có kênh ngực mẹ và không có điện cực da đầu. Vì không phương pháp nào thắng tuyệt đối, tác giả đề xuất FUSE: chạy song song một nhóm phương pháp rồi chọn đầu ra có chuỗi nhịp tim thai mượt nhất — ý tưởng đơn giản này đánh bại mọi phương pháp đơn lẻ và giành điểm cao nhất ở 2/4 hạng mục của PhysioNet/CinC Challenge 2013 trên 53 đội quốc tế.\par\vspace{3pt}
\noindent \textbf{Phương pháp.} Pipeline gồm 8 bước: tiền lọc bằng hai bộ Butterworth zero-phase nối tiếp (thông thấp bậc 5, thông cao bậc 3) với notch 50/60 Hz áp dụng có điều kiện, chuẩn hoá theo biên độ đo trong khoảng 1–5 giây rồi lấy tanh để khống chế ngoại lai; dò QRS mẹ trên từng kênh bụng đã lọc và trên kênh sau biến đổi ICA (thời gian trơ 250 ms cho MQRS và 150 ms cho FQRS), sau đó hiệu chỉnh lại vị trí đỉnh R mẹ riêng cho từng kênh trong cửa sổ ±30 ms. Tách nguồn chia làm bốn lớp: trừ mẫu trong miền quan sát (TS, TSc, TSm, TSlp, TSpca, TSEKF), BSS trực tiếp trên tín hiệu bụng (ICA, PCA), trừ mẫu trước rồi BSS trên phần dư, và BSS trước rồi trừ mẫu trong miền nguồn. Chuỗi fQRS được chọn theo hai tiêu chí SMI (số lần biến thiên nhịp tức thời vượt 29 bpm, chọn kênh có SMI nhỏ nhất) và BCM (loại kênh nếu tỷ lệ nhịp khớp chuỗi MQRS trong ±50 ms đạt ≥ 0,4), rồi làm mượt chuỗi RR bằng trung vị 5 nhịp gần nhất; tham số tối ưu bằng quét lưới và quét ngẫu nhiên cho nbC = 20, nbPC = 2, GR = 100, GQ = 5, fb = 10 Hz, fh = 99 Hz.\par\vspace{3pt}
\noindent \textbf{Mô hình.} Không có mạng nơ ron và không có học sâu — toàn bộ là xử lý tín hiệu cổ điển kết hợp tối ưu siêu tham số, nên số tham số huấn luyện được: không áp dụng. Hệ thống chỉ có 6 con số toàn cục (nbC = 20, nbPC = 2, GR = 100, GQ = 5, fb = 10 Hz, fh = 99 Hz) cộng các ngưỡng SMI 29 bpm, BCM 0,4, SMI > 26, nhịp hằng số 143 bpm và cửa sổ hiệu chỉnh ±30 ms; kích thước mô hình và độ trễ đều bài không nêu, nhưng mã nguồn được mở trên PhysioNet nên khả năng tái lập cao.\par\vspace{3pt}
\noindent \textbf{Dữ liệu.} Bộ CinC 2013 gồm 447 bản ghi, mỗi bản ghi 1 phút, 4 kênh bụng, fs = 1 kHz, gộp từ 5 nguồn (ADFECGDB 25 bản ghi công khai từ 5 sản phụ chuyển dạ 38–41 tuần × 5 phút, NIFECGDB 14, AFECGDB 20 bản ghi nhân tạo sinh bởi fecgsyn, PNIFECGDB 340 riêng tư, PSFECGDB 48 riêng tư và không công khai), chia thành set-a 75 bản ghi công khai cả tín hiệu lẫn nhãn, set-b 100 chỉ công khai tín hiệu và set-c 272 giấu hoàn toàn. Nhãn chuẩn không đồng nhất: lấy từ điện cực da đầu với PSFECGDB và ADFECGDB, gán tay hoặc suy ra từ các kênh bụng mẹ phụ mà thí sinh không được thấy với các bộ khác, và biết chính xác từ tham số mô phỏng với AFECGDB; tác giả thừa nhận một số lỗi nhãn vẫn còn tồn tại dù đã kiểm lại bằng mắt, và có 53 đội quốc tế tham gia với 93 bài nộp mã nguồn mở.\par\vspace{3pt}
\noindent \textbf{Giao thức đánh giá.} Dung sai khớp nhịp là ±50 ms, và hai hệ thống chấm điểm chạy song song không trùng nhau: F1 kèm Se và PPV được tính trên set-a, tức trong mẫu vì toàn bộ tham số và ngưỡng đều quét trên chính set-a, còn điểm Challenge E1, E2, E4, E5 mới chấm trên tập giấu nhưng chỉ là điểm kiểu RMS chứ không phải F1. Trước khi tính F1, bài cắt 2 giây đầu và 2 giây cuối mỗi đoạn 1 phút và loại 7 bản ghi (a33, a38, a47, a52, a54, a71, a74) vì nhãn tham chiếu không chính xác, còn 68/75 bản ghi; thống kê tổng là trung bình thô của F1, Se, PPV tính trên từng bản ghi. Bài dùng 4 kênh bụng đồng thời cộng cơ chế chạy nhiều phương pháp tách rồi chọn kênh phần dư tốt nhất nên không phải đơn kênh; không tinh chỉnh trên tập đích, dải lọc fb = 10 Hz và fh = 99 Hz tìm bằng quét lưới trên set-a, và chi phí tính toán bài không nêu.\par\vspace{3pt}
\noindent \textbf{Kết quả chính.} Trên set-a với fb = 10 Hz, thứ hạng F1 là FUSE-SMOOTH 96,0\% (Se 95,9, PPV 96,0, HRE 19,1, RRE 6,3), FUSE 95,0, TSpca-ICA 93,0, TSc-ICA 92,6, ICA-TSpca-ICA 92,4, TS-ICA 92,0, TSm-ICA 91,2, ICA-TSpca 89,2, TSlp-ICA 88,4, TSpca 86,1, TSm 82,1, TSEKF 81,9, TSlp 81,8, TS 81,6, TSc 81,5, ICA 63,7 và PCA 51,6; trên tập giấu, FUSE-CHALL đạt E1 = 179,4 bpm², E2 = 20,8 ms, E4 = 29,6 bpm², E5 = 4,7 ms, tốt nhất về E1 và E2 trong 53 đội. Phát hiện quan trọng nhất là đường cơ sở CONST-HR chỉ xuất nhịp hằng số 143 bpm mà không nhìn tín hiệu: F1 chỉ 23,0\% nhưng HRE = 172,2 và RRE = 8,9, tốt hơn mọi phương pháp đơn lẻ, chứng tỏ hệ chấm điểm kiểu RMS của Challenge có khuyết tật nghiêm trọng vì chỉ xét vị trí tương đối chứ không xét vị trí tuyệt đối của đỉnh. Hạ tần số cắt về 10 Hz (so với 2 Hz) làm tăng F1 ở 13/17 phương pháp với mức tăng lớn nhất +3,5 điểm (TSEKF), thêm một bước ICA sau trừ mẫu nâng F1 thêm 2–10 điểm (TS 81,6 lên TS-ICA 92,0), và làm mượt RR nâng FUSE từ 95,0 lên 96,0 đồng thời giảm HRE từ 132,9 xuống 19,1.\par\vspace{3pt}
\noindent \textbf{Điểm yếu.} Con số F1 = 96,0\% là trong mẫu vì mọi tham số quét trên set-a rồi F1 cũng báo cáo trên set-a, tập giấu chỉ có điểm Challenge kiểu RMS; thêm vào đó bài loại 7/75 bản ghi vì nhãn tham chiếu không chính xác xác định bằng quan sát mắt thường — quyết định chủ quan và làm sau khi đã nhìn dữ liệu — rồi còn cắt 2 giây đầu và cuối mỗi đoạn. Bốn ngưỡng SMI 29 bpm, BCM 0,4, SMI > 26 và nhịp hằng số 143 bpm đều đo bằng thực nghiệm trên set-a mà không có cơ sở lý thuyết, còn FUSE-CHALL được chính tác giả gọi là làm lệch về phía hệ thống chấm điểm của Challenge nên điểm E1/E2 hạng nhất có được một phần nhờ mẹo này. Nhãn chuẩn không đồng nhất và 20/447 bản ghi hoàn toàn nhân tạo, dữ liệu chỉ có 4 kênh và 1 phút mỗi bản ghi mà không có ca bệnh lý nào, không có bất kỳ con số chi phí tính toán nào dù FUSE phải chạy song song 5 thuật toán tách nguồn trên 4 kênh, dải lọc tối ưu không nhất quán với Chương 6 (fb = 20 Hz so với 10 Hz), và chính tác giả thừa nhận cách gộp thông tin của FUSE rất thô vì chỉ chọn một đầu ra chứ không gộp các nhãn lại.\par\vspace{3pt}
\noindent \textbf{\color{tdtblue}So với mô hình của nhóm.} Bài cùng thang đo với nhóm ở một điểm quan trọng là dung sai ±50 ms và công thức F1 giống hệt, nhưng chỉ so sánh được một phần vì Behar dùng 4 kênh bụng cộng 5 phương pháp tách chạy song song và cơ chế chọn kênh trong khi nhóm dùng 1 kênh và 1 mô hình duy nhất — đó là lợi thế cấu trúc chứ không phải lợi thế thuật toán — và vì F1 = 96,0\% của ông là trong mẫu trên chính tập quét tham số còn 97,43 ± 4,37 tách theo bản ghi của nhóm là ngoài mẫu nên khắt khe hơn về phương pháp luận. Đối chiếu khả dĩ nhất là con số xuyên bộ dữ liệu 77,34 ± 28,54 của nhóm trên CinC 2013 so với FUSE-SMOOTH 96,0 trên set-a: Behar cao hơn rõ ràng, nhưng phải giải thích đủ năm lý do khiến khoảng cách này không phải khoảng cách năng lực mô hình — 4 kênh so với 1 kênh, Behar tối ưu toàn bộ tham số trên chính set-a trong khi nhóm chạy zero-shot, Behar loại 7 bản ghi nhãn xấu, Behar cắt 2 giây đầu và cuối, và Behar có hậu xử lý làm mượt RR (+1,0 điểm F1) mà nhóm hiện chưa có. Cảnh báo rò rỉ cần kiểm chứng: Table 4.2 ghi CinC 2013 chứa 25 bản ghi từ ADFECGDB, khớp chính xác 5 sản phụ × 5 phút, nên nếu bản ghi ADFECGDB rơi vào set-a thì con số 77,34 bị nhiễm và điều đó giải thích chính xác hiện tượng lưỡng cực nhóm quan sát; ngoài ra F1 = 81,6\% của TS thuần tuý (4 kênh) chính là đường cơ sở cổ điển công bằng cho khối trừ mẫu mà nhóm đang dùng.\par\vspace{3pt}
\noindent \textbf{\color{tdtblue}Nhóm rút ra gì.} Áp dụng ngay với chi phí thấp: hậu xử lý làm mượt RR của Behar (lấy trung vị RR của 5 nhịp trước, chỉ làm mượt khi trung vị nằm trong 110–170 bpm, xoá nhịp cách nhịp trước < 70\% trung vị RR và chèn nhịp khi > 175\% trung vị RR), cơ chế BCM để loại chuỗi nhịp mẹ bị nhận nhầm thành nhịp thai, notch có điều kiện cho dữ liệu trộn nguồn 50/60 Hz, và hiệu chỉnh đỉnh R mẹ theo từng kênh trong ±30 ms trước khi trừ mẫu trung vị. Khi viết bài, dùng Bảng 7.2 làm bảng so sánh gốc, đưa CONST-HR vào bảng để bảo vệ việc báo cáo F1 thay vì RMSE, đối chiếu phát hiện dải lọc (Behar tăng tối đa +3,5 điểm với phương pháp cổ điển còn nhóm đo được +11,00 điểm với học sâu), giữ lợi thế đơn kênh làm luận điểm chính, và tuyệt đối không viết "vượt Behar 96,0\%" mà không nói rõ đó là con số trong mẫu trên set-a với 4 kênh sau khi loại 7 bản ghi.\par\vspace{3pt}
\vspace{6pt}
\subsection{Andreotti và cs. 2016, Physiol Meas \normalfont\small (p15)}
\noindent{\footnotesize\itshape Andreotti, Fernando; Behar, Joachim; Zaunseder, Sebastian; Oster, Julien; Clifford, Gari D. "An open-source framework for stress-testing non-invasive foetal ECG extraction algorithms." Physiological Measurement, tap 37, so 5, nam 2016, trang 627-648. DOI: 10.1088/0967-3334/37/5/627. (Ban PDF trong thu muc la ban luu tru CC BY-NC-ND 4.0 cua TU Dresden.)}\par\vspace{2pt}
\noindent{\small\color{tdtblue}\textbf{Kết luận: so được một phần}}\par\vspace{4pt}
\noindent \textbf{Họ làm gì.} Bài này không đề xuất mô hình mới mà xây dựng một bộ khung kiểm thử chuẩn (stress-testing) cho lĩnh vực đo ECG thai không xâm lấn, xuất phát từ nhận định rằng phần lớn công trình trước đó không thể đánh giá một cách có ý nghĩa do thiếu dữ liệu công khai, thước đo không thống nhất và không ai công bố mã nguồn. Nhóm tác giả tạo bộ dữ liệu mô phỏng lớn FECGSYNDB bằng công cụ fecgsyn, trong đó có chủ đích đưa vào các tình huống phi dừng (thai cử động, co thắt tử cung, nhịp ngoại tâm thu, song thai), đem 8 thuật toán tách fECG kinh điển ra so trên cùng một mặt bằng, rồi công bố toàn bộ dữ liệu, mã nguồn và quy trình đánh giá theo giấy phép GNU GPL trên PhysioNet.\par\vspace{3pt}
\noindent \textbf{Phương pháp.} Dữ liệu được sinh bằng fecgsyn theo mô hình lưỡng cực (mở rộng từ McSharry 2003 và Sameni 2007), với bảy kịch bản (Baseline; Case 0 chỉ thêm nhiễu; Case 1 thai cử động; Case 2 tăng/giảm tốc nhịp mẹ-thai; Case 3 co thắt tử cung; Case 4 nhịp ngoại tâm thu; Case 5 song thai), mỗi kịch bản sinh 10 lần × 5 mức nhiễu (0, 3, 6, 9, 12 dB) × 5 lần lặp, tổng 1.750 tín hiệu; tiền xử lý bằng bộ lọc Butterworth pha không, dải 3-100 Hz cho hai thí nghiệm đầu và dải 0,5-100 Hz cho thí nghiệm hình thái theo khuyến cáo của Hiệp hội Tim mạch Hoa Kỳ dành cho ECG người lớn. Ba nhóm thuật toán được so sánh: tách nguồn mù (BSSpca, BSSica với JADE và FAST-ICA), trừ mẫu (TSc, TSpca, TSekf) và lọc thích nghi (AMlms, AMrls, AMesn — cả ba đều cần một kênh tham chiếu ECG mẹ). Đỉnh fQRS được dò bằng bản Pan-Tompkins đã hiệu chỉnh theo Behar 2014a, còn phân tích hình thái dùng nhãn chuẩn R thai để dựng mẫu theo phương pháp Oster 2015 rồi phân đoạn bằng ecgpuwave nhằm đo FQT và tỉ số T/QRS.\par\vspace{3pt}
\noindent \textbf{Mô hình.} Không có mạng nơ-ron nào được huấn luyện (ngoại trừ AMesn là mạng echo state, nhưng bài không nêu số nơ-ron hay số tham số của nó); đây là 8 thuật toán xử lý tín hiệu cổ điển và bài không nêu số tham số, kích thước mô hình hay độ trễ của bất kỳ phương pháp nào. Các siêu tham số được nêu gồm: mẫu ECG mẹ dựng từ 20 chu kỳ và cập nhật mỗi nhịp, TSekf khởi tạo bằng 30 chu kỳ, FAST-ICA tối đa 1000 × Nch vòng lặp với hàm tương phản tanh, PCA giảm chiều giữ 99,9\% phương sai, cửa sổ khởi tạo tối đa 60 giây và không dùng bộ lọc làm mượt ngoại tuyến.\par\vspace{3pt}
\noindent \textbf{Dữ liệu.} Toàn bộ phần đánh giá dùng dữ liệu mô phỏng, không có một bản ghi lâm sàng thật nào: bộ FECGSYNDB gồm 1.750 tín hiệu, mỗi tín hiệu dài 5 phút và được chiếu lên 34 kênh (32 kênh bụng cùng 2 kênh tham chiếu ECG mẹ) ở tần số lấy mẫu 250 Hz, tổng cộng 145,8 giờ dữ liệu đa kênh và khoảng 1,1 triệu đỉnh R thai. Tham số mô phỏng: SNR thai/mẹ \textasciitilde{} N(-9, 2) dB; SNR mẹ/nhiễu thuộc \{0, 3, 6, 9, 12\} dB; nhịp thai \textasciitilde{} N(135, 25) bpm; nhịp mẹ \textasciitilde{} N(80, 20) bpm; tần số hô hấp thai \textasciitilde{} N(0,90; 0,05) Hz; tần số hô hấp mẹ \textasciitilde{} N(0,25; 0,05) Hz.\par\vspace{3pt}
\noindent \textbf{Giao thức đánh giá.} Không có chia huấn luyện/kiểm thử vì cả 8 phương pháp đều là xử lý tín hiệu không giám sát chạy trực tiếp trên từng bản ghi, nên không tồn tại khái niệm tách theo chủ thể hay tách theo bản ghi và cũng không có rò rỉ dữ liệu kiểu chia ngẫu nhiên. Dung sai khớp nhịp là 50 ms — bài nêu rõ chuẩn người lớn ANSI/AAMI/ISO EC57 là 150 ms nhưng cố ý thu hẹp xuống 50 ms để tính đến nhịp tim thai cao hơn (dẫn Andreotti 2014, Behar 2014b, Zaunseder 2013) — với các chỉ số Se, PPV, F1 = 2TP/(2TP+FN+FP) và MAE chỉ tính trên các TP; đơn vị phân tích là từng epoch 1 phút của từng tín hiệu. Cảnh báo lớn nhất: với mọi phương pháp, nhóm tác giả chọn kênh (hoặc thành phần độc lập) có F1 cao nhất cho từng epoch 1 phút rồi mới thống kê, tức chọn kênh kiểu oracle dựa trên nhãn chuẩn, nên mọi con số đều bị thổi phồng và bước chọn thành phần vốn khó nhất của tách nguồn mù bị bỏ qua hoàn toàn; ngoài ra nhóm BSS dùng 8 kênh bụng đồng thời còn nhóm AM tuy gọi là đơn kênh nhưng vẫn cần thêm kênh tham chiếu ECG mẹ (kênh 33).\par\vspace{3pt}
\noindent \textbf{Kết quả chính.} Thí nghiệm 1 (số kênh cho BSS): với 8 kênh, F1 trung bình đạt 97,46\% (BSSica-JADE), 97,40\% (BSSpca) và 97,22\% (FAST-ICA); khi Nch = 16, bước giảm chiều bằng PCA nâng F1 trung bình của ICA từ 95,3\% lên 97,3\%, và F1 gần như bão hòa khi Nch ≥ 8. Thí nghiệm 2 (F1 trung vị tổng thể): BSSica 99,9\% (IQR 6,4), BSSpca 98,7\%, AMesn 97,9\% (IQR 12,5), AMrls 97,3\%, AMlms 97,1\%, TSpca 96,0\% (IQR 18,5), TSc 95,0\%, TSekf 94,5\%; kịch bản khó nhất là co thắt tử cung (Case 3, F1 tụt còn 77,4-94,5\%) và nhịp ngoại tâm thu (Case 4, 83,4-88,4\%), trong khi MAE trung vị tốt nhất là 3,8 ms của TSekf, tiếp đến AMesn 4,0 ms và BSSica 4,1 ms. Thí nghiệm 3 (hình thái): số nhịp dùng được giảm đơn điệu từ 69,7\% ở mức nền xuống 11,6\% (Case 0 tại 0 dB), hệ số xác định r² giữa FQT và tỉ số T/QRS đo được với tham chiếu cho thấy BSSica rất kém (0,0067 đến 0,1886) trong khi TSc rất tốt cho tỉ số T/QRS (0,8361 đến 0,9155) và AMesn ở mức trung gian (0,3569 đến 0,8124), dẫn đến kết luận tuyệt đối không nên phân tích hình thái trong miền nguồn sau khi chạy BSS; kiểm định Friedman cho thấy khác biệt giữa các kịch bản là cực kỳ có ý nghĩa (p < 0,001) ở SNR thấp \{0, 3\} dB nhưng không có ý nghĩa (p > 0,05) ở 12 dB.\par\vspace{3pt}
\noindent \textbf{Điểm yếu.} Toàn bộ đánh giá thực hiện trên dữ liệu mô phỏng và chính tác giả thừa nhận F1 cao một phần là do sự đơn giản của mô hình lưỡng cực trong fecgsyn, với khoảng FQT mô phỏng thấp hơn khoảng sinh lý thật (200-340 ms), tức dữ liệu tổng hợp chưa đúng về mặt hình thái. Việc chọn kênh kiểu oracle khiến mọi con số chỉ là cận trên lý thuyết chứ không phải hiệu năng của một hệ thống tự động; nhóm AM tuy được tuyên là đơn kênh nhưng trên lâm sàng vẫn phải dán thêm điện cực ngực, còn nhóm BSS cần 8 kênh đồng thời nên hoàn toàn không phải bài toán đơn kênh. Bài không báo cáo bất kỳ chi phí tính toán nào (số tham số, độ trễ, kích thước mô hình, thời gian chạy), MAE được tính với nhãn chuẩn R thai chưa căn chỉnh theo từng kênh nên có sai số hệ thống nhỏ, ecgpuwave vốn thiết kế cho ECG người lớn thường không trả về được điểm mốc trên ECG thai, và các dải lọc 3-100 Hz cùng 0,5-100 Hz đều lấy theo chuẩn ECG người lớn chứ chưa tối ưu cho fQRS — chính tác giả viết rằng cần một thử nghiệm lâm sàng để xác nhận liệu chuẩn này có áp dụng được cho phân tích fECG hay không.\par\vspace{3pt}
\noindent \textbf{\color{tdtblue}So với mô hình của nhóm.} Không so sánh trực tiếp được về con số, và đây là điều phải nói rõ trong bài của nhóm: họ đo trên dữ liệu mô phỏng còn nhóm đo trên ADFECGDB lâm sàng thật, F1 99,9\% của BSSica là cận trên oracle sau khi chọn kênh tốt nhất theo nhãn chuẩn trong 8 kênh, và BSS cần 8 kênh còn AM cần thêm kênh tham chiếu mẹ trong khi FetalQRS-TCN chỉ cần 1 kênh bụng. Tuy vậy vẫn có những trục so sánh được và có lợi cho nhóm: cùng dung sai 50 ms và cùng công thức F1 nên nhóm được phép trích dẫn họ để biện minh cho lựa chọn 50 ms; về jitter, MAE tốt nhất của họ là 3,8 ms (TSekf, trên dữ liệu mô phỏng sạch) trong khi nhóm đo được 3,05 ms trên dữ liệu lâm sàng thật theo giao thức DPSS. Về dải lọc, họ dùng 3-100 Hz cho phát hiện đỉnh, trong khi thực nghiệm của nhóm cho thấy dải 3-90 Hz chỉ đạt 84,17 điểm F1 còn dải 10-60 Hz đạt 91,06 điểm — chênh khoảng 7 điểm, tức bộ khung chuẩn này đang dùng một dải lọc dưới tối ưu cho bài toán dò đỉnh; ngoài ra họ không báo cáo chi phí tính toán trong khi nhóm có 113.481 tham số, 0,48 MB và 4,35 ms cho mỗi cửa sổ 4 giây trên CPU.\par\vspace{3pt}
\noindent \textbf{\color{tdtblue}Nhóm rút ra gì.} Nên tải FECGSYNDB (miễn phí trên PhysioNet, tần số 250 Hz trùng khớp với tần số đích của pipeline nên nạp thẳng không cần lấy mẫu lại) để chạy FetalQRS-TCN trên 7 kịch bản phi dừng × 5 mức SNR, luôn báo cáo F1 kèm MAE/jitter thành cặp đúng như tác giả khuyến nghị, báo cáo tỉ lệ đoạn bị loại theo tiền lệ 69,7\% xuống 11,6\% của họ, và trích dẫn bài này vừa để biện minh cho dung sai 50 ms vừa làm căn cứ cho ý tưởng fSQI (họ viết rõ về tầm quan trọng của việc loại bỏ các đoạn ECG chất lượng kém bằng các thước đo chất lượng tín hiệu). Phải tránh việc báo cáo kết quả kênh tốt nhất như kết quả chính mà không ghi chú: nhóm cần đặt 97,43 (tách theo bản ghi trên đủ 4 kênh) làm kết quả chính và 99,18 (kênh lâm sàng tốt nhất) làm kết quả phụ có chú thích rõ, đồng thời không đặt con số của nhóm trên dữ liệu thật cạnh con số 99,9\% của họ mà thiếu cột ghi chú dữ liệu mô phỏng và chọn kênh oracle.\par\vspace{3pt}
\vspace{6pt}
\subsection{Clifford và cs. 2014, Physiol Meas \normalfont\small (p20)}
\noindent{\footnotesize\itshape Gari D. Clifford, Ikaro Silva, Joachim Behar, George B. Moody. "Non-invasive fetal ECG analysis" (tieu de trong metadata ghi la "Noninvasive Fetal ECG analysis"). Physiological Measurement, Volume 35, Issue 8, trang 1521-1536, 2014. DOI: 10.1088/0967-3334/35/8/1521. Day la bai xa luan (editorial) mo dau so dac biet ve cuoc thi PhysioNet/Computing in Cardiology Challenge 2013. File doc: ...}\par\vspace{2pt}
\noindent{\small\color{tdtblue}\textbf{Kết luận: so được một phần}}\par\vspace{4pt}
\noindent \textbf{Họ làm gì.} Đây là bài xã luận mở đầu số đặc biệt về cuộc thi PhysioNet/Computing in Cardiology Challenge 2013, không phải bài đề xuất thuật toán: bài trình bày vì sao điện tim thai không xâm lấn lại khó (chồng lấp cả về thời gian lẫn tần số giữa ECG mẹ và ECG thai), tình trạng chỉ có 3 bộ dữ liệu công khai và đều quá nhỏ trước năm 2013, thiết kế và cách chấm điểm cuộc thi, kỹ thuật của các đội dẫn đầu và tóm tắt 13 bài trong số đặc biệt. Với nhóm, giá trị lớn nhất là đây chính là tài liệu gốc định nghĩa bộ dữ liệu CinC 2013 mà nhóm đang dùng làm tập xuyên bộ dữ liệu, và nó chứa cảnh báo trực tiếp về rò rỉ dữ liệu giữa ADFECGDB và CinC 2013.\par\vspace{3pt}
\noindent \textbf{Phương pháp.} Ban tổ chức gộp dữ liệu từ 5 nguồn thành 447 bản ghi, chuẩn hóa tất cả về tần số lấy mẫu 1 kHz, thời lượng 1 phút và 4 kênh ECG bụng không xâm lấn, rồi chia thành Set A gồm 75 bản ghi công khai cả tín hiệu lẫn nhãn, Set B gồm 100 bản ghi chỉ công khai tín hiệu, và Set C gồm 272 bản ghi giấu cả tín hiệu lẫn nhãn, chạy trên máy chủ riêng của PhysioNet. Nhãn tham chiếu giai đoạn đầu lấy từ nhãn gốc của từng bộ dữ liệu, với các bộ không xâm lấn thì suy ra thủ công hoặc từ các kênh ECG mẹ bổ sung mà thí sinh không được thấy; nhãn giai đoạn cuối dùng phương pháp crowd-sourcing Bayes (Zhu 2013, 2014) hợp nhất nhãn gốc với nhãn của tất cả bài dự thi mã nguồn mở ở giai đoạn 1, có kiểm tra thủ công một tập con và tác giả thừa nhận "vẫn còn một số lỗi nhỏ trong nhãn". Việc chấm điểm dùng phần mềm WFDB phiên bản 10.5.19 qua ba pha nộp bài (25/04-01/06/2013 với 3 lần nộp; 01/06-25/08/2013 với 5 lần; 25/08-05/09/2013 với 1 lần), và bài tổng hợp pipeline 5 bước chung mà hầu hết các đội đều dùng: tiền xử lý, ước lượng thành phần mẹ, loại bỏ thành phần mẹ, ước lượng FHR/RR và hậu xử lý.\par\vspace{3pt}
\noindent \textbf{Mô hình.} Bài không đề xuất mô hình nào vì đây là xã luận, và không nêu số tham số hay kích thước mô hình của bất kỳ thuật toán nào. Bài chỉ mô tả vắn tắt năm thuật toán của các đội dẫn đầu: Andreotti (ICA cộng bộ lọc Kalman mở rộng cộng thích nghi mẫu mẹ và simulated annealing, hậu xử lý loại khoảng RR ngắn hơn 300 ms và biến thiên nhịp lớn hơn 70 ms); Lipponen và Tarvainen (Butterworth thông cao bậc 6 tần số cắt 2 Hz, xóa thành phần Fourier 50 Hz, phân rã riêng cho sóng P/QRS/T, zero 6 trị riêng lớn nhất cho QRS và 4 trị riêng lớn nhất cho P và T, trung bình trượt 100 ms để ước lượng chất lượng, mẫu tạo từ 20 phức bộ QRS lớn nhất, lọc trung bình trượt 30 ms); Varanini (lọc trung vị 60 ms, Butterworth thông thấp bậc 1 tần số cắt 3,17 Hz để tách xu thế, lọc trung vị 260 ms, notch tự động 3 hài bậc đầu, FastICA, SVD lấy 3 trị kỳ dị lớn nhất để mô hình nhịp mẹ); Podziemski và Gieratowski (độ dốc RS cộng ngưỡng biên độ, RR trong phạm vi 75 ms quanh giá trị trung bình); và Behar (lọc thông cao cộng notch, hợp nhất PCA, trừ mẫu và ICA, Pan-Tompkins sửa đổi, chọn kênh tốt nhất theo số lần biến thiên nhịp tức thời vượt 30 bpm, trả về chuỗi hằng số 143 bpm nếu không phát hiện được).\par\vspace{3pt}
\noindent \textbf{Dữ liệu.} Tổng 447 bản ghi từ 5 nguồn: ADFECGDB (Matonia 2006) 25 bản ghi, Simulated FECG tức FECGSYN (Behar 2014b) 20 bản ghi, NIFECGDB 14 bản ghi, "Non-Invasive FECG" 340 bản ghi và Scalp FECG Database 48 bản ghi (không công khai, chỉ dùng để chấm điểm mã nguồn mở trên Set C); tất cả chuẩn hóa về 1 kHz, 1 phút, 4 kênh bụng, chia Set A 75 / Set B 100 / Set C 272, trong đó 1 bản ghi ở Set A và 2 bản ghi ở Set B không chấm được do thiếu nhãn tham chiếu. Bài cũng mô tả ba bộ dữ liệu công khai trước đó: Daisy (8 kênh gồm 4 kênh bụng và 3 kênh ngực, 10 giây, 250 Hz), NIFECGDB (55 bản ghi đa kênh từ một đối tượng duy nhất, tuần 21-40, 1 kHz, không có nhãn tham chiếu) và ADFECGDB (5 phút, 4 kênh bụng, 5 sản phụ chuyển dạ tuần 38-41, có ECG da đầu làm tham chiếu).\par\vspace{3pt}
\noindent \textbf{Giao thức đánh giá.} Đây là giao thức thi đấu chứ không phải kiểm chứng chéo: thí sinh huấn luyện trên Set A có nhãn công khai, ban tổ chức chấm trên Set B và Set C mà thí sinh không có nhãn, tức tập giữ lại ở cấp bản ghi, không phải tách theo chủ thể và cũng không phải chia ngẫu nhiên trong cùng một bệnh nhân. Về dung sai khớp nhịp thì bài không nêu bất kỳ con số ms nào, và điểm cực kỳ quan trọng là cách chấm của CinC 2013 không dùng F1 và không dùng dung sai 50 ms. Chấm điểm chính thức gồm 5 sự kiện: E1 và E4 là chuỗi FHR chấm bằng lỗi phân loại từng nhịp với đơn vị (nhịp/phút)\textasciicircum{}2, đo tại 12 thời điểm cách nhau 5 giây; E2 và E5 là chuỗi RR chấm bằng căn bậc hai lỗi trung bình với đơn vị mili giây; E3 là khoảng QT; E1 và E2 chạy trên Set C còn E4 và E5 chạy trên Set B, bản ghi không được thí sinh gán nhãn bị phạt điểm rất cao và điểm cuối cùng là trung bình của tất cả bản ghi trong tập, còn F1 chỉ là hàm mà riêng Behar tự chọn để tối ưu thuật toán chứ không phải tiêu chí chấm của cuộc thi.\par\vspace{3pt}
\noindent \textbf{Kết quả chính.} 53 đội quốc tế tham gia, sinh ra 208 bộ nhãn và 93 bài dự thi mã nguồn mở, với điểm tốt nhất từng sự kiện là E1 = 179,439 (nhịp/phút)\textasciicircum{}2, E2 = 20,793 ms, E4 = 18,083 (nhịp/phút)\textasciicircum{}2 và E5 = 4,337 ms; Andreotti thắng E2 và E5, còn Behar (không chính thức vì đồng tổ chức nhưng bị mù với tập kiểm tra) đạt hạng 1 ở E1 và E4, hạng 2 ở E3, hạng 3 và hạng 2 ở E2 và E5. Phát hiện quan trọng nhất cho nhóm là bài ghi rõ nhận xét của Behar rằng chọn tần số cắt thông cao cao, ví dụ 10 Hz, cải thiện kết quả phát hiện, "có lẽ vì loại bỏ được các thành phần sóng P và T lớn của mẹ", tức bằng chứng độc lập từ y văn cho phát hiện về dải 10-60 Hz của nhóm. Bài cũng tự chỉ trích Rodrigues (2014) vì đã dùng ADFECGDB để huấn luyện thuật toán, việc này "có thể dẫn đến thiên lệch trong kết quả vì bộ dữ liệu này nằm trong set-a, set-b (và có thể một vài bản ghi trong set-c)".\par\vspace{3pt}
\noindent \textbf{Điểm yếu.} Nhãn tham chiếu giai đoạn cuối được tạo bằng cách hợp nhất nhãn gốc với nhãn của chính các thuật toán dự thi nên một phần nhãn là sản phẩm của thuật toán chứ không phải nhãn chuẩn độc lập, và tác giả thừa nhận vẫn còn lỗi; bài cũng không nêu dung sai khớp nhịp nên không thể tái lập. Bản ghi chỉ dài 1 phút nên không đánh giá được độ ổn định dài hạn, các tiêu chí chấm điểm theo (nhịp/phút)\textasciicircum{}2 và ms không đối chiếu được với F1/Se/PPV của đa số bài báo hiện đại khiến kết quả cuộc thi gần như không so sánh được với y văn sau này, còn sự kiện E3 kém tin cậy vì không có nhãn cho đoạn ST và khoảng QT ở cả điều kiện bình thường lẫn bất thường. Nghiêm trọng nhất là chồng lấn giữa các bộ dữ liệu, khi ADFECGDB (25 bản ghi) và FECGSYN (20 bản ghi) nằm bên trong 447 bản ghi của cuộc thi, nên bất kỳ ai huấn luyện trên ADFECGDB rồi kiểm tra trên CinC 2013 đều có nguy cơ rò rỉ dữ liệu và chính bài báo đã chỉ đích danh một đội mắc lỗi này; bài cũng thừa nhận cần bộ dữ liệu lớn hơn, nhiều bệnh nhân hơn, bản ghi dài hơn, nhiều chuyển đạo hơn kể cả ECG mẹ, và có bệnh lý như loạn nhịp, chậm phát triển trong tử cung, nhiễm toan thai.\par\vspace{3pt}
\noindent \textbf{\color{tdtblue}So với mô hình của nhóm.} So sánh được một phần và theo hướng ngoài ý muốn: bộ dữ liệu thì so sánh được vì đây chính là tập nhóm dùng để đo xuyên bộ dữ liệu (77,34 ± 28,54 F1), nhưng con số thì không so sánh được vì tiêu chí chính thức của CinC 2013 là lỗi FHR tính bằng (nhịp/phút)\textasciicircum{}2 và lỗi RR tính bằng ms chứ không phải F1 với dung sai 50 ms. Điểm phải sửa ngay trong bài của nhóm là giả định "50 ms là chuẩn CinC 2013" không được bài xã luận gốc này chứng minh, vì bài không nêu bất kỳ dung sai ms nào, nên nếu muốn viện dẫn mức 50 ms thì phải trích dẫn nguồn khác (nhiều khả năng là các bài benchmark của Behar, tức p13/p14), nếu không phản biện Q1 có thể bắt lỗi. Nguy cơ rò rỉ dữ liệu cũng nghiêm trọng: ADFECGDB (25 bản ghi) nằm trong 447 bản ghi của CinC 2013 và phân bố vào Set A, Set B và có thể cả Set C, trong khi nhóm huấn luyện trên ADFECGDB rồi kiểm tra trên CinC 2013, đúng lỗi mà Clifford đã chỉ đích danh ở Rodrigues (2014); kết quả lưỡng cực của nhóm (4 bản ghi F1 = 100 và 3 bản ghi 24-49\%) khiến rất đáng ngờ rằng 4 bản ghi đạt 100 chính là các bản ghi có nguồn gốc ADFECGDB, đây là giả thuyết phải kiểm chứng bằng cách đối chiếu danh sách bản ghi, và nếu đúng thì con số 77,34 hiện tại là thiên lệch lên, không dùng được làm bằng chứng tổng quát hóa.\par\vspace{3pt}
\noindent \textbf{\color{tdtblue}Nhóm rút ra gì.} Việc đầu tiên là kiểm tra chồng lấn giữa ADFECGDB và CinC 2013 bằng cách đối chiếu từng bản ghi trong tập đang test với 25 bản ghi ADFECGDB có trong cuộc thi, loại bỏ mọi bản ghi trùng rồi tính lại con số 77,34 ± 28,54, và nếu không loại được thì phải ghi rõ cảnh báo trong bài kèm trích dẫn chính câu của Clifford về trường hợp Rodrigues; đồng thời sửa cách trích dẫn dung sai 50 ms sang đúng nguồn (benchmark của Behar) hoặc ghi rõ đây là lựa chọn của nhóm có tham khảo AAMI EC57 (150 ms) và các benchmark fECG. Tiếp theo, dùng nhận xét về ngưỡng 10 Hz của Behar trong phần Discussion làm cơ sở vật lý độc lập cho mức tăng +11,00 điểm của dải 10-60 Hz, báo cáo thêm lỗi RR (ms) và lỗi FHR (nhịp/phút)\textasciicircum{}2 trên Set A nếu muốn tuyên bố tốt hơn các đội CinC, thận trọng với nhãn CinC 2013 (được tạo bằng crowd-sourcing từ chính các thuật toán dự thi nên có thể sai đúng ở các bản ghi khó) và kiểm tra trực quan 3 bản ghi đạt 24-49\% trước khi kết luận mô hình yếu, còn các con số nền của cuộc thi (447 bản ghi, 1 kHz, 1 phút, 4 kênh, Set A 75 / Set B 100 / Set C 272, 53 đội, điểm tốt nhất E5 = 4,337 ms) thì dùng cho phần Related Work.\par\vspace{3pt}
\vspace{6pt}
\subsection{Matonia và cs. 2020, Scientific Data \normalfont\small (p21)}
\noindent{\footnotesize\itshape Adam Matonia, Janusz Jezewski, Tomasz Kupka, Michal Jezewski, Krzysztof Horoba, Janusz Wrobel, Robert Czabanski, Radana Kahankova. "Fetal electrocardiograms, direct and abdominal with reference heartbeat annotations". Scientific Data (Nature), volume 7, so bai 200, nam 2020. DOI: 10.1038/s41597-020-0538-z. Nhan bai 22/11/2019, chap nhan 20/05/2020. Du lieu tren figshare, DOI: ...}\par\vspace{2pt}
\noindent{\small\color{tdtblue}\textbf{Kết luận: so được một phần}}\par\vspace{4pt}
\noindent \textbf{Họ làm gì.} Đây là bài mô tả dữ liệu (Data Descriptor) trong đó nhóm Ba Lan công bố công khai hai bộ tín hiệu điện tim thai nhi do chính họ thu bằng hệ thống KOMPOREL: B1 gồm 10 bản ghi thai kỳ mỗi bản 20 phút và B2 gồm 12 bản ghi chuyển dạ mỗi bản 5 phút, kèm bộ chỉ số đo quan hệ biên độ giữa ECG mẹ, ECG thai và nhiễu (WM, WF, WMF, WEM, WEF) cùng phân tích các mẫu biến thiên nhịp tim thai lâm sàng (acceleration, deceleration, LTV, STV). Điểm quan trọng nhất đối với nhóm là bộ ADFECGDB trên PhysioNet mà nhóm đang dùng chính là 5 bản ghi (r01, r04, r07, r08, r10) lấy ra từ bộ B2 này, tức chỉ chiếm 9,6\% tổng thời lượng tín hiệu bụng mà bài công bố.\par\vspace{3pt}
\noindent \textbf{Phương pháp.} Hệ thống KOMPOREL có nhiễu nội tại dưới 1 uV, hệ số khử nhiễu chung CMRR = 115 dB và dải thông FECG từ 0,05 đến 150 Hz (tần số cắt thông cao được nâng lên 1 Hz khi có nhiễu tần số thấp mạnh), dùng 4 điện cực đo A1-A4 đặt đều quanh đường rốn, điện cực tham chiếu chung V0 trên xương mu, điện cực nối đất N ở chân trái, điện cực Ag/AgCl 3M Red Dot 2271 với keo dẫn điện 3M Red Dot Trace Prep 2236 và điện cực xoắn vô trùng CETRO AB 15133C cho tín hiệu trực tiếp từ đầu thai nhi; số hóa 16 bit ở 500 Hz cho tín hiệu bụng và 1 kHz cho FECG trực tiếp. Dữ liệu công bố đã qua bộ lọc lược khử nhiễu răng cưa H(z) = (1/36)(1 - z\textasciicircum{}-50)/(1 - z\textasciicircum{}-1) x (1 - z\textasciicircum{}-10)... với tần số cắt đầu tiên khoảng 5 Hz và dải chặn 45-55 Hz để diệt nhiễu 50 Hz, sau đó khử ECG mẹ bằng trừ mẫu đầy đủ chu kỳ PQRST với mẫu được cập nhật liên tục, có trừ thêm đạo hàm bậc nhất của mẫu QRS để giảm ảnh hưởng của tần số lấy mẫu 500 Hz lên việc định vị điểm fiducial, và lọc không gian sơ bộ để tạo tín hiệu MECG phụ dùng định vị sóng R mẹ. Nhãn tham chiếu của B1 được tạo bằng cách chạy K hàm phát hiện, tính phiếm hàm Delta\_k đo biến thiên độ dài chuỗi RR rồi chọn hàm có Delta\_k nhỏ nhất, sau đó chuyên gia lâm sàng sửa tay và gắn cờ tin cậy bằng 0 cho các điểm không xác minh được; nhãn của B2 lấy từ FECG trực tiếp bằng lọc phối hợp chuẩn hóa cộng quy tắc quyết định cực tiểu hóa hàm chi phí dự đoán độ dài chu kỳ tim rồi chuyên gia xác minh, kênh tốt nhất được chọn bằng chỉ số chất lượng dựa trên hàm tự tương quan trong cửa sổ điều chỉnh theo nhịp tim thai, và FHR được phân tích bằng hệ thống MONAKO ở 4 Hz.\par\vspace{3pt}
\noindent \textbf{Mô hình.} Bài không đề xuất mô hình học máy nào và không nêu số tham số của mô hình nào, vì đây là bài mô tả dữ liệu. Các thuật toán được nhắc đến (trừ mẫu PQRST, lọc phối hợp chuẩn hóa, Pan-Tompkins sửa đổi, PCA, ICA, AICF, PFTAB) đều là thuật toán xử lý tín hiệu cổ điển không có tham số học được, và đều chỉ được mô tả bằng cách trích dẫn công trình trước của chính nhóm chứ không trình bày đầy đủ trong bài này.\par\vspace{3pt}
\noindent \textbf{Dữ liệu.} Bộ B1\_Pregnancy\_dataset gồm 10 bản ghi, mỗi bản 20 phút (tổng 200 phút), 500 Hz, 16 bit, 4 kênh bụng, tuổi thai 32 đến 42 tuần, với tổng 18.936 phức bộ QRS mẹ và 28.405 phức bộ QRS thai (130 phức bộ thai, tức 0,46\%, không xác minh được), chu kỳ tim thai trung bình 421,1 ms tương ứng 142,5 bpm, số phức bộ không trùng lặp là 8.106 mẹ và 17.575 thai (lấy bề rộng mQRS 100 ms và fQRS 40 ms), chỉ số biên độ WM trung bình 14,3 dB, WF trung bình 3,4 dB (có bản ghi âm: -1,6 và -0,7 dB), WMF gần 11 dB (8,5 đến 17,7 dB), WEM 0,09, WEF 0,52, FHR nền 142,5 ± 9,0 bpm, LTV\_Dawes 38,2 ± 11,2 ms và STV\_Dawes 6,4 ± 3,2 ms. Bộ B2\_Labour\_dataset gồm 12 bản ghi, mỗi bản 5 phút, 4 kênh bụng ở 500 Hz cộng 1 kênh FECG trực tiếp từ đầu thai nhi ở 1 kHz, tuổi thai 38 đến 42 tuần trong giai đoạn chuyển dạ tích cực, với 7.903 phức bộ QRS thai tham chiếu và 5.177 QRS mẹ (không trùng lặp: 5.142 thai và 2.415 mẹ), vị trí sóng R thai được xác định với độ chính xác 1 ms, WM trung bình 10,4 dB, WF trung bình 3,6 dB (thấp nhất -2,7, cao nhất 7,0), WMF trung bình 6,8 dB (1,8 đến 10,9), WEM 0,20, WEF 0,33, FHR nền 133,0 ± 7,2 bpm, LTV\_Dawes 46,1 ± 13,0 ms, STV\_Dawes 7,0 ± 2,7 ms và tỷ lệ mất tín hiệu từng bản ghi từ 0,0\% đến 27,5\% (bản ghi 03).\par\vspace{3pt}
\noindent \textbf{Giao thức đánh giá.} Bài không định nghĩa giao thức train/test nào, không tách theo chủ thể, không tách theo bản ghi và không chia ngẫu nhiên, vì đây là bài công bố dữ liệu không có mô hình để đánh giá. Về dung sai khớp nhịp, phần Technical Validation nêu rõ mức ±40 ms: một phức bộ fQRS được coi là phát hiện đúng nếu cách điểm tham chiếu sóng R không quá ±40 ms, giá trị này tương ứng bề rộng trung bình của phức bộ QRS thai, tức chặt hơn mức ±50 ms mà nhóm đang dùng. Các chỉ số họ định nghĩa gồm PI = (N - FP - FN)/N x 100, Acc = TP/(TP + FP + FN) x 100 (không có TN, thực chất tương đương Jaccard chứ không phải accuracy thông thường), Se = TP/(TP+FN), PPV = TP/(TP+FP), F1 = 2TP/(2TP+FN+FP), cùng sai số tuyệt đối trung bình của khoảng RR ký hiệu |Delta RR| tính bằng cả ms lẫn bpm; chính tác giả cảnh báo nhãn tham chiếu của B1 chỉ "chính xác ở mức độ giới hạn" vì được suy ra từ chính tín hiệu bụng nhiều nhiễu, khác hẳn B2 có nhãn chuẩn là FECG trực tiếp, và khuyến nghị bỏ qua các điểm có cờ tin cậy bằng 0 khi đánh giá hiệu năng phát hiện.\par\vspace{3pt}
\noindent \textbf{Kết quả chính.} Phần Technical Validation trích lại kết quả của chính nhóm trên bộ dữ liệu này, tất cả với dung sai ±40 ms: toàn bộ chuỗi xử lý tín hiệu bụng đạt PI tốt nhất 81,4\% nhưng khác biệt rất lớn giữa các chuyển đạo bụng (từ 70,4\% đến 99,8\%), còn đánh giá trên 10 tín hiệu của B2 với tham chiếu là FECG trực tiếp, tức chính công trình sinh ra ADFECGDB, đạt Acc = 98\%, Se = 98,97\%, PPV = 98,99\%, PI = 97,97\% và sai số tuyệt đối trung bình khoảng RR bằng 2,14 ms (tương đương 0,63 bpm). Trên toàn bộ bộ B2, phương pháp tách nguồn mù đạt F1 = 98,56\% với PCA (Se = 98,27\%, PPV = 98,86\%, Acc = 97,17\%) và F1 = 98,55\% với ICA (Se = 98,23\%, PPV = 98,87\%, Acc = 97,13\%). Bài nêu rõ 5 bản ghi r01, r04, r07, r08, r10 của ADFECGDB trên PhysioNet chính là lấy từ bộ B2 và đã được lấy mẫu lại lên 1 kHz để tiện so sánh với FECG trực tiếp, chỉ chiếm 9,6\% thời lượng tín hiệu bụng công bố; các bộ dữ liệu tham chiếu khác được đối chiếu gồm Daisy (10 giây, 250 Hz, không nhãn), NIFECGDB (55 bản ghi 1 kHz, chỉ có nhãn sóng R mẹ, tuần 21-40) và FECGSYNDB (1.750 mô phỏng 5 phút, 32 kênh bụng và 2 kênh ngực, tổng 145,8 giờ, 7 sự kiện sinh lý, 10 cấu hình lưỡng cực, 5 mức nhiễu).\par\vspace{3pt}
\noindent \textbf{Điểm yếu.} Số đối tượng rất nhỏ, chỉ 10 thai phụ ở B1 và 12 sản phụ ở B2, nên phương sai giữa các lần gấp khi tách theo bản ghi sẽ rất lớn (đây chính là lý do sai số chuẩn của nhóm là ±4,37 điểm), còn bản ghi B2 chỉ dài 5 phút và bản ghi B2\_Labour\_03 mất tới 27,5\% tín hiệu nên không đánh giá được độ ổn định dài hạn. Nhãn tham chiếu của B1 mang tính vòng tròn vì được suy ra từ chính tín hiệu bụng bằng thuật toán tự động cộng chuyên gia sửa tay, không có nhãn chuẩn độc lập và chính tác giả thừa nhận nhãn này chỉ "chính xác ở mức độ giới hạn", thêm 130 phức bộ (0,46\%) không xác minh được phải lọc bỏ bằng cờ tin cậy, nên chỉ bộ B2 với FECG trực tiếp mới là nhãn chuẩn thật. Bài không công bố mã nguồn của bất kỳ thuật toán nào và mọi con số hiệu năng đều là trích dẫn lại từ công trình cũ chứ không được tính lại theo một giao thức thống nhất, nên các mức 98,56\% hay 81,4\% không tái lập được; ngoài ra dữ liệu đã bị tiền lọc bằng bộ lọc lược với tần số cắt khoảng 5 Hz trước khi công bố nên người dùng không nhận được tín hiệu thô, và bài không báo cáo bất kỳ chi phí tính toán nào vì không có mô hình.\par\vspace{3pt}
\noindent \textbf{\color{tdtblue}So với mô hình của nhóm.} So sánh được một phần và đây là bài quan trọng nhất trong ba bài đối với nhóm, vì bộ B2\_Labour chứa chính 5 bản ghi ADFECGDB mà nhóm đang dùng (r01, r04, r07, r08, r10), cùng điều kiện chuyển dạ và cùng nhãn chuẩn là FECG trực tiếp từ đầu thai nhi: con số gần nhất để đối chiếu là PCA đạt F1 = 98,56\% và ICA đạt F1 = 98,55\% trên toàn bộ 12 bản ghi B2 với dung sai ±40 ms, trong khi mô hình của nhóm đạt macro F1 = 97,43 ± 4,37 (tách theo bản ghi, 20 đánh giá, toàn bộ 4 kênh bụng) và 99,18 ở kênh lâm sàng tốt nhất trên 5 bản ghi ADFECGDB với dung sai ±50 ms. Chỉ so sánh được một phần và phải nói rõ trong bài vì bốn lý do: dung sai của họ chặt hơn (40 ms so với 50 ms) nên con số của họ đo ở điều kiện khó hơn và không được nói nhóm hơn họ khi chưa đo lại ở 40 ms; PCA và ICA là kỹ thuật tách nguồn mù bắt buộc gộp cả 4 kênh bụng trong khi mô hình của nhóm là đơn kênh thực sự, đây là lợi thế rõ ràng cần nhấn mạnh; tập bản ghi khác nhau (12 bản ghi của họ so với 5 bản ghi của nhóm, tức có 7 bản ghi nhóm chưa từng thấy); và họ báo cáo một lần trên toàn bộ bản ghi chứ không tách theo bản ghi như 20 lần đánh giá của nhóm nên độ biến thiên khác nhau. Chỉ số |Delta RR| = 2,14 ms của họ và jitter 3,05 ms của nhóm là hai đại lượng khác nhau, vì |Delta RR| là sai số của khoảng RR còn jitter là độ lệch vị trí đỉnh so với nhãn, mà về toán học sai số RR thường lớn hơn khoảng căn 2 lần sai số vị trí nếu độc lập, nên không được đặt cạnh nhau như thể họ tốt hơn; về chi phí tính toán thì không so sánh được vì họ không báo cáo, và 113.481 tham số, checkpoint 0,48 MB cùng độ trễ 4,35 ms của nhóm là thông tin mà cả ba bài đều thiếu, tức khoảng trống nhóm nên khai thác.\par\vspace{3pt}
\noindent \textbf{\color{tdtblue}Nhóm rút ra gì.} Ưu tiên cao nhất là tải bộ B2 đầy đủ 12 bản ghi từ figshare (DOI 10.6084/m9.figshare.c.4740794, giấy phép CC BY 4.0) để lấy 7 bản ghi chuyển dạ chưa từng nằm trong ADFECGDB làm tập kiểm tra ngoài lý tưởng (cùng phân bố lâm sàng nên loại bỏ được yếu tố khác bộ dữ liệu, nhưng mô hình chưa từng thấy), và tải bộ B1 gồm 10 bản ghi thai kỳ 20 phút tuần 32-42 để mở rộng phạm vi ứng dụng từ chuyển dạ sang trước sinh, với điều kiện phải ghi rõ nhãn của B1 được suy ra từ chính tín hiệu bụng và phải lọc bỏ 130 điểm có cờ tin cậy bằng 0. Đồng thời cần bổ sung bảng độ nhạy dung sai ở 20 / 30 / 40 / 50 ms để đối chiếu công bằng với PCA 98,56\% và ICA 98,55\% đo ở ±40 ms, lập bảng so sánh trực tiếp có cột số kênh và cột dung sai, trích dẫn WF trung bình 3,4 dB ở B1 và 3,6 dB ở B2 cùng các giá trị âm (-1,6 và -0,7 dB ở B1; -2,7 dB ở B2) làm bằng chứng định lượng cho hướng chỉ số chất lượng tín hiệu fSQI được quyền từ chối trả lời, dùng bề rộng fQRS 40 ms và mQRS 100 ms để biện minh cho sigma = 12 ms của nhãn Gaussian, ghi rõ trong phần Methods rằng dữ liệu đã được tiền lọc ở khoảng 5 Hz và chặn 45-55 Hz (đồng thời kiểm tra lại xem tín hiệu ADFECGDB trên PhysioNet còn thành phần dưới 5 Hz không), và tránh nhầm công thức Acc = TP/(TP+FP+FN) của họ, thực chất là Jaccard, với accuracy thông thường.\par\vspace{3pt}
\vspace{6pt}
\clearpage
\section{Lý thuyết phân tích dữ liệu topo}
Nền tảng lý thuyết mà bản đề cương đầu tiên của nhóm dựa vào. Nhóm đọc kỹ nhóm này để hiểu vì sao giả thuyết ban đầu sai.

\subsection{Perea và Harer 2015, Found Comput Math \normalfont\small (p22)}
\noindent{\footnotesize\itshape Jose A. Perea, John Harer. "Sliding Windows and Persistence: An Application of Topological Methods to Signal Analysis". File PDF trong thu muc la BAN TIEN AN arXiv:1307.6188v2 [math.AT], nop 25/11/2013, ghi ngay 22/11/2013, 34 trang; PDF KHONG ghi ten tap chi va KHONG co DOI. Ban xuat ban chinh thuc (trich duoc tu danh muc tham khao cua p24 trong cung bo tai lieu): Foundations of Computational ...}\par\vspace{2pt}
\noindent{\small\color{warnred}\textbf{Kết luận: không so trực tiếp được}}\par\vspace{4pt}
\noindent \textbf{Họ làm gì.} Bài xây dựng khung lý thuyết đo tính tuần hoàn của chuỗi thời gian bằng cách kết hợp nhúng cửa sổ trượt (time-delay embedding) với đồng điều bền vững chiều 1: trải một đoạn tín hiệu thành đám mây điểm nhiều chiều, nếu tín hiệu tuần hoàn thì đám mây có dạng vòng tròn và "độ tròn" của vòng chính là maximum persistence của sơ đồ bền vững chiều 1. Đóng góp chính là chứng minh toán học rằng độ tròn này đạt cực đại khi kích thước cửa sổ Mτ xấp xỉ độ dài chu kỳ của tín hiệu, cùng các định lý về sự phụ thuộc vào số chiều nhúng, kích thước cửa sổ và trường hệ số; phương pháp được đặt tên SW1PerS, đây không phải bài về ECG và không có dữ liệu y sinh thật.\par\vspace{3pt}
\noindent \textbf{Phương pháp.} Pipeline thực nghiệm (mục 7): nội suy spline bậc 3 tín hiệu lấy mẫu thành hàm liên tục; khử nhiễu bắt buộc bằng trung bình trượt cửa sổ 7 điểm rồi 1 vòng mean-shift trên đám mây đã chuẩn hóa với ngưỡng ε = cos(π/16) đo bằng độ tương tự cosine; nhúng cửa sổ trượt SW\_\{M,τ\}f(t) = (f(t), f(t+τ), ..., f(t+Mτ)); căn giữa và chuẩn hóa từng điểm; lọc Vietoris-Rips và tính đồng điều bền vững chiều 1 trên trường hữu hạn F\_p; điểm tuần hoàn Score(S) = mp(dgm)/√3 với mp là maximum persistence. Phần lý thuyết: Prop 5.1 khẳng định các vector u\_n, v\_n độc lập tuyến tính khi và chỉ khi M ≥ 2N (chặt hơn điều kiện đủ của Takens/Nyquist), giả định chuẩn của bài là M = 2N và Mτ < 2π; Prop 5.2/5.5 cho τ = 2π/(L(M+1)); Thm 4.5 chặn d\_B(dgm(X),dgm(Y)) ≤ 2√(4k−2)·||R\_N f\textasciicircum{}(k)||\_2·√(M+1)/(N+1)\textasciicircum{}(k−1/2); Thm 5.6 khẳng định đám mây chuẩn hóa nằm trên đường cong trong N-torus; Thm 6.8 cho mp(dgm) ≥ √3·max(r̂\_n) − δ·κ\_N và Corollary 6.10 cho mp(dgm\_vô cực(f,w)) ≥ 2√3·max|f̂(n)|; Thm 6.6/6.7 chứng minh hội tụ khi N → vô cực và khi lưới lấy mẫu mịn dần. Cài đặt đồng điều bền vững dùng thuật toán Union-Find dựa trên Mischaikow-Nanda nhưng chi tiết được hẹn sang tài liệu [23] chưa có sẵn; mục 6.3 chứng minh sơ đồ bền vững phụ thuộc trường hệ số (ví dụ g1(t) = 0,6cos(t) + 0,8cos(2t) và g2 = 0,8cos(t) + 0,6cos(2t) với M = 4, τ = 2π/5, T gồm 151 điểm cho kết quả F\_2 khác F\_3 vì xuất hiện dải Möbius) nên khuyên chọn số nguyên tố p > N.\par\vspace{3pt}
\noindent \textbf{Mô hình.} Bài không nêu số tham số vì đây không phải mô hình học máy và không có tham số huấn luyện; chỉ có siêu tham số hình học gồm số chiều nhúng M+1, bậc cắt cụt Fourier N, trường hệ số F\_p và số chu kỳ giả định L, với thực nghiệm dùng N = 10 (suy ra M = 2N = 20, cửa sổ 21 chiều), p = 11, thí nghiệm 1 thử L = 2, 3, 4 rồi lấy điểm tốt nhất còn thí nghiệm 2 cố định L = 2. Bài không nêu thời gian chạy, không nêu bộ nhớ và không nêu kích thước mô hình.\par\vspace{3pt}
\noindent \textbf{Dữ liệu.} Toàn bộ là dữ liệu tổng hợp, không có dữ liệu sinh lý thật: thí nghiệm 1 (Shape Independence) dùng 10 dạng sóng khác nhau gồm cosin 2 chu kỳ thuần, cosin cộng nhiễu Gauss ở 3 mức (phương sai 25\%, 50\%, 75\% biên độ), răng cưa có nhiễu 25\%, cos(φ(t)) với φ(t) = e\textasciicircum{}(at+b), cosin tắt dần có nhiễu 3 chu kỳ, tín hiệu nhọn 3 chu kỳ, sóng vuông có nhiễu 2 chu kỳ và 1 hàm ngẫu nhiên từ 5 hệ số Fourier phức, mỗi hàm lấy 50 mẫu cách đều. Thí nghiệm 2 (Classification Rates) dùng 6 dạng gồm 4 đường tuần hoàn (cosin, cosin có xu thế, cosin tắt dần, cosin đỉnh nhọn) và 2 không tuần hoàn (hằng số, tuyến tính), mỗi dạng sinh 100 profile bằng cách đổi pha nên tổng cộng 600 profile, mỗi profile lấy 50 mẫu cách đều trên [0, 2π] và thêm nhiễu Gauss độ lệch chuẩn 0\%, 25\%, 50\% biên độ; bài không nêu tần số lấy mẫu vật lý vì tín hiệu là vô thứ nguyên, trục thời gian tính theo radian.\par\vspace{3pt}
\noindent \textbf{Giao thức đánh giá.} Không có chia train/test, không có tách theo chủ thể, không có tách theo bản ghi, vì SW1PerS không học tham số nào mà chỉ tính một điểm số rồi vẽ đường ROC trên toàn bộ 600 profile; do đó không tồn tại rò rỉ dữ liệu theo nghĩa học máy nhưng cũng không thể gọi là độ chính xác trên tập kiểm thử độc lập. Chỉ số đánh giá là AUC của ROC (TPR so với FPR) khi quét ngưỡng trên điểm số, đối thủ so sánh là JTK\_CYCLE, Lomb-Scargle và Total Persistent Homology (Cohen-Steiner và cộng sự) đều đặt ở tham số mặc định hoặc gợi ý; bài không nêu dung sai khớp nhịp (ms) và không có Se/PPV/F1 vì không làm bài toán dò đỉnh. Cần lưu ý về phương pháp luận: ở thí nghiệm 1, bài thử L = 2, 3, 4 rồi báo cáo điểm tốt nhất trên chính tập đánh giá, tức là một dạng tối ưu hóa siêu tham số trên tập kiểm thử.\par\vspace{3pt}
\noindent \textbf{Kết quả chính.} Thí nghiệm 1 xếp hạng 10 tín hiệu theo điểm tuần hoàn: SW1PerS cho điểm từ 0,9488 (cosin thuần) giảm dần qua 0,8810 (ba tín hiệu), 0,8132, 0,7455, 0,6777, 0,6770 xuống 0,4066 (hai tín hiệu cuối), trong khi JTK\_CYCLE trả p-value từ 1,77E-30 đến 1,00E+00 và Lomb-Scargle từ 1,46E-07 đến 9,98E-01; tác giả kết luận chỉ SW1PerS là không thiên vị một hình dạng sóng cụ thể. Thí nghiệm 2 đo AUC theo từng dạng sóng ở 3 mức nhiễu: ở SD = 0\% cả bốn phương pháp (JTK, LS, Total PH, SW1PerS) đều đạt AUC = 1 cho cả 4 dạng; ở SD = 25\% lần lượt là JTK 1/0,94/0,99/1, LS 1/0,98/1/1, Total PH 0,97/0,65/0,92/0,98 và SW1PerS 1/0,90/0,99/0,99; ở SD = 50\% là JTK 0,96/0,56/0,72/0,93, LS 0,99/0,66/0,86/0,97, Total PH 0,78/0,59/0,77/0,81 và SW1PerS 0,98/0,71/0,87/0,93, nghĩa là ở mức nhiễu cao SW1PerS vượt Lomb-Scargle ở dạng tắt dần (0,71 so với 0,66) và đỉnh nhọn (0,87 so với 0,86) nhưng thua ở cosin (0,98 so với 0,99) và có xu thế (0,93 so với 0,97). Hình 7 cho thấy nếu bỏ khử nhiễu (bỏ trung bình trượt và mean-shift) thì kết quả suy giảm đáng kể, và tác giả gọi khử nhiễu là yếu tố then chốt của pipeline SW1PerS.\par\vspace{3pt}
\noindent \textbf{Điểm yếu.} Toàn bộ là dữ liệu tổng hợp với mỗi tín hiệu chỉ 50 mẫu và không có bất kỳ kiểm chứng nào trên tín hiệu sinh lý thật, trong khi chi tiết cài đặt đồng điều bền vững được hẹn sang một bài chưa công bố ([23]) nên không tái lập được nguyên vẹn. Điểm số phụ thuộc việc đoán đúng số chu kỳ L trong cửa sổ và thí nghiệm 1 thử ba giá trị L rồi lấy kết quả tốt nhất ngay trên tập đánh giá; ngoài ra bài thua Lomb-Scargle ở dạng cosin và có xu thế, tức với tín hiệu gần hình sin thì phương pháp cổ điển vẫn tốt hơn. Ràng buộc M ≥ 2N cùng cửa sổ khoảng một chu kỳ làm kích thước phức Rips bùng nổ khi muốn giữ nhiều hài bậc cao mà bài tự thừa nhận là điểm nghẽn tính toán nhưng không báo cáo thời gian chạy hay bộ nhớ; giả thuyết trung tâm (maximum persistence cực đại đúng tại cửa sổ bằng chu kỳ) chưa được chứng minh mà chỉ ghi trong phần hướng phát triển, và hàm hằng số bị coi là không tuần hoàn theo quy ước riêng của bài (Remark 7.1), ảnh hưởng trực tiếp đến AUC.\par\vspace{3pt}
\noindent \textbf{\color{tdtblue}So với mô hình của nhóm.} Không so sánh trực tiếp được vì ba lý do: bài toán khác hẳn (họ chấm điểm tuần hoàn toàn cục cho cả đoạn tín hiệu, còn nhóm làm định vị đỉnh từng mẫu theo kiểu sequence-to-sequence), dữ liệu khác hẳn (tổng hợp 50 mẫu vô thứ nguyên, không phải ADFECGDB hay CinC 2013), và không có Se/PPV/F1 theo nhịp, không có dung sai ±50 ms, không có giao thức tách theo bản ghi. Tuy nhiên đây là cơ sở lý thuyết mạnh nhất để bảo vệ phát hiện của nhóm: định lý của bài nói maximum persistence chỉ lớn khi kích thước cửa sổ Mτ xấp xỉ chu kỳ, nên cửa sổ 300 ms ngắn hơn chu kỳ tim thai khoảng 430 ms không thể tạo ra vòng H1 — đó là lý do toán học chứ không phải lỗi cài đặt. Ngoài ra kết quả Hình 7 (bỏ khử nhiễu thì AUC sụt mạnh) trùng khớp với kết luận của nhóm rằng tiền xử lý quan trọng hơn kiến trúc.\par\vspace{3pt}
\noindent \textbf{\color{tdtblue}Nhóm rút ra gì.} Nên áp dụng: nếu dùng TDA làm chỉ số chất lượng tín hiệu thai (fSQI) thì phải đặt kích thước cửa sổ trượt xấp xỉ một chu kỳ RR thai (khoảng 400-500 ms, tương ứng khoảng 100-125 mẫu ở 250 Hz) chứ không được ngắn hơn — căn cứ là Corollary 6.10 và Prop 5.5 — đồng thời chọn số chiều nhúng M ≥ 2N với N là bậc hài cao nhất muốn giữ, chọn số nguyên tố p > N khi tính đồng điều bền vững để tránh hiện tượng xoắn kiểu Möbius, bắt buộc khử nhiễu trước khi nhúng (với fECG chính là dải thông 10-60 Hz của nhóm), và dùng điểm chuẩn hóa mp(dgm)/√3 thuộc [0, 1] làm fSQI vì nó có ý nghĩa hình học rõ ràng thay vì p-value khó diễn giải. Nên tránh: không lấy điểm tốt nhất qua nhiều giá trị L rồi báo cáo như kết quả cuối mà phải cố định siêu tham số trên tập huấn luyện, và không kỳ vọng SW1PerS thắng các phương pháp phổ cổ điển trên tín hiệu gần hình sin.\par\vspace{3pt}
\vspace{6pt}
\subsection{Adams và cs. 2017, JMLR \normalfont\small (p23)}
\noindent{\footnotesize\itshape Henry Adams, Tegan Emerson, Michael Kirby, Rachel Neville, Chris Peterson, Patrick Shipman, Sofya Chepushtanova, Eric Hanson, Francis Motta, Lori Ziegelmeier. "Persistence Images: A Stable Vector Representation of Persistent Homology". Journal of Machine Learning Research 18 (2017) 1-35. Nop 7/2016, xuat ban 2/2017. Bien tap vien: Michael Mahoney. Giay phep CC-BY 4.0, trang bai: ...}\par\vspace{2pt}
\noindent{\small\color{warnred}\textbf{Kết luận: không so trực tiếp được}}\par\vspace{4pt}
\noindent \textbf{Họ làm gì.} Sơ đồ bền vững là một đa tập điểm trong mặt phẳng với số điểm thay đổi tùy mẫu, không cộng trừ được nên không dùng trực tiếp làm đầu vào cho SVM, cây quyết định, mạng nơ-ron hay giảm chiều; bài biến sơ đồ này thành một vector hữu hạn chiều gọi là Persistence Image (PI) bằng cách rải một Gauss có trọng số lên từng điểm rồi tích phân trên lưới ô vuông, mỗi pixel là một số thực. Bài chứng minh phép biến đổi này ổn định theo khoảng cách 1-Wasserstein và thực nghiệm cho thấy PI phân loại tốt hơn persistence landscape đồng thời nhanh hơn sơ đồ gốc hàng nghìn lần; đây chính là biểu diễn mà nhóm đã thử (cnn2d\_PersistenceImage) và bị chốt bảng.\par\vspace{3pt}
\noindent \textbf{Phương pháp.} Bước 1 đổi tọa độ T(x,y) = (x, y−x) từ (sinh, chết) sang (sinh, độ dài bền vững); bước 2 đặt lên mỗi điểm u một phân bố xác suất khả vi (bài dùng Gauss đối xứng chuẩn hóa g\_u(x,y) = 1/(2πσ²)·exp(−[(x−u\_x)² + (y−u\_y)²]/(2σ²))) nhân với hàm trọng số f không âm, liên tục, khả vi từng khúc và bằng 0 trên trục hoành — điều kiện bắt buộc để có tính ổn định — cho ra mặt persistence ρ\_B(z) = Σ f(u)·φ\_u(z); bước 3 chia lưới n ô, mỗi pixel là tích phân của ρ\_B trên ô đó. Hàm trọng số dùng trong thực nghiệm là tuyến tính từng khúc w\_b(t) = 0 nếu t ≤ 0, t/b nếu 0 < t < b, 1 nếu t ≥ b, với b là độ dài bền vững lớn nhất trong tất cả các lần thử của thí nghiệm; PI của nhiều chiều đồng điều (H0, H1, ...) được nối thành một vector duy nhất, và khi tọa độ sinh bằng 0 với mọi điểm (thường gặp ở H0) thì dùng PI một chiều thay vì hai chiều. Các định lý ổn định: Thm 4 chặn ||ρ\_B − ρ\_B'||\_∞ ≤ √10·(||f||\_∞·|∇φ| + ||φ||\_∞·|∇f|)·W1(B,B'); Thm 5 cho ảnh với hệ số √(10n)·A (hệ số này tiến ra vô cực khi tăng độ phân giải); Thm 9 và 10 riêng cho Gauss có hệ số tốt hơn và không phụ thuộc n, ||I(ρ\_B) − I(ρ\_B')||\_1 ≤ (√5·|∇f| + √(10/π)·||f||\_∞/σ)·W1(B,B'); và Remark 6 rất quan trọng khi khẳng định nhân k(B,B') sinh bởi tích vô hướng hai PI là không ổn định theo W\_p với mọi 1 < p ≤ ∞, tức không ổn định theo khoảng cách bottleneck.\par\vspace{3pt}
\noindent \textbf{Mô hình.} Đây không phải mạng nơ-ron mà là một phép trích xuất đặc trưng; các bộ phân loại đặt bên trên PI gồm K-medoids với K = 6 và 1.000 lần khởi tạo ngẫu nhiên (cải thiện bằng lặp Voronoi), SVM tuyến tính chính quy chuẩn 1 (sparse SVM, SSVM) kiểu one-against-all, và subspace discriminant ensemble kiểu Ho 1998. Bài không nêu số tham số mô hình, không nêu kích thước mô hình, không nêu độ trễ suy luận; chỉ nêu chiều vector đặc trưng (độ phân giải 20×20 cho vector 400 chiều ở mục 6.3, 10×10 cho 100 chiều ở mục 6.4.2) và phần cứng đo thời gian là laptop Intel Core i5 1,3 GHz, 4 GB RAM.\par\vspace{3pt}
\noindent \textbf{Dữ liệu.} Toàn bộ là dữ liệu tổng hợp, không có dữ liệu y sinh, không có ECG, không có chuỗi thời gian sinh lý: (1) bộ hình học 6 lớp gồm hình lập phương đơn vị, đường tròn đường kính 1, mặt cầu đường kính 1, 3 cụm tâm ngẫu nhiên trong lập phương đơn vị, 3 cụm lồng trong 3 cụm (nhiễu loạn < 0,1) và hình xuyến đường kính lớn 1 đường kính nhỏ 0,5, mỗi lớp sinh 25 đám mây và mỗi đám mây 500 điểm lấy ngẫu nhiên đều cộng nhiễu Gauss, tổng 150 đám mây ở hai mức nhiễu η = 0,05 và η = 0,1, tính sơ đồ H0 và H1 bằng lọc Vietoris-Rips với metric Euclid trong R³; (2) hệ động lực rời rạc linked twist map x\_\{n+1\} = x\_n + r·y\_n(1−y\_n) mod 1, y\_\{n+1\} = y\_n + r·x\_n(1−x\_n) mod 1 với 5 giá trị r = 2,5; 3,5; 4,0; 4,1; 4,3, mỗi giá trị 50 điều kiện ban đầu ngẫu nhiên và mỗi quỹ đạo 1.000 vòng lặp, tổng 250 đám mây 2 chiều. (3) Phương trình đạo hàm riêng Kuramoto-Sivashinsky bất đẳng hướng (aKS) mô phỏng phổ Fourier trên lưới không gian 512×512, bước thời gian Runge-Kutta ETD bậc 4, điều kiện biên tuần hoàn và điều kiện đầu là nhiễu trắng biên độ thấp, với 5 giá trị r = 1; 1,25; 1,5; 1,75; 2 và mỗi giá trị 30 lần chạy cho tổng 150 mặt, xét các thời điểm t = 3, 5, 10 (t = 15 và 20 cho kết quả tương tự t = 10); bài không nêu tần số lấy mẫu vật lý vì không có tín hiệu thời gian thực.\par\vspace{3pt}
\noindent \textbf{Giao thức đánh giá.} Đây là điểm yếu nhất của bài: thí nghiệm chính (mục 6.1, Bảng 1) dùng K-medoids, tức phân cụm không giám sát chạy trên cả 150 đám mây mà không có tập huấn luyện và tập kiểm thử tách rời, trong đó "độ chính xác" là tỷ lệ phần trăm đám mây được gán vào medoid cùng lớp và bài lấy kết quả của lần phân cụm có điểm thấp nhất trong 1.000 lần khởi tạo, tức con số lạc quan nhất chứ không phải kỳ vọng; mục 6.2 quét 20 độ phân giải (5×5 đến 100×100, bước 5) và 20 giá trị phương sai (0,01 đến 0,2, bước 0,01) rồi đo độ chính xác trên chính tập đó, không có tập giữ riêng. Chỉ mục 6.3 có giao thức chia dữ liệu rõ ràng (SSVM one-against-all với kiểm định chéo 5 fold), còn mục 6.4.1 chỉ viết là thực hiện phân loại và kiểm định chéo bằng subspace discriminant ensemble rồi chạy 10 lần lấy trung bình mà không nêu số fold cũng không nêu tỷ lệ chia train/test, và ở mục 6.4.2 chỉ riêng bộ phân loại dựa trên phương sai chiều cao có nói rõ khớp phân bố chuẩn trên 2/3 dữ liệu và lấy trung bình 100 lần chia, còn phần PI kết hợp subspace ensemble thì bài không nêu tỷ lệ chia. Không có khái niệm bệnh nhân hay bản ghi nên không có tách theo chủ thể và tách theo bản ghi, cũng không có dung sai ms vì bài không làm bài toán phát hiện sự kiện theo thời gian; ngoài ra tham số b của hàm trọng số được định nghĩa là độ dài bền vững lớn nhất trong tất cả các lần thử, tức lấy cả từ dữ liệu kiểm thử, là một dạng rò rỉ dữ liệu nhẹ.\par\vspace{3pt}
\noindent \textbf{Kết quả chính.} Bảng 1 (ma trận khoảng cách 150×150, phân cụm K-medoids, độ chính xác và thời gian, mức nhiễu 0,05 và 0,1): sơ đồ gốc PD H0 với L1 đạt 96,0\% trong 37.346 s và 96,0\% trong 42.613 s, PD H0 L2 đạt 91,3\% trong 24.656 s, PD H0 L∞ đạt 60,7\% trong 1.133 s, PD H1 L1 đạt 100\% trong 657 s và 96,0\% trong 703 s, PD H1 L2 đạt 100\% trong 984 s, PD H1 L∞ đạt 81,3\% trong 527 s; persistence landscape PL H0 L1 đạt 92,7\% trong 29 s, PL H1 L1 đạt 83,3\% trong 36 s, PL H1 L2 đạt 83,3\% trong 50 s; còn Persistence Image PI H0 đạt 93,3\% (L1), 92,7\% (L2) và 94,0\% (L∞) đều trong 9 s, PI H1 đạt 100\% với cả ba chuẩn trong 17-18 s ở nhiễu 0,05 và 95,3-96,0\% ở nhiễu 0,1, nghĩa là PI đạt 100\% trong 17 giây trong khi sơ đồ gốc cần 657 giây cho cùng kết quả và PI H0 nhanh hơn PD H0 khoảng 4000 lần (9 s so với 37.346 s), với tham số σ = 0,1 và độ phân giải 20×20. Mục 6.3 cho thấy SSVM 5 fold trên PI 20×20 với phương sai 0,0001 ở nhiễu 0,05 đạt 100\% độ chính xác và chỉ cần 10 pixel phân biệt (ví dụ lớp đường tròn chỉ cần pixel số 21 và 22 trong 400, lớp hình xuyến cần pixel 59 và 98). Mục 6.4.1 (linked twist map, subspace discriminant ensemble, trung bình 10 lần chạy) cho thấy nối H0 + H1 đạt 82,5\% so với 49,8\% khi chỉ dùng H0 và 65,7\% khi chỉ dùng H1; mục 6.4.2 (Bảng 2, aKS) tại t = 3/5/10 lần lượt cho mặt được thu nhỏ rồi phân loại 26,0\%/19,3\%/19,3\%, bộ phân loại phân bố chuẩn trên phương sai chiều cao 20,74\%/75,2\%/77,62\%, PI H0 58,3\%/96,0\%/94,7\%, PI H1 67,7\%/87,3\%/93,3\% và PI H0 + H1 72,7\%/95,3\%/97,3\%; về độ nhạy tham số, khi đổi phương sai Gauss từ 0,0001 đến 0,1 thì độ chính xác H0 đổi dưới 1 điểm phần trăm còn H1 dao động khoảng 5 điểm, và khi đổi độ phân giải từ 5 lên 20 thì H0 đổi hơn 3 điểm rồi tụt 6 điểm ở độ phân giải 25.\par\vspace{3pt}
\noindent \textbf{Điểm yếu.} Thí nghiệm chính (Bảng 1) không phải phân loại có giám sát mà là phân cụm K-medoids không giám sát trên toàn bộ 150 mẫu, con số báo cáo lại là lần tốt nhất trong 1.000 khởi tạo nên không có tập kiểm thử độc lập và con số mang tính lạc quan; mục 6.4.1 cùng phần PI của mục 6.4.2 không nêu giao thức chia dữ liệu (số fold, tỷ lệ train/test) nên không tái lập được; và tham số b của hàm trọng số lấy từ độ dài bền vững lớn nhất của tất cả các lần thử kể cả dữ liệu đánh giá nên có rò rỉ dữ liệu nhẹ. Tính ổn định chỉ đúng với 1-Wasserstein, còn Remark 6 khẳng định rõ nhân sinh bởi PI không ổn định với W\_p khi p > 1 kể cả bottleneck, nghĩa là một điểm trong sơ đồ bị dịch xa có thể làm PI thay đổi không kiểm soát; ba lựa chọn về độ phân giải, phân bố cùng phương sai và hàm trọng số đều không chuẩn tắc, điều bài tự nhận là điểm yếu. Chi phí tính toán chỉ báo cáo bằng giây trên một laptop i5 1,3 GHz cho ma trận khoảng cách, còn bài không nêu số tham số, không nêu độ trễ suy luận một mẫu, không nêu kích thước mô hình; toàn bộ dữ liệu là tổng hợp và không có ứng dụng y sinh nào do chính nhóm tác giả thực hiện (ứng dụng khảo cổ học chỉ là trích dẫn từ Zeppelzauer 2016); ngoài ra bài tự mâu thuẫn ở giá trị tham số r của linked twist map khi thân bài ghi 2,5 còn chú thích Hình 4 ghi 2.\par\vspace{3pt}
\noindent \textbf{\color{tdtblue}So với mô hình của nhóm.} Không so sánh trực tiếp được với FetalQRS-TCN vì miền dữ liệu hoàn toàn khác (đám mây điểm tổng hợp và mặt PDE, không phải ECG bụng), không có Se/PPV/F1 theo nhịp, không có dung sai ms, không có giao thức tách theo bản ghi hay tách theo chủ thể. Tuy nhiên đây là bài báo gốc của biểu diễn mà nhóm đã thử nghiệm và chốt bảng (cnn2d\_PersistenceImage đạt 27,42 so với 97,14 của CNN 1D cùng ngân sách tham số), và điểm mấu chốt để nhóm viện dẫn là toàn bộ thí nghiệm của Adams và cộng sự đều là phân loại toàn cục — một nhãn cho cả một đám mây điểm, một quỹ đạo, một mặt — trong khi bài không hề khẳng định và cũng không kiểm chứng rằng PI phù hợp cho định vị sự kiện theo thời gian ở độ phân giải từng mẫu. Vì vậy kết quả 27,42 của nhóm không mâu thuẫn với bài báo mà cho thấy PI bị dùng sai mục đích chứ không phải PI sai; ngoài ra Remark 6 (không ổn định theo bottleneck) và việc PI đạt 100\% chỉ với 10 pixel cho thấy PI mã hóa thông tin rất thưa, phù hợp với việc CNN 2D trên PI học kém.\par\vspace{3pt}
\noindent \textbf{\color{tdtblue}Nhóm rút ra gì.} Nên áp dụng: nếu vẫn dùng PI cho fSQI thì phải đóng khung lại thành bài toán một nhãn cho một cửa sổ đúng thiết kế gốc của PI chứ không dùng cho sequence-to-sequence, đồng thời nối PI của H0 và H1 lại với nhau vì bài chứng minh lợi ích lớn (82,5\% so với 49,8\% khi chỉ H0 và 65,7\% khi chỉ H1 trên linked twist map; 97,3\% so với 94,7\% và 93,3\% trên aKS), không cần quét siêu tham số lớn mà có thể cố định 20×20 và σ = 0,1 do độ chính xác gần như không đổi khi σ từ 0,0001 đến 0,1 và độ phân giải từ 5×5 đến 100×100, và dùng SSVM tuyến tính để chọn pixel nhằm diễn giải vùng nào của sơ đồ bền vững mang thông tin — điều rất có giá trị cho một bài báo Q1 về fSQI. Nên tránh: tuyệt đối không tính tham số chuẩn hóa b (max persistence) trên toàn bộ dữ liệu mà phải tính chỉ trên tập huấn luyện nếu không sẽ là rò rỉ dữ liệu, không báo cáo kết quả phân cụm tốt nhất trong N lần khởi tạo như độ chính xác mà phải báo cáo trung bình và độ lệch chuẩn, và phải nêu rõ số fold cùng tỷ lệ chia — đây chính là thứ bài này thiếu khiến nó không tái lập được.\par\vspace{3pt}
\vspace{6pt}
\subsection{Carrière và cs. 2020, AISTATS \normalfont\small (p24)}
\noindent{\footnotesize\itshape Mathieu Carriere, Frederic Chazal, Yuichi Ike, Theo Lacombe, Martin Royer, Yuhei Umeda. "PersLay: A Neural Network Layer for Persistence Diagrams and New Graph Topological Signatures". Proceedings of the 23rd International Conference on Artificial Intelligence and Statistics (AISTATS) 2020, Palermo, Italy. PMLR Volume 108. Ban quyen 2020. FILE KHONG CO DOI va KHONG CO so trang. Don vi: Columbia ...}\par\vspace{2pt}
\noindent{\small\color{warnred}\textbf{Kết luận: không so trực tiếp được}}\par\vspace{4pt}
\noindent \textbf{Họ làm gì.} Bài có hai đóng góp; thứ nhất là định nghĩa PersLay, một lớp mạng nơ-ron khả vi tự động nhận đầu vào là sơ đồ bền vững với công thức PersLay(Dg) = op(\{w(p)·φ(p)\} với p thuộc Dg), trong đó op là phép gộp bất biến hoán vị (tổng, max, min, giá trị lớn thứ k) còn hàm trọng số w và hàm biến đổi điểm φ đều học được, nên thay vì cố định cách vector hóa như Persistence Image hay Persistence Landscape thì PersLay học cách vector hóa theo đúng tác vụ, và bài chỉ rõ hầu hết cách vector hóa cổ điển (landscape, silhouette, persistence surface/PI của Adams 2017, Sliced Wasserstein) chỉ là trường hợp riêng khi chọn φ và op cụ thể. Đóng góp thứ hai là một họ chữ ký tô-pô mới cho đồ thị dựa trên sơ đồ bền vững mở rộng xây trên Heat Kernel Signature, có chứng minh ổn định; đây không phải bài về tín hiệu, không có ECG và không có chuỗi thời gian.\par\vspace{3pt}
\noindent \textbf{Phương pháp.} Chữ ký đồ thị: tính Laplacian chuẩn hóa L\_w = I − D\textasciicircum{}(−1/2) A D\textasciicircum{}(−1/2) với trị riêng 0 ≤ λ\_1 ≤ ... ≤ λ\_n ≤ 2, dùng Heat Kernel Signature hks\_\{G,t\}(v) = Σ\_k exp(−t·λ\_k)·ψ\_k(v)² làm hàm lọc trên đỉnh rồi tính bền vững mở rộng (xét cả tập mức dưới và tập mức trên) cho ra 4 loại điểm Ord0, Rel1, Ext+0, Ext−1, ưu điểm quan trọng là mọi điểm đều có tọa độ hữu hạn nên không cần xử lý đặc biệt như Hofer 2017; Thm 2.2 chứng minh ổn định theo nhiễu loạn đồ thị với d\_B(Dg(G,t), Dg(G',t)) ≤ C(G,t)·||W||\_F và Thm 2.3 chứng minh ánh xạ t → Dg(G,t) là 2-Lipschitz, tức d\_B(Dg(G,t), Dg(G,t')) ≤ 2|t − t'|. PersLay cài đặt ba hàm biến đổi điểm là tam giác (Λ\_p(t) = max\{0, y − |t − x|\}), Gauss (Γ\_p(t) = exp(−||p − t||²/(2σ²))) và đường thẳng (L\_δ(p) là tích vô hướng của p với vector chỉ phương cộng bias); hàm trọng số w là một lưới N×N trên hình vuông đơn vị với mỗi ô là một tham số học được (N thường bằng 10 hoặc 20), phép gộp op là tổng, kiến trúc chỉ hai lớp (PersLay → chuẩn hóa batch → một lớp fully-connected) có kết hợp thêm đặc trưng đồ thị truyền thống (trị riêng Laplacian chuẩn hóa và các decile của HKS), tiền xử lý prom(500) giữ 500 điểm xa đường chéo nhất, tối ưu bằng Adam với tốc độ học 0,01, Exponential Moving Average decay 0 / 0,9 / 0,99 tùy bộ, 70 đến 1.000 epoch và mini-batch 128. Tham số khuếch tán thường dùng t = 0,1 và t = 10 (có thử học luôn t nhưng không cải thiện đáng kể mà làm tăng mạnh thời gian chạy vì phải tính lại toàn bộ sơ đồ sau mỗi epoch nên không đưa vào kết quả chính); về tính liên tục, bài thừa nhận công thức với op = tổng nói chung không liên tục theo các metric thông dụng của sơ đồ, theo Divol và Lacombe 2019 nó liên tục theo d\_s khi và chỉ khi φ(p) = φ\_ngoài(p)·||p − đường chéo||\_s, và khi s = 1 với φ\_ngoài 1-Lipschitz thì ổn định với chặn d\_1.\par\vspace{3pt}
\noindent \textbf{Mô hình.} Mạng rất nhỏ, chỉ dùng 2 lớp (PersLay cộng một lớp fully-connected sau chuẩn hóa batch), và tác giả nói rõ chọn kiến trúc đơn giản để tạo ra hiểu biết chứ không nhằm đạt điểm cao nhất; cấu hình từng kênh được ghi ở Bảng 5 phụ lục theo ký hiệu Im(p,(a,b),q,op) nghĩa là lưới p×p, tích chập a bộ lọc kích thước b×b, hàm trọng số học trên lưới q×q, và Pm(d1,d2,q,op) là biến đổi đường thẳng với d1 đường, hàm đồng biến hoán vị chiều d2, trọng số lưới q×q. Bài không nêu số tham số huấn luyện, không nêu kích thước checkpoint và không nêu độ trễ suy luận cho một mẫu; chỉ có thời gian chạy (Bảng 6 phụ lục, trung bình 100 lần chạy) là 2,30 đến 5,91 giây trên MUTAG bằng CPU Intel i5-8350U 1,70 GHz và 26,0 đến 61,4 giây trên COLLAB bằng GPU NVIDIA P100.\par\vspace{3pt}
\noindent \textbf{Dữ liệu.} Không có dữ liệu tín hiệu sinh lý, không có ECG và không có tần số lấy mẫu; dữ liệu hệ động lực tổng hợp dùng chính linked twist map như bài Adams 2017 với 5 lớp r = 2,5; 3,5; 4,0; 4,1; 4,3 và mỗi quỹ đạo 1.000 điểm, tạo ORBIT5K gồm 1.000 quỹ đạo mỗi lớp (5.000 đám mây) và ORBIT100K gồm 20.000 quỹ đạo mỗi lớp (100.000 đám mây, quy mô mà các phương pháp nhân không xử lý nổi), sơ đồ sinh bằng lọc AlphaComplex chiều 0 và 1 của Gudhi với các lớp cân bằng. Ngoài ra có 11 bộ đồ thị chuẩn (Bảng 4 phụ lục, ghi số đồ thị / số lớp / số đỉnh trung bình / số cạnh trung bình): REDDIT5K 5.000/5/508,5/594,9; REDDIT12K 12.000/11/391,4/456,9; COLLAB 5.000/3/74,5/2457,5; IMDB-B 1.000/2/19,77/96,53; IMDB-M 1.500/3/13,00/65,94; COX2 467/2/41,22/43,45; DHFR 756/2/42,43/44,54; MUTAG 188/2/17,93/19,79; PROTEINS 1.113/2/39,06/72,82; NCI1 4.110/2/29,87/32,30; NCI109 4.127/2/29,68/32,13, trong đó COX2, DHFR, MUTAG, PROTEINS, NCI1 và NCI109 có nhãn thuộc tính trên đỉnh đồ thị.\par\vspace{3pt}
\noindent \textbf{Giao thức đánh giá.} ORBIT5K và ORBIT100K chia ngẫu nhiên 70\% huấn luyện và 30\% kiểm thử rồi báo cáo độ chính xác kiểm thử trung bình trên 100 lần chạy (làm giống Le và Yamada 2018 để so sánh được), còn các bộ đồ thị dùng 10 lần chạy kiểm định chéo 10 fold và báo cáo hai cột là "Mean" (trung bình của 10 lần 10-fold, nhất quán với Zhang và cộng sự 2018) và "Max" (kết quả 10-fold tốt nhất trong 10 lần, được đối chiếu với con số một lần 10-fold của các đối thủ). Không có tách theo chủ thể và không có tách theo bản ghi vì các bộ dữ liệu không có khái niệm bệnh nhân hay bản ghi lặp lại — mỗi đồ thị là một quan sát độc lập nên chia ngẫu nhiên ở đây không phải rò rỉ dữ liệu, khác hẳn trường hợp nhiều cửa sổ ECG từ cùng một sản phụ — và cũng không có dung sai ms vì bài không làm bài toán phát hiện sự kiện theo thời gian. Điểm cần cảnh báo: Bảng 6 phụ lục (ảnh hưởng của kích thước lưới trọng số, lựa chọn φ và lựa chọn op) báo cáo trực tiếp độ chính xác trên tập kiểm thử cho từng lựa chọn và đánh dấu cấu hình được dùng ở phần chính, tức siêu tham số được chọn dựa trên kết quả kiểm thử; tương tự Hình 8 quét tham số t từ 10\textasciicircum{}−2 đến 10\textasciicircum{}2 và vẽ độ chính xác kiểm thử.\par\vspace{3pt}
\noindent \textbf{Kết quả chính.} Bảng 1 (ORBIT, độ chính xác ± độ lệch chuẩn): trên ORBIT5K, Persistence Scale Space Kernel đạt 72,38 (±2,4), Persistence Weighted Gaussian Kernel 76,63 (±0,7), Sliced Wasserstein Kernel 83,6 (±0,9), Persistence Fisher Kernel 85,9 (±0,8) và PersLay 87,7 (±1,0); trên ORBIT100K chỉ PersLay chạy được và đạt 89,2 (±0,3) vì các phương pháp nhân không xử lý nổi quy mô này. Bảng 2 (đồ thị, cột PersLay Mean/Max so với đối thủ tốt nhất): REDDIT5K 55,6/56,5 (RetGK 56,1; GIN 57,0; GCNN 52,9; FGSD 47,8), REDDIT12K 47,7/49,1 (RetGK 48,7; GCNN 46,6), COLLAB 76,4/78,0 (RetGK 81,0; FGSD 80,0; GIN 80,1), IMDB-B 71,2/72,6 (GIN 74,3; FGSD 73,6), IMDB-M 48,8/52,2 (FGSD 52,4; GIN 52,1), COX2 80,9/81,6 (RetGK 80,1), DHFR 80,3/80,9 (RetGK 81,5), MUTAG 89,8/91,5 (FGSD 92,1; RetGK 90,3), PROTEINS 74,8/75,9 (GCNN 76,3; GIN 75,9), NCI1 73,5/74,0 (RetGK 84,5 — thua xa) và NCI109 69,5/70,1 (FGSD 78,8); để so sánh, Hofer và cộng sự 2017 cũng dùng sơ đồ bền vững nuôi mạng chỉ đạt 54,5\% trên REDDIT5K và 44,5\% trên REDDIT12K. Nghiên cứu loại bỏ ở Bảng 7 phụ lục (chỉ đặc trưng phổ / chỉ sơ đồ mở rộng / chỉ sơ đồ thường / PersLay đầy đủ) cho REDDIT5K 49,7/55,0/52,5/55,6; REDDIT12K 39,7/44,2/40,1/47,7; COLLAB 67,8/71,6/69,2/76,4; IMDB-B 67,6/68,8/64,7/71,2; MUTAG 85,8/85,1/70,2/89,8; DHFR 69,5/78,2/71,8/80,3; NCI1 65,3/72,3/68,9/73,5, tức bền vững mở rộng hơn hẳn bền vững thường, rõ nhất ở MUTAG (85,1 so với 70,2, chênh 14,9 điểm); còn Bảng 6 cho thấy lưới trọng số càng lớn càng quá khớp (MUTAG train/test: không lưới 92,3/88,9, lưới 2×2 91,1/88,8, 5×5 91,7/89,6, 10×10 92,3/89,9, 20×20 93,7/88,3, 50×50 94,1/87,7; COLLAB: không lưới 76,5/75,3, 10×10 80,0/76,5, 20×20 83,5/73,9, 50×50 94,0/71,3), trong khi trên MUTAG hàm biến đổi Gauss đạt 92,5/89,7 so với đường thẳng 89,2/84,2 và tam giác 91,5/85,0, phép gộp tổng đạt 92,3/89,5 so với max 91,9/87,4.\par\vspace{3pt}
\noindent \textbf{Điểm yếu.} Bài báo cáo cột "Max" là kết quả 10-fold tốt nhất trong 10 lần lặp rồi đối chiếu với con số một lần của đối thủ nên đây là so sánh không đối xứng và thiên lệch lạc quan (con số trung thực là cột "Mean"), đồng thời Bảng 6 và Hình 8 chọn siêu tham số (kích thước lưới, φ, op, tham số t) dựa trên độ chính xác trên tập kiểm thử. Bài không nêu số tham số huấn luyện, không nêu kích thước mô hình và không nêu độ trễ suy luận một mẫu mà chỉ có thời gian chạy tổng của một lần thí nghiệm; đồng thời thua xa các phương pháp phi-tô-pô trên các bộ có nhãn đỉnh (NCI1 73,5 so với RetGK 84,5, kém 11 điểm; NCI109 69,5 so với FGSD 78,8; COLLAB 76,4 so với RetGK 81,0) và bài tự thừa nhận tô-pô không phân biệt được trên các bộ NCI. PersLay với op = tổng nói chung không liên tục theo metric của sơ đồ bền vững, muốn có liên tục thì phải ép hàm trọng số nhỏ gần đường chéo trong khi bài lại lập luận rằng không có lý do gì coi điểm gần đường chéo là kém quan trọng, tức có mâu thuẫn giữa tính ổn định lý thuyết và tự do học tập; ngoài ra tiền xử lý prom(500) cắt bỏ điểm tùy ý mà không khảo sát ảnh hưởng, mọi tác vụ đều là phân loại toàn cục (một nhãn cho một đồ thị hoặc một quỹ đạo) chứ không có bài toán định vị sự kiện theo thời gian nào, và toàn bộ là dữ liệu benchmark công khai không có ứng dụng y sinh thật.\par\vspace{3pt}
\noindent \textbf{\color{tdtblue}So với mô hình của nhóm.} Không so sánh trực tiếp được với FetalQRS-TCN vì miền dữ liệu hoàn toàn khác (đồ thị mạng xã hội, phân tử hóa học, quỹ đạo hệ động lực — không có tín hiệu sinh lý), không có Se/PPV/F1 theo nhịp, không có dung sai ±50 ms, không có tách theo bản ghi hay tách theo chủ thể, và không có số tham số để đối chiếu với 113.481 tham số của nhóm. Ý nghĩa gián tiếp lại rất lớn cho lập luận của nhóm: PersLay là nỗ lực mạnh nhất hiện nay để bắt TDA làm việc trong một pipeline học sâu vì nó học luôn cách vector hóa thay vì cố định như Persistence Image mà nhóm đã thử, vậy mà ngay cả khi học đầu-cuối trên các tác vụ phân loại toàn cục thuận lợi nhất cho TDA thì nó cũng chỉ ngang bằng các phương pháp hiện đại và thua rõ rệt ở nhiều bộ (NCI1 kém 11 điểm so với RetGK). Điều này ủng hộ mạnh mẽ quyết định chuyển hướng của nhóm: nếu cả PersLay học được cũng chỉ ngang bằng trên tác vụ dễ nhất của nó thì kỳ vọng Persistence Image kết hợp CNN 2D thắng CNN 1D ở tác vụ khó hơn nhiều (định vị đỉnh từng mẫu) là không có cơ sở, và đây là trích dẫn tốt để biến 27,42\% của nhóm từ kết quả thất bại thành kết quả âm tính có kỳ vọng trước.\par\vspace{3pt}
\noindent \textbf{\color{tdtblue}Nhóm rút ra gì.} Nên áp dụng: nếu nhóm quay lại TDA cho fSQI thì PersLay là lựa chọn đúng đắn hơn Persistence Image cố định vì lưới trọng số học được còn cho phép diễn giải vùng nào của sơ đồ quan trọng sau huấn luyện (Hình 5 phụ lục) — rất có giá trị cho bài báo Q1 — đồng thời nên xét dùng bền vững mở rộng thay vì bền vững thường vì nó bắt được cả nhánh hướng lên chứ không chỉ đỉnh và mọi điểm có tọa độ hữu hạn (chênh lệch 14,9 điểm trên MUTAG cho thấy điều này quan trọng), nên giữ lưới trọng số N = 10 vì lưới lớn gây quá khớp (COLLAB 50×50: huấn luyện 94,0 nhưng kiểm thử 71,3), nên bắt chước cách làm nghiên cứu loại bỏ của Bảng 7 (báo cáo riêng "chỉ đặc trưng phổ", "chỉ sơ đồ bền vững" và "kết hợp") để chứng minh TDA đóng góp thêm bao nhiêu so với các chỉ số SQI cổ điển, và nên chọn kiến trúc tối giản kèm nói rõ mục đích là để hiểu chứ không để đạt điểm cao nhất. Nên tránh: tuyệt đối không báo cáo "Max" (lần tốt nhất trong nhiều lần lặp) như chỉ số chính mà phải báo cáo trung bình và độ lệch chuẩn trên toàn bộ giao thức tách theo bản ghi đúng như nhóm đang làm (97,43 ± 4,37), không chọn siêu tham số dựa trên độ chính xác tập kiểm thử như Bảng 6 của họ, và phải báo cáo số tham số, độ trễ cùng kích thước checkpoint — thứ bài này thiếu và nhóm đang làm tốt hơn (113.481 tham số, 0,48 MB, 4,35 ms).\par\vspace{3pt}
\vspace{6pt}
\subsection{Som và cs. 2020, CVPRW \normalfont\small (p25)}
\noindent{\footnotesize\itshape Anirudh Som, Hongjun Choi, Karthikeyan Natesan Ramamurthy, Matthew P. Buman, Pavan Turaga. "PI-Net: A Deep Learning Approach to Extract Topological Persistence Images". Metadata cua file PDF ghi subject = "IEEE Conference on Computer Vision and Pattern Recognition Workshops" (CVPRW), nam 2020. Ma nguon: https://github.com/anirudhsom/PI-Net. Kinh phi: NIH R01GM135927, NSF CAREER 1452163, IRB ...}\par\vspace{2pt}
\noindent{\small\color{warnred}\textbf{Kết luận: không so trực tiếp được}}\par\vspace{4pt}
\noindent \textbf{Họ làm gì.} Bài giải quyết nút thắt lớn nhất của TDA khi dùng trong học sâu: việc tính Persistence Diagram rồi chuyển thành Persistence Image (PI) bằng công cụ giải tích truyền thống rất chậm và không khả vi nên không gắn được vào mạng neural để lan truyền ngược. Nhóm tác giả huấn luyện một CNN học ánh xạ trực tiếp từ dữ liệu thô sang PI, bỏ qua toàn bộ bước tính PD, với hai kiến trúc Signal PI-Net cho chuỗi thời gian đa biến và Image PI-Net cho ảnh nhiều kênh, rồi kiểm chứng trên nhận dạng hoạt động người và phân loại ảnh.\par\vspace{3pt}
\noindent \textbf{Phương pháp.} Nhãn chuẩn PI được tạo bằng Scikit-TDA (Ripser và Persim): với chuỗi thời gian dùng lọc theo tập mức, trong đó cực tiểu địa phương sinh ra đặc trưng H0 và cực đại địa phương làm nó chết; với ảnh thì mỗi kênh được biến thành đám mây điểm 3D rồi tính đồng điều bền vững bậc 1. Tham số PI dùng lưới 50 x 50 và hạt nhân Gaussian trọng số theo life-time, với chuỗi thời gian đặt độ lệch chuẩn 0,25 và dải birth-time [-10, 10], với CIFAR10/CIFAR100 đặt [0; 0,3] và 0,01, với SVHN đặt [0; 0,2] và 0,005. Mạng hồi quy PI được huấn luyện bằng MSE (Signal PI-Net, batch 128, 1000 epoch) hoặc Binary Cross-Entropy (Image PI-Net, batch 32, 415 epoch); khi dùng cho phân loại, PI được giảm chiều từ 7500 xuống 32 bằng PCA rồi nối vào đầu ra lớp global-average-pooling của CNN 1D.\par\vspace{3pt}
\noindent \textbf{Mô hình.} Signal PI-Net gồm 4 lớp Conv1D với 128, 256, 512 và 1024 bộ lọc, hạt nhân 3, kèm BatchNorm, ReLU và MaxPool, kết thúc bằng global-average-pooling và lớp Dense kích thước 2500 x n rồi reshape thành 50 x 50 x n; Image PI-Net dùng Conv2D với tất cả các lớp đều 128 bộ lọc, lớp biến tiềm ẩn Dense 2500 và một lớp Deconv2D. Bộ phân loại đi kèm gồm MLP 8 lớp Dense mỗi lớp 128 đơn vị, CNN 1D 10 lớp với 32 bộ lọc, và DenseNet depth 16 với growth rate 12; bài không nêu số tham số huấn luyện được của bất kỳ mạng nào, cũng không nêu dung lượng checkpoint.\par\vspace{3pt}
\noindent \textbf{Dữ liệu.} GENEactiv gồm 29 lớp hoạt động, 152 đối tượng, gia tốc kế 3 trục đeo cổ tay lấy mẫu 100 Hz và hạ mẫu xuống 50 Hz, khung không chồng lấn dài 250 bước (5 giây) và 500 bước (10 giây); USC-HAD gồm 12 lớp hoạt động, 14 đối tượng, 100 Hz hạ mẫu xuống 50 Hz, chỉ trích khung 250 bước; bài không nêu tổng số khung cụ thể. Phần ảnh dùng CIFAR10 (50.000 huấn luyện / 10.000 kiểm tra), CIFAR100 (50.000/10.000) và SVHN (73.257/26.032), mỗi ảnh 32 x 32 x 3.\par\vspace{3pt}
\noindent \textbf{Giao thức đánh giá.} Giao thức không đồng nhất giữa hai bộ dữ liệu và đây là điểm yếu nghiêm trọng nhất: với GENEactiv, khoảng 75\% khung được dùng để huấn luyện và phần còn lại để kiểm tra, tức chia theo khung chứ không theo đối tượng, và bài không nêu có tách đối tượng hay không, nên phải coi đây là chia ngẫu nhiên trên cùng đối tượng, tức rò rỉ dữ liệu trên 152 người. Với USC-HAD, 8 đối tượng đầu dùng để huấn luyện và 6 đối tượng còn lại để kiểm tra, tương đương tách theo chủ thể một lần chia chứ không xoay vòng. Đơn vị phân tích là khung tín hiệu; kết quả nhận dạng hoạt động báo cáo trung bình ± độ lệch chuẩn trên 5 lần chạy và kết quả ảnh trên 3 lần chạy; p-value (t-test hai phía so với DenseNet gốc) chỉ có cho thí nghiệm ảnh, không có kiểm định thống kê nào cho thí nghiệm chuỗi thời gian; đây không phải bài toán dò đỉnh nên không có khái niệm dung sai ms.\par\vspace{3pt}
\noindent \textbf{Kết quả chính.} F1 có trọng số trên GENEactiv với khung 250 bước: MLP-PI 46,45 ± 0,32; MLP-Signal PI-Net 49,47 ± 0,69; MLP-SF 41,70 ± 0,41; CNN 1D 53,56 ± 0,31; CNN 1D + PI 54,28 ± 0,23; CNN 1D + Signal PI-Net 54,41 ± 0,21; với khung 500 bước, CNN 1D 54,97 ± 1,35 lên 56,41 ± 0,22 khi thêm Signal PI-Net; trên USC-HAD khung 250 bước, CNN 1D 54,58 ± 0,64 so với CNN 1D + Signal PI-Net 57,82 ± 0,78. Phân loại ảnh: CIFAR10 DenseNet 83,80 ± 0,12 lên 84,82 ± 0,19 khi thêm Image PI-Net (p = 0,0160), SVHN 95,65 ± 0,00 lên 95,95 ± 0,08 (p = 0,0038); hàm mất cuối cùng của Signal PI-Net là 0,00158 trên tập kiểm tra với khung 250 bước. Thời gian tính PI cho một ảnh: TDA truyền thống trên CPU 3567,29 ± 867,74 ms so với PI-Net trên CPU 125,45 ± 5,30 ms và trên GPU 2,52 ± 0,02 ms, tức nhanh hơn khoảng 1400 lần, đo trên NVIDIA GTX Titan Xp 12 GB.\par\vspace{3pt}
\noindent \textbf{Điểm yếu.} Rò rỉ dữ liệu ở GENEactiv (chia 75/25 theo khung trên 152 người) cộng với số tuyệt đối rất thấp (F1 chỉ 46–58\% cho nhận dạng hoạt động) khiến mọi chênh lệch vài điểm trên nền 54\% đều kém thuyết phục. Lợi ích của PI là biên: CNN 1D tăng từ 53,56 lên 54,41, tức 0,85 điểm, trong khi độ lệch chuẩn là 0,21–0,31 và bài không nêu p-value cho phần chuỗi thời gian; chính tác giả viết rằng việc hợp nhất PI chỉ cải thiện "marginally". Bài không báo cáo số tham số của bất kỳ mạng nào nên không thể biết phần cải thiện đến từ thông tin tô-pô hay chỉ từ việc thêm tham số, không báo cáo thời gian suy luận của Signal PI-Net, không nêu bộ lọc dải thông cho tín hiệu gia tốc, có mâu thuẫn nhỏ giữa khung 5 giây và mô tả đặc trưng "tính trên mỗi khung 10 giây", và chỉ đo sai số xấp xỉ PI bằng MSE chứ không bằng một chỉ số tô-pô như khoảng cách Wasserstein.\par\vspace{3pt}
\noindent \textbf{\color{tdtblue}So với mô hình của nhóm.} Không so sánh trực tiếp được vì khác bài toán (nhận dạng hoạt động và phân loại ảnh chứ không phải dò đỉnh fQRS), khác dữ liệu (gia tốc kế đeo tay và ảnh CIFAR/SVHN, không có ECG bụng), khác chỉ số (F1 có trọng số đa lớp so với macro F1 khớp nhịp ±50 ms) và không có khái niệm dung sai ms. Tuy vậy bài củng cố gián tiếp kết quả của nhóm: ngay trong một công trình chủ trương dùng TDA, PI dùng một mình (MLP-PI 46,45) vẫn thua xa CNN 1D học từ dữ liệu thô (53,56) trên cùng bộ dữ liệu, còn ghép PI vào CNN 1D chỉ tăng 0,85 điểm mà không có kiểm định thống kê. Kết quả của nhóm (PI + CNN 2D đạt 27,42 so với CNN 1D đạt 97,14 cùng ngân sách tham số) do đó không mâu thuẫn với tài liệu mà là trường hợp cực đoan của cùng một xu thế, vì Persistence Diagram là một đa tập điểm không có thứ tự nên mất thông tin vị trí thời gian tuyệt đối, đúng thứ mà bài toán dò đỉnh cần nhất.\par\vspace{3pt}
\noindent \textbf{\color{tdtblue}Nhóm rút ra gì.} Nếu nhóm vẫn giữ một nhánh TDA cho fSQI thì không nên tính PI bằng thư viện giải tích trong vòng lặp huấn luyện mà tính trực tiếp peak prominence (tương đương H0 lọc theo tập mức), đồng thời có thể mượn cách họ giảm chiều PI bằng PCA xuống 32 chiều rồi nối vào đầu ra global-average-pooling để đưa điểm fSQI vào FetalQRS-TCN mà không làm phồng số tham số. Bảng thời gian của họ (3567 ms cho một ảnh 32 x 32) là bằng chứng định lượng để biện minh cho quyết định loại TDA khỏi đường ống thời gian thực của nhóm (độ trễ 4,35 ms cho cửa sổ 4 giây); ngược lại phải tránh kiểu chia 75/25 theo khung trên cùng một chủ thể và tránh báo cáo cải thiện nhỏ mà không kèm kiểm định thống kê.\par\vspace{3pt}
\vspace{6pt}
\subsection{Turkes và cs. 2022, NeurIPS \normalfont\small (p26)}
\noindent{\footnotesize\itshape Renata Turkes (University of Antwerp), Guido Montufar (UCLA), Nina Otter (Queen Mary University of London). "On the Effectiveness of Persistent Homology". arXiv:2206.10551v3 [math.AT], ban v3 ngay 16 thang 1 nam 2023 (ban PDF day du kem phu luc, 39 trang). Ban rut gon 17 trang trong cung thu muc (p26\_Turkes2022-EffectivenessPH.pdf) co dong chan trang ghi ro "36th Conference on Neural Information ...}\par\vspace{2pt}
\noindent{\small\color{warnred}\textbf{Kết luận: không so trực tiếp được}}\par\vspace{4pt}
\noindent \textbf{Họ làm gì.} Đây là bài tự phản biện mạnh nhất trong cộng đồng TDA: đồng điều bền vững (PH) đã được dùng thành công ở hàng trăm ứng dụng nhưng không ai biết rõ vì sao và với loại bài toán nào thì thực sự hiệu quả, nên các tác giả phân biệt rõ "usefulness" (có ích) với "effectiveness" (thực sự tạo ra kết quả mong muốn). Họ đo PH trên ba bài toán hình học cơ bản — đếm số lỗ hổng, ước lượng độ cong, phát hiện tính lồi — so với các đường cơ sở học máy và học sâu, chứng minh một định lý mới về tính lồi, định nghĩa lọc tubular, và rút ra một bộ hướng dẫn: PH hữu ích khi và chỉ khi tín hiệu của bài toán là số chu trình k chiều, độ cong, hoặc tính lồi.\par\vspace{3pt}
\noindent \textbf{Phương pháp.} Bài đếm số lỗ hổng dùng phức alpha với hàm lọc Distance-to-Measure (m = 0,03, tức 30 láng giềng gần nhất) để chống điểm ngoại lai, trích PD bậc 1 rồi chuyển thành PI hoặc PL hoặc chỉ lấy life-time của 10 chu trình bền nhất, phân loại bằng SVM. Bài độ cong tính ma trận khoảng cách hyperbolic, Euclid hoặc cầu tương ứng với độ cong âm, bằng 0 và dương, dựng phức Vietoris-Rips qua Ripser, trích PD bậc 0 và bậc 1 rồi hồi quy bằng SVR. Bài tính lồi biến đám mây điểm thành ảnh nhị phân trên lưới 20 x 20, dùng phức lập phương thay vì Vietoris-Rips (vì Vietoris-Rips nối các thành phần rời rạc lại và làm mất dấu hiệu lõm), lọc bằng hàm tubular đo khoảng cách tới một đường thẳng, thử 9 đường thẳng, mỗi đường lấy life-time của thành phần liên thông bền thứ hai rồi lấy giá trị lớn nhất qua 9 đường làm một số vô hướng duy nhất đưa vào SVM.\par\vspace{3pt}
\noindent \textbf{Mô hình.} PH không phải mạng neural nên không có số tham số; bộ phân loại và hồi quy là SVM/SVR của sklearn với C thuộc \{0,001; 1; 100\}, tìm tham số bằng GridSearchCV 3 fold. Các đường cơ sở gồm SVM trên ma trận khoảng cách đôi một (lấy mẫu con 100 điểm), MLP nông và sâu (độ sâu 1–5, bề rộng 64–4096, dropout 0,5, mỗi cấu hình chỉ huấn luyện 2 epoch trong lưới tìm tham số) và PointNet theo kiến trúc gốc không tăng cường dữ liệu; bài không nêu số tham số của PointNet hay của các MLP.\par\vspace{3pt}
\noindent \textbf{Dữ liệu.} Toàn bộ thí nghiệm chính dùng đám mây điểm tổng hợp trong R\textasciicircum{}2 và R\textasciicircum{}3: 1.000 đám mây mỗi cái 1.000 điểm từ 20 hình dạng cho bài đếm lỗ hổng (nhãn thuộc \{0, 1, 2, 4, 9\} lỗ), 1.010 đám mây mỗi cái 500 điểm cho bài độ cong (kappa thuộc [-2, 2]), và 960 đám mây mỗi cái 5.000 điểm cho bài tính lồi (480 từ 8 hình quy chuẩn và 480 từ hình ngẫu nhiên). Dữ liệu thật duy nhất là FLAVIA gồm 1.907 ảnh lá cây kích thước 1200 x 1600, nhãn là độ lồi tính bằng tỷ số diện tích hình trên diện tích bao lồi, xử lý bằng phức lập phương 30 x 30 với 8 trong 9 đường tubular đi qua gốc của chiếc lá.\par\vspace{3pt}
\noindent \textbf{Giao thức đánh giá.} Vì là dữ liệu tổng hợp nên không có khái niệm chủ thể hay bệnh nhân và không có nguy cơ rò rỉ dữ liệu kiểu tách theo chủ thể; họ chia ngẫu nhiên (80\% đám mây để huấn luyện ở bài lỗ hổng, 400 huấn luyện và 80 kiểm tra cho mỗi kịch bản ở bài tính lồi, 70\% ảnh ở FLAVIA) và điều đó hợp lệ với bài toán của họ. Đơn vị phân tích là một đám mây điểm hoặc một ảnh; họ kiểm tra độ vững qua tịnh tiến, xoay [-20, 20] độ, kéo giãn hệ số [0,8–1,2], trượt [-0,2–0,2], nhiễu Gauss và điểm ngoại lai (thay thế 0–10\% số điểm), cùng bốn kịch bản ngoài phân phối cho tính lồi. Mỗi thí nghiệm lặp 3 lần và báo cáo trung bình kèm độ lệch chuẩn, trong đó độ lệch chuẩn của PH luôn bằng 0,00 vì SVM tất định; bài không nêu kiểm định thống kê nào và không có khái niệm dung sai ms vì đây không phải bài toán dò đỉnh.\par\vspace{3pt}
\noindent \textbf{Kết quả chính.} Đếm số lỗ hổng, độ chính xác trên dữ liệu gốc: PH 1,00; PH simple 0,94; ML 0,67; NN nông 0,53; NN sâu 0,50; PointNet 1,00 — nhưng khi tịnh tiến, PointNet sụp còn 0,20 trong khi PH giữ 1,00, khi xoay PointNet còn 0,74, và với điểm ngoại lai PH đạt 0,93 so với PointNet 0,55. Ước lượng độ cong theo MSE: PH bậc 0 rút gọn 0,06; PH bậc 0 đầy đủ 0,08; PH bậc 1 0,18; ML 0,34; NN nông 0,50; NN sâu 0,45; PointNet 196,94 ± 330,31 với ba lần chạy lần lượt 12,28 / 0,25 / 578,28, tức cực kỳ không ổn định; đáng chú ý là chỉ lấy 10 khoảng dài nhất làm MSE tăng từ 0,06 lên 0,21, nghĩa là nhiều khoảng ngắn mới mang thông tin độ cong. Tính lồi: PH đạt 1,00 ở kịch bản quy chuẩn, 0,85 ở ngẫu nhiên, 0,78 và 0,96 ở hai kịch bản chéo, so với PointNet 1,00 / 0,57 / 0,52 / 0,54 và ML 0,74 / 0,56 / 0,59 / 0,54; trên FLAVIA hồi quy tuyến tính trên vector 9 chiều đạt MSE 0,00065; định lý 1 khẳng định X là lồi khi và chỉ khi với mọi đường thẳng trong R\textasciicircum{}d, PD bậc 0 trên lọc tubular chứa đúng một khoảng.\par\vspace{3pt}
\noindent \textbf{Điểm yếu.} Toàn bộ thí nghiệm chính là dữ liệu tổng hợp trong R\textasciicircum{}2 và R\textasciicircum{}3, dữ liệu thật duy nhất là ảnh lá cây FLAVIA; không có chuỗi thời gian, không có tín hiệu sinh lý, không có dữ liệu 1 chiều. Các đường cơ sở bị làm yếu một cách hệ thống: MLP chỉ huấn luyện 2 epoch cho mỗi cấu hình, ML và NN chỉ nhận ma trận khoảng cách của 100 điểm lấy mẫu con trong khi PH nhìn toàn bộ 1.000 hoặc 5.000 điểm, và PointNet không được tăng cường dữ liệu dù bài gốc của nó có dùng; chính tác giả thừa nhận họ không nhằm chứng minh PH vượt trội so với state-of-the-art. Kết quả tính lồi phụ thuộc mạnh vào việc chọn thủ công 9 đường thẳng tubular (các đám mây bị phân loại sai đều có chỗ lõm nằm ngoài 9 đường đó), mỗi cấu hình chỉ chạy 3 lần và không có kiểm định thống kê, quy mô bài toán rất nhỏ (tối đa 9 lỗ, 20 hình dạng), còn chi phí tính toán chỉ được vẽ biểu đồ mà không có con số cụ thể trong văn bản.\par\vspace{3pt}
\noindent \textbf{\color{tdtblue}So với mô hình của nhóm.} Không so sánh trực tiếp được về con số vì khác hoàn toàn bài toán (phân loại hình dạng từ đám mây điểm, không có ECG, không có đỉnh, không có dung sai ms, không có bệnh nhân). Nhưng đây là căn cứ quan trọng nhất để biện minh cho quyết định chuyển hướng của nhóm: bộ hướng dẫn của họ liệt kê đúng ba loại tín hiệu mà PH làm tốt — số chu trình k chiều (lọc alpha, khoảng dài nhất), độ cong (Vietoris-Rips, nhiều khoảng ngắn) và tính lồi (phức lập phương, tubular, khoảng bậc hai) — còn bài toán xác định vị trí mẫu của phức bộ QRS thai nhi trong tín hiệu 1 chiều với dung sai 50 ms không thuộc bất kỳ loại nào trong ba loại đó. Vì Persistence Diagram là một đa tập điểm không có thứ tự nên về nguyên tắc không mã hóa được vị trí thời gian tuyệt đối, do đó phát hiện của nhóm (PI + CNN 2D đạt 27,42 so với 97,14 của CNN 1D) hoàn toàn nhất quán với khung lý thuyết của bài này chứ không phải một lỗi cài đặt.\par\vspace{3pt}
\noindent \textbf{\color{tdtblue}Nhóm rút ra gì.} Trích bộ hướng dẫn và phần kết luận của Turkes và cs. làm căn cứ học thuật chính thức cho luận điểm "TDA không phù hợp với bài toán dò vị trí đỉnh", đồng thời áp dụng đúng khung ba bước đó cho nhánh fSQI: xác định tín hiệu (có tồn tại vòng lặp tuần hoàn ứng với nhịp tim thai hay không), chọn phép lọc (nhúng trễ để biến tính chu kỳ thành vòng lặp H1), chọn đặc trưng (max persistence hoặc tổng life-time của H1), và dùng lọc DTM với m = 0,03 nếu nhánh này phải chịu nhiễu và điểm ngoại lai — đây cũng là cơ sở cho yêu cầu cửa sổ phải dài hơn một chu kỳ tim thai (khoảng 430 ms) thì H1 mới có thể có vòng lặp. Ngược lại phải tránh hai lỗi của họ là làm yếu đường cơ sở (huấn luyện 2 epoch, lấy mẫu con 100 điểm) trong khi nhóm đã so 16 kiến trúc cùng ngân sách tham số, và báo cáo chi phí tính toán chỉ bằng biểu đồ thay vì con số (nhóm đã có 113.481 tham số, 0,48 MB, 4,35 ms).\par\vspace{3pt}
\vspace{6pt}
\clearpage
\section{Nhúng trễ và chọn tham số}
Hai công trình về cách chọn tham số trễ. Quan trọng vì chúng đi trước phần đóng góp mà bản đề cương đầu tiên định nhận.

\subsection{Trần và cs. 2019, arXiv \normalfont\small (p27)}
\noindent{\footnotesize\itshape Quoc Hoan Tran, Yoshihiko Hasegawa (Department of Information and Communication Engineering, Graduate School of Information Science and Technology, The University of Tokyo). "Topological time-series analysis with delay-variant embedding". arXiv:1803.00208v4 [physics.data-an], ban v4 ghi ngay 10-11 thang 1 nam 2019. Ma nguon: https://github.com/OminiaVincit/delay-variant-embed. FILE PDF KHONG GHI ...}\par\vspace{2pt}
\noindent{\small\color{tdtblue}\textbf{Kết luận: so được một phần}}\par\vspace{4pt}
\noindent \textbf{Họ làm gì.} Đây là bài TDA cho chuỗi thời gian 1 chiều, gần với miền dữ liệu của nhóm nhất trong ba bài: để dùng TDA cho chuỗi thời gian phải nhúng chuỗi vào không gian nhiều chiều bằng tọa độ trễ, nhưng việc chọn độ trễ tau là bài toán không có lời giải lý thuyết, mọi cách chọn đều theo kinh nghiệm và giá trị tốt cho ứng dụng này lại tệ cho ứng dụng khác trên cùng dữ liệu. Giải pháp của họ là không cố định tau mà coi tau là một biến: tính persistence diagram cho nhiều giá trị tau rồi xếp chồng thành một persistence diagram ba chiều (b, d, tau), chứng minh đặc trưng này ổn định với nhiễu, định nghĩa hàm hạt nhân xác định dương trên không gian đó rồi đưa vào SVM để phân loại chuỗi thời gian.\par\vspace{3pt}
\noindent \textbf{Phương pháp.} Với mỗi tau trong tập K giá trị, chuỗi được nhúng vào không gian m chiều bằng tọa độ trễ, dựng phức Vietoris-Rips bằng cách tăng dần bán kính epsilon để được filtration, rồi trích PD bậc 0 (thành phần liên thông) và bậc 1 (vòng lặp), sau đó xếp chồng thành PD ba chiều PD3 = \{(b, d, tau)\}. Hạt nhân được định nghĩa bằng tổng các cặp Gaussian trừ đi ảnh đối xứng qua mặt phẳng chéo, với d\textasciicircum{}2 = |b1 − b2|\textasciicircum{}2 + |d1 − d2|\textasciicircum{}2 + xi\textasciicircum{}2·|tau1 − tau2|\textasciicircum{}2, được chứng minh xác định dương bằng cách viết thành tích vô hướng trong L2, đặt xi = sigma và sigma\textasciicircum{}2 = 0,5 nhân trung vị của các sigma\_s\textasciicircum{}2, rồi chuẩn hóa. Ma trận Gram kết hợp hai bậc theo K = alpha\_0·K0 + alpha\_1·K1 với alpha\_0 + alpha\_1 = 1 và alpha\_0 chọn bằng kiểm định chéo từ 14 giá trị trong khoảng 0 đến 1, cuối cùng phân loại bằng SVM trong không gian hạt nhân.\par\vspace{3pt}
\noindent \textbf{Mô hình.} Không phải mạng neural mà là SVM trên hạt nhân PD ba chiều, nên không có tham số học được theo nghĩa mạng, và bài không nêu số tham số cũng không nêu kích thước mô hình. Siêu tham số gồm K = 10 giá trị độ trễ tau = 1, 2, ..., 10 (riêng thí nghiệm điều chế tần số dùng K = 20), chiều nhúng m thuộc 2–10 chọn bằng kiểm định chéo, và chiều nhúng hiệu dụng M = m x K (ví dụ m = 6 và K = 20 cho M = 120).\par\vspace{3pt}
\noindent \textbf{Dữ liệu.} Dữ liệu tổng hợp Hes1 mô phỏng bộ dao động gen tự điều hòa âm có trễ bằng thuật toán Gillespie, với 1.000 tế bào cho mỗi chế độ (dao động và không dao động), đo mỗi nu = 64, 32, 16 hoặc 8 phút trong 4.096 phút, chuẩn hóa về trung bình 0 phương sai 1 và thêm nhiễu trắng phương sai 0,1. Tám bộ dữ liệu thật từ kho UCR: ECG200 (200 chuỗi dài 96), ECG5000 (5.000 chuỗi dài 140), ECGFiveDays (884 chuỗi dài 136, hai ngày của cùng một bệnh nhân), TwoLeadECG (1.162 chuỗi dài 82), NonInvasiveFetalECGThorax1 và Thorax2 (mỗi bộ 3.765 chuỗi dài 750, 42 lớp, ECG thai nhi không xâm lấn), EigenWorms (259 chuỗi dài 900) và FordB (4.446 chuỗi dài 500).\par\vspace{3pt}
\noindent \textbf{Giao thức đánh giá.} Với Hes1, họ chia ngẫu nhiên 50/50, lặp 100 lần chia và báo cáo độ chính xác trung bình — hợp lệ vì 1.000 tế bào là các mô phỏng độc lập; với tám bộ thật, họ dùng đúng chia train/test có sẵn của kho UCR, tức không phải tách theo chủ thể cũng không phải tách theo bản ghi, và đơn vị phân tích là cả chuỗi. Cần lưu ý ECGFiveDays là bản ghi của cùng một bệnh nhân ở hai ngày khác nhau nên bản chất là nhận dạng phiên ghi, và kết quả 99,9\% trên bộ đó không thể hiểu là khả năng tổng quát hóa lâm sàng. Lỗi giao thức rõ nhất là đường cơ sở single-delay được chọn tau bằng cách lấy độ chính xác lớn nhất trên tập kiểm tra, dù chính tác giả thừa nhận đáng lẽ phải dùng kiểm định chéo — may là lỗi này làm lợi cho đường cơ sở nên kết luận vẫn đứng vững; bài không nêu dung sai ms, không nêu tần số lấy mẫu của các bộ ECG, không nêu bộ lọc dải thông và không nêu kiểm định thống kê nào cho bảng kết quả chính.\par\vspace{3pt}
\noindent \textbf{Kết quả chính.} Độ chính xác phân loại (\%) của delay-variant so với single-delay: ECG200 90,0 so với 87,0 (1NN Euclid và 1NN DTW cùng 88,0); ECG5000 93,6 so với 92,1; ECGFiveDays 99,9 so với 92,0; TwoLeadECG 99,4 so với 94,0; EigenWorms 83,1 so với 83,1; FordB 90,8 so với 78,2. Trên hai bộ ECG thai nhi, Thorax1 đạt 91,8 (tăng 13,6 điểm so với single-delay 78,2 và 65,9 điểm so với Learned Shapelets 25,9 vốn sụp vì có 42 lớp) và Thorax2 đạt 93,0 so với 83,6; delay-variant tốt nhất ở 4 trong 8 bộ, đứng nhì ở 4 bộ còn lại và tốt nhất theo trung bình. Trên Hes1 chỉ có biểu đồ chứ không có bảng số: delay-variant luôn không kém single-delay ở mọi chiều nhúng, và khi so theo chiều nhúng hiệu dụng M thì độ chính xác gần như không đổi trên dải M từ 40 đến 100 trong khi single-delay suy giảm vì điểm nhúng bị gấp và thưa thớt quá mức; giới hạn ổn định lý thuyết là d\_B ≤ 2·√m·max|x(t) − y(t)|; bài không báo cáo thời gian tính, số tham số, kích thước mô hình hay độ trễ.\par\vspace{3pt}
\noindent \textbf{Điểm yếu.} Bài không báo cáo chi phí tính toán dưới bất kỳ dạng nào — không có số giây, số tham số hay dung lượng bộ nhớ — dù phải tính 10 filtration Vietoris-Rips cho mỗi chuỗi, chắc chắn rất chậm nhưng không biết chậm bao nhiêu. Các đường cơ sở 1NN, EE và LS không được chạy lại mà lấy số từ tài liệu tham khảo, trong khi họ có hạ mẫu với hệ số 2, 2, 3 và 2 cho bốn bộ mà không giải thích lý do, nên so sánh không hoàn toàn đồng nhất; thêm vào đó tau của single-delay được chọn bằng giá trị lớn nhất trên tập kiểm tra. Không có kiểm định thống kê nên chênh lệch 90,0 so với 88,0 trên 200 chuỗi hoàn toàn có thể là nhiễu ngẫu nhiên, không nêu tần số lấy mẫu nên không quy được độ dài chuỗi ra giây, và toàn bộ thí nghiệm là phân loại cả chuỗi chứ không có bài toán định vị nào.\par\vspace{3pt}
\noindent \textbf{\color{tdtblue}So với mô hình của nhóm.} Chỉ so sánh được một phần: họ có dùng hai bộ ECG thai nhi (NonInvasiveFetalECGThorax1 và Thorax2, mỗi bộ 3.765 chuỗi dài 750 mẫu) và đạt 91,8\% và 93,0\%, nhưng đó là phân loại 42 lớp hình thái do chuyên gia gán nhãn trên toàn bộ đoạn chuỗi chứ không phải dò vị trí phức bộ QRS thai nhi, không có dung sai ms, không có Se/PPV theo nhịp, không có jitter và không có khái niệm khử ECG mẹ, nên 91,8\% của họ và 99,18\% macro F1 của FetalQRS-TCN không thể đặt cạnh nhau. Về giao thức, họ dùng chia train/test có sẵn của UCR với ít nhất một bộ là cùng một bệnh nhân ở hai ngày, yếu hơn cách tách theo bản ghi và tách theo chủ thể xoay vòng của nhóm; về chi phí tính toán thì nhóm hơn hẳn vì họ không báo cáo gì trong khi nhóm có 113.481 tham số, 0,48 MB và 4,35 ms cho một cửa sổ 4 giây. Về giá trị khoa học, đây là minh chứng tích cực cho hướng fSQI mới: TDA trên tín hiệu 1 chiều chỉ phát huy tác dụng khi đi qua nhúng trễ để biến tính chu kỳ thành vòng lặp H1 và khi mục tiêu là phân loại trạng thái động lực học của cả đoạn chứ không phải định vị sự kiện, khớp với phát hiện của nhóm rằng cửa sổ 300 ms ngắn hơn một chu kỳ tim thai (khoảng 430 ms) nên không thể chứa vòng lặp.\par\vspace{3pt}
\noindent \textbf{\color{tdtblue}Nhóm rút ra gì.} Áp dụng trực tiếp cho fSQI: thay vì chọn một độ trễ, hãy quét tau qua một dải rồi xếp chồng thành PD ba chiều — với ECG bụng 250 Hz và nhịp tim thai 120–160 lần/phút (chu kỳ khoảng 375–500 ms, tương ứng 94–125 mẫu), dải tau nên phủ từ khoảng 1/10 đến 1/2 chu kỳ này, cửa sổ phải bao ít nhất 2 chu kỳ (khớp với cửa sổ 4 giây đang dùng), và nên dùng kết hợp hạt nhân K = alpha\_0·K0 + alpha\_1·K1 với alpha\_0 chọn bằng kiểm định chéo để dữ liệu tự quyết định phần đóng góp của H0 (tức peak prominence) so với H1 (tức tính chu kỳ). Đồng thời trích bất đẳng thức ổn định d\_B ≤ 2·√m·max|x − y| làm cơ sở lý thuyết cho tuyên bố fSQI bền với nhiễu, và có thể thêm Thorax1/Thorax2 làm bộ dữ liệu phụ với mốc so sánh 91,8\% / 93,0\% (ghi rõ đó là phân loại chuỗi, không phải dò đỉnh), nhưng phải tránh hai lỗi của họ là chọn siêu tham số trên tập kiểm tra và công bố bảng kết quả không kèm kiểm định thống kê.\par\vspace{3pt}
\vspace{6pt}
\subsection{Tan và cs. 2023, arXiv:2302.03447 \normalfont\small (p28)}
\noindent{\footnotesize\itshape Eugene Tan, Shannon Algar, Débora Corrêa, Michael Small, Thomas Stemler, David Walker. "Selecting embedding delays: An overview of embedding techniques and a new method using persistent homology." arXiv:2302.03447v1 [math.DS], nộp 7-2-2023, trang đầu ghi ngày 8-2-2023. Complex Systems Group, Department of Mathematics and Statistics, The University of Western Australia. LƯU Ý: file PDF trong repo ...}\par\vspace{2pt}
\noindent{\small\color{warnred}\textbf{Kết luận: không so trực tiếp được}}\par\vspace{4pt}
\noindent \textbf{Họ làm gì.} Đây là bài tổng quan kèm đề xuất phương pháp thuộc lĩnh vực hệ động lực phi tuyến, không phải công trình y sinh: khi chỉ quan sát được một chuỗi thời gian vô hướng, định lý Takens cho phép dựng lại không gian pha bằng nhúng trễ, nhưng chọn sai độ trễ τ thì bản dựng lại vô nghĩa, trong khi hàng chục phương pháp chọn τ hiện có (thông tin tương hỗ, tự tương quan, fill-factor, khuếch đại nhiễu, L-statistic, PECUZAL, MDOP) thường mâu thuẫn nhau và không giải thích được vì sao một độ trễ là tốt. Nhóm tác giả đề xuất SToPS (Significant Times on Persistent Strands), dùng đồng điều bền vững để chấm điểm độ quan trọng động lực của từng độ trễ τ, tạo ra phổ thời gian đặc trưng S(τ) mà mỗi đỉnh đều giải thích được theo chu kỳ thật của tín hiệu.\par\vspace{3pt}
\noindent \textbf{Phương pháp.} Với mỗi ứng viên độ trễ τ, chuỗi được chuẩn hóa về trung bình 0 và độ lệch chuẩn 1 rồi dựng nhúng trễ 2 chiều; sau đó lấy mẫu N = 250 đoạn quỹ đạo ngắn (strand) có độ dài l = 4τ theo heuristic một phần tư chu kỳ, vị trí đầu chọn bằng bộ lấy mẫu ngẫu nhiên k-means++ để phủ đều không gian pha. Mỗi strand được tính đồng điều bền vững bằng lọc Vietoris-Rips và chỉ được coi là có vòng (lỗ H1) nếu độ bền lớn nhất vượt ngưỡng ρ bằng khoảng cách trung bình giữa hai điểm liên tiếp trong không gian pha, đồng thời biên lỗ phải có ít nhất Nhole = 8 điểm nhằm loại các lỗ giả do đoạn quá ngắn không khép kín được. Điểm của mỗi strand là Si(τ) = αi(τ)·γi(τ) với α là độ tròn (tỉ số trị riêng λ2/λ1 của PCA trên các điểm biên lỗ) và γ là hiệu suất (diện tích lỗ chia diện tích đa giác lồi nhỏ nhất bao strand), lấy trung bình và độ lệch chuẩn qua 250 strand để được phổ Sμ(τ) và Sσ(τ); các đỉnh của phổ được chọn bằng mắt (nguyên văn "the peaks of each profile are selected by observation"), rồi hai độ trễ trội nhất được dùng dựng nhúng 3 chiều để huấn luyện mạng dự báo một bước.\par\vspace{3pt}
\noindent \textbf{Mô hình.} SToPS bản thân không phải mô hình học máy mà là thuật toán chọn tham số, không có tham số huấn luyện; bộ dự báo dùng để kiểm chứng là mạng truyền thẳng 4 lớp với kiến trúc 3:128:128:3, kích hoạt ReLU, learning rate 0,001, batch size 512, 30 epoch. Bài không nêu số tham số (theo suy luận của người tổng hợp, khoảng 17.411 tham số), không nêu bộ tối ưu, không nêu hàm mất mát và không nêu thư viện tính đồng điều bền vững, chỉ nêu dùng gói DynamicalSystems.jl cho PECUZAL và MDOP.\par\vspace{3pt}
\noindent \textbf{Dữ liệu.} Ba lớp chuỗi thời gian, hoàn toàn không có ECG, không có dữ liệu y sinh và không có bệnh nhân: chuỗi nhân tạo tổng ba sóng sin (ω = \{1, 5, 30\}, φ = \{0, 0.25, 0.75\}, A = \{1, 0.5, 0.2\}, dt = 0,001, ba thời gian đặc trưng lý thuyết τ = (250, 50, 8) mẫu); thành phần x của hệ Lorenz (Runge-Kutta bậc 4 bước 0,0004, lấy mẫu lại về dt = 0,004, chạy 25.000 bước); và điện thế nơ-ron LP trong hạch somatogastric của tôm hùm với hai thang thời gian, độ trễ kỳ vọng khoảng τ = (12, 400). Mọi chuỗi bị giới hạn 25.000 điểm khi tính độ trễ, riêng chuỗi tôm hùm khi huấn luyện dự báo được lấy mẫu thưa mỗi 3 điểm còn khoảng 200.000 bước; bài không nêu tần số lấy mẫu vật lý (Hz) của dữ liệu tôm hùm nên mọi độ trễ chỉ tính bằng số mẫu, không quy ra ms.\par\vspace{3pt}
\noindent \textbf{Giao thức đánh giá.} Không có khái niệm tách theo chủ thể hay tách theo bản ghi vì không có đối tượng, chỉ có ba chuỗi tín hiệu; giao thức là dùng nửa đầu chuỗi để huấn luyện và nửa sau để kiểm định, đơn vị phân tích là chuỗi thời gian. Kiểm định bằng dự báo freerun n bước (Lorenz 25 bước trên 25.000 bước; tôm hùm 10 bước trên tập đã lấy mẫu thưa, tương đương 30 bước dữ liệu gốc) với lỗi E(t) = ||NN\textasciicircum{}(n)(x(t)) − x(t+n)||, báo cáo dạng phân phối cùng trung bình và trung vị; bài không nêu dung sai khớp nhịp vì đây không phải bài toán dò đỉnh. Bài không chạy nhiều seed ngẫu nhiên, không có kiểm định thống kê (không Wilcoxon, không t-test, không khoảng tin cậy), và việc chọn đỉnh phổ là thủ công bằng mắt — chính tác giả thừa nhận điều này "đưa vào một mức chủ quan trong diễn giải phổ thời gian đặc trưng".\par\vspace{3pt}
\noindent \textbf{Kết quả chính.} Về phát hiện thang thời gian: với chuỗi tổng sóng sin (kỳ vọng τ = 250, 50, 8), thông tin tương hỗ chỉ tìm được thành phần tần số cao nhất, PECUZAL chỉ ra một độ trễ sai τ = 341 và MDOP ra τ = (241, 42, 271, 180) với chỉ hai giá trị gần đúng, còn SToPS bắt được cả ba thang; ở Lorenz, SToPS ra (41, 130) trong đó τ = 130 là artefact gây quá-nhúng do chính tác giả thừa nhận, còn với nơ-ron tôm hùm thì PECUZAL và MDOP thất bại hoàn toàn ở thang nhanh trong khi SToPS lấy được cả hai (τ ≈ (5, 500)) dù có đỉnh phụ giả ở τ = (40, 60). Về lỗi dự báo freerun, trung vị trên Lorenz 25 bước là SToPS 0,106 | PECUZAL 0,122 | MDOP 0,112, nhưng bài ghi y hệt ba con số đó cho tôm hùm nên gần như chắc chắn là lỗi sao chép và số liệu này không đáng tin; chính tác giả viết cải thiện "không có ý nghĩa và không nên coi là lợi ích nhắm tới của SToPS". Về chi phí tính toán (Bảng I, Ryzen 7 4800HS, 16 GB RAM, một luồng, chuỗi 25.000 bước), ở Lorenz với τmax = 200 SToPS tốn 2.396 giây so với MDOP 0,42 giây (theo suy luận, chậm hơn khoảng 5.700 lần) và ở tôm hùm với τmax = 600 tốn 7.266 giây; phổ SToPS giữ nguyên vị trí đỉnh ở mọi mức nhiễu thử, nhưng dưới SNR 17 dB cả SToPS lẫn PECUZAL đều không cho ra kết quả nào, chỉ MDOP còn chạy được.\par\vspace{3pt}
\noindent \textbf{Điểm yếu.} Bài hoàn toàn không có dữ liệu ECG, không có bệnh nhân và không có chỉ số y sinh nên mọi con số không chuyển thẳng sang bài toán fQRS; việc chọn đỉnh phổ làm thủ công bằng mắt và chưa được tự động hóa. Bộ trung vị lỗi dự báo của tôm hùm trùng khít với bộ của Lorenz (0,106 | 0,122 | 0,112) khiến phần kết quả định lượng chính không kiểm chứng được, và phần khảo sát nhiễu mâu thuẫn nội bộ (phần chữ nói năm mức nhưng liệt kê bốn, còn Hình 14 hiện sáu mức 40/37/30/27/20/17 dB). Không có kiểm định thống kê, không lặp seed, không khoảng tin cậy nên khẳng định SToPS tốt hơn không có cơ sở thống kê; chi phí lên tới 7.266 giây cho một chuỗi 25.000 điểm nên không dùng được thời gian thực, cài đặt chưa song song hóa, không công bố mã nguồn trong bản PDF, chưa thử trên chuỗi đa biến, và SToPS sụp đổ khi SNR dưới 17 dB.\par\vspace{3pt}
\noindent \textbf{\color{tdtblue}So với mô hình của nhóm.} Không so sánh được về số liệu với mô hình của nhóm vì bài không làm dò đỉnh, không có ECG, không có F1/Se/PPV, không có dung sai theo ms và không có chủ thể để tách theo bản ghi, nhưng lại rất có giá trị làm cơ sở lý thuyết cho quyết định chuyển hướng. Bài chứng minh bằng chính cấu trúc thuật toán rằng muốn có lỗ H1 trong nhúng trễ 2 chiều thì đoạn tín hiệu phải dài ít nhất l = 4τ, tức khoảng một chu kỳ đầy đủ, và còn phải thêm ràng buộc Nhole ≥ 8 điểm biên để loại lỗ giả sinh ra khi đoạn quá ngắn không khép kín được — đúng hiện tượng nhóm đã quan sát khi cửa sổ 300 ms ngắn hơn chu kỳ tim thai khoảng 430 ms. Chi phí 2.396 giây cho một chuỗi 25.000 điểm với τmax = 200 trên một luồng CPU so với 4,35 ms cho một cửa sổ 4 giây của FetalQRS-TCN, cùng với việc SToPS mất tác dụng dưới SNR 17 dB, biện minh trực tiếp cho việc không đưa đồng điều bền vững nặng vào đường dẫn thời gian thực trên ECG bụng đơn kênh vốn có tỉ số tín hiệu thai trên tạp rất thấp.\par\vspace{3pt}
\noindent \textbf{\color{tdtblue}Nhóm rút ra gì.} Nên trích dẫn ràng buộc độ dài strand l = 4τ và ràng buộc Nhole ≥ 8 làm căn cứ học thuật cho lập luận "cửa sổ ngắn hơn một chu kỳ tim thì không thể chứa vòng nhịp" khi giải thích vì sao Persistence Image kết hợp CNN 2D đứng chót bảng (27,42 so với 97,14), đồng thời dùng Bảng I về thời gian tính toán làm dẫn chứng số thật cho phần chi phí; theo heuristic một phần tư chu kỳ, chu kỳ tim thai khoảng 430 ms cho τ khoảng 107 ms, tức khoảng 27 mẫu ở 250 Hz và cửa sổ tối thiểu 4τ khoảng 108 mẫu (đây là suy luận, bài không tính cho ECG). Nên tránh cách chọn tham số bằng mắt và lỗi lặp số liệu giữa hai thí nghiệm như bài này: nhóm phải tự động hóa, kiểm tra chéo bảng kết quả trước khi nộp và báo cáo kèm kiểm định thống kê (nhóm đã có Wilcoxon p = 0,0000 cho quét trường tiếp nhận), đồng thời không đưa Vietoris-Rips hay đặc trưng H1 vào bất kỳ nhánh chạy thời gian thực nào.\par\vspace{3pt}
\vspace{6pt}
\clearpage
\section{Topo cho tín hiệu tim}
Các công trình đã áp dụng topo cho tín hiệu tim người lớn. Đây là chỗ nhóm tìm được tiền lệ cho đóng góp C3.

\subsection{Ren và cs. 2023, Frontiers in Neuroscience \normalfont\small (p29)}
\noindent{\footnotesize\itshape Yonglian Ren, Feifei Liu, Shengxiang Xia, Shuhua Shi, Lei Chen, Ziyu Wang. "Dynamic ECG signal quality evaluation based on persistent homology and GoogLeNet method." Frontiers in Neuroscience, tập 17, số bài 1153386, 2023. DOI: 10.3389/fnins.2023.1153386. Nhận 29-1-2023, chấp nhận 20-2-2023, xuất bản 8-3-2023. Đơn vị: School of Science, Shandong Jianzhu University, Tế Nam, Trung Quốc. Giấy phép ...}\par\vspace{2pt}
\noindent{\small\color{tdtblue}\textbf{Kết luận: so được một phần}}\par\vspace{4pt}
\noindent \textbf{Họ làm gì.} Bài giải bài toán đánh giá chất lượng tín hiệu ECG cho thiết bị đeo, tức phân loại nhị phân một đoạn ECG là "chấp nhận được" hay "không chấp nhận được" — chính là bài toán mà nhóm vừa chuyển hướng sang với fSQI có quyền từ chối trả lời, nên đây là bài gần với hướng mới nhất trong ba bài. Thay vì trích đặc trưng dạng sóng và miền thời gian - tần số như các chỉ số chất lượng tín hiệu truyền thống, nhóm tác giả dùng đồng điều bền vững để bắt thông tin tô-pô, vẽ kết quả ra ảnh (persistence barcode hoặc persistence diagram) rồi đưa vào GoogLeNet đã tiền huấn luyện, so sánh hai kiểu lọc là Vietoris-Rips và SubLevel-Set.\par\vspace{3pt}
\noindent \textbf{Phương pháp.} Mỗi đoạn ECG được chuẩn hóa về khoảng (−1, 1) bằng hàm mapminmax của MATLAB, và bài không áp dụng bộ lọc dải thông riêng nào. Nhánh Vietoris-Rips dựng đám mây điểm bằng cửa sổ trượt theo công thức (3) của bài (cắt chuỗi thành ma trận d hàng, mỗi hàng t phần tử liên tiếp, ràng buộc d·t ≤ n) rồi tăng dần bán kính ε để ghi lại lúc sinh và lúc chết của các lỗ, xuất barcode 0 chiều và 1 chiều; nhánh SubLevel-Set coi ECG là hàm liên tục và quét ngưỡng α từ âm vô cùng lên dương vô cùng để xuất persistence diagram. Bước mấu chốt và cũng là điểm yếu là barcode hoặc diagram không được véc-tơ hóa thành đặc trưng số mà được vẽ thành ảnh RGB 224 × 224 × 3 rồi đưa vào GoogLeNet bằng học chuyển tiếp; siêu tham số duy nhất được quét là độ dài cửa sổ trượt (bộ 12 chuyển đạo thử 1, 2, 3, 4, 5 giây với số chiều 120, 60, 36, 24, 24; bộ đơn chuyển đạo thử 0,1; 0,2; 0,3; 0,4; 0,5 giây với số chiều 100, 50, 33, 25, 20), đánh giá bằng kiểm định chéo 10 mảnh.\par\vspace{3pt}
\noindent \textbf{Mô hình.} GoogLeNet (Inception v1) dùng học chuyển tiếp; bài không nêu số tham số. Bài cũng không nêu số epoch, learning rate, batch size, bộ tối ưu, cách thay lớp đầu ra, số lớp đóng băng, thư viện tính đồng điều bền vững, phần cứng, thời gian huấn luyện, độ trễ suy luận hay kích thước mô hình, và về chi phí tính toán chỉ nói định tính rằng số chiều đám mây điểm quá lớn sẽ gây tốn kém — đây là công trình không tái lập được.\par\vspace{3pt}
\noindent \textbf{Dữ liệu.} Nguồn dữ liệu là PhysioNet/CinC Challenge 2011 với 1.000 bản ghi ECG 12 chuyển đạo, mỗi bản 10 giây, ghi đồng thời, tần số lấy mẫu 500 Hz, độ phân giải 16 bit, băng thông chẩn đoán 0,05–100 Hz, nhãn gốc gồm 773 "acceptable", 225 "unacceptable" và 2 "indeterminate"; bộ 12 chuyển đạo dùng 998 tín hiệu sau khi bỏ 2 tín hiệu indeterminate. Bộ đơn chuyển đạo do nhóm tự gán nhãn lại 12.000 tín hiệu theo thang 5 mức của Silva 2011 và Liu 2018 (9.941 "acceptable", 2.059 "unacceptable"), nhưng sau khi dò chuyển đạo rỗng loại đi 1.071 đoạn thì chỉ còn 988 đoạn xấu thật, nên họ cộng thêm nhiễu từ MIT-BIH Noise Stress Test Database (2.252 bản bw, 2.252 bản em, 2.252 bản ma) và 1.126 bản nhiễu Gauss ở −10 dB vào 7.882 tín hiệu "acceptable", thu được bộ cuối 9.941 "acceptable" và 8.870 "unacceptable" — theo suy luận, 7.882 trên 8.870 tức 88,9\% lớp âm là nhiễu nhân tạo, chỉ 11,1\% là tín hiệu xấu thật.\par\vspace{3pt}
\noindent \textbf{Giao thức đánh giá.} Kiểm định chéo 10 mảnh với nguyên văn "All the segments were randomly divided into 10 groups", tức chia ngẫu nhiên theo đoạn, không tách theo chủ thể và không tách theo bản ghi; đơn vị phân tích là đoạn 10 giây. Đây là rò rỉ dữ liệu theo hai cơ chế độc lập: 12 chuyển đạo của cùng một bản ghi được ghi đồng thời trên cùng một bệnh nhân nên chuyển đạo I của bệnh nhân X có thể nằm ở tập huấn luyện còn chuyển đạo V2 của chính bệnh nhân đó nằm ở tập kiểm tra; nghiêm trọng hơn, 7.882 mẫu "unacceptable" được tạo bằng cách cộng nhiễu vào chính các tín hiệu "acceptable" có trong bộ dữ liệu nên bản sạch và bản nhiễu của cùng một tín hiệu có thể rơi vào hai phía huấn luyện và kiểm tra. Bài không nêu dung sai khớp nhịp vì đây là bài toán phân loại chất lượng chứ không phải dò đỉnh QRS, đồng thời không nêu cách chia mảnh có phân tầng theo lớp hay không, không nêu số lần lặp lại và không có kiểm định thống kê khi so với đối thủ, chỉ nêu trung bình ± độ lệch chuẩn.\par\vspace{3pt}
\noindent \textbf{Kết quả chính.} Trên bộ 12 chuyển đạo, SubLevel-Set tốt nhất với mAcc 98,04 ± 0,11 | F1 98,40 ± 0,06 | Se 97,15 ± 0,11 | Sp 98,93 ± 0,23 (Acc 97,70), còn Vietoris-Rips giảm dần theo cửa sổ từ mAcc 95,16 ± 0,20 ở 1 giây xuống 93,11 ± 0,19; 91,75 ± 0,21; 90,22 ± 0,29 và 89,78 ± 0,26 ở 5 giây; so với nhóm đối chứng, phương pháp tốt nhất là Zhao và Zhang 2018 đạt mAcc 91,67 (Acc/Se/Sp 94,67 / 90,33 / 93,00). Trên bộ đơn chuyển đạo, Vietoris-Rips với cửa sổ 0,1 giây tốt nhất (mAcc 98,55 ± 0,13 | F1 98,62 ± 0,12 | Se 98,37 ± 0,18 | Sp 98,85 ± 0,20), giảm còn 89,76 ± 0,56 ở 0,5 giây trong khi SubLevel-Set đạt mAcc 97,25 ± 0,39, và các đối thủ chạy lại trên cùng bộ dữ liệu đạt picaSQI (Liu 2020) 93,11 ± 0,73; bSQI\_4 (Liu 2018) 89,28 ± 0,55; bSQI\_2 84,08 ± 0,75; MS-QI (Falk và Maier 2014) 85,61 ± 1,05. Theo suy luận tự tính từ hai bảng, trên bộ đơn chuyển đạo việc đổi độ dài cửa sổ từ 0,1 giây lên 0,5 giây làm mAcc rơi 8,79 điểm trong khi đổi kiểu lọc từ Vietoris-Rips sang SubLevel-Set chỉ đổi 1,30 điểm; trên bộ 12 chuyển đạo, quét cửa sổ đổi 5,38 điểm còn đổi kiểu lọc đổi 2,88 điểm.\par\vspace{3pt}
\noindent \textbf{Điểm yếu.} Rò rỉ dữ liệu nghiêm trọng ở hai tầng — chia ngẫu nhiên 10 mảnh không tách bệnh nhân hay bản ghi, cộng với việc lớp "unacceptable" phần lớn là bản sao có nhiễu của chính các tín hiệu "acceptable" trong cùng bộ dữ liệu — nên con số 98,55\% gần như chắc chắn lạc quan quá mức; bài toán còn bị làm dễ đi vì 88,9\% lớp âm là nhiễu tổng hợp cộng ở −10 dB, khiến mạng chỉ cần học có nhiễu bw/em/ma/Gauss hay không chứ không phải học chất lượng lâm sàng thật, với chỉ 988 mẫu xấu thật. Nhãn 12 chuyển đạo không nhất quán theo thừa nhận của chính tác giả, quy trình tự gán nhãn lại 12.000 đoạn không được mô tả (mấy người gán, có độc lập không, độ đồng thuận bao nhiêu), và tuyên bố "đơn chuyển đạo" chỉ đúng cho bộ 2 vì bộ 1 thực chất gộp cả 12 kênh vào một đám mây điểm, kèm mâu thuẫn nội bộ giữa phần chữ và bảng về số chiều ứng với cửa sổ 0,1 giây hay 1 giây cũng như giữa năm và sáu giá trị cửa sổ. Bài không tái lập được (không nêu siêu tham số huấn luyện, không nêu thư viện TDA, không công bố mã nguồn), không có kiểm định thống kê nào khi tuyên bố vượt đối thủ, không có phân tích theo mức nhiễu, không có đường cong ROC hay ngưỡng vận hành, và lãng phí thông tin khi vẽ barcode ra ảnh 224×224×3 rồi cho CNN 2D đọc.\par\vspace{3pt}
\noindent \textbf{\color{tdtblue}So với mô hình của nhóm.} So sánh được một phần và là bài tham chiếu quan trọng nhất cho hướng fSQI, nhưng không so được về con số với FetalQRS-TCN vì nhiệm vụ khác hẳn (họ phân loại chất lượng đoạn 10 giây, nhóm định vị từng đỉnh fQRS với dung sai ±50 ms), đối tượng khác (ECG người lớn, không có thai nhi và không có bài toán khử ECG mẹ), chỉ số khác (F1 phân loại nhị phân trên đoạn so với F1 dò nhịp có dung sai thời gian), và họ có rò rỉ dữ liệu trong khi nhóm dùng tách theo bản ghi nghiêm ngặt — đặt 98,55 của họ cạnh 97,43 ± 4,37 của nhóm là so sánh sai lệch có lợi cho họ. Họ kém hơn ở ba điểm có thể nêu rõ trong bài báo: giao thức chia ngẫu nhiên theo đoạn thay vì tách theo bản ghi; tính tái lập (nhóm có 113.481 tham số, checkpoint 0,48 MB, độ trễ 4,35 ms còn họ không nêu một siêu tham số nào); và cơ sở dữ liệu âm tính với 88,9\% là nhiễu tổng hợp. Họ hơn nhóm ở quy mô dữ liệu (998 bản ghi 12 chuyển đạo và 18.811 đoạn đơn chuyển đạo so với 5 sản phụ của ADFECGDB) và ở việc có bảng đối chiếu với 6 phương pháp chỉ số chất lượng tín hiệu kinh điển; điểm trùng khớp đáng chú ý là kết quả của họ củng cố phát hiện của nhóm rằng siêu tham số tiền xử lý quan trọng hơn lựa chọn thuật toán — với họ độ dài cửa sổ đổi 8,79 điểm còn kiểu lọc tô-pô chỉ đổi 1,30 điểm, với nhóm dải lọc đổi 11,00 điểm còn kiến trúc mạng chỉ đổi 0,41 điểm (p = 0,7012).\par\vspace{3pt}
\noindent \textbf{\color{tdtblue}Nhóm rút ra gì.} Nên trích dẫn bài này làm nền cho hướng fSQI vì nó cho thấy đồng điều bền vững có giá trị thật cho bài toán đánh giá chất lượng (mAcc 98,04 bằng SubLevel-Set) trong khi thất bại cho bài toán định vị đỉnh, đồng thời bắt chước cấu trúc bảng đối chiếu với các chỉ số chất lượng kinh điển (bSQI, MS-QI, picaSQI, kSQI, pSQI) và ghi nhận rằng SubLevel-Set — loại lọc nhóm đã chứng minh tương đương peak prominence — vẫn đủ mạnh (97,25 trên đơn chuyển đạo), mở đường triển khai fSQI bằng peak prominence rẻ tiền mà tương đương. Nên tránh tuyệt đối việc chia ngẫu nhiên 10 mảnh theo đoạn thay vì tách theo bản ghi, việc cân bằng lớp bằng cách cộng nhiễu vào chính các mẫu tốt rồi trộn ngẫu nhiên (mọi biến thể sinh ra từ cùng một bản ghi gốc phải nằm cùng một phía huấn luyện hoặc kiểm tra và phải báo cáo riêng kết quả trên tập xấu thật), việc vẽ barcode hay persistence diagram ra ảnh cho CNN 2D, và việc bỏ trống siêu tham số như epoch, learning rate, batch size, seed và phần cứng.\par\vspace{3pt}
\vspace{6pt}
\subsection{Dindin và cs. 2019, arXiv:1906.05795 \normalfont\small (p30)}
\noindent{\footnotesize\itshape Meryll Dindin, Yuhei Umeda, Frédéric Chazal. "Topological Data Analysis for Arrhythmia Detection through Modular Neural Networks." arXiv:1906.05795v1 [cs.LG], 13-6-2019. Đơn vị: École Centrale Paris (Dindin), Fujitsu Nhật Bản (Umeda), INRIA (Chazal). File PDF này KHÔNG ghi DOI và KHÔNG ghi tên hội nghị. LƯU Ý: trong danh mục tài liệu tham khảo của bài p29 (Ren 2023) có ghi bản hội nghị của đúng ...}\par\vspace{2pt}
\noindent{\small\color{tdtblue}\textbf{Kết luận: so được một phần}}\par\vspace{4pt}
\noindent \textbf{Họ làm gì.} Bài giải bài toán phát hiện và phân loại loạn nhịp tim từ ECG người lớn, với trọng tâm không phải điểm số mà là khả năng tổng quát hóa sang bệnh nhân mới, vì các mô hình ECG thường tụt điểm thảm hại trên bệnh nhân chưa từng thấy do khác biệt cá thể về nhịp nhanh, nhịp chậm và hình dạng sóng. Ý tưởng cốt lõi là dùng đường cong Betti từ đồng điều bền vững làm một kênh đầu vào riêng trong mạng nơ-ron đa kênh — đường cong Betti bất biến với co giãn trục thời gian và co giãn biên độ nên mô tả hình dạng nhịp tim độc lập với nhịp nhanh hay chậm của từng người — kết hợp với kênh tín hiệu thô, kênh đặc trưng thủ công, kênh FFT và một bộ tự mã hóa huấn luyện không giám sát trên nhịp bình thường để xử lý mất cân bằng lớp.\par\vspace{3pt}
\noindent \textbf{Phương pháp.} Tiền xử lý gồm lấy mẫu lại về 200 Hz, khử trôi đường nền bằng wavelet Daubechies, lọc FIR đóng vai trò dải thông với tần số cắt 0,05 Hz và 50 Hz cộng một bộ lọc Kalman, rồi co tín hiệu về [0, 1] và tịnh tiến cho trung bình gần 0; ECG được cắt thành các đoạn chồng lấn một phần gồm một số cố định nhịp liên tiếp (nhãn lấy từ đỉnh trung tâm) rồi nội suy chuẩn hóa độ dài, trong đó vị trí nhịp lấy từ chú thích có sẵn nên bài nói rõ "peak detection thus unnecessary". Đặc trưng gồm đồng điều bền vững của lọc sublevel-set và upper-level-set trên chuỗi 1 chiều (tính bằng thư viện GUDHI, độ phức tạp O(n log n)) chuyển thành đường cong Betti kèm định lý bất biến rằng đường cong Betti của f(t) và f(at) là như nhau với mọi a > 0; đặc trưng thủ công gồm biến đổi Fourier, quan hệ tuyến tính giữa các thành phần P, Q, R, S, T, các thống kê cực trị, trung bình, độ lệch chuẩn, kurtosis, skewness, entropy, số lần cắt trục và PCA rút về 10 thành phần; cùng bộ tự mã hóa 6 lớp ẩn (đầu vào 400 chiều, không gian ẩn 20 chiều, mất mát MSE) huấn luyện trên toàn bộ nhịp bình thường rồi đóng băng. Các kênh tích chập và kết nối đầy đủ được nối vào một mạng chung, huấn luyện với PReLU khởi tạo he-normal, dropout ủ giảm từ 0,5 xuống 0,0 sau 100 epoch và bộ tối ưu Adadelta learning rate ban đầu 1,0, theo chiến lược hai tầng: trước tiên phân biệt nhịp bình thường với bất thường, sau đó mới phân loại 13 lớp chỉ trên các nhịp loạn nhịp.\par\vspace{3pt}
\noindent \textbf{Mô hình.} Mạng nơ-ron mô-đun đa kênh gồm các kênh CNN 1 chiều và MLP nối vào một mạng kết nối đầy đủ chung, cộng một bộ tự mã hóa 6 lớp ẩn (đầu vào 400, không gian ẩn 20); bài không nêu số tham số. Bài cũng không nêu số lớp tích chập, kích thước kernel, số kênh của từng khối, kích thước mô hình, độ trễ suy luận, số epoch tổng, batch size hay thư viện học sâu và không công bố mã nguồn trong bản PDF này, chỉ nêu thời gian huấn luyện khoảng 10 giờ trên một GPU GeForce GTX không rõ model.\par\vspace{3pt}
\noindent \textbf{Dữ liệu.} Năm cơ sở dữ liệu PhysioNet MIT-BIH theo Bảng 1 (số bệnh nhân / số nhãn / số giờ): Arrhythmia 48 / 109.494 / 24; LongTerm 7 / 668.735 / 147; NormalSinus 18 / 1.729.629 / 437; SupraVentricular 78 / 182.165 / 38; TWave 90 / 790.565 / 180, tất cả lấy mẫu 360 Hz, độ phân giải 11 bit trên dải 10 mV, mô tả là ECG đơn chuyển đạo và mỗi nhịp được ít nhất hai bác sĩ tim mạch chú thích độc lập rồi giải quyết bất đồng; sau tiền xử lý mọi tín hiệu được lấy mẫu lại về 200 Hz. Theo suy luận tự cộng, tổng cộng khoảng 3.480.588 nhịp, 826 giờ và 241 bệnh nhân, nhưng bài nhiều lần viết "240 patients" nên có mâu thuẫn một bệnh nhân giữa bảng và phần chữ.\par\vspace{3pt}
\noindent \textbf{Giao thức đánh giá.} Điểm mạnh nhất của bài là kiểm định chéo theo bệnh nhân: tập huấn luyện, kiểm định và kiểm tra là các bệnh nhân khác nhau và được hoán vị qua toàn bộ bệnh nhân, đồng thời điểm kiểm định (trên bệnh nhân đã biết) được báo cáo tách biệt khỏi điểm kiểm tra (trên bệnh nhân chưa từng thấy) đúng tinh thần tách theo chủ thể; đơn vị phân tích là một đoạn gồm số nhịp cố định, mỗi vòng có lấy mẫu thiếu ngẫu nhiên lớp đa số và điểm số được trọng số hóa để bù mất cân bằng. Tuy vậy cỡ tập kiểm tra mâu thuẫn nội bộ: phần Testing Methodology ghi giữ 5 bệnh nhân mỗi vòng, Mục 4.1 ghi 10 thí nghiệm với 225 bệnh nhân dùng cho huấn luyện (70\%) và kiểm định (30\%) cùng 15 bệnh nhân kiểm tra không chồng lấn giữa các thí nghiệm, còn Mục 4.2 lại quay về 5 bệnh nhân; riêng hai benchmark PVC theo giao thức bài gốc chia tách bệnh nhân rõ ràng (22 trên 22 bệnh nhân với MIT-BIH sau khi bỏ 4 trên 48, và 120 trên 120 bệnh nhân với 5 cơ sở gộp), còn benchmark 8 lớp không nêu rõ cách chia. Bài không nêu dung sai khớp nhịp theo ms vì họ không dò đỉnh mà nhận vị trí nhịp từ chú thích của bác sĩ, đồng thời không có kiểm định thống kê nào (không Wilcoxon, không t-test) kể cả cho so sánh có hay không kênh TDA, và không báo cáo độ lệch chuẩn mà chỉ liệt kê 10 giá trị thô trong Bảng 2.\par\vspace{3pt}
\noindent \textbf{Kết quả chính.} Mọi con số đều là độ chính xác có trọng số chứ không phải F1: phát hiện loạn nhịp nhị phân đạt 98\% trên tập kiểm định và 90\% trên bệnh nhân mới, còn phân loại 13 lớp đạt 97,3\% kiểm định và 80,5\% kiểm tra — khoảng cách 8 điểm và 16,8 điểm chính là thước đo khác biệt cá thể. Thí nghiệm bật và tắt kênh TDA qua 10 lần chạy cho trung bình (theo suy luận tự tính, bài không tính) 0,868 so với 0,820 ở nhiệm vụ phát hiện, chênh +4,8 điểm, và 0,841 so với 0,786 ở nhiệm vụ phân loại, chênh +5,5 điểm; tuy nhiên TDA làm xấu đi ở 4 trên 10 thí nghiệm phát hiện và 1 trên 10 thí nghiệm phân loại, kèm 1 trường hợp hòa. Trên các benchmark (Acc/PPV/Se): phát hiện ngoại tâm thu thất PVC trên MIT-BIH đạt 99,2 / 95,1 / 95,5 so với Li 2018 đạt 99,4 / 93,9 / 98,2, tức Li hơn về accuracy và độ nhạy còn bài này chỉ hơn về PPV; PVC trên 5 cơ sở gộp đạt 96,1 / 97,0 / 96,0 so với Li 2018 đạt 95,6 / 94,1 / 92,7; phân loại 8 lớp trên MIT-BIH đạt 99,0 / 99,0 / 98,5 so với Jun 2017 đạt 99,0 / 98,5 / 97,8, Kiranyaz 2017 (5 lớp) 96,4 / 79,2 / 68,8 và Melgani 2008 (6 lớp) 91,7 / không có / 93,8; chi phí khoảng 10 giờ huấn luyện trên GPU.\par\vspace{3pt}
\noindent \textbf{Điểm yếu.} Bài có mâu thuẫn nội bộ về số bệnh nhân giữ lại để kiểm tra (5 hay 15) và về tổng số bệnh nhân (240 trong phần chữ, 241 khi cộng Bảng 1), đồng thời không có kiểm định thống kê cho tuyên bố TDA cải thiện dù chỉ đưa ra 10 cặp số thô, không tính trung bình, không tính độ lệch chuẩn và TDA thua ở 4 trên 10 thí nghiệm phát hiện. Bài không tái lập được vì không nêu số tham số, kích thước mô hình, độ trễ suy luận, kiến trúc chi tiết từng kênh, số epoch tổng, batch size hay thư viện học sâu, và cũng không nêu độ dài cửa sổ theo giây do cắt theo số nhịp ("as many consecutive heartbeats as wanted") nên không quy đổi được 400 chiều đầu vào bộ tự mã hóa ra giây; lấy mẫu thiếu ngẫu nhiên lớp đa số làm mất thông tin và khiến số liệu phụ thuộc seed nhưng bài không lặp nhiều seed. Bài không dò đỉnh mà nhận vị trí nhịp từ chú thích chuyên gia nên không có dung sai theo ms, không thể coi là hệ thống đầu-cuối và đã giả định giải xong phần khó nhất, trong khi cửa sổ cắt theo nhịp còn phụ thuộc bệnh nhân; benchmark 8 lớp không nêu rõ giao thức chia dữ liệu nên con số 99,0\% không kiểm chứng được và đáng ngờ vì cao hơn hẳn 80,5\% mà chính họ báo cáo với kiểm định chéo theo bệnh nhân, và chính tác giả thừa nhận không công trình nào khác dùng cùng thiết lập kiểm tra nên 90\% và 80,5\% không đối chiếu được với ai.\par\vspace{3pt}
\noindent \textbf{\color{tdtblue}So với mô hình của nhóm.} So sánh được một phần vì cùng lĩnh vực ECG và cùng dùng đồng điều bền vững lọc sublevel-set, nhưng không so được về con số: họ phân loại nhãn của một nhịp đã biết vị trí còn nhóm định vị vị trí đỉnh chưa biết, họ không dò đỉnh nên không có dung sai ±50 ms và không có Se hay PPV theo nghĩa dò nhịp, chỉ số của họ là độ chính xác có trọng số còn của nhóm là macro F1 dò nhịp (hai đại lượng không quy đổi được), và đối tượng là ECG người lớn đơn chuyển đạo không có bài toán khử ECG mẹ. Nhóm hơn họ ở tính tái lập (113.481 tham số, checkpoint 0,48 MB, độ trễ 4,35 ms cho một cửa sổ 4 giây trên CPU, so với chỉ "khoảng 10 giờ trên GeForce GTX"), ở kiểm định thống kê (Wilcoxon bắt cặp p = 0,0000 cho quét trường tiếp nhận và p = 0,7012 cho so sánh kiến trúc, họ không có kiểm định nào), ở tính đầu-cuối vì nhóm dò đỉnh từ tín hiệu thô, và ở tính nhất quán số liệu; ngược lại họ hơn nhóm về quy mô (240 bệnh nhân và khoảng 3,48 triệu nhịp so với 5 sản phụ của ADFECGDB, đây là hạn chế lớn nhất cần nói thẳng trong phần Limitations) và về việc có bảng benchmark đối chiếu với 5 công trình khác theo đúng giao thức gốc. Quan trọng nhất về mặt học thuật, định lý bất biến của họ rằng đường cong Betti của f(t) và f(at) trùng nhau là bằng chứng toán học do chính phe ủng hộ TDA công bố cho thấy TDA cố ý vứt bỏ thông tin thời gian: với họ đó là ưu điểm chống nhịp nhanh và nhịp chậm, còn với bài toán định vị đỉnh của nhóm đó là nhược điểm chí tử vì thứ nhóm cần chính là thông tin thời gian mà bất biến đó xóa đi.\par\vspace{3pt}
\noindent \textbf{\color{tdtblue}Nhóm rút ra gì.} Nên trích dẫn định lý bất biến đường cong Betti làm lập luận trung tâm giải thích vì sao Persistence Image kết hợp CNN 2D đứng chót bảng 16 kiến trúc (27,42 so với 97,14), học cách trình bày điểm kiểm định tách biệt điểm kiểm tra theo bệnh nhân để phơi bày đúng độ lớn khác biệt cá thể (tương ứng với nhóm là 97,43 ± 4,37 khi tách theo bản ghi so với 77,34 ± 28,54 khi chạy xuyên bộ dữ liệu CinC 2013), bắt chước ý tưởng kênh song song bằng cách giữ FetalQRS-TCN làm xương sống và gắn nhánh fSQI dựa trên peak prominence chạy song song, áp dụng mẫu thiết kế tự mã hóa huấn luyện trên cửa sổ chất lượng tốt rồi dùng sai số tái tạo làm cơ sở từ chối trả lời, và dùng dải lọc FIR 0,05–50 Hz của họ làm điểm đối chiếu cho khảo sát 8 dải lọc của nhóm (dải 0,5–100 Hz chỉ đạt 72,17 và dải 1–45 Hz đạt 80,06, cho thấy dải rộng kiểu chẩn đoán không phù hợp cho fQRS thai). Nên tránh việc tuyên bố cải thiện mà không có kiểm định thống kê, việc để mâu thuẫn số liệu giữa các mục (5 hay 15 bệnh nhân, 240 hay 241), việc bỏ qua số tham số và độ trễ suy luận, và việc giả định vị trí đỉnh cho sẵn mà không khai báo minh bạch.\par\vspace{3pt}
\vspace{6pt}
\subsection{Chung và cs. 2020 \normalfont\small (p31)}
\noindent{\footnotesize\itshape Yu-Min Chung, Chuan-Shen Hu, Yu-Lun Lo, Hau-Tieng Wu. "A Persistent Homology Approach to Heart Rate Variability Analysis with an Application to Sleep-Wake Classification". Ban trong file la tien an pham arXiv:1908.06856v2 [eess.SP], 02/05/2020, 38 trang. FILE PDF NAY KHONG IN TEN TAP CHI VA KHONG IN DOI. Ban xuat ban chinh thuc duoc chinh bai p32 trich dan (tai lieu tham khao [37] cua p32): ...}\par\vspace{2pt}
\noindent{\small\color{tdtblue}\textbf{Kết luận: so được một phần}}\par\vspace{4pt}
\noindent \textbf{Họ làm gì.} Bài đặt câu hỏi liệu có thể dùng đồng điều bền vững để định lượng biến thiên nhịp tim (HRV) hay không: nhóm tác giả không đo đỉnh QRS mà lấy chuỗi nhịp tim tức thời (IHR, đơn vị nhịp/phút) đã có sẵn rồi trích đặc trưng tô-pô từ đó. Đóng góp chính là "persistence statistics" — quy mỗi sơ đồ bền vững về 16 con số thống kê rẻ tiền thay cho khoảng cách bottleneck/Wasserstein rất đắt — với ứng dụng minh hoạ là phân loại giai đoạn ngủ (thức/ngủ, REM/NREM, thức-REM-NREM) chỉ từ ECG một đạo trình; đây chính là bài gốc của hướng "TDA cho HRV" mà nhóm đang đối chiếu.\par\vspace{3pt}
\noindent \textbf{Phương pháp.} Dò đỉnh R bằng thuật toán tự động chuẩn (tài liệu [32] của bài), lọc trung vị 5 nhịp để loại artifact (nhịp quá gần hoặc quá xa nhịp trước bị loại hoặc nội suy), tính IHR = 60/(r\_i − r\_\{i−1\}), nội suy cubic bảo toàn hình dạng và lấy mẫu lại ở 4 Hz; sau đó cắt thành epoch 30 giây theo đúng phân đoạn của chuyên gia (loại epoch có dưới 5 đỉnh R), ghép thêm 2 epoch liền trước để được chuỗi 90 giây = 360 mẫu rồi trừ trung vị của chính đoạn 90 giây đó để khử khác biệt giữa các cá thể. Hai phép lọc được chạy song song: sublevel-set filtration trên chính tín hiệu một chiều (chỉ cho P0) và Vietoris-Rips filtration trên ảnh trễ Takens với (p, τ) = (120, 1), tức nhúng vào R\textasciicircum{}120 với p = 120 tương đương 30 giây (cho P0 và P1). Với mỗi sơ đồ bền vững, nhóm tác giả tạo hai tập M = \{(d+b)/2\} và L = \{d−b\}, trên mỗi tập tính 8 thống kê (trung bình, độ lệch chuẩn, độ lệch, độ nhọn, phân vị 25/50/75, entropy bền vững) thành 16 con số, rồi ghép 3 sơ đồ (P0 sublevel, P0 VR, P1 VR) thành vector đặc trưng — theo công thức (17) là 3 × 16 = 48 đặc trưng, nhưng bài không in rõ con số tổng này — cuối cùng chuẩn hoá z-score theo từng đối tượng và huấn luyện SVM (thư viện tính đồng điều bền vững: DIPHA và Ripser).\par\vspace{3pt}
\noindent \textbf{Mô hình.} SVM nhân tuyến tính, dùng hàm fitcsvm mặc định của Matlab 2014b; khi có hơn 2 lớp thì dùng ECOC một-đối-một qua hàm fitcecoc mặc định. Đây không phải mạng nơ-ron nên không có khái niệm "số tham số huấn luyện được" theo nghĩa của nhóm; bài cũng không nêu số vector hỗ trợ và không nêu kích thước mô hình.\par\vspace{3pt}
\noindent \textbf{Dữ liệu.} Bốn cơ sở dữ liệu đa ký giấc ngủ, tất cả đều là người lớn chứ không phải thai nhi: CGMH-training (Bệnh viện Chang Gung, Đài Loan, IRB No. 101-4968A3) gồm 90 người không ngừng thở khi ngủ (AHI < 5) nhưng chỉ giữ 56 người có ít nhất 10\% epoch là thức để tránh mất cân bằng lớp, ECG đạo trình 2 lấy mẫu 200 Hz, mỗi bản ghi tối thiểu 5 giờ, tổng 9.150 epoch thức và 54.547 epoch ngủ; CGMH-validation 27 người thu độc lập cùng phòng lab; DREAMS Subjects Database (Bỉ, DOI 10.5281/zenodo.2650142) 20 người khoẻ mạnh, máy Brainnet, 200 Hz, tối thiểu 7 giờ; và UCDSADB trên PhysioNet (St. Vincent's University Hospital, Dublin) 25 người có ngừng thở khi ngủ ở nhiều mức độ, máy Holter, 128 Hz, tối thiểu 6 giờ, dùng đạo trình ECG thứ nhất. Nhãn giai đoạn ngủ do 2 chuyên gia đọc theo chuẩn AASM 2007, mỗi epoch 30 giây không chồng lấn.\par\vspace{3pt}
\noindent \textbf{Giao thức đánh giá.} Kiểm định chéo cơ sở dữ liệu (cross-database validation) — mạnh hơn cả tách theo chủ thể: huấn luyện trên một cơ sở dữ liệu và kiểm tra trên toàn bộ các cơ sở dữ liệu còn lại, với tuyên bố rõ ràng rằng các tập kiểm định không được dùng để tinh chỉnh tham số mô hình và không loại bỏ đối tượng nào; vì tập huấn luyện và tập kiểm tra là những người khác nhau ở bệnh viện khác nhau nên không có rò rỉ dữ liệu. Đơn vị phân tích là epoch 30 giây; để xử lý mất cân bằng lớp, nhóm tác giả hạ mẫu lớp đa số chỉ trong tập huấn luyện nhưng kiểm tra trên toàn bộ tập kiểm tra. Bảng chính dùng 1 random seed (seed = 1), phụ lục SI lặp lại với 20 seed (1-20) và báo cáo độ lệch chuẩn dưới 1\%; trong phần thảo luận họ cũng phân biệt rõ rằng tách theo chủ thể là nghiêm túc còn giao thức không tách theo chủ thể thì có xu hướng đánh giá cao quá mức, còn về dung sai khớp nhịp thì bài không nêu vì đây không phải bài toán dò đỉnh nên không có bước đối chiếu đỉnh dự đoán với đỉnh thật.\par\vspace{3pt}
\noindent \textbf{Kết quả chính.} Bài toán Thức vs Ngủ, huấn luyện trên CGMH-training (Bảng 1): trên CGMH-validation đạt SE 70,9 ± 16,0\%, SP 78,9 ± 5,4\%, Acc 75,8 ± 4,4\%, PR 38,1 ± 19,6\%, F1 0,452 ± 0,140, AUC 0,824 ± 0,090, Kappa 0,322 ± 0,123; trên DREAMS đạt SE 66,9 ± 16,1\%, Acc 74,7 ± 5,0\%, F1 0,445 ± 0,146, AUC 0,789 ± 0,090; trên UCDSADB (bệnh nhân có ngừng thở khi ngủ) chỉ còn SE 57,6 ± 15,5\%, Acc 70,6 ± 5,4\%, F1 0,407 ± 0,140, AUC 0,702 ± 0,094, Kappa 0,238 ± 0,133; ngay trên chính tập huấn luyện cũng chỉ đạt Acc 75,2 ± 5,4\% và F1 0,438. REM vs NREM trên CGMH-validation (Bảng 2) đạt SE 78,1 ± 17,4\%, SP 79,6 ± 6,5\%, Acc 77,8 ± 8,3\%, F1 0,510 ± 0,160, AUC 0,849 ± 0,108, Kappa 0,393 ± 0,175; bài toán ba lớp Thức-REM-NREM (Bảng 3) đạt Acc 68,3 ± 6,4\% (CGMH-training), 67,6 ± 9,2\% (CGMH-validation), 66,3 ± 6,4\% (DREAMS) và 57,1 ± 7,1\% (UCDSADB), trong đó +P của NREM cao nhất 85,6 ± 16,0\% nhưng +P của Thức chỉ 43,7 ± 16,6\% và +P của REM chỉ 40,0 ± 17,7\%. Chi phí tính toán chỉ được báo cáo một phần: huấn luyện một mô hình SVM mất 5-7 phút trên Windows 7, CPU i5-4570, 32 GB RAM, còn Ripser tính cho một epoch mất 0,06 giây với P1, 1,7 giây với P2 và 106 giây với P3, tức tăng gần 1.800 lần từ chiều 1 lên chiều 3.\par\vspace{3pt}
\noindent \textbf{Điểm yếu.} Bài không nêu dải lọc thông dải cho ECG — toàn bộ khâu tiền xử lý chỉ được dẫn chiếu sang bài [50] — nên khâu quan trọng nhất theo thực nghiệm của nhóm là không tái lập được; bài cũng không nêu số tham số mô hình, độ trễ suy luận hay kích thước mô hình, chỉ có thời gian huấn luyện 5-7 phút và thời gian chạy Ripser. Hiệu năng tuyệt đối rất thấp: F1 phát hiện trạng thái Thức chỉ 0,41-0,45, Kappa 0,24-0,32 (mức đồng thuận yếu), độ chính xác dương tính chỉ 34-38\% nghĩa là cứ 3 epoch báo "Thức" thì 2 sai, và bảng số thô cho thấy FP dao động 126-175 epoch mỗi đối tượng trong khi TP chỉ 76-101; nguyên nhân một phần là họ hạ mẫu để cân bằng lớp khi huấn luyện nhưng kiểm tra trên tập thật mất cân bằng (9.150 so với 54.547) — một lựa chọn trung thực nhưng khiến F1 không thể đặt cạnh các bài dùng tập kiểm tra cân bằng. Bảng chính chỉ dùng 1 random seed; toàn bộ phần tuyên bố "tốt hơn state-of-the-art" là so sánh với số in trong bài khác chứ không chạy lại thuật toán đối thủ trên cùng dữ liệu nên không kiểm soát được điều kiện so sánh (chính họ thừa nhận khó có so sánh trực tiếp); và họ chỉ dùng được đến P1 vì P2/P3 quá đắt (106 giây mỗi epoch).\par\vspace{3pt}
\noindent \textbf{\color{tdtblue}So với mô hình của nhóm.} Không thể so sánh trực tiếp về con số vì đây là bài toán khác hẳn — phân loại giai đoạn ngủ của người lớn từ chuỗi HRV đã có sẵn, đo bằng Acc/Kappa/AUC trên epoch 30 giây — trong khi nhóm dò đỉnh fQRS từ ECG bụng thai nhi và đo bằng F1 khớp nhịp với dung sai ±50 ms; nhưng nếu đọc về mức độ thành công thì F1 cao nhất của họ là 0,510 (REM/NREM) và 0,452 (Thức/Ngủ), so với 97,43 ± 4,37 macro F1 của FetalQRS-TCN, cho thấy TDA trên HRV là công cụ mô tả yếu chứ không phải công cụ đo lường mạnh. Điểm trùng khớp quan trọng nhất là họ dùng đúng sublevel-set filtration trên tín hiệu một chiều mà nhóm đã chứng minh trùng khớp 100\% với peak prominence cổ điển (định lý Elder) trên 2.000 cửa sổ, nghĩa là phát hiện của nhóm áp dụng trực tiếp lên bài này; khác biệt chiến lược là họ tóm tắt sơ đồ bền vững bằng 16 thống kê vô hướng rẻ tiền, còn thí nghiệm thất bại của nhóm dùng Persistence Image + CNN 2D (27,42\%) — cùng một lý thuyết, hai cách vector hoá, và cách của họ rõ ràng hợp lý hơn. Về giao thức thì đối chiếu được một phần: kiểm định chéo cơ sở dữ liệu của họ tương đương phép thử xuyên bộ dữ liệu không tinh chỉnh của nhóm trên CinC 2013 (77,34 ± 28,54), và cả hai đều tụt hiệu năng khi đổi sang tập lạ — họ tụt từ Acc 75,8\% xuống 70,6\% khi gặp bệnh nhân ngừng thở khi ngủ.\par\vspace{3pt}
\noindent \textbf{\color{tdtblue}Nhóm rút ra gì.} Nên áp dụng: trích dẫn bài này (Frontiers in Physiology 2021) làm tiền lệ cho giao thức kiểm định chéo cơ sở dữ liệu không tinh chỉnh mà nhóm đã chạy trên CinC 2013, nhấn mạnh trong phần Methods rằng chính họ tuyên bố trừ bài [50] ra thì không ai báo cáo cross-database; mượn nguyên 16 persistence statistics (trung bình, độ lệch chuẩn, độ lệch, độ nhọn, phân vị 25/50/75, entropy bền vững trên cả hai tập M và L) làm ứng viên đặc trưng cho chỉ số chất lượng tín hiệu thai fSQI vì rẻ, ổn định và không cần CNN; ghi lại chi phí Ripser 0,06 / 1,7 / 106 giây cho P1/P2/P3 làm bằng chứng định lượng biện minh vì sao fSQI chỉ nên dùng H0 sublevel-set và tránh chiều đồng điều từ 2 trở lên trong hệ thời gian thực; và học cách báo cáo mean ± std theo từng đối tượng kèm Kappa và AUC chứ không chỉ Acc. Nên tránh: không được để trống khâu tiền xử lý bằng cách dẫn chiếu sang bài khác — báo cáo của nhóm phải in rõ dải lọc 10-60 Hz, notch 50 Hz và tần số lấy mẫu 250 Hz — đồng thời không lấy số của đối thủ từ bài báo rồi tự tuyên bố "tốt hơn", nếu cần so sánh thì phải chạy lại trên cùng dữ liệu.\par\vspace{3pt}
\vspace{6pt}
\subsection{Dominguez-Monterroza và cs. 2025, PLoS One \normalfont\small (p32)}
\noindent{\footnotesize\itshape Andy Dominguez-Monterroza, Alfonso Mateos Caballero, Antonio Jimenez-Martin (Department of Artificial Intelligence, Universidad Politecnica de Madrid, Tay Ban Nha). "Age-dependent patterns of cardiac complexity unveiled by topological data analysis of pediatric heart rate variability". PLoS One. 2025;20(12):e0337620. doi: 10.1371/journal.pone.0337620. Nhan bai 02/09/2025, chap nhan 11/11/2025, ...}\par\vspace{2pt}
\noindent{\small\color{warnred}\textbf{Kết luận: không so trực tiếp được}}\par\vspace{4pt}
\noindent \textbf{Họ làm gì.} Đây là một nghiên cứu mô tả chứ không phải bài toán học máy: nhóm tác giả nhúng chuỗi khoảng RR của trẻ em từ 1 tháng tuổi đến 17 tuổi vào không gian pha bằng phép nhúng trễ Takens, tính đồng điều bền vững bậc 1 (H1, tức các "vòng lặp" trong quỹ đạo), rút 5 chỉ số mô tả rồi kiểm định thống kê xem chúng có khác nhau giữa 7 nhóm tuổi và có tương quan với các chỉ số HRV cổ điển (SDNN, RMSSD, pNN50) hay không. Kết luận của họ là trẻ càng lớn thì cấu trúc tô-pô của nhịp tim càng gọn và đồng nhất hơn; bài đáng chú ý với nhóm vì trong phần Thảo luận họ tự nêu hướng mở rộng sang HRV thai nhi, tức đang mở cửa cho đúng chủ đề của nhóm.\par\vspace{3pt}
\noindent \textbf{Phương pháp.} Chuỗi RR được lấy sẵn từ PhysioNet (không tự dò đỉnh R) với kiểm soát chất lượng: chỉ nhận bản ghi có tổng artifact dưới 8\% và không có đoạn artifact đơn lẻ nào vượt 20 giây, artifact do chuyên gia tim mạch xác định và sửa, và bài ghi rõ không áp dụng bất kỳ phép lọc số nào ngoài việc sửa artifact. Với mỗi người, họ chọn một đoạn liên tục sạch gồm 3.000 khoảng RR, chuẩn hoá z-score theo từng đối tượng, rồi nhúng trễ Takens v\_i = (RR\_i, RR\_\{i+τ\}, ..., RR\_\{i+(d−1)τ\}) với τ chọn bằng thông tin tương hỗ trung bình (trung vị tối ưu trả về 7-12 nhịp, họ chốt τ = 10 nhịp), còn d thì chốt bằng 3 mặc dù chính họ đo được False Nearest Neighbors tụt dưới 5\% ở d ≈ 6 và tỉ số Cao bão hoà ở d ≈ 9, với lý do "thực hành phổ biến trong HRV phi tuyến" và để dễ trích vòng H1 trong không gian ba chiều. Trên đám mây điểm họ chạy lọc Vietoris-Rips và chỉ lấy chiều đồng điều 1 (thư viện Ripser cho Python), rút 5 chỉ số gồm số đặc trưng bền vững N1, tổng bền vững TP1, bền vững cực đại MP1, bền vững trung bình μ1 và entropy bền vững PE1, đồng thời tính persistence landscape lớp đầu tiên (k = 1) rời rạc hoá trên lưới N = 100 giá trị lọc cách đều thành vector trong R\textasciicircum{}100 để lập ma trận khoảng cách Euclid từng đôi một giữa và trong các nhóm tuổi.\par\vspace{3pt}
\noindent \textbf{Mô hình.} Bài không có mô hình học máy nào: không có mạng nơ-ron, không có bộ phân loại, không có số tham số. Toàn bộ "mô hình" chỉ là Ripser (Python) tính sơ đồ bền vững cùng persistence landscape và các kiểm định thống kê phi tham số, chạy trên Python 3.12 với scipy, scikit-posthocs và statsmodels.\par\vspace{3pt}
\noindent \textbf{Dữ liệu.} Bộ dữ liệu công khai "RR interval time series from healthy subjects (version 1.0.0)" của Irurzun IM, Garavaglia L, Defeo MM, Thomas Mailland J trên PhysioNet 2021, chỉ giữ đối tượng từ 1 tháng tuổi đến 17 tuổi nên còn lại 127 người; ghi ECG Holter 24 giờ, 3 đạo trình, máy DMS300-7 / DMS300-3A / Galix, tần số lấy mẫu 512-1024 Hz, thu với phê duyệt đạo đức của National University of La Plata, Argentina, người tham gia không dùng thuốc và có điện tâm đồ bình thường. Đoạn phân tích là 3.000 khoảng RR mỗi người, tương đương khoảng 25-43 phút tuỳ nhịp tim (từ khoảng 120 nhịp/phút ở trẻ nhỏ đến khoảng 70 nhịp/phút ở thiếu niên), phân tầng thành 7 nhóm tuổi (Bảng 1): sơ sinh 0-1 tháng n = 8; nhũ nhi sớm 1-5 tháng n = 33; nhũ nhi muộn 6-11 tháng n = 29; tập đi 1-2 tuổi n = 22; mẫu giáo 3-5 tuổi n = 10; học đường 6-11 tuổi n = 15; thiếu niên 12-17 tuổi n = 10; tổng n = 127.\par\vspace{3pt}
\noindent \textbf{Giao thức đánh giá.} Bài không có giao thức đánh giá học máy: không chia train/test, không tách theo chủ thể, không tách theo bản ghi, không chia ngẫu nhiên — đơn giản vì không có mô hình dự đoán nào để đánh giá. Thay vào đó là kiểm định thống kê mức nhóm: Kruskal-Wallis H test cho khác biệt giữa các nhóm tuổi, hậu kiểm Dunn có hiệu chỉnh Holm-Bonferroni cho đa so sánh, và hệ số tương quan Pearson giữa chỉ số tô-pô với chỉ số HRV cổ điển (tất cả đều được log-hoá trước khi tính tương quan), ngưỡng ý nghĩa hai phía p < 0,05. Đơn vị phân tích là một người bằng một điểm dữ liệu (mỗi người cho đúng một đoạn 3.000 nhịp và một bộ 5 chỉ số), do đó không có rò rỉ dữ liệu kiểu chia ngẫu nhiên trên cùng bệnh nhân nhưng cũng không có bất kỳ bằng chứng nào về khả năng khái quát sang đối tượng mới; về dung sai khớp nhịp thì bài không nêu, vì đây không phải bài toán dò đỉnh và chuỗi RR được lấy sẵn với artifact do chuyên gia tim mạch sửa tay.\par\vspace{3pt}
\noindent \textbf{Kết quả chính.} Chỉ số HRV cổ điển (Bảng 2, trung vị [khoảng tứ phân vị]): Mean RR tăng đơn điệu từ 404 ms [395-417] ở sơ sinh lên 674 ms [669-680] ở thiếu niên, SDNN từ 62,9 ms [57-66] ở sơ sinh xuống thấp nhất 36,7 ms ở nhũ nhi muộn rồi lên 81,4 ms [78-85] ở thiếu niên, pNN50 từ 4,2\% lên 14,1\% (cả ba đều p < 0,001), riêng RMSSD không có ý nghĩa thống kê với p = 0,115 (bảng ghi 0,11). Chỉ số tô-pô theo Kruskal-Wallis (Bảng 3): N1 p < 0,001 (khác biệt giữa sơ sinh với mẫu giáo và tập đi), TP1 p < 0,001 (thiếu niên so với nhũ nhi sớm; sơ sinh so với học đường), MP1 p = 0,0048 (sơ sinh so với thiếu niên và học đường), μ1 p < 0,001 (sơ sinh so với mọi nhóm lớn hơn trừ tập đi) và PE1 p < 0,001 (sơ sinh so với mẫu giáo, tập đi). Tương quan Pearson mạnh nhất (log-hoá, p đã hiệu chỉnh Holm): μ1 với ln(SDNN) r = −0,75, μ1 với ln(Tuổi) r = −0,67, μ1 với ln(MeanRR) r = −0,65 (cả ba đều p < 10\textasciicircum{}−10), TP1 với ln(SDNN) r = −0,57, N1 với ln(pNN50) r = 0,49, N1 với ln(Tuổi) r = 0,41, N1 với ln(SDNN) r = 0,40 và PE1 với các chỉ số HRV cổ điển r từ 0,32 đến 0,43, ngưỡng lọc của hình là |r| > 0,52 cho 6 tương quan mạnh nhất; chi phí tính toán thì bài không nêu — không có thời gian chạy, cấu hình máy, số tham số hay độ trễ.\par\vspace{3pt}
\noindent \textbf{Điểm yếu.} Bài không có bài toán dự đoán nào — không train/test, không mô hình, không độ chính xác/F1/AUC — nên chỉ chứng minh các chỉ số tô-pô khác nhau giữa các nhóm tuổi chứ không chứng minh chúng dự đoán được tuổi hay bệnh lý của một cá thể mới, đó là khoảng cách rất lớn giữa "có ý nghĩa thống kê" và "dùng được trong lâm sàng". Nghiêm trọng hơn là mâu thuẫn nội tại: phần Tóm tắt và Thảo luận viết rằng sơ sinh và nhũ nhi có số lượng và cường độ đặc trưng bền vững cao hơn (N1 cao), trong khi mục Kết quả 4.2 viết ngược lại rằng sơ sinh có N1 và PE1 thấp hơn có ý nghĩa nhưng MP1 và μ1 cao hơn — vì vậy phải đọc Bảng 3 và mục 4.2, tuyệt đối không trích dẫn theo Tóm tắt. Các hạn chế còn lại: mất cân bằng nhóm nghiêm trọng (sơ sinh n = 8, mẫu giáo n = 10, thiếu niên n = 10 so với nhũ nhi sớm n = 33) khiến Kruskal-Wallis và Dunn có lực thống kê rất yếu và dễ ra kết quả không lặp lại được; tự mâu thuẫn về tham số nhúng khi đo được d ≈ 6 (FNN) và d ≈ 9 (tiêu chí Cao) nhưng vẫn chọn d = 3, nếu chiều nhúng sai thì cấu trúc tô-pô tìm được có thể chỉ là ảnh giả của phép chiếu; đoạn phân tích được định nghĩa theo số nhịp (3.000 nhịp) chứ không theo thời gian nên thời lượng thật biến thiên 25-43 phút giữa các nhóm tuổi — mà tuổi lại chính là biến kết quả và độ dài chuỗi ảnh hưởng trực tiếp đến số vòng H1 tìm được — nhưng bài không kiểm soát nhiễu gây rối này; ngoài ra bài chỉ lấy H1 và bỏ qua H0, trong khi H0 sublevel-set mới là thứ rẻ và ổn định nhất (việc không nêu chi phí tính toán, dung sai hay số tham số thì chấp nhận được vì bài không có mô hình).\par\vspace{3pt}
\noindent \textbf{\color{tdtblue}So với mô hình của nhóm.} Không thể so sánh trực tiếp, lý do rất đơn giản: bài không có một con số hiệu năng nào để đặt cạnh macro F1 97,43 ± 4,37 của FetalQRS-TCN vì họ không dò đỉnh, không phân loại, không có tập kiểm tra, toàn bộ kết quả là p-value và hệ số tương quan mức nhóm; nhiệm vụ cũng khác hẳn — mô tả sự trưởng thành hệ thần kinh tự chủ ở trẻ em từ chuỗi RR sạch đã được chuyên gia sửa tay, trong khi nhóm làm việc với ECG bụng đơn kênh đầy nhiễu và phải tự tìm đỉnh fQRS bị ECG mẹ lấn át. Điểm đối chiếu có ích nhất và cũng quan trọng nhất là ngữ cảnh thời gian: họ cần 3.000 nhịp (25-43 phút) và trễ τ = 10 nhịp (khoảng 5-10 giây) để TDA có ý nghĩa, còn nhóm đã phát hiện cửa sổ 300 ms ngắn hơn một chu kỳ tim thai (\textasciitilde{}430 ms) nên không thể chứa "vòng nhịp" — hai kết quả độc lập cùng nói rằng TDA trên nhịp tim cần ngữ cảnh hàng chục giây đến hàng chục phút chứ không phải hàng trăm mili giây, đây là bằng chứng bên ngoài ủng hộ quyết định chuyển hướng của nhóm (dùng TDA làm fSQI chứ không làm bộ dò đỉnh). Một chi tiết chiến lược: cuối phần Thảo luận họ viết rõ ý định mở rộng sang HRV thai nhi và thừa nhận phải điều chỉnh tiền xử lý cùng tham số nhúng vì nhịp tim thai cao hơn và đặc tính nhiễu khác, tức chính họ nhận mình chưa làm được điều đó.\par\vspace{3pt}
\noindent \textbf{\color{tdtblue}Nhóm rút ra gì.} Nên áp dụng: trích dẫn bài này (PLoS One tháng 12/2025, Universidad Politecnica de Madrid) làm khoảng trống nghiên cứu cho phần Giới thiệu vì chính họ tuyên bố mở rộng TDA sang HRV thai nhi là việc chưa ai làm — đúng chỗ đứng của nhóm, với dữ liệu ADFECGDB sẵn có; mượn cách trình bày minh bạch về tham số nhúng (báo cáo cả kết quả AMI với τ tối ưu 7-12 nhịp, FNN d ≈ 6, tiêu chí Cao d ≈ 9 rồi mới nói chọn gì và vì sao, dù lựa chọn cuối cùng mâu thuẫn với số đo); ghi nhận τ = 10 nhịp và 3.000 nhịp làm mốc ngữ cảnh tối thiểu của TDA trên nhịp tim để biện minh vì sao TDA không dùng được ở quy mô 300 ms; và học cách dùng kiểm định phi tham số Kruskal-Wallis kèm hậu kiểm Dunn có hiệu chỉnh Holm-Bonferroni khi so sánh nhiều nhóm. Nên tránh: không viết bài kiểu này nếu nhắm tạp chí kỹ thuật Q1 vì không tập kiểm tra, không mô hình, không con số hiệu năng thì người phản biện sẽ từ chối ngay; không để Tóm tắt mâu thuẫn với Bảng kết quả mà phải đối chiếu từng con số trong Tóm tắt với từng ô trong Bảng trước khi nộp; và không định nghĩa độ dài đoạn phân tích theo số nhịp khi biến kết quả liên quan đến nhịp tim (với thai nhi nhịp khoảng 140/phút thì sai lầm này còn nghiêm trọng hơn), cửa sổ phải cố định theo thời gian (4 giây như hiện tại) và phải nói rõ.\par\vspace{3pt}
\vspace{6pt}
\subsection{Karan và Kaygun 2021 \normalfont\small (p33)}
\noindent{\footnotesize\itshape Alperen Karan, Atabey Kaygun (Department of Mathematical Engineering, Istanbul Technical University, 34467 Istanbul, Tho Nhi Ky). "Time Series Classification via Topological Data Analysis". arXiv:2102.01956v2 [stat.ML], 11/06/2021, 31 trang. Email: karana@itu.edu.tr, kaygun@itu.edu.tr. FILE PDF NAY KHONG IN TEN TAP CHI VA KHONG IN DOI - chi la ban tien an pham arXiv.}\par\vspace{2pt}
\noindent{\small\color{tdtblue}\textbf{Kết luận: so được một phần}}\par\vspace{4pt}
\noindent \textbf{Họ làm gì.} Bài xây dựng một quy trình TDA tổng quát để phân loại chuỗi thời gian đơn biến, rồi áp dụng vào bài toán nhận diện căng thẳng từ tín hiệu sinh lý. Ý tưởng kỹ thuật trung tâm là subwindowing (cửa sổ con): thay vì tính đồng điều bền vững trực tiếp trên một cửa sổ 60 giây vốn rất đắt và rất nhạy nhiễu, họ cắt cửa sổ 60 giây thành nhiều cửa sổ con 4 giây, tính đặc trưng tô-pô trên từng cửa sổ con rồi lấy trung bình và độ lệch chuẩn làm đặc trưng của cả cửa sổ, nhờ vậy nhiễu bị "nhốt" trong vài cửa sổ con và bị làm loãng, đồng thời chi phí tính toán trở nên độc lập với độ dài cửa sổ (họ chứng minh bằng công thức rằng khi cửa sổ trượt một bước thì trung bình và độ lệch chuẩn mới chỉ cần cập nhật từ giá trị cũ cộng một phần tử vào và một phần tử ra).\par\vspace{3pt}
\noindent \textbf{Phương pháp.} Tiền xử lý: tín hiệu trên 100 Hz (kênh ngực WESAD 700 Hz) được hạ mẫu về 100 Hz bằng cách lấy mỗi phần tử thứ 7, DriveDB được nâng từ 15,5 Hz lên 496 Hz bằng nội suy tuyến tính rồi lấy mỗi phần tử thứ 31 để được 16 Hz, ba trục gia tốc x, y, z được lấy trung bình thành một chuỗi, và bài không nêu bất kỳ phép lọc thông dải nào. Cửa sổ trượt 60 giây với bước dịch 2 giây, trong mỗi cửa sổ có các cửa sổ con 4 giây bước dịch 2 giây — chọn 4 giây vì đủ dài để chứa ít nhất một chu kỳ của tín hiệu tựa tuần hoàn (để vòng lặp hiện ra trong ảnh nhúng trễ) nhưng đủ ngắn để tính đồng điều bền vững còn khả thi — và mỗi cửa sổ con được nhúng trễ với 4 chiều khác nhau là 0,5·fs, 1·fs, 1,5·fs, 2·fs (ví dụ fs = 100 Hz thì chiều nhúng là \{50, 100, 150, 200\}). Trên mỗi đám mây điểm họ chạy lọc Rips đến chiều tối đa 1 để được H0 và H1 (thư viện Ripser), cộng thêm 2 sơ đồ chiều 0 từ upper level set và lower level set của chính cửa sổ con (thư viện Dionysus-2), tổng cộng 6 sơ đồ mỗi cửa sổ con, và trên mỗi lớp đồng điều rút 7 đặc trưng (khoảng cách Wasserstein W1 đến đường chéo, khoảng cách bottleneck W∞ đến đường chéo, entropy bền vững, chuẩn L1 và L2 của đường cong Betti, chuẩn L1 và L2 của persistence landscape lớp đầu) — bài ghi rõ tổng cộng 70 đặc trưng mỗi cửa sổ con; đặc trưng của cửa sổ là trung bình cộng độ lệch chuẩn của các cửa sổ con, sau khi loại đặc trưng hằng số và đặc trưng có tương quan trên 0,9 rồi chuẩn hoá về [0, 1] trên cả tập huấn luyện lẫn tập kiểm tra.\par\vspace{3pt}
\noindent \textbf{Mô hình.} Bốn bộ phân loại cổ điển từ scikit-learn: Random Forest và AdaBoost Decision Tree với 100 cây và độ sâu tối đa 5; Support Vector Classifier nhân tuyến tính với tham số điều chuẩn C = 0,1; và Linear Discriminant Analysis, tất cả đặt random\_state = 0 để tái lập. Bài không nêu tổng số tham số của bất kỳ mô hình nào, không nêu kích thước mô hình và không nêu độ trễ suy luận.\par\vspace{3pt}
\noindent \textbf{Dữ liệu.} Ba bộ dữ liệu: bộ tổng hợp mô phỏng bằng thư viện NeuroKit2 cho 20 người (RESP và ECG sinh ở 50 Hz, mỗi người 120 giây không stress cộng 120 giây stress, với ba thí nghiệm là tăng nhịp thở từ nền 15 rpm lên 16-20 rpm, tăng nhịp tim từ nền 70 bpm lên 71-75 bpm, hoặc giữ nhịp tim 70 bpm nhưng tăng độ lệch chuẩn từ 1 lên 2, 3, 4, 5, lặp lại với 2 mức nhiễu 0,1 và 0,3); WESAD gồm 15 người với 3 trạng thái (baseline khoảng 20 phút, stress bằng Trier Social Stress Test khoảng 10 phút, amusement khoảng 7 phút), thiết bị ngực RespiBAN ghi ACC, RESP, ECG, EDA, EMG, TEMP đều ở 700 Hz và thiết bị cổ tay Empatica E4 ghi ACC 32 Hz, BVP 64 Hz, EDA 4 Hz, TEMP 4 Hz. DriveDB trên PhysioNet ghi 17 người lái xe (nghỉ - cao tốc - trong phố, mỗi phiên 60-90 phút) với ECG, EMG, GSR chân và tay, HR, RESP tất cả ở 15,5 Hz, nhưng chỉ 9 người được phân tích vì 8 người còn lại thiếu hoặc không đọc được mốc đánh dấu chuyển trạng thái.\par\vspace{3pt}
\noindent \textbf{Giao thức đánh giá.} Tách theo chủ thể (leave-one-subject-out cross-validation) cho tất cả các bộ dữ liệu — huấn luyện trên tất cả mọi người trừ một, kiểm tra trên người còn lại, rồi lấy trung bình — với lý do được nêu nguyên văn: vì cửa sổ trượt chồng lấn rất nhiều (60 giây, dịch 2 giây) nên chia ngẫu nhiên sẽ gây rò rỉ dữ liệu, và họ khẳng định cách chia này không gây rò rỉ ngay cả khi cửa sổ chồng lấn cao. Quan trọng nhất, họ cố ý chạy thêm một giao thức thứ hai để đo lường mức rò rỉ: kiểm định chéo nội bộ từng đối tượng, tức với mỗi người thì chia mỗi trạng thái làm đôi, dùng nửa đầu dự đoán nửa sau và ngược lại — đúng kiểu chia ngẫu nhiên trên cùng bệnh nhân — rồi so sánh hai giao thức bằng kiểm định t bắt cặp lặp lại; nhãn là trạng thái của cả cửa sổ 60 giây, còn dung sai khớp nhịp thì bài không nêu vì đây không phải bài toán dò đỉnh. Để đối chiếu với bản gốc: bài WESAD gốc cũng dùng cửa sổ 60 giây nên họ chọn 60 giây, còn bài DriveDB gốc dùng cửa sổ 300 giây không chồng lấn với leave-one-out trộn lẫn cửa sổ của mọi người (tức có rò rỉ theo tiêu chuẩn của nhóm) trên 24 người và đạt khoảng 97\%.\par\vspace{3pt}
\noindent \textbf{Kết quả chính.} WESAD ba lớp (Bảng 2, tách theo chủ thể, cửa sổ 60 giây): gộp tất cả kênh với SVC và dùng tất cả sơ đồ đạt 81,35\% độ chính xác, F1 73,44\% (bản gốc WESAD 79,57\%); riêng từng kênh thì ECG 76,27\% (nhúng trễ) và 72,72\% (tất cả sơ đồ) so với bản gốc 66,29\%, ACC ngực 74,47\% so với bản gốc 56,56\% (hơn 17 điểm), RESP 75,57\% so với 72,37\%, còn TEMP ngực chỉ 48,92-54,75\% (kém bản gốc 55,68\%); WESAD hai lớp (Bảng 3) gộp tất cả đạt 94,46\% với F1 93,26\% (bản gốc 92,28\%), riêng ECG 88,70\% so với 85,44\% và riêng ACC ngực 87,69\% so với 73,87\%. DriveDB ba lớp (Bảng 4) tốt nhất là RESP 85,81\% với F1 79,68\%, trong khi ECG chỉ đạt 53,32-58,49\% (gần mức ngẫu nhiên của bài ba lớp) và HR 59,30-60,61\%; DriveDB hai lớp (Bảng 5) tốt nhất là RESP với level-set 98,07\% và F1 97,97\%, gộp tất cả 96,24\%, ECG chỉ 65,74-71,11\%; quét độ dài cửa sổ trên WESAD với SVC và cửa sổ con 4 giây cho 75,17\% ở 10 giây, 76,63\% ở 20 giây, 78,36\% ở 30 giây, 81,35\% ở 60 giây và 85,01\% ở 120 giây (cửa sổ con 3 giây: 74,38 / 77,92 / 79,08 / 80,67 / 85,18; cửa sổ con 5 giây: 75,39 / 77,95 / 80,07 / 81,25 / 84,29), mở rộng đến 300 giây đạt 89,86\% ba lớp và 96,42\% hai lớp trên WESAD, 91,19\% và 96,96\% trên DriveDB. Con số quan trọng nhất với nhóm là về rò rỉ dữ liệu: trên WESAD ba lớp, mô hình nội bộ từng đối tượng đạt trung bình 96\% trong khi mô hình tách theo chủ thể chỉ đạt 81\%, khác biệt có ý nghĩa với p < 0,05 (kiểm định t bắt cặp lặp lại), tức chia ngẫu nhiên trên cùng đối tượng thổi phồng kết quả lên khoảng 15 điểm phần trăm — với bài hai lớp khác biệt không đạt p < 0,05 và họ quy cho hiệu ứng trần, còn trên DriveDB cả hai bài toán đều không đạt ý nghĩa và họ quy cho cỡ mẫu nhỏ 9 người; ghi chú thêm là nếu chỉ dùng upper level set của tất cả tín hiệu thì còn 61 đặc trưng mà vẫn đạt 79,61\% (ba lớp) và 92,47\% (hai lớp), còn chi phí tính toán thì bài không nêu con số nào (không có số tham số, độ trễ mili giây, kích thước mô hình hay cấu hình máy) mà chỉ có lập luận định tính rằng subwindowing làm chi phí độc lập với độ dài cửa sổ.\par\vspace{3pt}
\noindent \textbf{Điểm yếu.} Con số đẹp nhất là kết quả gộp đa kênh chứ không phải đơn kênh: 94,46\% và 81,35\% trên WESAD là gộp toàn bộ 10 kênh ngực và cổ tay (ACC, RESP, ECG, EDA, EMG, TEMP, BVP), trong khi riêng ECG một kênh chỉ đạt 88,70\% (hai lớp) và 76,27\% (ba lớp), còn trên DriveDB ECG một kênh chỉ 65,74\% và 53,32\% tức gần như vô dụng — ai trích dẫn "98,07\%" mà không nói rõ đó là tín hiệu hô hấp (RESP) chứ không phải ECG là trích dẫn sai. Có một lỗi phương pháp xử lý tín hiệu thực sự: hạ mẫu 700 Hz xuống 100 Hz bằng cách lấy mỗi phần tử thứ 7 mà không lọc chống chồng phổ trước, khiến mọi thành phần trên 50 Hz bị gấp phổ vào dải quan tâm — nghiêm trọng với ECG vốn có năng lượng đáng kể trên 50 Hz — và bài cũng không nêu bất kỳ dải lọc thông dải nào. Các hạn chế còn lại: thiên lệch chọn mẫu trên DriveDB (ghi 17 người nhưng chỉ phân tích 9 vì mốc chuyển trạng thái thiếu hoặc không đọc được, và chính họ thừa nhận kết quả không đạt ý nghĩa thống kê là do cỡ mẫu nhỏ); cỡ mẫu nhỏ nói chung (WESAD 15 người, DriveDB 9 người) khiến tách theo chủ thể trên 9 fold có độ tin cậy thống kê rất thấp; không nêu số tham số, độ trễ hay kích thước mô hình nên không đánh giá được khả năng triển khai thời gian thực dù họ tuyên bố ưu thế về chi phí; bài chỉ là tiền ấn phẩm arXiv, không có tên tạp chí và không có DOI nên không rõ đã qua bình duyệt hay chưa; phần so sánh với "Original Findings" là đọc số từ bài gốc chứ không chạy lại thuật toán gốc trên cùng phân chia, lặp lại đúng vấn đề của p31; và cửa sổ càng dài thì độ chính xác càng cao (10 giây 75\% lên 300 giây 96\%) nhưng số cửa sổ độc lập càng ít và độ trễ phát hiện càng lớn, điều họ có nhắc đến nhưng không định lượng.\par\vspace{3pt}
\noindent \textbf{\color{tdtblue}So với mô hình của nhóm.} Không thể so sánh trực tiếp về con số: đây là bài toán phân loại trạng thái căng thẳng trên cửa sổ 60 giây của người lớn, đo bằng độ chính xác và F1 ở mức cửa sổ, còn FetalQRS-TCN là bài toán dò đỉnh fQRS trên ECG bụng thai nhi đo bằng F1 khớp nhịp với dung sai ±50 ms, không có điểm chung nào về nhãn hay dung sai. Nhưng về giao thức thì đối chiếu được một phần và rất có giá trị: tách theo chủ thể của họ tương đương tách theo bản ghi và tách theo chủ thể của nhóm, và đây là bài duy nhất trong ba bài chủ động đo tác hại của rò rỉ dữ liệu bằng cách chạy song song hai giao thức trên cùng dữ liệu — 96\% nội bộ đối tượng so với 81\% tách theo chủ thể ở bài ba lớp, chênh 15 điểm, p < 0,05 — một bằng chứng định lượng từ bên thứ ba hoàn toàn phù hợp với phát hiện nội bộ của nhóm rằng repo fECG-TDA cũ báo 91\% trong khi tách theo chủ thể thật chỉ đạt 59\% (chênh 32 điểm). Về hiệu năng tuyệt đối thì không thể đặt cạnh nhau, nhưng đáng chú ý là ngay cả khi gộp 10 kênh và dùng cửa sổ 300 giây họ vẫn chỉ đạt 96,42\% trên bài hai lớp, trong khi FetalQRS-TCN đạt 99,18 macro F1 trên kênh lâm sàng tốt nhất với cửa sổ chỉ 4 giây; ngoài ra lập luận của họ rằng cửa sổ con phải đủ dài để chứa ít nhất một chu kỳ của tín hiệu tựa tuần hoàn chính là lập luận nhóm đã dùng để bác bỏ cửa sổ 300 ms (ngắn hơn chu kỳ tim thai khoảng 430 ms).\par\vspace{3pt}
\noindent \textbf{\color{tdtblue}Nhóm rút ra gì.} Nên áp dụng: trích dẫn cặp số 96\% (nội bộ đối tượng) so với 81\% (tách theo chủ thể), p < 0,05, WESAD ba lớp làm viện dẫn tốt nhất trong cả ba bài để biện minh vì sao nhóm bắt buộc dùng tách theo bản ghi và tách theo chủ thể, và vì sao con số 91\% trong README cũ là không hợp lệ — nên đưa thẳng vào phần Methods hoặc Discussion; mượn kỹ thuật subwindowing kèm công thức cập nhật tăng dần μ\_\{N+1\} = μ\_N + (f\_\{M+1\} − f\_1)/M cho fSQI để chi phí không tăng theo độ dài cửa sổ, phù hợp với ràng buộc thời gian thực 4,35 ms của nhóm; lấy nguyên tắc cửa sổ con phải chứa ít nhất một chu kỳ tín hiệu (thai nhi chu kỳ khoảng 430 ms nên cửa sổ con tối thiểu nên từ 0,5 đến 1 giây trở lên) để chính thức bác bỏ cửa sổ 300 ms; và học cách trình bày bảng quét độ dài cửa sổ (10/20/30/60/120/300 giây) song song với bảng quét cửa sổ con (3/4/5 giây) như nhóm đã có với trường tiếp nhận (172/748/1.516/3.052 ms), đồng thời tách rõ kết quả đơn kênh và kết quả gộp kênh trong cùng một bảng. Nên tránh: tuyệt đối không hạ mẫu bằng cách lấy mỗi phần tử thứ N mà không lọc chống chồng phổ (nhóm đang lấy mẫu lại về 250 Hz nên phải dùng bộ lọc chống chồng phổ đúng chuẩn và ghi rõ trong bài); không để tiêu đề hay tóm tắt ngụ ý "đơn kênh" trong khi con số tốt nhất là gộp nhiều kênh, và nếu báo cáo 99,18 trên kênh tốt nhất thì phải nói rõ đó là kết quả chọn kênh chứ không phải kết quả trung bình (nhóm cũng nên tách rõ tách theo bản ghi toàn bộ 4 kênh đạt 97,43 với tách theo bản ghi trên kênh lâm sàng tốt nhất đạt 99,18); và không tuyên bố ưu thế về chi phí tính toán mà không đưa số, vì nhóm đã có 113.481 tham số, 0,48 MB, 4,35 ms — lợi thế rõ ràng so với cả ba bài, không bài nào nêu đủ ba con số này.\par\vspace{3pt}
\vspace{6pt}
\clearpage
\end{document}